% Things to do:
% - fix chapter-section references "ch. N, §M", as when you click on "N" you are redirected to the beginning of chapter N, but it should be better to have it behave as the usual "sec. N.M" links
% - make the table of contents and the chapter/section fonts the same as Dieudonné's
% - the asterisk of a footnote is usually before a dot/comma, not after
% - add space before and after arguments of operators like cos/sin as in the original book
% - in sections 2.3, 4.1 there are multiple footnotes with the same (*) mark on the same page
% - the integrals in text mode are not exactly the same as those in math mode, while Diuedonné wrote them with the same symbol at the same height
% - substitute all divergences with \divergence
% - add space between the end of a section and the beginning of another
% - the tilde in the equations that use \widetilde sould only cover the top of the elongated "SS" of the integral, it should not cover the white space on the left

\documentclass[12pt,a4paper,leqno,twoside]{extbook}

\usepackage{adjustbox}
\usepackage{amsmath,amssymb}
\usepackage[english]{babel}
\usepackage{bookmark}
\usepackage[autostyle=true]{csquotes} % otherwise babel complains
\usepackage[shortlabels]{enumitem}
\usepackage{etoolbox}
\usepackage[default]{fontsetup} % New Computer Modern book weight
\usepackage[bottom,hang,splitrule]{footmisc}
\usepackage{footnotebackref}
\usepackage[hmargin=2.5cm,vmargin=2cm,twoside=false]{geometry}
\usepackage{indentfirst}
\usepackage{mathtools}
\usepackage[protrusion=true,expansion=true]{microtype}
\usepackage[a-1a]{pdfx}
\usepackage{realscripts} % for OpenType super/subscripts
\usepackage[tiny,compact,explicit]{titlesec}
\usepackage{todonotes}

\usepackage[backend=biber,maxnames=2,mcite,sorting=custom]{biblatex}
\addbibresource{bibliography.bib}
\DeclareFieldFormat{pages}{p.~#1} % Dieudonné used "p." instead of "pp." even for multiple pages
\DeclareSortingScheme{custom}{\sort{\field{entrykey}}}
\DefineBibliographyStrings{english}{
    andothers = {\mkbibemph{et\addabbrvspace al\adddot}} % "et al." in italics
}
\renewcommand*{\mkbibnamefamily}[1]{\textsc{#1}} % surnames in small caps
\renewcommand*{\newunitpunct}{\addcomma\space} % comma instead of dot between authors and title

\usepackage{chngcntr}
\counterwithout{equation}{chapter}

\usepackage{hyperref}
\hypersetup{colorlinks,allcolors=blue}

\usepackage{fancyhdr}
\pagestyle{fancy}
\fancyhf{} % clear all header and footer fields
\fancyhead[EL]{\thepage} % odd pages, left header
\fancyhead[OR]{\thepage} % even pages, right header
\titleformat{\chapter}[display]{\centering\scshape}{\bfseries\chaptertitlename~\thechapter}{0pt}{#1}
\titleformat{\section}[block]{\centering\bfseries}{\S\arabic{section}. #1.}{0em}{}
\renewcommand{\chaptermark}[1]{\markboth{#1}{#1}} % a hack needed for fancyhdr to display just the chapter name
\renewcommand{\headrulewidth}{0pt} % remove all header rules

\usepackage{perpage}
\MakePerPage{footnote}
\renewcommand*{\thefootnote}{\tiny{($\ast$)}}

\usepackage{tocloft}
\makeatletter
    \renewcommand{\cftchappresnum}{Chapter~} % prefix before number
    \renewcommand{\cftchapaftersnum}{:~} % text after number
    \setlength{\cftchapnumwidth}{8em} % wrapped lines align with the chapter title
\makeatother
\renewcommand{\cftchapfont}{\bfseries\scshape}
\renewcommand{\cftchapleader}{\cftdotfill{\cftdotsep}}

\addtolength{\footnotesep}{1em} % add vertical space between footnotes
\renewcommand\footnotemargin{\parindent} % footnote text aligned with text margin
\renewcommand*{\footnotelayout}{\normalsize}

\renewcommand{\arraystretch}{1.2}
\renewcommand{\thechapter}{\Roman{chapter}}
\renewcommand{\thesection}{§\arabic{section}}

\newcommand{\mytodo}[1]{\todo[inline]{TODO: #1}}
\newcommand{\supth}{\texorpdfstring{\textsuperscript{th}}{th}}
\newcommand{\sub}[1]{\texorpdfstring{\textsubscript{#1}}{#1}}

\DeclareMathOperator*{\Liminf}{lim. \! inf \,}
\DeclareMathOperator*{\Limsup}{lim. \! sup \,}
\DeclarePairedDelimiter{\norm}{\lVert}{\rVert}
\DeclarePairedDelimiter{\ip}{\langle}{\rangle}
\newcommand{\B}{\mathcal{B}}
\newcommand{\Cont}{\mathcal{C}}
\newcommand{\C}{\mathbb{C}}
\newcommand{\R}{\mathbb{R}}
\newcommand{\N}{\mathbb{N}}
\newcommand{\Sph}{\mathbb{S}}
\newcommand{\Card}{\operatorname{Card}}
\newcommand{\Dch}{\operatorname{Dch}}
\newcommand{\diam}{\operatorname{diam}}
\newcommand{\dom}{\operatorname{dom}}
\newcommand{\grad}{\operatorname{grad}}
\newcommand{\spec}{\operatorname{Sp}}
\newcommand{\vol}{\operatorname{vol}}
\newcommand{\Ker}{\operatorname{Ker}}
\newcommand{\der}[2]{\frac{d #1}{d #2}}
\newcommand{\pder}[2]{\frac{\partial #1}{\partial #2}}
\newcommand{\divergence}[1]{\left(\pder{#1}{x}\right)^2 + \left(\pder{#1}{y}\right)^2 + \left(\pder{#1}{z}\right)^2}
\newcommand{\Iint}{\displaystyle\iint}
\newcommand{\Iiint}{\displaystyle\iiint}

% for Pringsheim's notation in equation (F); "pfrac" stands for "Pringsheim's fraction"
\newcommand{\pfrac}[2]{%
    \begin{array}{@{} c @{}}%
        \multicolumn{1}{c|}{#1} \\%
        \hline%
        \multicolumn{1}{|c}{#2}%
    \end{array}%
}

\allowdisplaybreaks
\frenchspacing
\raggedbottom % prevent huge vertical gaps between paragraphs
\sloppy % do not overfill

\title{History of Functional Analysis}
\author{Jean Dieudonné}
\date{1981}

\begin{document}

\maketitle

\tableofcontents

\fancyhead[C]{\textsc{\leftmark}} % header for the introduction (it's different from the headers of the chapters)

\chapter*{Introduction}
\markboth{Introduction}{Introduction} % otherwise it shows "CONTENTS" and the names of the each chapter becomes the name of its previous chapter
\addcontentsline{toc}{chapter}{\protect\numberline{}Introduction}

One may give different definitions of ``Functional Analysis''. Its name name might suggest that it contains all parts of mathematics which deal with functions, but that would practically mean \emph{all} mathematical Analysis. We shall adopt a narrower definition: for us, it will be the study of topological vector spaces and of mappings \(u: \Omega \to F\) from a part of \(\Omega\) of a topological vector space \(E\) into a topological vector space \(F\), these mappings being assumed to satisfy various algebraic and topological conditions. A moment of reflection shows that this already covers a large part of modern Analysis, in particular the theory of partial differential equations.

Functional analysis thus appears as a rather complex blend of Algebra and Topology, and it should therefore surprise no one that the development of these two branches of mathematics had a strong influence on its own evolution. As a matter of fact, it is almost impossible to dissociate the early history of General Topology (and even of the set-theoretic language) from the beginnings of Functional Analysis, since the sets and spaces which (after the subsets of \(\R^n\)) attracted most attention consisted of \emph{functions}.

With regard to Algebra, as the most frequently studied mappings between topological vector spaces are \emph{linear}, it is quite natural that linear Algebra should have greatly influenced Functional Analysis. In fact, at the end on the XIX\supth{} century, the old idea that infinitesimal Calculus was derived from the algebraic ``Calculus of differences'' by a ``limit process'' began to acquire a more precise and more influential form when Volterra applied a similar idea to an integral equation
\begin{equation}
    \label{eq:1}
    \int_a^y \varphi(x) H(x,y) dx = f(y)
\end{equation}
for an unknown function \(\varphi\), the functions \(f\) and \(H\) being continuous in \([a,b]\) and \([a,b] \times [a,b]\) respectively, with \(f(a) = 0\). He divides \([a,b]\) into \(n\) subintervals by the points \(y_k = a + k \frac{b-a}{n} \; (1 \leq k \leq n)\), replaces \(y\) in \eqref{eq:1} by these \(n\) values, and the integral by the corresponding Riemann sums, which gives him a system of \(n\) linear equations
\begin{equation}
    \label{eq:2}
    \begin{dcases}
        h_{11} z_1               \hfill = b_1 \\
        h_{21} z_1 + h_{22} z_2  \hfill = b_2 \\
        \dotfill \\
        h_{n1} z_1 + h_{n2} z_2 + \ldots + h_{nn} z_n = b_n
    \end{dcases}
\end{equation}
with \(j_{jk} = H(y_j, y_k), z_k = \varphi(y_k)\) and \(b_k = f(y_k)\); the integral equation \eqref{eq:1} was thus considered as obtained from systems \eqref{eq:2} by a limit process when the number of unknowns became infinite.

Unfortunately, linear Algebra, as it was understood in the XIX\supth{} century (and even much later) did not readily lend itself to affording a good guidance to such generalizations. Its own evolution had been very slow and painful, stretching over 130 years, and in a succession of stages which, to our eyes, is exactly the \emph{reverse} of the \emph{logical} sequence of notions, namely

\[
    \begin{array}{c}
        \text{linear equations} \\
        \downarrow \\
        \text{determinants}\\
        \downarrow \\
        \text{linear and bilinear forms}\\
        \downarrow \\
        \text{matrices} \\
        \downarrow \\
        \text{vector spaces and linear maps}
    \end{array}
\]

In spite of the unsuccessful efforts of Grassmann and Peano, the intrinsic aspects and the geometric point of view in linear Algebra remained in the background until 1900; one would readily speak with Cayley (1843) of vectors and linear subspaces, but they were invariably considered as parts of some \(\R^n\); in other words, everything in a vector space was always referred to a \emph{fixed basis}, and linear maps were only handled through their \emph{matrices} corresponding to these bases. The various ``reduction'' theorems were known in 1880, but only through complicated computations of determinants, and without any geometric interpretation. Furthermore, Frobenius, who had been the most influential mathematician in building up a synthesis of the linear Algebra of his time, had unfortunately taken a step backward (even with respect to Cayley) by electing to work systematically with bilinear forms \(\sum_{p,q} a_{pq} x_p y_q\) instead of working with matrices \((a_{pq})\). Finally, before 1930 nobody had a correct conception of duality between finite dimensional \emph{vector spaces}; even in van der Waerden's book (1931), such a vector space and its dual are still \emph{identified}.

All this was to weigh heavily on the evolution of linear Functional Analysis; in particular it followed (over a shorted span of years) the same unfortunate succession of stages through which linear Algebra had to go; and it is only after it was realized that the current conception of vectors as ``\(n\)-tuples'' could not possibly be extended to infinite dimensional function spaces, that this conception was finally abandoned and that genuinely geometrical notions won the day.

The diagram at the end of this Introduction tries to depict graphically in some detail the successive stages of the history of Functional Analysis, by mentioning the actions and reactions of the various parts of mathematics which took part in it. If one were to reduce this complicated history to a few key words, I think the emphasis should fall on the evolution of two concepts: \emph{spectral theory} and \emph{duality}. Both of course stem from the very concrete problems encountered in the solution of linear equations (or systems of linear equations), where the unknowns are \emph{functions}. The basic concepts of spectral theory: eigenvalues, eigenfunctions and expansions in series of such functions were already known at the beginning of the XIX\supth{} century, in the theory of Fourier series; they would form the model on which all further advances were patterned. But it took 60 years of strenuous efforts to extend the theory from the Sturm--Liouville problem in ordinary differential equations to the partial differential equation of the vibrating membrane. It was gradually realized that the heart of the matter lay, not in the differential (or partial differential) equations themselves, but in \emph{integral equations} associated to them; at first they were not explicitly written down, so that one can only speak of ``crypto-integral'' equations, to designate the use of methods resting on evaluations of integrals, and which only later emerged as standard methods in the theory of integral equations.

The remarkable feature of this history is that, after such a slow incubation period, so to speak, spectral theory, in the span of a few years, reached complete maturity, giving birth in the process to the concept of linear duality, which began at last to be understood by analysts, before becoming later familiar to all mathematicians by a kind of backlash effect. What is interesting in this rapid advance is that it was accomplished in a series of what one may call discrete \emph{jumps}, in each of which the decisive step was to ignore the special features of the problem under consideration, and to make it accessible by inserting it into a more general context.

The first of these ``discontinuities'' occurred in 1896--1900, when Le Roux, Volterra and Fredholm, instead of working on the \emph{special} integral equations studied by their predecessors (Abel, Liouville, Beer--Neumann), elected to use \emph{minimal} assumptions on the kernels, and in so doing discovered that the theory was far simpler that it was generally thought.

The second step was taken by Hilbert in his 1906 papers, subordinating the too special theory of symmetric integral equations to the much more general concept of infinite ``bounded'' quadratic forms, which turned out to provide the frame needed for all subsequent progress in ordinary and partial differential equations.

The contemporary discovery of the Lebesgue integral, and the geometric and topological concepts introduced by Fréchet in Analysis immediately led Hilbert's successors to translate his results into the language of what we now call Hilbert space, linking the euclidean geometry to integration theory, and making possible the discussion of the most general systems of linear equations in such a space.

This in turn led F. Riesz in 1910--1913 to introduce \(L^p\) and \(l^p\) spaces for any exponent \(p\) such that \(1 < p < +\infty\), and to discover the natural duality between the \emph{different} spaces \(L^p\) and \(L^q\) with \(\frac{1}{p} + \frac{1}{q} = 1\), in sharp distinction from the muddleheaded ideas on the matter, which the accidental self-duality of Hilbert space had failed to dispel.

But although F. Riesz, in the treatment of systems of linear equations in \(l^p\) spaces, was the first to obtain a condition which later was seen to consist in a particular application of the Hahn--Banach theorem, he failed to visualize that condition as amounting to an extension property of a continuous linear form defined on a subspace. This fourth ``jump'' was only accomplished by Helly in 1921, again by generalizing the theory of systems of linear equations from the special \(l^p\) spaces to \emph{any} normed subspace of \(\C^{\mathbb{N}}\). After that, only two more steps were needed to reach the present status of the theory, with the passage to general normed spaces (together with the use of transfinite induction) by Hahn and Banach and a little later the extension of duality theory to locally convex spaces during the period 1935--1945.

This process of successive generalizations may thus have reached a point of diminishing returns around the middle of the century. Inasmuch as we are able to judge from events probably too recent to allow a proper perspective, the theory of topological vector spaces, after 1950, has stabilized as one of the standard tools of modern mathematics, together with linear and multilinear Algebra, General Topology and measure theory. The advances which have been achieved during the last 30 years mainly consist in new imaginative ways to use the fundamental tools of Functional Analysis, either in theories where they had not been applied before, such as differential geometry and differential topology (K-theory, theory of the Atiyah--Singer index, foliations), or in the construction of more powerful methods to handle functional equations (distributions, Sobolev spaces, pseudo-differential operators and their generalizations)·

This volume grew out of a series of lectures which I gave in Rio de Janeiro in 1979, at the invitation of Prof. Jorge Alberto Barroso of the Universidade Federal de Rio de Janeiro, to whom go my most heartfelt thanks. I am also very grateful to him for the pains he took in supervising the preparation of the manuscript for publication.

\clearpage

% headers for the other chapters
\fancyhead[EC]{\textsc{chapter} \thechapter}
\fancyhead[OC]{\textsc{\leftmark}}

\chapter{Linear differential equations and the Sturm--Liouville problem}
\label{ch:1}
\setcounter{equation}{0}

\section[Differential equations and partial differential equations in the XVIII\supth{} century]{Differential equations and partial differential equations\\in the XVIII\supth{} century}
\label{sec:1.1}

Until around 1750, the notion of \emph{function} of one variable was a very hazy one. The domain where it was defined was seldom described with precision; it was tacitly assumed that around each point \(x_0\), the function was equal to a power series in \(x - x_0\) and its derivatives were obtained by taking the derivatives of each term of the series. To solve a differential equation of order \(n\)
\begin{equation}
    \label{eq:1.1}
    y^{(n)} = F(x, y, y', y'', \ldots, y^{(n-1)})
\end{equation}
one would therefore substitute in \eqref{eq:1.1} for \(y\) and its derivatives a power series \(\sum\limits_{k=0}^\infty c_k (x - x_0)^k\) and its derivatives, and identify the series on both sides, which would determine each  \(c_k\) for \(k \leq n\) as a function of \(c_0, c_1, \ldots, c_{k-1}\); the solution thus depended on \(n\) arbitrary parameters \(c_0, c_1, \ldots, c_{n-1}\). The very few cases in which it was possible to write explicitly the solution by means of primitives of known functions (such as the linear equation \(y' = a(x)y + b(x)\) of order 1) were already known at the end of the XVII\supth{} century.

After 1760 began the first general study of linear equations of arbitrary order
\begin{equation}
    \label{eq:1.2}
    L(y) \equiv y^{(n)} + a_1(x) y^{(n-1)} + \ldots + a_n(x) y = b(x).
\end{equation}
D'Alembert observed that the knowledge of a particular solution of the equation and of all solutions of the homogeneous equation \(L(y) = 0\) yields by addition all solutions of \eqref{eq:1.2}. A little later, Lagrange \cite[vol.~I, p.~474]{135} showed that the general solution of \(L(y) = 0\) may be written \(\sum\limits_{k=1}^n c_k y_k\) where the \(c_k\) are arbitrary constants, and the \(y_k (1 \leq k \leq n)\) particular solutions (which he tacitly assumed to be linearly independent). Then, by his famous method of  ``variation of constants'' \cite[vol.~IV, p.~159]{135}, he showed how to obtain also the solutions of \eqref{eq:1.2} when the \(y_k\) were known: the solution is written in the form \(y = \sum\limits_{k=1}^n z_k y_k\), where the \(z_k\) are unknown functions, subject to \(n-1\) linear relations
\begin{equation}
    \label{eq:1.3}
    \sum_{k=1}^n z'_k y_k^{(\nu)} = 0 \quad (0 \leq \nu \leq n-2)
\end{equation}
These conditions imply that \(y^{(\nu)} = \sum\limits_{k=1}^n z_k y_k^{(\nu)}\) for \(0 \leq \nu \leq n-1\); replacing \(y\) by \(\sum\limits_{k=1}^n z_k y_k\) in \eqref{eq:1.2} and using the fact that the \(y_k\) satisfy \(L(y_k) = 0\), one obtains for the \(z'_k\) another linear equation
\begin{equation}
    \label{eq:1.4}
    \sum_{k=1}^n z'_k y_k^{(n-1)}  = b(x)
\end{equation}
from which, by the Cramer formulas, one can compute the \(z_k' (1 \leq k \leq n)\) and the problem is thus reduced to computing their primitives,

Lagrange also introduced \cite[vol.~I, p.~471]{135} the notion of \emph{adjoint} of a linear differential operator \(L\), which was to acquire great importance later: he showed that there exists a linear differential operator \(M\) satisfying an identity
\begin{equation}
    \label{eq:1.5}
    z L(y) - y M(z) = \frac{d}{dx}\left(B(y, z)\right)
\end{equation}
where \(B\) is bilinear in \((y, y', \ldots, y^{(n-1)})\) and \((z, z', \ldots, z^{(n-1)})\), constituting a generalization of the classical ``integration by parts''; he deduced from that formula that if a solution \(z\) of \(M(z) = 0\) was known, solutions of \(L(y) = 0\) could be obtained by solving an equation \(B(y,z) = \text{Const}\), of order \(n-1\).

Partial differential equations were not considered until the middle of the XVIII\supth{} century, in connection with problems of Mechanics or Physics and then they were of order 2 at least (see \ref{sec:1.2}). The study of partial differential equations of first order was only begun by Euler and Lagrange after 1770. Euler was able to solve a few particular equations, and then Lagrange found general methods which enabled his followers, Charpit and Monge, to reduce the solution of a general equation of first order
\begin{equation}
    \label{eq:1.6}
    F\left(x,y,z, \pder{z}{x}, \pder{z}{y}\right) = 0
\end{equation}
to the solution of a system of ordinary differential equations, an idea which was developed later by Cauchy in his concept of ``characteristic curves''.

\section{Fourier expansions}
\label{sec:1.2}

In 1747, d'Alembert gave the first mathematical treatment of the general problem of the small vibrations of a string of length \(a\), fixed at each extremity; the string moves in a plane where the axis \(0x\) is along the position of the string at rest, the segment \(0 \leq x \leq a\); if \(y = u(x,t)\) is the equation of the string at time \(t\), d'Alembert shows that, if \(u(x,t)\) remains small, it satisfies the equation
\begin{equation}
    \label{eq:1.7}
    \frac{\partial^2 u}{\partial t^2} = c^2 \frac{\partial^2 u}{\partial x^2}
\end{equation}
where \(c\) is a known function of \(x\) alone, and is constant if the density of the string is constant. When \(c\) is constant, taking \(X = x-ct\) and \(Y = x+ct\) as new variables reduces the equation to \(\frac{\partial^2 u}{\partial X \partial Y} = 0\), and d'Alembert concluded that the solution of \eqref{eq:1.7} is given by
\begin{equation}
    \label{eq:1.8}
    u(x,t) = f(x - ct) + g(x + ct)
\end{equation}
where \(f\) and \(g\) are ``arbitrary'' functions. A year later, Euler interpreted this result as meaning that (for \(c=1\)) \(u(x,t)\) was known once the two functions of \(x\),
\begin{equation}
    \label{eq:1.9}
    u(x,0) = \varphi(x), \quad \pder{u}{t}(x,0) = \psi(x)
\end{equation}
were prescribed, the value of \(u(x,t)\) being explicitly given by
\begin{equation}
    \label{eq:1.10}
    u(x,t) = \frac{1}{2} \left(\varphi(x-t) + \varphi(x+t)\right) + \frac{1}{2} \int_{-t}^t \psi(x - \xi) d\xi
\end{equation}
(Euler only gives a geometric construction equivalent to this formula). Now it was well known experimentally that \(\varphi(x)\) could be quite different from an analytic function, for instance it could have no derivative at some points, and this led Euler to introduce, in addition to what he called ``continuous'' functions (i.e. analytic functions in our sense) more general ones which he baptized ``mechanical'' without giving their precise definition (from the context they seem to be piecewise twice differentiable functions in our terminology).

On the other hand, already in 1715, B. Taylor, by a direct argument which did not use equation \eqref{eq:1.7}, had concluded that (when \(c\) is constant) for any integer $n \geq 1$, the function
\begin{equation}
    \label{eq:1.11}
    u_n(x,t) = \sin\frac{n\pi x}{a} \cos\frac{n\pi ct}{a}
\end{equation}
represented vibrations of the string, namely for \(n = 1\) the ``fundamental'' tone, and for \(n = 2,3, \ldots\), its ``harmonics''. As it was well known that the sound emitted by a vibrating string was in general a mixture of several ``harmonics'', Daniel Bernoulli, in 1750, proposed that the general solution \eqref{eq:1.10} could also be written as a series
\begin{equation}
    \label{eq:1.12}
    u(x,t) = \sum_{n=1}^\infty a_n \sin\frac{n\pi x}{a} \cos\frac{n\pi c}{a}\left(t - b_n\right)
\end{equation}
for suitable values of the \(a_n\) and \(b_n\). However, in 1753, Euler observed that this would imply that an arbitrary ``mechanical'' function defined in an interval \(-a \leq x \leq a\) could be written as a series
\begin{equation}
    \label{eq:1.13}
    \frac{a_0}{2} + a_1 \cos\frac{\pi x}{a} b_1 \sin\frac{\pi x}{a} + a_2 \cos\frac{2\pi x}{a} + b_2 \sin\frac{2\pi x}{a} + \ldots
\end{equation}
and he believed that such a series of analytic functions could only represent an analytic function. His opinion was shared (with some variations) by almost all other mathematicians of his time, and no progress was made on this question until the beginning of Fourier's work on the theory of heat (see \cite[(2), t.~XI\sub{2}, pp.~273--300]{065}). Having to solve equations such as
\begin{equation}
    \label{eq:1.14}
    \frac{\partial^2 u}{\partial x^2} + \frac{\partial^2 u}{\partial y^2} = 0
\end{equation}
\begin{equation}
    \label{eq:1.15}
    \frac{\partial^2 u}{\partial x^2} - \frac{\partial^2 u}{\partial y} = 0
\end{equation}
for various boundary conditions, he systematically looks for solutions of the form \(u(x,y) = v(x)w(y)\) and, following D. Bernoulli, wants to obtain the most general solution as series whose terms are these particular ones. In so doing, he is brought back to the problem of expressing a function \(f\) as a series \eqref{eq:1.13}, but this time he adds to D. Bernoulli's argument the formulas giving actually the values of the coefficients \(a_n, b_n\)
\begin{equation}
    \label{eq:1.16}
    a_n = \frac{1}{\pi} \int_{-\pi}^\pi f(x)\cos{nx} dx, \quad b_n = \frac{1}{\pi} \int_{-\pi}^\pi f(x)\sin{nx} dx
\end{equation}
(when \(a = \pi\)) which as a matter of fact had already been obtained by Clairaut and Euler, without realizing their interest. Using these formulas Fourier was able to show on many examples of non analytic functions that the corresponding Fourier series converged to \(\frac{1}{2} \left(f(x_+) + f(x_-)\right)\), and expressed his conviction that this was true for ``arbitrary'' functions, although his attempts and those of Cauchy to prove that result were unsuccessful and the first proof for a piecewise monotonic and piecewise continuous function was only given by Dirichlet in 1829. One should also mention in that connection that in 1799, Parseval had given the formula
\begin{equation}
    \label{eq:1.17}
    \frac{a_0^2}{2} + \sum_{n=1}^\infty \left(a_n^2 + b_n^2\right) = \frac{1}{\pi} \int_{-\pi}^\pi \left(f(t)\right)^2 dt
\end{equation}
by a purely formal computation, without any proof of convergence.

These results gave the impetus to the vast theory of \emph{trigonometric series}, which was to be one of the main concerns of most analysts in the XIX\supth{} century, centered around the criteria of convergence of such series and the relations between its sum and its coefficients. The evolution of that theory was closely linked to a gradual precision and deepening of the notions of set of real numbers, of function and of integral. But before 1920 there was not much contact between that theory and the development of Functional Analysis as we understand it.

On the contrary, other results of Fourier in his \emph{Theory of heat} triggered the birth of \emph{spectral theory}. For instance \cite[vol.~I, p.~304]{067} he shows that the ``cooling off'' problem for a solid sphere of radius \(r\), when one assumes spherical symmetry for the problem, is governed by the partial differential equation
\begin{equation}
    \label{eq:1.18}
    \pder{u}{t} = k \left(\frac{\partial^2 u}{\partial x^2} + \frac{2}{x} \pder{u}{x}\right)
\end{equation}
with the ``boundary conditions'' that \(u(x,t)\) must remain finite when \(x\) tends to 0, and satisfy the relation
\begin{equation}
    \label{eq:1.19}
    \pder{h}{u} + hu = 0 \quad \text{for \(x=r\) and all \(t\)},
\end{equation}
where \(h\) and \(k\) are constants. Using his favorite method of ``separation of variables'', Fourier obtains solutions
\begin{equation}
    \label{eq:1.20}
    u(x,t) = \frac{1}{x} \exp{\left(-k\lambda^2 t\right)}\sin{\lambda x}
\end{equation}
provided the parameter \(\lambda\) is a solution of the transcendental equation
\begin{equation}
    \label{eq:1.21}
    \frac{\lambda r}{\operatorname{tg}{\lambda r}} = 1 - hr.
\end{equation}
He easily proves that the equation has an infinity of real roots \(\lambda_n\) tending to \(+\infty\). To obtain a solution of \eqref{eq:1.18} with boundary condition \eqref{eq:1.19} and such that \(u(x,0)\) is a given function \(f(x)\), he proceeds as before, writing \(xf(x)\) as a series \(\sum\limits_{n=1}^\infty c_n \sin{\lambda_n x}\); he shows that one has again the ``orthogonality'' relations (of course he does not use that word)
\begin{equation}
    \label{eq:1.22}
    \int_0^r \sin \lambda_n x \sin \lambda_m x dx = 0 \quad \text{for \(m \neq n\)}
\end{equation}
and from them deduces the relations
\begin{equation}
    \label{eq:1.23}
    % I use \Bigg here because they have to be the same size as those of eq. 1.34, where for some reason I don't need \Bigg for the parentheses
    c_n = \Bigg(\int_0^r xf(x)\sin{\lambda_n x} dx\Bigg) \Bigg/ \Bigg(\int_0^r \sin^2{\lambda_n x} dx\Bigg)
\end{equation}
without of course any rigorous justification, nor any proof of the fact that the series converges to \(xf(x)\).

\section{The Sturm--Liouville theory}
\label{sec:1.3}

The results of Fourier on the theory of heat were continued and expanded by Poisson. Their work led Ch. Sturm in 1836 and J. Liouville one year later to build a general theory which would include all cases considered by Fourier and Poisson, without assuming the possibility of explicit integration. They consider a second order differential equation
\begin{equation}
    \label{eq:1.24}
    y'' = q(x)y + \lambda y = 0
\end{equation}
where \(q\) is a real valued continuous function in a compact interval \([a,b]\) of \(\R\), and \(\lambda\) a complex parameter. The first problem is to consider boundary conditions of the form
\begin{equation}
    \label{eq:1.25}
    y(a)\cos\alpha - y'(a)\sin\alpha = 0, \quad y(b)\cos\beta - y'(b)\sin\beta = 0
\end{equation}
where \(alpha\) and \(\beta\) are two positive constants, and to determine for what values of \(\lambda\) the problem has a non trivial solution (an ``eigenfunction'' for the ``eigenvalue'' \(\lambda\) in our present day language).

A first remark, which had already essentially been made by Poisson, is that if \(\lambda, \mu\) are two different eigenvalues, and \(u, v\) two corresponding ``eigenfunctions'', then from the relations
\[
    u'' - qu + \lambda u = 0, \quad v'' - qv + \mu v = 0,
\]
one deduces
\[
    u'' v - v'' u + (\lambda - \mu) uv = 0
\]
and as \(\int\limits_a^b \left(u'' v - v'' u\right)dx = (u'v - v'u)\rvert_a^b = 0\) because of \eqref{eq:1.25}, one obtains
\begin{equation}
    \label{eq:1.26}
    (\lambda - \mu) \int_a^b u(x)v(x)dx = 0.
\end{equation}
A first consequence of this relation is that eigenvalues are necessarily \emph{real} numbers, Indeed, if \(\lambda\) was not real, then \(\overline\lambda\) would also be an eigenvalue with eigenfunction \(\overline{u}\), and substituting \(\overline\lambda\) and \(\overline{u}\) for \(\mu\) and \(v\) in \eqref{eq:1.26}, one obtains \(\int\limits_a^b |u(x)|^2 dx = 0\), contrary to assumption.

The main contribution of Sturm was the proof that there are infinitely many eigenvalues \(\lambda_1 < \lambda_2 < \ldots < \lambda_n \ldots\), tending to \(+\infty\). In his study of vibrating strings, d'Alembert had al ready considered an equation of the form \(y'' - \lambda\varphi(x)y = 0\) where \(\varphi\) is not constant, and had tried to prove that there is a single value of \(\lambda\) for which there is a solution in \([a,b]\) vanishing at \(a\) and \(b\) and nowhere else; his idea was to study the corresponding Riccati equation for \(y'/y\) when \(\lambda\) varies \cite[(2), vol.~XI\sub{2}, p.~311]{065}. Sturm elects a similar approach: he considers a solution \(u(x,\lambda)\) of \eqref{eq:1.24} satisfying the \emph{first} condition \eqref{eq:1.25}, and fixed for instance by the condition \(u(a,\lambda) = 1\) (or \(u'(a,\lambda) =1\) if \(\alpha = 0\)), and he studies the \emph{variation} of \(u(x,\lambda)\) as a function of \(\lambda\); the \(\lambda_n\) are therefore the solutions of the equation \(u(b,\lambda)\cos\beta - u' (b,\lambda)\sin\beta = 0\). He is thus led to \emph{compare} solutions of two equations
\begin{equation}
    \label{eq:1.27}
    y'' + q_1(x)y = 0, \quad y'' + q_2(x)y = 0
\end{equation}
when \(q_1(x) < q_2(x)\), and discovers many remarkable such ``comparison theorems'', of which we will only quote the one which leads to the existence of the eigenvalues. Sturm's paper is rather long-winded and not very clear (\cite{209}, \cite[p.~259--268]{S}) and there is a much simpler formulation of his result: an equation \(y'' + q(x)y = 0\) is written as a system of two first order equations by the usual introduction of two functions \(y_1 = y, y_2 = y'\), which gives \(y_1' = y_2, y_2' = -q(x)y_1\), and then one takes as new unknowns two functions \(r,\theta\) such that \(y_1 = r \sin\theta, y_2 = r \cos\theta\), which leads to the system
\begin{equation}
    \label{eq:1.28}
    r' = (1-q(x))r \sin\theta \cos\theta
\end{equation}
\begin{equation}
    \label{eq:1.29}
    \theta' = \cos^2\theta + q(x)\sin^2\theta
\end{equation}
where the second equation now is of the first order only\footnote{This device seems to have first been introduced by H. Prüfer \cite{180}.}. The comparison theorem which is needed is then the following one: consider solutions \(\varphi_0, \varphi_1\) in \([a,b]\) of the two equations
\begin{equation}
    \label{eq:1.30}
    \theta' = \cos^2\theta + q_1(x)\sin^2\theta, \quad \theta' = \cos^2\theta + q_2(x)\sin^2\theta
\end{equation}
and suppose that \(q_1(x) < q_2(x)\) in \([a,b]\). Then, if for a number \(\alpha \in ]a,b[\) one has \(\varphi_1(\alpha) \leq \varphi_2(\alpha)\), one also has \(\varphi_1(x) < \varphi_2(x)\) for \(\alpha < x < b\). The proof is very simple and consists in computing the derivative of the function \(w(x) = \varphi_2(X) - \varphi_1(x)\) and showing that there is a continuous function \(f\) in \([a,b]\) such that \(w'(x) - f(x)w(x) \geq 0\), which implies that \(w\) cannot change sign,

If now we apply the preceding change of variable to \eqref{eq:1.24}, we get the equation
\begin{equation}
    \label{eq:1.31}
    \theta' = \cos^2\theta + (\lambda - q(x))\sin^2\theta
\end{equation}
and we consider the solution \(\omega(x,\lambda)\) such that \(\omega(a,\lambda) = \alpha\); the eigenvalues \(\lambda\) are the solutions of the equations
\begin{equation}
    \label{eq:1.32}
    \omega(b, \lambda) = \beta + n\pi \quad \text{for \(n \in \mathbb{Z}\)}
\end{equation}
Sturm's comparison theorem then shows that for each \(x \in ]a,b]\) the function \(\lambda \mapsto \omega(t, \lambda)\) is \emph{strictly increasing}, and in addition, from \eqref{eq:1.31} it follows that if \(\omega(x,\lambda) = k\pi\) for an integer \(k\), then \(\pder{\omega}{x}(x, \lambda) = 1\). From these facts it is easy to show that each equation \eqref{eq:1.32} has one and only one solution \(\lambda_n\) for each \(n \geq 1\) and no solution for \(n \leq 0\); in addition, the corresponding eigenfunction \(u_n\) may be shown to have exactly \(n\) zeroes in the interval \(]a,b[\) \cite[p.~435--441]{052}.

Building on these results of Sturm, Liouville then proceeds to give a general formulation to the expansions of Fourier and Poisson. From relation \eqref{eq:1.26} where \(\lambda\) and \(\mu\) are replaced by \(\lambda_n\) and \(\lambda_m\) it follows that
\begin{equation}
    \label{eq:1.33}
    \int_a^b u_m(x)u_n(x)dx = 0 \quad \text{for \(m \neq n\)}.
\end{equation}
To each function \(f\), defined and continuous in \([a,b]\), Liouville associates its ``generalized Fourier coefficients''
\begin{equation}
    \label{eq:1.34}
    c_n = \left(\int_a^b f(x)u_n(x) dx\right) \Bigg/ \left(\int_a^b u_n^2(x)dx\right)
\end{equation}
and considers the ``generalized Fourier series'' \(\sum\limits_{n=1}^\infty c_n u_n(x)\). In order to study its convergence, he needs more information on the behavior of \(\lambda_n\) and \(u_n\) when \(n\) tends to \(+\infty\). He observes that, if \(\lambda = \rho^2 > 0\), any solution of \eqref{eq:1.24} satisfies a relation of the form
\begin{equation}
    \label{eq:1.35}
    y(x) = A \cos{\rho x} + B \sin{\rho x} + \frac{1}{\rho} \int_a^x q(t)y(t)\sin{\rho(x-t)}dt
\end{equation}
(which can be deduced from Lagrange's ``variation of constants'' method, by writing \eqref{eq:1.24} as \(y + \rho^2 y = q(x)y\), although this is not the way Liouville proves \eqref{eq:1.35}). Applying this to \(y = u_n\), so that \(\rho\) is replaced by \(\lambda_n^\frac{1}{2}\), he gives a sketchy proof that \(\rho_n = \frac{(n-1)\pi}{b-a} + O(1/n)\) and (if \(\cos\alpha \neq 0\)) \(u_n(x) = \sqrt\frac{2}{b-a} \cos{\rho_n x} + O(1/n)\) (when \(u_n\) is normalized by the condition \(\int\limits_a^b u_n^2(x)dx = 1\)). This allows him to prove that the series \(\sum\limits_{n=1}^\infty c_n u_n(x)\) converges, provided the usual Fourier series of \(f\) converges. He still has to show that, if \(f\) is continuous, the function \(F(x) = \sum\limits_{n=1}^\infty c_n u_n(x)\) is equal to \(f(x)\); he assumes (without proof) that \(F\) is continuous and that \(c_n = \int\limits_a^b F(x)u_n(x)dx\), and is reduced to proving that the relations \(\int\limits_a^b (F(x)-f(x))u_n(x)dx = 0\) for all \(n\) imply \(F = f\) (first appearance of the property of ``completeness'' of an orthonormal system); but this he can only do under the additional assumption that \(F-f\) has only a finite number of zeroes in \([a,b]\). The complete proof of the relation \(f(x) = \sum\limits_{n=1}^\infty c_n u_n(x)\) was only given (for \(f\) piecewise \(C^2\)) at the end of the XIX\supth{} century, as well as the relation \(\sum\limits_{n=1}^\infty c_n^2 = \int\limits_a^b f^2(x)dx\); Liouville had only proved the corresponding inequality \(c_1^2 + \ldots  + c_N^2 \leq \int\limits_a^b f^2(x)dx\) for all \(N\) (named after Bessel, who had proved it for the trigonometric system) (\cite{151}, \cite[p.~268--281]{S}).

These remarkable results were to form the pattern of \emph{spectral theory}, the main efforts of analysts in that direction being directed to a generalization of the Sturm--Liouville theory to some types of partial differential equations; but in the first half of the XIX\supth{} century, the theory of these equations was far less advanced than the theory of ordinary differential equations, and it is only after 1880 that progress became possible (see Chapter \ref{ch:3}).

\chapter{The ``crypto-integral'' equations}
\label{ch:2}
\setcounter{equation}{0}

\section{The method of successive approximations}
\label{sec:2.1}

The study of celestial mechanics during the XVIII\supth{} century by the method of perturbations consisted, for the theory of the movements of planets, to first neglect their mutual attraction, which gave for each planet a Keplerian orbit around the sun, and then to find the deviations of the actual orbits from the Keplerian ones by taking into account the attraction of other planets; due to the fact that the masses of the planets are much smaller than the mass of the sun, these deviations were expected to be small. Translated into mathematical terms, this amounted, in the simplest cases, to find good approximations for the solutions of a system of differential equations
\begin{equation}
    \label{eq:2.1}
    y'_i = \varepsilon f_{1i}(x, y_1, \ldots, y_n) + \varepsilon^2 f_{2i}(x, y_1, \ldots, y_n) + \ldots \quad (1 \leq i \leq n)
\end{equation}
where the parameter \(\varepsilon\) on the right-hand sides is ``small''. The general conception of function in XVIII\supth{} century mathematics naturally led to try to express the \(y_i\) as a power series in \(\varepsilon\)
\begin{equation}
    \label{eq:2.2}
    y_i = a_i + \varepsilon y_{1i} + \varepsilon^2 y_{2i} + \ldots \quad (1 \leq i \leq n),
\end{equation}
to substitute these expressions in and identify the coefficients of the successive powers of \(\varepsilon\) on both sides. This led to a succession of equations
\[
    \begin{array}{l}
        y'_{1i} = f_{1i}(x, a_1, \ldots, a_n) \\
        y'_{2i} = F_{2i}(x, y_{11}, \ldots, y_{1n}) \\
        y'_{3i} = F_{3i}(x, y_{11}, \ldots, y_{1n}, y_{21}, \ldots, y_{2n}) \\
        \dotfill
    \end{array}
\]
all of which had right-hand sides which were known functions, hence were reduced to mere ``quadratures''. No attempt was made to justify mathematically those procedures; the goal of these computations was merely to obtain a satisfactory agreement with observations.

It is well-known that Cauchy was the first mathematician who proved existence theorems for \emph{general} types of differential equations, for which no explicit solution is available. His strategy was to consider the various methods introduced earlier for the purpose of numerical computations, and to show that, under certain conditions, these methods actually gave \emph{convergent} approximation processes having a solution as limit. In particular, in a paper published in 1835 in Prag (\cite{040}, (2), vol.~XI, p.~399--465), he takes up the method outlined above, not for an ordinary differential equation, but for a linear partial differential equation of first order (which was known to be equivalent to a system of ordinary differential equations)
\begin{equation}
    \label{eq:2.3}
    \pder{U}{t} = \sum_{i=1}^p A_i(t, x_1, \ldots, x_p) \pder{U}{x_i};
\end{equation}
the problem is to find a solution which for \(t=0\) reduces to a given function \(u(x_1, \ldots,x_n)\), and Cauchy transforms \eqref{eq:2.3} into the equivalent ``integro-differential'' equation by considering \(x_1, \ldots, x_p\) as parameters:
\begin{equation}
    \label{eq:2.4}
    U(t, x_1, \ldots, x_p) = u(x_1, \ldots, x_p) + \int_0^t \left( \sum_{i=1}^p A_i(s, x_1, \ldots, x_p) \right) \pder{U}{x_i} ds
\end{equation}
which he solves by successive approximations, starting with \(U_0 = u\), and defining
\[
    U_n(t, x_1, \ldots, x_p) = u(x_1, \ldots, x_p) + \int_0^t \left( \sum_{i=1} A_i(s, x_1, \ldots, x_p) \right) \pder{U_{n-1}}{x_i} ds
\]
by induction; but he is only able to prove convergence towards a solution when the \(A_i\) are analytic functions.

In his 1837 papers on the Sturm--Liouville problem Liouville independently applied a similar method to the linear differential equation \(y'' = f(x)y\), for which he wants to find a solution in \([a,b]\) satisfying the boundary condition \(y'(a) - hy(a) = 0\). He starts from the function \(y_0(x) = 1+h(x-a)\) satisfying that condition, and considers the series
\begin{equation}
    \label{eq:2.5}
    y = y_0 + y_1 + \ldots + y_n + \ldots
\end{equation}
where the \(y_n\) are determined for \(n > 0\) by the recursive equations
\[
    y_{n+1}(x) = \int_a^x dt \int_a^t f(s)y_n(s)ds.
\]
It must be remembered that at that time the concept of uniform convergence had not yet been formulated, and no justification had been given for asserting the continuity of a convergent series of continuous functions, or differentiating or integrating such a series termwise, Liouville proves very easily that there is a constant \(C\) such that
\[
    \left| y_n(x) \right| \leq C^n(x-a)^{2n} / (2n)!
\]
from which he concludes that the series \eqref{eq:2.5} giving \(y(x)\) converges for every \(x\); but he tacitly takes for granted that \(y\) is a \(C^2\) function and a solution of his problem.

In addition, Liouville makes the interesting remark that the function \(y\) can also be defined by the relation
\begin{equation}
    \label{eq:2.6}
    y = y_0 + \int_a^x dt \int_a^t f(s) y(s) ds
\end{equation}
(which he could also have written \(y = y_0 + \int_a^x (x-t)f(t)y(t)dt\)), thus giving what is probably the first example of what will be called later a ``Volterra integral equation of the second kind'' (see chap. \ref{ch:4}); if one writes \(z_n = y_0 + y_1 + \ldots + y_n\), Liouville observes that the \(z_n\) are given by \(z_0 = y_0\), and the recursive equations
\begin{equation}
    \label{eq:2.7}
    z_{n+1}(x) = y_0 + \int_a^x dt \int_a^t f(s) z_n(s) ds
\end{equation}
which is the standard process of ``successive approximations'' for these equations (\cite{151}, \cite[p.~268--281]{S}).

We have already seen that a little later in his papers of 1837, Liouville gives another ``integral equation'' equivalent to an equation \(y' = f(x)y\) (chap. \ref{ch:1}, \ref{sec:1.3}, equation \eqref{eq:1.35}). This exemplifies a general idea: if a linear differential operator \(P\) is such that the equation \(P \cdot u = f\) can be solved by a formula \(u = y_0 + G \cdot f\), where \(G\) is a linear operator, then the equation \(P \cdot u + Q \cdot u = 0\), where \(Q\) is an operator, is equivalent to \(u - G \cdot (Q \cdot u) = y_0\); in the case of Liouville, \(P \cdot u = u'' + \rho^2 u\) and \(Q \cdot u = -qu\), and \(G\) is an integral operator (cf. chap. \ref{ch:9}, \ref{sec:9.5}).

The simplest application of this idea is to the proof of Cauchy's existence and uniqueness theorem for an ordinary differential equation \(y' = f(x,y)\), which, with the initial condition \(y(x_0) = y_0\), is equivalent to \(y = y_0 + \int_{x_0}^x f(t,y)dt\). In this general form it is given by E. Picard in his 1890 paper on successive approximations \cite[vol.~II, p.~197--200]{172}, where it comes as an afterthought, the bulk of the paper being concerned with applications of the method to partial differential equations. However, in these applications, Picard is directly influenced by the fundamental earlier works of C. Neumann on the Laplace equation and of H.A. Schwarz on the equation of vibrating membranes, which are the direct forerunners of the theory of integral equations; we will describe in detail C. Neumann's results in \ref{sec:2.4} of this chapter, and H.A. Schwarz's paper in chap. \ref{ch:3}, \ref{sec:3.1}.

\section{Partial differential equations in the XIX\supth{} century}
\label{sec:2.2}

During the whole XIX\supth{} century, the theory of partial differential equations (in contrast with the theory of ordinary differential equations) has remained in an embryonic stage. The only general theorem, patterned after the Cauchy theorem on local existence and uniqueness of solutions of ordinary differential equations, is the Cauchy--Kowalewska theorem: suppose we have a system of \(r\) equations in \(r\) unknown real functions \(v_1, \ldots, v_r\) of \(p+1\) real variables \(x_1, \ldots, x_p\), of type
\begin{equation}
    \label{eq:2.8}
    \pder{v_j}{x_{p+1}} = H_j\left(x_1, \ldots, x_{p+1}, v_1, \ldots, v_{r}, \pder{v_1}{x_1}, \pder{v_1}{x_2}, \ldots, \pder{v_r}{x_{p-1}}, \pder{v_r}{x_p}\right) \hfill (1 \leq j \leq r)
\end{equation}
where the right hand sides do not contain any derivative with respect to \(x_{p+1}\) and are supposed to be real and \emph{analytic} with respect to their \(p+1+r+rp\) variables, in a neighborhood \(V_0\) of 0 in \(\R^{p+1+r+rp}\); then there is a small neighborhood \(V\) of 0 in \(\R^{p+1}\) such that \eqref{eq:2.8} has in \(V\) a unique solution \((v_1,\ldots,v_r)\) consisting of \emph{analytic} functions in \(V\), such that \(v_j(x_1,\ldots,x_p,0) = 0\) in \(V \cap \R^p\) for \(1 \leq j \leq r\).

The tendency (inherited from the XVIII\supth{} century) to consider that the most interesting functions were analytic was still very strong during the whole XIX\supth{} century, and therefore at first the analyticity restrictions of the Cauchy--Kowalewska theorem did not worry mathematicians very much, However, as it was known that some special types of partial differential equations, such as the scalar equation of first order and some types of second order equations, had solutions under much less stringent restrictions, people began to wonder if some other method than Cauchy's ``method of majorants'' (which could only be applied to analytic functions) would not yield a generalization of the Cauchy--Kowalewska theorem, at least for \(C^\infty\) functions. The question remained unanswered until 1956, when H. Lewy gave the surprising example of a system of two linear equations in 3 variables, with \(C^\infty\) coefficients
\[
  \begin{dcases}
    \pder{v_1}{x_1} = \pder{v_2}{x_2} - 2 x_2 \pder{v_1}{x_3} - 2x_1 \pder{v_2}{x_3} - f(x_3) \\
    \pder{v_2}{x_1} = -\pder{v_1}{x_2} + 2 x_1 \pder{v_1}{x_3} - 2x_2 \pder{v_2}{x_3}
  \end{dcases}
\]
which, for a suitable choice of the real \(C^\infty\) function \(f\), has \emph{no solution} whatsoever around any point (even if one al lows solutions which are distributions).

We shall not discuss the numerous local studies of analytic systems of partial differential equations (not necessarily reducible to the form \eqref{eq:2.8}) which followed the Cauchy--Kowalewska theorem, since they had no influence on the development of Functional Analysis as we understand it.

The remainder of the theory of partial differential equations until 1890 was limited to very special scalar equations (mostly linear equations or order 2) generally derived from physical problems\footnote{See the interesting description of these problems given by Poincaré in the Introduction of his 1890 paper on the equations of mathematical physics (\cite{177}, vol.~IX, p.~28--32)}, such as the equation of vibrating strings and its generalizations to 3 and 4 variables (the ``wave equations''), the Laplace equation \(\Delta u = 0\) in 2 and 3 variables, the heat equation in 2, 3 and 4 variables. For these equations, the techniques of ``separation of variables'' or of Fourier transforms (see chapter \ref{ch:7}, \ref{sec:7.6}) gave special solutions or solutions depending on ``arbitrary'' functions. But until 1825 the determination of solutions by boundary conditions (of which we have seen a few examples in Chapter I) was always restricted to \emph{explicitly described and particular} such conditions.

A first attempt of classification of second order equations in 2 variables had been made by Laplace \cite[vol.~IX, p.~21--28]{137} He considered ``quasi-linear'' equations, i.e. those of the form
\begin{equation}
  \label{eq:2.9}
    A(x,y) \frac{\partial^2 z}{\partial x^2} + B(x,y) \frac{\partial^2 z}{\partial x \partial y} + C(x,y) \frac{\partial^2 z}{\partial y^2} + F\left( x,y,z, \pder{z}{x}, \pder{z}{y} \right) = 0
\end{equation}
linear in the second order derivatives. As he did not have a clear idea of the distinction between real and complex variables, and therefore did not hesitate to give complex values to \(x\) and \(y\), he asserted that a suitable change of variables could reduce the terms of \eqref{eq:2.9} containing second order derivatives either to \(\frac{\partial^2 z}{\partial x \partial y}\) or to \(\frac{\partial^2 z}{\partial x^2}\) when \(A, B, C\) are not all identically zero! With the development of the theory of functions of one complex variable, it was soon realized that, for real variables \(x, y\), equations \eqref{eq:2.9} where the second order derivatives enter by \(\frac{\partial^2 z}{\partial x^2} + \frac{\partial^2 z}{\partial y^2}\) (called \emph{elliptic} equations) had to be sharply distinguished from those (called \emph{hyperbolic} equations) where the second order derivatives enter by \(\frac{\partial^2 z}{\partial x \partial y}\) or \(\frac{\partial^2 z}{\partial x^2} - \frac{\partial^2 z}{\partial y^2}\). The study of general boundary conditions for hyperbolic equations only begins around 1860 and will have little contact with Functional Analysis until around 1925 (see chapter \ref{ch:9}, \ref{sec:9.5}). On the contrary, the various problems connected with the Laplace equation in 2 or 3 variables will be one of the main concerns of analysts from 1828 onwards, and will become the impetus leading to the theory of integral equations, and thence to our modern Functional Analysis.

\section{The beginnings of potential theory}
\label{sec:2.3}

In 1748, D. Bernoulli had introduced in the theory of newtonian attraction the function \(\Omega(M) = \sum\limits_i\left(m_i \mu/r_i\right)\) for a point \(M\) of mass \(\mu\) attracted by a finite number of punctual masses \(m_i\), where \(r_i\), is the distance of \(M\) to the mass \(m_i\); and in 1773 Lagrange observed that the knowledge of that function immediately gave the components of the attraction exerted on \(M\), by taking the derivatives of \(\Omega\) with respect to the coordinates \(x, y, z\) of \(M\). When the finite number of masses is replaced by a solid \(V\) of density \(\rho\) and the point \(M\) is outside \(V\), the function \(\Omega\) becomes
\begin{equation}
    \label{eq:2.10}
    \Omega(x,y,z) = \mu \iiint_V \frac{\rho(\xi, \eta, \zeta) d\xi d\eta d\zeta}{r(x, y, z, \xi, \eta, \zeta)}
\end{equation}
with \(r(x, y, z, \xi, \eta, \zeta) = \left( (x-\xi)^2 + (y-\eta)^2 + (z-\zeta)^2 \right)^{\frac{1}{2}}\) \cite[vol.~VI, p.~349]{135}.

In 1782 and 1785, Laplace showed that outside of \(V\) the function \(\Omega\) satisfied the equation
\begin{equation}
    \label{eq:2.11}
    \Delta \Omega \equiv \frac{\partial^2 \Omega}{\partial x^2} + \frac{\partial^2 \Omega}{\partial y^2} + \frac{\partial^2 \Omega}{\partial z^2} = 0
\end{equation}
\cite[vol.~X, p.~361--363 and vol.~XI, p.~276--280]{137}, and in 1813 Poisson completed that result by showing that if \(\rho\) discontinuous in \(V\), the integral \eqref{eq:2.10} is still meaningful inside \(V\), and \(\Omega\) satisfies the ``Poisson equation''
\begin{equation}
    \label{eq:2.12}
    \Delta\Omega + 4\pi \rho = 0
\end{equation}
(\cite{178}, \cite[p.~342--346]{S}). His idea is to consider the value of \(\Omega\) at a point \(M\) in \(V\) as the sum of the corresponding functions \(\Omega_1, \Omega_2\) relative to a small ball \(V_1\) of center \(M\) and to the complement \(V_2\) of \(V_1\) in \(V\); one has then \(\Delta\Omega_2 = 0\), and when the radius of \(V_1\) tends to 0 Poisson shows that \(\Delta\Omega_1\) tends to \(-4\pi\rho(M)\). (In fact his argument is not rigorous when one only assumes the continuity of \(\rho\), and the existence of \(\Delta\Omega\) is only guaranteed when \(\rho\) satisfies a Hölder condition; when \(\rho\) is merely continuous, equation \eqref{eq:2.12} is valid only if the second order derivatives are taken in the sense of the theory of distributions (chap. \ref{ch:8}, \ref{sec:8.3})).

After the discovery of Coulomb's laws (1785) the Laplace equation became of central importance in electrostatics; it also was found to govern ``stationary'' phenomena in hydrodynamics and the theory of heat. Finally the so-called ``Cauchy--Riemann'' equations for real functions \(P, Q\) of \(x, y\) such that \(P + iQ\) is an analytic function of \(x + iy\), were known since the middle of the XVIII\supth{} century, and they implied that \(P\) and \(Q\) were solutions of the Laplace equation in 2 variables. Very early in the XIX\supth{} century, Gauss was well aware of this connection and of the fact that one obtained solutions of the Laplace equation in 2 variables by replacing the function \eqref{eq:2.10} by
\begin{equation}
    \label{eq:2.13}
    \Omega(x, y) = \iint_D \rho(\xi, \eta) \log\frac{1}{r(x, y, \xi, \eta)} d\xi d\eta
\end{equation}
for a bounded domain \(D\) in the plane. The development by Cauchy of the theory of holomorphic functions of a complex variable could thus be used to yield properties of harmonic functions of 2 variables, such as for instance the non existence of relative extrema for such a function in its domain of definition; it was then natural to conjecture that similar properties were also valid for harmonic functions of 3 (and later for \(n \geq 4\)) variables, although they had to be proved by other means.

The first paper dealing with \emph{general} boundary conditions for a partial differential equation was written in 1828 by George Green, a self-taught English mathematician (1793--1841); it is concerned with electrostatics and the general study in that theory of what Green for the first time calls \emph{potential functions}. By that he not only means the functions of the form \eqref{eq:2.10}, but also what will later be called \emph{simple layer potentials}, namely functions of the type
\begin{equation}
    \label{eq:2.14}
    \Omega(M) = \iint_\Sigma \frac{\rho(P)}{MP} d\sigma(P)
\end{equation}
where \(\Sigma\) is a smooth surface, \(\rho\) (the ``density'') a continuous function on \(\Sigma\) and \(d\sigma\) the element of area on \(\Sigma\); he was naturally led to such functions by the known experimental fact that on conductors the electric charges are concentrated on their surface.

Green was interested in the relations between the surface density \(\rho\) and the potential it defines. He first establishes the famous theorem which, for the operator \(\Delta\), generalizes to 3 dimensions the relation between a differential operator and its adjoint (Chapter \ref{ch:1}, formula \eqref{eq:1.5}):
\begin{equation}
    \label{eq:2.15}
    \iiint_V \left( u\Delta v - v\Delta u \right) d\omega = \iint_\Sigma \left( v\pder{u}{n} - u\pder{v}{n} \right) d\sigma
\end{equation}
where \(\Sigma\) is a smooth surface limiting a bounded volume \(V\), \(u\) and \(v\) are \(C^2\) in a neighborhood of \(\overline{V}\), \(\pder{u}{n}\) is the derivative of \(u\) along the exterior normal of \(\Sigma\) \footnote{Lagrange \cite[vol.~I, p.~263]{135} and Gauss \cite[vol.~V, p.~22]{82} had already obtained more particular relations of that kind between volume and surface integrals.}. He then has the original idea \footnote{It is of course the same idea which leads to the Cauchy formula giving the value of a holomorphic function inside a domain \(D\) when it is known on the boundary of \(D\). However, it is unlikely that Green knew Cauchy's papers} of considering a function \(u\) which, still \(C^2\) for all points different from a point \(M\) in \(V\), becomes infinite at \(M\) in such a way that the difference \(u(P) - (1/MP)\) is bounded when \(P\) tends to \(M\); he applies \eqref{eq:2.15} to the volume \(V\) from which a small ball of center \(M\) has been excised, and by letting the radius of the ball tend to 0, he obtains the formula
\begin{equation}
    \label{eq:2.16}
    4\pi v(M) + \iiint_V \left( u\Delta v - v\Delta u \right) d\omega = \iint_\Sigma \left( v \pder{u}{n} - u \pder{v}{n} \right) d\sigma
\end{equation}
provided of course the triple integral exists. Taking in particular \(u(P) = 1/MP\) would give for a solution \(v\) of \(\Delta v = 0\)
\begin{equation}
    \label{eq:2.17}
    4\pi v(M) = \iint_\Sigma \left( v \pder{\left(\frac{1}{r}\right)}{n} - \frac{1}{r} \pder{v}{n} \right) d\sigma \qquad (\text{with } r(P)=MP)
\end{equation}
in other words, an integral formula which would solve the Laplace equation when \(v\) \emph{and} \(\pder{v}{n}\) were known on \(\Sigma\). This was in agreement with what was known at the time for partial differential equations of the second order, such as the equation of vibrating strings (Chapter \ref{ch:1}, \ref{sec:1.2}). However, experiments showed that \(v\) was entirely determined by its values on \(\Sigma\), and therefore it was not possible to take for both \(v\) and \(\pder{v}{n}\) on \(\Sigma\) \emph{arbitrary} continuous functions, so that the situation appeared quite different from the boundary conditions for hyperbolic equations. Furthermore, there was at least one case when an explicit formula gave \(v\) inside \(V\) by an integral extended to \(\Sigma\), namely the \emph{Poisson formula} for a ball \(V\) of center 0 and radius \(a\), published in 1820:
\begin{equation}
    \label{eq:2.18}
    v(M) = \frac{1}{4\pi} \iint_\Sigma \frac{a^2 - \rho^2}{ar^3} v(P) d\sigma
\end{equation}
with \(\rho = OM\). Green observed that one would have a similar formula for general domains \(V\):
\begin{equation}
    \label{eq:2.19}
    v(M) = \frac{1}{4\pi} \iint_\Sigma v(P) \pder{G}{n}(M, P) d\sigma
\end{equation}
by substituting in his formula \eqref{eq:2.16} for \(u\) a function \(G(M,P)\) such that: 1° in \(V \times V\), \(G\) is \(C^2\) provided \(M \neq P\) and \(\pder{G}{n}\) exists on \(\Sigma\); 2° \(G(M,P) - 1/MP\) remains bounded when \(P\) tends to \(M\); 3° \(G(M,P) = 0\) when \(M\) is in \(V\) and \(P\) on \(\Sigma\); 4° when \(M\) is fixed in \(V\), \(G\), as a function of \(P\), satisfies the Laplace equation in \(V\). He could not prove the existence of such a ``Green function'', but made it plausible by an appeal to experimental facts: when the surface \(\Sigma\) is connected to the ground, and an electric charge +1 is put at the point \(M\), it ``induces'' an electric charge on \(\Sigma\) such that the total potential of that charge and the punctual charge at \(M\) is 0 on \(\Sigma\); that potential should be the function \(G(M,P)\) (\cite{090}, \cite[p.~347--358]{S}).

Finally, by an ingenious use of his formula \eqref{eq:2.15}, Green could prove that in \(V \times V\), one had \(G(P,M) = G(M,P)\) for \(M \neq P\).

\section{The Dirichlet principle}
\label{sec:2.4}

Gauss had very early been interested in the Laplace equation, both in 2 variables in connection with his work on complex numbers, and in 3 variables in relation with his astronomical studies, and we have seen that in his 1813 paper on the attraction of spheroids, he had proved particular cases of the Green formula \eqref{eq:2.15}. After 1830, he devoted much of his time to the study of magnetism, both experimentally and theoretically, and thus was led to new research on potential theory, which he published in 1840 \cite[vol.~V, p.~197--242]{090}. In that paper, he quotes no other work on the subject, and it is very unlikely that he ever heard of Green (whose work was not widely known, even in England)\footnote{The fact that Gauss also uses the word ``potential'' with the same meaning may be attributed to the fact that the word (in its Latin form) was commonly used in the XVIII\supth{} century by ``natural philosophers''.}; he expands his 1813 formulas and obtains in this way some new particular cases of Green's formula \eqref{eq:2.15}, although he does not seem to have thought of formula \eqref{eq:2.16}. The closest approach to the latter is his famous ``mean value formula''
\begin{equation}
    \label{eq:2.20}
    v(0) = \frac{1}{4\pi} \iint_\Sigma v(P)d\sigma
\end{equation}
for a harmonic function \(v\) in a sphere \(\Sigma\) of center 0, for which it is quite surprising that he should not have observed that it was a special case of Poisson's formula \eqref{eq:2.18} which he cannot have failed to know.

As Green had done, Gauss was particularly interested in the behavior of simple layer potentials \eqref{eq:2.14} when \(M\) tends to a point on the surface \(\Sigma\); by a careful study, he shows that the potential \(\Omega\) is continuous everywhere, and that the normal derivatives at a point \(M_0\) of \(\Sigma\) exist on both sides of the surface, but have \emph{different values}, their difference being \(4\pi \rho(M_0)\); all this had been taken for granted without proof by Green.

Gauss attacked several problems related to potential theory, some of which were to become the focus of active research after 1930. One was the \emph{equilibrium problem}: find a distribution of electric charges on a closed surface \(\Sigma\) giving a potential which is \emph{constant} on \(\Sigma\); another consisted in replacing charges inside \(\Sigma\) by charges on \(\Sigma\) in such a way that the potential outside \(\Sigma\) remains the same (what would later be called a ``sweeping-out'' process), and Gauss showed that it could be solved if the equilibrium problem had a solution.

Regarding the latter, Gauss introduced a new idea which was to become quite central in potential theory: he observed that if the potential \(\Omega\) is given by \eqref{eq:2.14} with \(rho \geq 0\), and \(U\) is any continuous function on \(\Sigma\), then if \(\rho\) is chosen such that the integral \(\iint_\Sigma (\Omega - 2U)\rho d\sigma\) takes the \emph{smallest} possible value among all possible choices of \(\rho\), then \(\Omega - U\) is constant on \(\Sigma\), and he added that the existence of such a density \(\rho\) was obvious.

By adding to \(\Omega\) a suitable constant, this method of Gauss solved the problem of finding a harmonic function \(u\) in the volume \(V\), continuous in \(\overline{V} = V \cup \Sigma\), and equal on \(\Sigma\) to a given function \(U\) \footnote{If such a problem is solved, it implies the existence of the Green function: one considers the function \(u(M, P)\) harmonic in \(V\) (as a function of \(P\)) which takes the values \(-1/MP\) on \(\Sigma\); the Green function is then \(G(M, P) = u(M, P) + (1/MP)\), provided one shows that \(\pder{G}{n}\) exists and is continuous on \(\Sigma\).}. The same problem was considered a little later by W. Thompson (the future Lord Kelvin) in 1847 and by Dirichlet around the same time in his lectures (published long afterwards) \cite[p.~380--387]{S}; it became known as the \emph{Dirichlet problem}. Their idea is similar to Gauss's: they consider the volume integral
\begin{equation}
    \label{eq:2.21}
    \iiint_V \left( \left(\pder{v}{x}\right)^2 + \left(\pder{v}{y}\right)^2 + \left(\pder{v}{z}\right)^2 \right) d\omega
\end{equation}
and the function \(v\) continuous in \(\overline{V}\) with continuous and bounded first derivatives in \(V\) (\(v\) taking the given values on \(\Sigma\)), for which the integral \eqref{eq:2.21} takes its smallest value; applying the standard techniques of the Calculus of variations, they easily show that such a function is indeed harmonic in \(V\).

The great success of this idea is probably due to the imaginative use Riemann almost immediately made of it, in his epoch-making papers on holomorphic functions, Riemann surfaces and abelian integrals. By considering the real and imaginary parts of such functions, he was the first to realize that the existence theorems he needed could be derived from similar existence theorems for these harmonic functions, which he thought he could prove by adapting Dirichlet's argument to similar integrals in 2 variables, called by him ``Dirichlet principle'' \cite[p.~97]{182}.

His magnificent results attracted considerable attention, but soon mathematicians realized that they rested on three properties for which W. Thompson, Dirichlet and Riemann did not give any proof at all:

1) For a given continuous function \(g\) on \(\Sigma\), there exist continuous functions \(v\) in \(\overline{V}\) whose restriction to \(\Sigma\) is \(g\) and for which the integral \eqref{eq:2.21} is meaningful.

2) If such functions exist, there is one for which the smallest value of \eqref{eq:2.21} is attained.

3) For that function \(v\), the second order derivatives \(\frac{\partial^2 v}{\partial x^2}, {\partial^2 v}{\partial y^2}, {\partial^2 v}{\partial z^2}\) exist.

However, in 1871, F. Prym presented an example (for two variables and \(V\) the disk \(x^2 + y^2 < 1\)) where \emph{no function} \(v\) satisfying 1) existed \cite{181}\footnote{The discovery of that fact is usually attributed to Hadamard, who published a similar example in 1906 \cite[vol.~III, p.~1245--1248]{094}}. On the other hand, in 1870, Weierstrass observed that in all problems of the Calculus of variations which had been studied since the beginning of the XVIII\supth{} century, properties 2) and 3) had been taken for granted without any proof, and he gave a very simple example in which property 2) does not hold: the problem of minimizing the integral \(\int_{-1}^1 xy'^2dx\) among all \(C^1\) functions \(y\) defined in the interval [-1, 1] and satisfying the boundary conditions \(y(-1) = a, y(1) = b\), with \(a \neq b\) \cite[p.~390--391]{S}.

Spurred by these difficulties, Weierstrass and his pupils (P. Du Bois-Reymond, A. Kneser, S. Zaremba) undertook to put the Calculus of variations on sounder foundations and were able to rescue many classical results from the suspicion raised by such counterexamples. But the ``Dirichlet principle'' eluded their efforts, and it was only in 1899 that Hilbert, using new ideas in what was called his ``direct method'', was able to give a complete justification of the use Riemann had made of that ``principle'' \cite[vol.~III, p.~10--37]{111}.

\section{The Beer--Neumann method}
\label{sec:2.5}

We shall see in later chapters how the concepts and tools used by Hilbert and the Weierstrass school contributed to the birth of General Topology and later to the introduction of such notions as ``weak'' solutions of partial differential equations. Meanwhile, the challenge remained to prove the existence of a solution to the Dirichlet problem and similar boundary values problems for the Laplace equation, at least under conditions such as were used in Riemann's work. Between 1870 and 1890, that challenge was successfully taken up by three mathematicians: H.A. Schwarz around 1870, C. Neumann in 1877 and H. Poincaré in 1887.

We shall not discuss in detail the contributions of Schwarz and Poincaré, which did not influence directly the development of Functional Analysis. Both are based on the idea of \emph{approximation}: starting from known solutions of the Dirichlet problem for special kinds of domains, an approximation process enables one to get solutions for much more general domains. Schwarz limits himself to 2 variables; he first considers domains limited by a convex polygon, for which it is possible to prove directly (by explicit construction) the existence of a conformal mapping on the unit disk, hence the existence of a solution of the Dirichlet problem (by transferring the Poisson formula from the circle to the polygon). Using the maximum principle, it is then possible to prove the existence of the solution for a \emph{convex} domain by approximating it by a sequence of inscribed convex polygons. A little later, he invented an ingenious ``alternating process'' which enabled him to show that when one can solve the Dirichlet problem for two domains in the plane, it is also possible to solve it for their union, and from that result he finally showed that the Dirichlet problem in the plane is solvable for any domain limited by piecewise analytic curves \cite[vol.~II, p.~133--210]{196}.

Poincaré's famous ``sweeping-out method'' applies to any number of dimensions. To solve the Dirichlet problem for a bounded domain \(V\) limited by a surface \(\Sigma\), he shows (using the maximum principle) that it is enough to consider the case in which the function given on \(\Sigma\) is the restriction to \(\Sigma\) of a function \(\Phi\) defined in a neighborhood \(W\) of \(\overline{V}\), of class \(C^2\) and such that \(\Delta\Phi \geq 0\). By Poisson's equation \eqref{eq:2.12}, \(\Phi\) is the sum of a harmonic function and a potential \(\Phi_0\) of masses \(\geq 0\). The fundamental idea is that if \(B\) is a ball contained in \(W\), it is possible to use the Poisson integral \eqref{eq:2.18} extended to the surface of \(B\) in order to replace \(\Phi_0\) by another potential which coincides with \(\Phi_0\) outside \(B\) and is smaller than \(\Phi_0\) inside \(B\); the masses inside \(B\) have been ``swept out'' on the surface of \(B\). One then takes an infinite sequence of balls \(B_n\) whose union is \(V\), and one applies the ``sweeping-out'' process repeatedly to the \(B_n\) in the order \(B_1, B_2, B_1, B_2, B_3, B_1, B_2, B_3, B_4, \ldots\) (each \(B_n\), is ``swept-out'' in finitely many times). The corresponding sequence of potentials is decreasing, hence has a limit in \(V\); using Harnack's inequalities (consequences of the Poisson formula \eqref{eq:2.18}) and the maximum principle, Poincaré is able to show that this limit is a solution of the Dirichlet problem, provided the boundary \(\Sigma\) satisfies a ``regularity'' condition, namely, for any point \(M \in \Sigma\), there must be a small ball whose intersection with \(\overline{V}\) is reduced to \(M\) \cite[vol.~IX, p.~33--54]{177}; later, Zaremba could replace the small ball by a small cone of vertex \(M\) in that condition.

In contrast with Schwarz's and Poincaré's papers, the Beer--Neumann method was a landmark in Functional Analysis by introducing the first example of what was later to be called a ``Fredholm integral equation of the second kind''. Green's formula \eqref{eq:2.17} naturally introduced still another type of potential:
\begin{equation}
    \label{eq:2.22}
    u(M) = \iint_\Sigma \rho(P) \frac{\partial}{\partial n} \left( \frac{1}{MP} \right) d\sigma
\end{equation}
which was harmonic outside the surface \(\Sigma\). It also occurred in the theory of magnetism, from which it got its name of \emph{double layer potential}: it was there conceived as the limit of a difference of two simple layer potentials, one with density \(\mu\) on \(\Sigma\), the other with density \(\mu\) on a surface \(\Sigma'\) parallel to \(\Sigma\) and at an ``infinitely small'' distance \(\varepsilon\); when \(\varepsilon\) tends to 0, \(\mu\) was supposed to increase to \(+\infty\) in such a way that the product \(\mu \varepsilon\) tended to \(rho\).

Such a potential had been shown to have near \(\Sigma\) a behavior quite similar to the normal derivative of a simple layer potential, studied by Gauss: when \(M\) tends to a point \(M_0\) of \(\Sigma\) along the normal to \(\Sigma\) at \(M_0\), \(u(M)\) tends to a limit on each side of \(\Sigma\), but these limits are \emph{different} in general; however, \(\pder{u}{n}\) is the \emph{same} on both sides.

Formula \eqref{eq:2.22} also had a nice geometric interpretation; one has \(\frac{\partial}{\partial n} \left( \frac{1}{MP} \right) = \frac{\cos \varphi}{MP^2}\) where \(\varphi\) is the angle between \(MP\) with the normal to \(\Sigma\) at \(P\), and \(\frac{\cos \varphi}{MP^2} d\sigma\) is the infinitesimal ``solid angle'' from which \(d\sigma\) is ``seen'' from the point \(M\).

Around 1860, C. Beer proposed to obtain a solution to the Dirichlet problem by formula \eqref{eq:2.22} for a suitable density \(\rho\) on \(\Sigma\). From the continuity properties of double layer potentials, it follows that if \(\Sigma\) is a smooth surface, and \(g(M)\) is the function on \(\Sigma\) to which the solution \(u(M)\) must be equal, the unknown density must satisfy the equation
\begin{equation}
    \label{eq:2.23}
    2\pi \rho(M) + \iint_\Sigma \rho(P) \frac{\partial}{\partial n}\left(\frac{1}{MP}\right) d\sigma = g(M) \quad \text{for } M \in \Sigma.
\end{equation}
He then concluded that one could compute \(\rho\) by the usual device of ``successive approximations'' (\ref{sec:2.1}) starting with \(\rho_0(M) = \frac{1}{2\pi} g(M)\) and defining recursively \(\rho_n(M)\) by
\[
    2\pi \rho_n(M) + \iint_\Sigma \rho_{n-1}(P) \frac{\partial}{\partial n}\left(\frac{1}{MP}\right) d\sigma = 0 \quad \text{for } n \geq 1
\]
so that the series \(\rho(M) = \rho_0(M) + \rho_1(M) + \ldots + \rho_n(M) + \ldots\) would give the solution to \eqref{eq:2.23}; but he made no attempt to prove that the series converged.

In 1877, Carl Neumann attempted to give such a proof \cite{165}. He restricted himself to the case in which the domain \(V\) is bounded and \emph{convex}, but he allowed a non smooth boundary \(\Sigma\); equation \eqref{eq:2.23} must then be modified to
\begin{equation}
    \label{eq:2.24}
    4\pi \rho(M) = \iint_\Sigma \left( \rho(M) - \rho(P) \right) \frac{\cos \varphi}{MP^2} d\sigma + f(M)
\end{equation}
with \(f\) continuous on \(\Sigma\), and the successive approximations are given by \(4\pi\rho_0(M) = f(M)\) and, for \(n \geq 1\),
\begin{equation}
    \label{eq:2.25}
    4\pi \rho_n(M) = \iint_\Sigma \left( \rho_{n-1}(M) - \rho_{n-1}(P) \right) \frac{\cos \varphi}{MP^2} d\sigma
\end{equation}
Neumann's idea is to consider the maximum value \(L_n\) and minimum value \(l_n\) of \(\rho_n\), and to show that there is a number \(q\) such that \(0 < q < 1\) and
\begin{equation}
    \label{eq:2.26}
    L_n - l_n \leq \left( L_0 - l_0 \right) q^{n - 1}
\end{equation}
from which he majorizes \(|\rho_n(M)|\) by a multiple of \(q^n\) using \eqref{eq:2.25}, and he can conclude that the series \(\sum\limits_{n=0}^\infty \rho_n(M)\) converges to a continuous function.

To prove \eqref{eq:2.26}, Neumann divides \(\Sigma\) into two parts \(A_n, B_n\) respectively defined by the conditions
\[
    \begin{aligned}
        \frac{1}{2} \left( L_{n-1} + l_{n-1} \right) \leq \rho_{n-1}(P) \leq L_{n-1} \quad &\text{for } A_n \\
        l_{n-1} \leq \rho_{n-1}(P) < \frac{1}{2} \left( L_{n-1} + l_{n-1} \right) \quad &\text{for } B_n
    \end{aligned}
\]
and he deduces from \eqref{eq:2.25} that for all points \(M\) of \(\Sigma\)
\[
    \left( L_{n-1} - l_{n-1} \right) \left( A_n(M) + \frac{1}{2} B_n(M) \right) \leq 4\pi \rho_n(M) \leq \left( L_{n-1} - l_{n-1} \right) \left( \frac{1}{2} A_n(M) + B_n(M) \right)
\]
where \(A_n(M)\) and \(B_n(M)\) are the solid angles from which \(A_n\) and \(B_n\) are ``seen'' from \(M\). This implies
\[
    L_n - l_n \leq \left(L_{n-1} - l_{n-1}\right) q
\]
where \(q\) is the \emph{least upper bound} of the quantity
\begin{equation}
    \label{eq:2.27}
    \Lambda(M, M', A, B) = \frac{1}{4\pi} \left( \frac{1}{2} A(M) + B(M) + A(M') + \frac{1}{2} B(M') \right)
\end{equation}
when \(M\) and \(M'\) vary arbitrarily in \(\Sigma\), \(A\) is an arbitrary closed part of \(\Sigma\) and \(B\) its complement. One is thus faced with the \emph{purely geometric} problem of showing that \(q < 1\). The expression \eqref{eq:2.27} can be written
\[
    \frac{1}{4\pi} \left( A(M) + B(M) + A(M') + B(M') - \frac{1}{2} \left( A(M) + B(M') \right) \right)
\]
and also
\[
    \frac{1}{4\pi} \left( \frac{1}{2} \left( A(M) + B(M) \right) + \frac{1}{2} \left( A(M') + B(M') \right) + \frac{1}{2} \left( A(M') + B(M) \right) \right)
\]
and as one always has \(A(M) + B(M) \leq 2\pi\) (maximum value of the solid angle from which the whole of \(\Sigma\) is ``seen'' from one of its points), the problem can also be formulated in two equivalent ways:
\begin{equation}
    \label{eq:2.28}
    A(M) + B(M') \geq 4\pi r \quad \text{for an } r > 0,
\end{equation}
\begin{equation}
    \label{eq:2.29}
    A(M) + B(M') \leq 4\pi s \quad \text{for an } s > 1,
\end{equation}
for all points \(M, M'\) in \(\Sigma\) in two parts \(A, B\). The form \eqref{eq:2.28} of that condition immediately shows that there is an \emph{exceptional} type of convex set for which it cannot be satisfied, namely the case in which \(V\) is the \emph{intersection of two convex cones} (``double cone''): indeed we then have \(A(M) = B(M') = 0\) if \(A\) is the surface of one of the cones, \(B\) the surface of the other, \(M\) the vertex of \(A\) and \(M\) the vertex of \(B\). Furthermore, this particular choice of \(A, B, M\) and \(M'\) is the \emph{only one} for which \(A(M) + B(M')\) may be 0. However, when the exceptional case is excluded, Neumann concludes, from the fact that \(A(M) + B(M') > 0\) for all choices of \(A, B, M\) and \(M'\), that there is an \(r > 0\) for which \eqref{eq:2.28} is satisfied for all these choices, and does not give a proof of that assertion valid for all convex sets other than double cones. This gap in Neumann's proof seems to have remained undetected until Lebesgue drew attention to it is 1937 (\cite{138}, vol.~IV, p.~151--166). He shows in addition how one can fill in that gap by a compactness argument: there are two points \(M_0, M_0'\) in \(\Sigma\), limits of sequences \((M_k), (M'_k)\) such that for each \(k\) there is a splitting of \(\Sigma\) in two parts \(A_k, B_k\) such that \(A_k(M_k) + B_k(M'_k)\) tends to the l.u.b. \(4\pi s\) of \(A(M) + B(M')\) for all choices of \(A, B, M, M'\). On the other hand there are a point \(N\) of \(\Sigma\) and neighborhoods \(V(M_0), V(M'_0), V(N)\) of \(M_0, M'_0, N\) respectively in \(\Sigma\) such that the planes of support at all points of \(V(N)\) do not intersect \(V(M_0)\) nor \(V(M'_0)\) (it is here that the assumption that \(V\) is not a double cone is used); an elementary geometrical argument then gives an upper bound \(< 1\) for \(s\). Historically, such an argument would have been barely possible in the late 1870's, but I strongly doubt that C. Neumann was familiar enough with the use of the ``Bolzano--Weierstrass'' theorem (as it was called at that time) to have thought of it. He was apparently satisfied with the fact that for simple convex sets, such as ellipsoids, it was possible to compute explicitly an upper bound \(< 1\) for \(s\).

C. Neumann dealt in the same way with the Dirichlet problem in the plane, with a similar gap in his proof.

For a long time, the restrictions on the surface \(\Sigma\) in all the existence proofs of the Dirichlet problem were thought to be imperfections of the methods of proof; but in 1912, Lebesgue gave an example (in 3 dimensions) of a bounded open set \(V\) (homeomorphic to a ball) such that there is a continuous function on the boundary \(\Sigma\) of \(V\), for which the Dirichlet problem has no solution (\cite{138}, vol.~IV, p.~131). This was the starting point of modern Potential theory, where, on one hand, the initial formulation of the Dirichlet problem is modified in such a way that it always has a unique ``solution'' for any bounded domain, the word ``solution'' being interpreted in some ``weak'' sense; on the other hand, the behavior of these ``weak'' solutions on the boundary of the domain is investigated under various conditions \cite{030}. The detailed history of that extensive theory is outside the scope of this book.

\chapter{The equation of vibrating membranes}
\label{ch:3}
\setcounter{equation}{0}

\section{H.A. Schwarz's 1885 paper}
\label{sec:3.1}

The same physical arguments which lead to the equation of vibrating strings (Chap. \ref{ch:1}, \ref{sec:1.2}, equation \eqref{eq:1.7}) apply to the small vibrations of a membrane which at rest is in the plane \(Oxy\), and has a constant density: if \(z = u(x, y, t)\) is the equation of its surface at time \(t\), the function \(u\) satisfies the equation
\begin{equation}
    \label{eq:3.1}
    \frac{\partial^2 u}{\partial x^2} + \frac{\partial^2 u}{\partial y^2} = \frac{\partial^2 u}{\partial t^2}
\end{equation}
(for suitable units of length and time). The usual method of ``separation of variables'' consists here in looking for solutions \(u(x, y, t) = v(x, y)w(t)\) and one finds for \(v\) the equation (also called ``Helmholtz's equation'')
\begin{equation}
    \label{eq:3.2}
    \frac{\partial^2 u}{\partial x^2} + \frac{\partial^2 u}{\partial y^2} + \lambda v = 0
\end{equation}
for a constant \(\lambda\). If in addition the membrane at rest is a bounded portion \(\Omega\) of the plane and is \emph{fixed} at its boundary \(\Sigma\) (which means that \(u(x, y, t) = 0\) for all \(t\) if \((x,y) \in \Sigma\)), \(\lambda\) must be \(> 0, w(t) = \sin \sqrt{\lambda t}\), and one has to find a solution \(v\) of \eqref{eq:3.2} which vanishes on \(\Sigma\) and is not identically 0. Contrasting with the easy solution of the corresponding problem for the vibrating string, the elucidation of that problem was going to challenge the ingenuity of mathematicians during the whole second half of the XIX\supth{} century.

Experimental evidence, as well as the explicit solution of the problem for very special domains \(\Omega\), such as a rectangle or a disk, showed that, just as in the case of the vibrating string, solutions of \eqref{eq:3.2} vanishing on \(\Sigma\) and not identically 0 could only exist when \(\lambda\) was equal to one of an infinite sequence \((\lambda_n)\) of real numbers \(> 0\) (the ``eigenvalues'' of the problem), tending to \(+\infty\).

The first attempt to prove such a result for general domains \(\Omega\) was made by H. Weber in 1869 \cite{247}, by an adaptation of the variational method used by Riemann for the Dirichlet problem. Using Green's formula (Chapter \ref{ch:2}, formula \eqref{eq:2.15}) he first shows that if \(\mu_1, \mu_2\) are two distinct eigenvalues, \(v_1, v_2\) corresponding ``eigenfunctions'', then
\begin{equation}
    \label{eq:3.3}
    \left(\mu_1 - \mu_2\right) \iint_\Omega v_1(x, y) v_2(x, y) dxdy = 0
\end{equation}
from which he deduces, as Poisson had done for ordinary differential equations (Chap. \ref{ch:1}, \ref{sec:1.3}), that the eigenvalues are necessarily real numbers. To determine the smallest eigenvalue \(\lambda_1\), he considers the Dirichlet integral
\begin{equation}
    \label{eq:3.4}
    F(v) = \iint_\Omega \left( \left(\pder{v}{x}\right)^2 + \left(\pder{v}{y}\right)^2 \right) dxdy
\end{equation}
for \(C^2\) functions \(v\) in \(\overline\Omega\), equal to 0 on \(\Sigma\) and subject to the additional constraint
\begin{equation}
    \label{eq:3.5}
    \iint_\Omega v^2 dxdy = 1.
\end{equation}

He assumes, as Riemann, that in this set \(\mathcal{F}_1\) of functions, there is one for which \(F(v)\) is equal to its greatest lower bound \(\lambda_1\) and by the usual methods of the Calculus of variations, he shows that this function \(v_1\) is a solution of \eqref{eq:3.2} for \(\lambda = \lambda_1\).

He next considers the subset \(\mathcal{F}_2\) of \(\mathcal{F}_1\) defined by the additional condition
\begin{equation}
    \label{eq:3.6}
    \iint_\Omega v(x, y)v_1(x, y)dxdy = 0,
\end{equation}
takes the function \(v_2 \in \mathcal{F}_2\) for which \(F(v_2)\) is equal to its greatest lower bound \(\lambda_2\) and shows that \(v_2\) is a solution of \eqref{eq:3.2} for \(\lambda = \lambda_2\). The induction process is then obvious, and Weber concludes that he has proved the existence of an increasing infinite sequence \(\left(\lambda_n\right)\) of positive eigenvalues to each of which there corresponds an eigenfunction \(v_n\) normalized by condition \eqref{eq:3.5}, and orthogonal to each other. But he does not try to prove that \(\lim\limits_{n \to \infty} \lambda_n = +\infty\), nor that functions in \(\mathcal{F}_1\) possess a ``Fourier expansion'' \(\sum\limits_n c_n v_n\) defined in the same manner as in the Sturm--Liouville problem (Chap. \ref{ch:1}, \ref{sec:1.3}, formula \eqref{eq:1.33}) (a result which he states however, without proof).

Weber's proofs were of course subject to the general criticisms of Weierstrass against the Calculus of variations, but no one seems to have tried to find more rigorous ones until 1885, In that year, H.A. Schwarz published a long paper on the theory of minimal surfaces, in which he had to consider a type of equation slightly more general than \eqref{eq:3.2}:
\begin{equation}
    \label{eq:3.7}
    \frac{\partial^2 v}{\partial x^2} + \frac{\partial^2 v}{\partial y^2} + \lambda^2 pv = 0
\end{equation}
where \(p\) is a continuous function in a domain \(D\), with values \(> 0\); his arguments apply for any such function, but in fact he is only interested in the particular case \(p(x, y) = - 8/(1 + x^2 + y^2)^2\) \cite[vol.~I, p.~223--269]{195}.

Schwarz's paper is extremely remarkable by the originality of its methods, which do not seem to have been inspired by any previous work; it may be that the study of the Sturm--Liouville problem led him to arguments which later could be transferred almost verbatim to general integral equations with symmetric kernels (see Chap. \ref{ch:5}, \ref{sec:5.2}), but there is no hint in his paper of such an influence, and in fact he quotes nobody, not even Weber.

His starting point is not the problem of existence of eigenvalues \(\lambda^2\) for equation \eqref{eq:3.7}, but a ``Dirichlet problem'' for the equation
\begin{equation}
    \label{eq:3.8}
    \Delta w + \xi pw = 0
\end{equation}
depending on a parameter \(\xi\); he limits himself to the case where \(w\) is subject to the condition of being equal to 1 on the boundary \(\Gamma\) of \(D\). Using the time honored method of representing the solution as a power series in \(\xi\) (Chap. \ref{ch:2}, \ref{sec:2.1})
\begin{equation}
    \label{eq:3.9}
    w = w_0 + \xi w_1 + \ldots + \xi^n w_n + \ldots
\end{equation}
he takes for \(w_0\) the constant function equal to 1, and imposes on the \(w_n\) for \(n \geq 1\) to vanish on \(\Gamma\); they are then determined inductively by the equations
\begin{equation}
    \label{eq:3.10}
    \Delta w_n + pw_{n-1} = 0 \quad \text{for } n \geq 1.
\end{equation}
He assumes that the Green function \(G(M, P)\) for the domain \(D\) exists (remember that he himself had proved that existence in extensive cases (Chap. \ref{ch:2}, \ref{sec:2.4})); the properties of that function implied that for any function \(f\) continuous in \(\overline{D}\) the equation
\begin{equation}
    \label{eq:3.11}
    \Delta w + f = 0
\end{equation}
has a unique solution vanishing on \(\Gamma\), given by the formula
\begin{equation}
    \label{eq:3.12}
    w(M) = \frac{1}{2\pi} \iint_D f(P) G(M, P) d\omega \qquad (\text{with } d\omega = dxdy)
\end{equation}
Therefore his functions are given explicitly by
\begin{equation}
    \label{eq:3.13}
    w_n(M) = \frac{1}{2\pi} \iint_D p(P) w_{n-1}(P) G(M, P) d\omega.
\end{equation}

One must now investigate the convergence of \eqref{eq:3.9} for small enough values of \(|\xi|\), and it is here that Schwarz's original contributions begin. His main tool is the inequality named after him\footnote{That inequality had been discovered by Buniakowsky in 1859, but does not seem to have been noticed nor used by many mathematicians before 1885. If is of course a direct generalization of the corresponding inequality for finite sums, which goes back at least to Cauchy.}
\[
    \left( \iint_D fg d\omega \right)^2 \leq \left( \iint_D f^2 d\omega \right) \left( \iint_D g^2 d\omega \right)
\]
for any two functions \(f, g\) continuous in \(\Gamma\); this gives from \eqref{eq:3.13}
\begin{equation}
    \label{eq:3.14}
    \begin{gathered}
    4\pi^2\left(w_n(M)\right)^2 \leq \left(\iint_D p^2(P)G^2(M, P)d\omega\right) \left(\iint_D w^2_{n-1}(P)d\omega\right) \\
    \leq A \left( \iint_D w^2_{n-1}(P)d\omega \right)
    \end{gathered}
\end{equation}
where \(A\) is a constant independent of \(n\) (due to the properties of the Green function of a bounded domain). Schwarz is thus led to study the numbers
\begin{equation}
    \label{eq:3.15}
    W_{n, k} = \iint_D p w_k w_{n-k} d\omega
\end{equation}
which, using the \emph{symmetry} of the Green function, he shows are independent of \(k\), so that \(W_{n, k} = W_{n, 0}\), which he writes \(W_n\). He also proves that
\begin{equation}
    \label{eq:3.16}
    W_{n, k} = \iint_D \left(\pder{w_{k+1}}{x} \pder{w_{n-k}}{x} + \pder{w_{k+1}}{y} \pder{w_{n-k}}{y}\right) dxdy.
\end{equation}
Finally, using the Schwarz inequality, he obtains the relation
\begin{equation}
    \label{eq:3.17}
    W_n^2 \leq W_{n-1} W_{n+1},
\end{equation}
hence the sequence of numbers \(W_n/W_{n-1}\) is increasing; on the other hand, integrating \eqref{eq:3.14} gives \(W_{2n} \leq BW_{2n-2}\) for a constant \(B\) independent of \(n\), and therefore the limit of the sequence \((W_n/W_{n-1})\) is a finite number \(c > 0\). It follows then from \eqref{eq:3.14} that the series \eqref{eq:3.9} is absolutely and uniformly convergent in \(\overline{D}\) for \(|\xi| < 1/ \sqrt{c}\); the properties of the Green function enable one to show that the derivatives of \(w\) are also given by convergent series obtained by differentiating \eqref{eq:3.9} termwise, and that \(w\) then satisfies \eqref{eq:3.8} and is equal to 1 on the boundary \(\Gamma\).

But Schwarz goes one step further. He proves that when \(\xi = 1/\sqrt{c}\), the general term of the series \eqref{eq:3.9} tends uniformly to a limit \(U_1\) which is not identically 0 in \(\overline{D}\) but vanishes on the boundary and is solution of
\begin{equation}
    \label{eq:3.18}
    \Delta w + (1/c) pw = 0.
\end{equation}
He has thus proved the existence of the \emph{smallest eigenvalue} \(\lambda_1^2 = 1/c\) of the equation \eqref{eq:3.8} for functions vanishing on the boundary, and of the corresponding eigenfunction.

It should be observed here that these developments in fact are just another treatment of a ``crypto-integral'' equation (which Schwarz does not write, however). If one writes \(w = w_0 + \xi v\) and ``solves'' equation \eqref{eq:3.8} by formula \eqref{eq:3.12} (using the same idea as Liouville in 1837 to obtain his ``Volterra integral equation'' (Chap. \ref{ch:2}, \ref{sec:2.1} and Chap. \ref{ch:1}, \ref{sec:1.3}, equation \eqref{eq:1.35})), one gets for \(v\) this time a ``Fredholm integral equation''
\begin{equation}
    \label{eq:3.19}
    v(M) = g(M) + \frac{\xi}{2\pi} \iint_D p(P) G(M, P) v(P) d\omega
\end{equation}
with
\[
    g(M) = \frac{1}{2\pi} \iint_D G(M, P) p(P) d\omega.
\]
Schwarz's procedure is therefore essentially the same as C. Neumann's for the Dirichlet problem (Chap. \ref{ch:2}, \ref{sec:2.4}), at least as a starting point; the main difference is in the emphasis put by Schwarz on the dependence on the parameter \(\xi\).

To appreciate the originality and power of Schwarz's method, it is perhaps not superfluous to show how it can be translated, almost without change, in the theory of self-adjoint compact operators in a separable Hilbert space \(E\). Suppose \(U\) is such an operator in \(E\), which in addition we suppose \emph{positive}, i.e. \((U\cdot f| f) \geq 0\) for all \(f \in E\). The spectrum of \(U\) then consists in a decreasing sequence (finite or infinite) \(\mu_1 \geq \mu_2 \geq \ldots \geq \mu_n \geq \ldots 0\), where each \(\mu_n\) is an eigenvalue counted a number of times equal to its multiplicity; 0 is always in the spectrum but \(\Ker(U)\) may be reduced to 0 or have infinite dimension; for each \(\mu_n\) there is an eigenvector \(\varphi_n\) of norm 1, such that \(E\) is the \emph{Hilbert sum} of the one-dimensional spaces \(\C \varphi_n\) and of \(\Ker(U)\). Let
\[
    w_0 = \sum_n d_n \varphi_n + w'_0
\]
with \(w'_0 \in \Ker(U)\), be the expression of a vector \(w_0 \in E\) for that decomposition. Then, for any \(m \geq 1\), we have
\[
    U^m \cdot w_0 = \sum_n \mu^m_n d_n \varphi_n
\]
and therefore the Schwarz series \eqref{eq:3.9} is equal to
\[
    w = \sum_{m=0}^\infty \xi^m U^m \cdot w_0 = \sum_n \left( \sum_{m=0}^\infty \xi^m \mu^m_n \right) d_n \varphi_n = \sum_n \frac{d_n}{1 - \xi \mu_n} \varphi_n
\]
provided \(|\xi| < \mu_1^{-1}\), and we have \(w = \xi U \cdot w + w_0 - w'_0\). For \(\xi = 1/\mu_1\),
\[
    \xi^m U^m \cdot w_0 = \sum_n \xi^m \mu^m_n d_n \varphi_n = \sum_n \left( \mu_n / \mu_1 \right)^m d_m \varphi_n
\]
tends to \(d_1 \varphi_1\) if \(\mu_1\) is a simple eigenvalue, to the sum of the \(d_n \varphi_n\) such that \(\mu_n = \mu_1\) in general. Finally, we have
\[
    W_m = \left( U^m \cdot w_0 | w_0 \right) = \left( U^{m-k} \cdot w_0 | U^k \cdot w_0 \right) = \sum_n \mu^m_n d^2_n
\]
from which it follows that if the \(d_n\) such that \(\mu_n = \mu_1\) are not all 0, the ratio \(W_m / W_{m-1}\) tends to \(\mu_1\); furthermore, one has
\[
    \begin{gathered}
        W_{2m} = \norm{U^m \cdot w_0}^2 = |\left( U^{m+1} \cdot w_0 | U^{m-1} \cdot w_0 \right)| \leq \norm{U^{m+1} \cdot w_0} \cdot \norm{U^{m-1} \cdot w_0} \\
        = \left( W_{2m-2} W_{2m+2} \right)^\frac{1}{2}.
    \end{gathered}
\]
To get the inequality \(W_m \leq \left( W_{2m-2} W_{2m+2} \right)^\frac{1}{2}\) for all integers \(m\), it is enough to consider the unique compact positive operator \(V\) such that \(V^2 = U\), and apply to \(V\) the preceding argument. Of course the concept of the ``square root'' of a positive self-adjoint operator was not available to Schwarz, and this is why he had to use the expression \eqref{eq:3.16} for his numbers \(W_n\).

In 1893, E. Picard published a short Comptes-Rendus Note (\cite{172}, vol.~II, p.~545--550) in which he went one step further. For any point \(M \in D\), the function \(w(M, \xi)\) given by Schwarz's series \eqref{eq:3.9} is holomorphic in the circle \(|\xi| < \xi_1 = 1/\sqrt{c}\), and Picard investigates the \emph{analytic continuation} of \(\xi \mapsto w(M, \xi)\) beyond that circle: he shows that such a continuation exists in a circle \(|\xi| < \xi_2\), of radius independent of \(M \in D\), and that it has a \emph{simple pole} with residue \(-\xi_1 U_1(M)\) at the point \(\xi_1\). He limits himself to the case in which \(p = 1\) and \(\Gamma\) is convex and smooth, and his idea is to adapt the method of C. Neumann (Chap. \ref{ch:2}, \ref{sec:2.4}) to evaluate the differences \(\left| \xi^n_1 w_n - \xi^{n-1}_1 w_{n-1} \right|\); with apparently the same gap as in Neumann's argument (the details are not given in the Note) he ``proves'' that there are constants \(C\) and \(q < 1\) independent of \(M\), such that \(\left| \xi^n_1 w_n - \xi^{n-1}_1 w_{n-1} \right| \leq Cq^n\), hence \(\left| \xi^n_1 w_n - U_1 \right| \leq C'q^n\) for another constant \(C'\), hence his result. writing
\[
    w = \frac{U}{1 - (\xi / \xi_1)} + v
\]
he looks for a power series development
\begin{equation}
    \label{eq:3.20}
    v = v_0 + \xi v_1 + \ldots \xi^n v_n + \ldots
\end{equation}
similar to \eqref{eq:3.9} but which should converge in a circle \(|\xi| < \xi_2\) with \(\xi_2 > \xi_1\). He determines the \(v_n\) by the successive approximations
\[
    \Delta v_0 - \xi_1 U_1 = 0, \qquad \Delta v_n + v_{n-1} = 0 \quad \text{for } n \geq 1
\]
with the boundary conditions: \(v_0 = 1\) on \(\Gamma\) and \(v_n = 0\) on \(\Gamma\) for \(n \geq 1\). Introducing numbers similar to the \(W_{n, k}\) of Schwarz, he is able to prove that the radius of convergence \(\xi_2\) of \eqref{eq:3.20} is finite, but he cannot show that there is an eigenfunction corresponding to \(\xi_2\) and vanishing on \(\Gamma\).

\section{The contributions of Poincaré}
\label{sec:3.2}

In 1890, H. Poincaré published in the American Journal of Mathematics a long paper developing some of his research done since 1887, which had been announced in three Comptes-Rendus Notes (\cite{177}, vol.~IX, p.~15--113). The paper consists of two completely independent parts; in the first, he describes in detail his ``sweeping-out'' method for the solution of the Dirichlet problem (Chap. \ref{ch:2}, \ref{sec:2.4}). The second part is devoted to the \emph{cooling off} problem in the theory of heat, which had been treated by Fourier in some particular cases, for instance the cooling off of a sphere when the temperature is a function of the distance to the center (Chap. \ref{ch:1}, \ref{sec:1.2}). The \emph{general} cooling off problem had been presented by Fourier in the following form: given a solid body \(V\) of constant density, isotropic for the propagation of radiations, one has to find the temperature \(u(x,y,z,t)\) inside \(V\), as a function of the coordinates \(x, y, z\) and the time \(t\), when the outside temperature is 0. Fourier shows that the function \(u\) must satisfy inside \(V\) an equation (where \(a\) is constant)
\begin{equation}
    \label{eq:3.21}
    \pder{u}{t} = a^2 \Delta u
\end{equation}
and in addition is subject to the boundary condition on the surface \(\Sigma\) of \(V\)
\begin{equation}
    \label{eq:3.22}
    \pder{u}{n} + hu = 0
\end{equation}
where \(\pder{u}{n}\) is the normal derivative (towards the exterior) and \(h\) is a constant \(\geq 0\) (see Chap. \ref{ch:4}, \ref{sec:4.4}). The usual method of ``separation of variables'' led to solutions of the form \(u(x,y,z,t) = e^{-\lambda a^2 t} v(x,y,z)\), where \(v\) should be a solution of the Helmholtz equation
\begin{equation}
    \label{eq:3.23}
    \Delta v + \lambda v = 0
\end{equation}
with a \emph{different} boundary condition from the one deriving from the equation of vibrating membranes, namely
\begin{equation}
    \label{eq:3.24}
    \pder{v}{n} + hv = 0 \quad \text{on } \Sigma.
\end{equation}
In his 1869 paper, H. Weber had also considered that problem, but he had only described his variational method to obtain eigenvalues and eigenfunctions for the particular case \(h=0\). Poincaré apparently was unaware of Weber's paper and never mentioned it in his own work; what he does in 1890 is first to repeat Weber's arguments for the general boundary condition \eqref{eq:3.24}, replacing the Dirichlet integral by the function
\begin{equation}
    \label{eq:3.25}
    F(v) = h\iint_\Sigma v^2 d\sigma + \iiint_V \left(\divergence{v}\right)d\omega.
\end{equation}
Having thus obtained an increasing infinite sequence \(\left(\lambda_n\right)\) of eigenvalues and the corresponding sequence \(\left(v_n\right)\) of eigenfunctions, Poincaré is of course aware of the non rigorous character of his ``proof''; however, having for the time being no better arguments at his disposal, he takes for granted the existence of \(\lambda_n\) and \(v_n\) and proceeds to study them in more detail, and in the first place to prove that the sequence \(\left(\lambda_n\right)\) tends to \(+\infty\), a question which Weber had not been able to answer. In his attack on that problem, it is quite remarkable to see Poincaré introducing a whole batch of \emph{completely new ideas}. In the first place, he considers the eigenvalues as functions \(\lambda_n(h, V)\) of the constant \(h\) and the domain \(V\) and begins to study the way in which they \emph{depend} on \(h\) and \(V\), a trend of thought which will later blossom in the work of H. Weyl and R. Courant, and even now has not entirely lost its interest. Poincaré first shows that, for \(V\) fixed, \(\lambda_n(h, V)\) is increasing with \(h\), by an application of Green's formula to the eigenfunctions \(v_n(h, V), v_n(h', V)\) corresponding to two values of \(h\); as he wants to prove that \(\lambda_n(h, V)\) tends to \(+\infty\), he can assume that \(h = 0\), which implies that \(\lambda_1(0, V) = 0\) and \(v_1(0, V)\) is a constant.

The second idea is to \emph{decompose} \(V\) into a union of smaller solids \(V_1, V_2, \ldots, V_p\); the variational definition of \(\lambda_n\) enables him to prove that if \(p \leq n-1, \lambda_n(0, V)\) is \emph{at least equal} to the smallest of the numbers \(\lambda_2(0, V_1), \ldots, \lambda_2(0, V_p)\). Poincaré is thus led to \emph{minorize} \(\lambda_2(0, V)\) by a number depending only on the geometry of \(V\); by definition (since \(h = 0\)), this means finding a lower bound of the expression
\begin{equation}
    \label{eq:3.26}
    \frac{\iiint_V \left(\divergence{v}\right) d\omega}{\iiint_V v^2 d\omega}
\end{equation}
where \(v\) is a \(C^2\) function in \(\overline{V}\), subject to the condition
\begin{equation}
    \label{eq:3.27}
    \iiint_V v d\omega = 0.
\end{equation}
He assumes \(V\) is \emph{convex}; using polar coordinates and the standard methods of the Calculus of variations, he obtains as lower bound
\begin{equation}
    \label{eq:3.28}
    C \cdot \vol(V)/(\diam(V))^5
\end{equation}
where \(C\) is an absolute constant; one should here stress the fact that this Poincaré inequality is the first example of what we now call \emph{``a priori'' inequalities} (cf. Chap. \ref{ch:9}, \ref{sec:9.4}). Returning to the minoration of \(\lambda_n(0, V)\), he takes \(p = n-1\), assumes that \(V\) can be decomposed in \(n-1\) solids \(V_j\) which are convex and have a diameter tending to 0 with \(1/n\) and such that the ratio of their volume to the fifth power of their diameter tends to \(+\infty\) with \(n\); this gives him his conclusion.

Poincaré's next step is to investigate how the knowledge of the \(\lambda_n\) and \(v_n\) gives the solution of the cooling off problem, when the temperature \(u(x,y,z,0)\) is a known function \(f(x,y,z)\) in \(V\) at time \(t = 0\). Fourier's method consists in writing
\begin{equation}
    \label{eq:3.29}
    u(x,y,z,t) = \sum_{n=1}^\infty c_n \exp\left( -\lambda_n a^2 t \right) v_n(x,y,z)
\end{equation}
which gives for the unknown coefficients \(c_n\) the condition that \(f = \sum\limits_{n=1}^\infty c_n v_n\), hence, from the orthogonality relations, \(c_n = \iiint_V fv_n d\omega\). But Poincaré, no more than Weber, is not at that time able to prove that this Fourier expansion converges to the function \(f\) in \(V\). However, taking his cue from Tchebychef's results in approximation theory, he shows (by a clever use of Green's formula) that the integral
\[
    S_n = \iiint_V \left( u - \sum_{k=1}^n c_k \exp\left( -\lambda_k a^2 t \right) v_k \right)^2 d\omega
\]
satisfies an inequality \(S_n \leq C \cdot \exp\left(-\lambda_{n+1} a^2 t\right)\) where \(C\) is independent of \(n\) and \(t\); in other words, for \(t > 0\), he proves the convergence of the series in \eqref{eq:3.29} in what we now call the topology of Hilbert space\footnote{The method of least squares of Legendre--Gauss had led Tchebychef to define a ``best approximation'' to a function \(F\), by a linear combination \(\sum\limits_{j=1}^N a_j \psi_j\) of given functions \(\psi_j (1 \leq j \leq N)\), by the condition that
\[
    \sum_{k=1}^n \rho(x_k)\left( F(x_k) - \sum_{j=1}^N a_j \psi_j(x_k) \right)^2
\]
be minimum, for given points \(x_k (1 \leq k \leq n)\) and given ``weight'' \(\rho\). Gram, in 1883, generalized the problem by considering instead of a finite sum, an integral
\begin{equation}
    \tag{+}
    \label{eq:3.+}
    \int_a^b \rho(x) \left( F(x) - \sum_{j=1}^N a_j \psi_j(x) \right)^2 dx
\end{equation}
and he solved the problem in an original way, by applying to the \(\psi_j\) the ``orthogonalization process'' usually attributed to E. Schmidt \cite{089}. He was thus reduced to the case in which the \(\psi_j\) form an orthonormal system (for the measure \(\rho dx\)), where he showed that the \(a_j\) giving the best approximation \(a\) are the ``Fourier coefficients'' \(\int_a^b \rho(x)F(x)\psi_j(x)dx\). He went on to consider an \emph{infinite} orthonormal system \((\psi_n)\) and investigated under which conditions the minimum value \(\mu_n\) of the integral \eqref{eq:3.+} tends to 0 when \(n\) increases to \(+\infty\); he was able to see that this was linked to the ``completeness'' of the system \((\psi_n)\), i.e. the fact that no function other than the constant 0 is orthogonal to all \(\psi_n\). It is unlikely that Poincaré had any knowledge of Gram's paper.}.

The final section of Poincaré's paper (if we except a kind of postscript which we will discuss later in Chapter \ref{ch:4}, \ref{sec:4.4}) is devoted to the general study of the eigenfunctions \(v_n\) (their existence being admitted). In general, if \(v\) satisfies \eqref{eq:3.23} and \eqref{eq:3.24}, use of Green's formula shows that there is a formula similar to Green's expression of the potential (Chap. \ref{ch:2}, \ref{sec:2.3}, formula \eqref{eq:2.17})
\begin{equation}
    \label{eq:3.30}
    -4\pi v(M) = \iint_\Sigma v\left( \pder{T}{n} + hT \right)d\sigma
\end{equation}
where \(T\) (replacing the function \(1/r\)) is now \(\exp(i/\sqrt{\lambda} r)/r\). Using that formula, he is able to show, after a rather long discussion (patterned on the study of double layer potentials but more difficult), that \(v\) is continuous in \(\overline{V}\), and to obtain bounds for its derivatives in \(V\).

The second paper devoted by Poincaré to the equation of vibrating membranes (\cite{177}, vol.~IX, p.~123--196) is even more original. It is likely that in 1890, he was not aware of Schwarz's paper of 1885. The publication of Picard's note in 1893 immediately attracted his attention, and in a few months he had seen that by combining Schwarz's method and his ``a priori'' inequality of 1890, he could go beyond Picard and prove the analytic continuation of the function \(\xi \mapsto w(M, \xi)\) as a \emph{meromorphic function in the whole complex plane}, obtaining at the same time the existence of the long sought eigenvalues and eigenfunctions for the Helmholtz equation (with the same boundary condition as Schwarz).

Poincaré starts with a simplification and an improvement of his inequality for the expression \eqref{eq:3.26}; using Schwarz's inequality, he is able to replace his lower bound \eqref{eq:3.28} by \(C/(\diam(V))^2\) for a convex solid \(V\). He then only assumes that for a general solid \(V\) it is possible to decompose it in convex solids having arbitrary small diameters, and uses this idea of decomposition to prove the following crucial lemma: given \(p\) arbitrary \(C^2\) functions \(F_1, F_2, \ldots, F_p\) in \(\overline{V}\), it is possible to choose \(p\) numbers \(\alpha_1, \ldots, \alpha_p\) in such a way that, for \(v = \alpha_1 F_1 + \ldots \alpha_p F_p\), one has \(\iiint_V v d\omega = 0\) and the ratio \eqref{eq:3.26} is \emph{at least} \(L_p\), where \(L_p\) is a number which only depends on \(V\) and \(p\) (and not on the \(F_j\)) and \emph{tends to} \(+\infty\) with \(p\). This is simply done by decomposing \(V\) in the union of \(p-1\) convex subsets \(V_j\), and choosing the coefficients \(\alpha_j\) by the \(p-1\) conditions \(\iiint_{V_j} v d\omega = 0 \; (1 \leq j \leq p - 1)\).

Poincaré, as Picard, limits himself to the case in which the function \(p\) in equation \eqref{eq:3.8} is the constant 1, but considers a problem which slightly generalizes Schwarz's, namely he looks for a function \(v\) solution of
\begin{equation}
    \label{eq:3.31}
    \Delta v + \xi v + f = 0
\end{equation}
and vanishing on the boundary \(\Sigma\), with \(f\) an \emph{arbitrary} \(C^\infty\) function (if in Schwarz's equation \eqref{eq:3.8} with \(p = 1\), one writes \(w = w_0 + \xi v\), the equation for \(v\) is \eqref{eq:3.31} with \(f=w_0\)); he will make a very clever use of this arbitrariness. He starts by observing that Schwarz's method works just as well for arbitrary \(f\) as for \(f = 1\), and proves the existence of the solution of \eqref{eq:3.31} vanishing on \(\Sigma\) for \emph{small enough} \(|\xi|\); he writes it
\[
    v = [f, \xi] = v_0 + \xi v_1 + \ldots + \xi^n v_n + \ldots
\]
with \(\Delta v_0 + f = 0, \Delta v_n + v_{n-1}=0\) for \(n \geq 1\), the \(v_n\) all vanishing on \(\Sigma\).

For any given integer \(p\), he introduces \(p\) arbitrary coefficients \(\alpha_1, \ldots, \alpha_p\) and forms the function (defined at least for small \(\xi\))
\begin{equation}
    \label{eq:3.32}
    w = [\alpha_1 f + \alpha_2 v_0 + \ldots + \alpha_p v_{p-2}, \xi] = w_0 + w_1 \xi + \ldots + w_n \xi^n + \ldots
\end{equation}
Next, applying his lemma for the evaluation of the Schwarz integrals \(W_n\) \emph{corresponding to} \(w\), he is able to show that, for a suitable \emph{choice} of the \(\alpha_j\) the series \eqref{eq:3.32} \emph{converges for} \(|\xi| \leq L_p\) (uniformly in \(\overline{V}\)). But if one writes \(u_j = \left[v_{j-2}, \xi\right]\), one has
\begin{equation}
    \label{eq:3.33}
    \begin{dcases}
        \alpha_1 v + \alpha_2 u_2 + \ldots \alpha_p u_p = w \\
        v - u_2 \xi = v_0 \\
        u_2 - u_3 \xi = v_1 \\
        \dotfill \\
        u_{p-1} - u_p \xi = v_{p-2}
    \end{dcases}
\end{equation}
a linear system from which Cramer's formulas give
\begin{equation}
    \label{eq:3.34}
    v = P / D
\end{equation}
with
\[
    D =
    \left|
    \begin{matrix}
        \alpha_1 & \alpha_2 & \alpha_3 & \ldots & \alpha_{p-1} & \alpha_p \\
        1 & -\xi & 0 & \ldots & 0 & 0 \\
        0 & 1 & -\xi & \ldots & 0 & 0 \\
        \hdotsfor{6} \\
        0 & 0 & 0 & & 1 & -\xi
    \end{matrix}
    \right|
    = \alpha_p - \alpha_{p-1}\xi + \alpha_{p-2}\xi^2 - \ldots + (-1)^{p-1} \alpha_1 \xi^p
\]
and
\[
    D =
    \left|
    \begin{matrix}
        w & \alpha_2 & \alpha_3 & \ldots & \alpha_{p-1} & \alpha_p \\
        v_0 & -\xi & 0 & \ldots & 0 & 0 \\
        v_1 & 1 & -\xi & \ldots & 0 & 0 \\
        \hdotsfor{6} \\
        v_{p-2} & 0 & 0 & \ldots & 1 & -\xi
    \end{matrix}
    \right|
\]
which shows that \(P\), as \(w\), is equal to a series
\begin{equation}
    \label{eq:3.35}
    P = P_0 + P_1 \xi + \ldots + P_n \xi^n + \ldots
\end{equation}
where the \(P_n\) are \(C^\infty\) functions in \(\overline{V}\) vanishing on \(\Sigma\), the series being uniformly convergent in \(\overline{V}\) for \(|\xi| < L_p\), and all derivatives of \(P\) (with respect to \(\xi\) or to \(x,y,z\)) being obtained by derivating termwise the series. This shows that \(\xi \mapsto v(M, \xi)\) extends to a \emph{meromorphic} function in \(|\xi|<L_p\) with a finite number of poles, which are roots of \(D(\xi) = 0\) and \emph{independent} of \(M\); it is easy to show, using Green's formula, that these poles are all \emph{simple}. As this is true for any \(p, (\xi \mapsto v(M, \xi)\) extends to a meromorphic function in the \emph{whole complex plane}, with simple real and positive poles independent of \(M\); furthermore, for each one of these poles \(\lambda_n\), the function \(P(M, \lambda_n)\) satisfies \(\Delta P + \lambda_n P = 0\); in other words, one has found for each \(\lambda_n\) an eigenfunction \(u_n\) corresponding to that eigenvalue. In addition, Poincaré's \emph{a priori} inequality enables him to show that \(\lambda_n \geq c \cdot n^{2/3}\), where \(c\) is a constant,

The remainder of Poincaré's 1894 paper is devoted to two questions:

A) In the last 4 sections of the paper, he takes up again the problem of Fourier expansions (when the boundary condition is \(v=0\)). Attaching to the function \(f\) its ``Fourier coefficients'' \(c_n = \iiint_V f u_n d\omega\) (where the eigenfunctions \(u_n\) have been normalized by \(\iiint_V u^2_n d\omega = 1\)), he first deduces from the relations \(\Delta u_n + \lambda_n u_n = 0\) and Schwarz's inequality, that \(|u_n| \leq A \lambda_n\) in \(\overline{V}\) (\(A\) constant), and that the \(c_n\) are uniformly bounded. From that it follows that for \(\xi\) different from the eigenvalues the unique solution of \eqref{eq:3.31} vanishing on \(\Sigma\) is given by the absolutely and uniformly convergent series
\begin{equation}
    \label{eq:3.36}
    v = -\sum_n \frac{c_n u_n \xi^2}{\lambda^2_n \left( \xi - \lambda_n \right)} + v_0 + v_1 \xi;
\end{equation}
in addition, Poincaré shows that if the series \(\sum\limits_n c_n u_n\) is absolutely convergent, its sum is equal to \(f\); he cannot prove that for ``arbitrary'' functions \(f\) (probably at least \(C^4\)), vanishing on \(\Sigma\), the series converges, but he proves absolute convergence when in addition \(\Delta f\) and \(\Delta^2 f\) also vanish on \(\Sigma\).

B) Before returning to the question of Fourier expansions, Poincaré had tried to extend his results on the existence of eigenvalues and eigenfunctions for the boundary condition \eqref{eq:3.24} of the cooling off problem. He realizes that Schwarz's method would work, and therefore also his own existence theorem, \emph{provided} one could prove the existence of a ``Green function'' for the Laplace equation with that new boundary condition, i.e. a function \(G(M, P)\) having the same properties as the usual Green function, with the exception that, for \(M \in V, P \mapsto G(M, P)\) satisfies \eqref{eq:3.24} on the boundary.

\emph{In the special case} \(h=0\), C. Neumann, in his work on the Dirichlet problem, had shown how to obtain such a ``Green function'' (also named ``Neumann function'') when \(V\) is convex and not a double cone. He had observed that by changing the sign before the integral in the Beer--Neumann equation (Chap. \ref{ch:2}, \ref{sec:2.4}, formula \eqref{eq:2.23}), the solution of that new equation gave a density \(\rho\) such that the corresponding double layer potential, in the \emph{exterior} of \(V\) (complement of \(\overline{V}\)) is harmonic, tends to 0 at infinity and to \(-g\) on the boundary \(\Sigma\) (a solution to what is called the ``exterior Dirichlet problem''). From this result, he had shown how to obtain a solution of what is now called the \emph{Neumann problem} for the Laplace equation: find in \(V\) a harmonic function \(u\) such that \(u\) is continuous in \(\overline{V}\) and has on \(\Sigma\) a normal derivative \(\pder{n}{u}\) equal to a \emph{given} continuous function \(g\); a necessary condition for the existence of the solution (deduced from Green's formula applied to \(u\) and the constant 1) is that \(\iint_\Sigma g d\sigma = 0\). Neumann proves that this condition is sufficient (the solutions being determined up to an additive constant): he considers the \emph{simple layer} potential \(w\) defined by the density \(\frac{1}{4\pi} g\); it is continuous on \(\Sigma\) and its normal derivative jumps by \(-g\) when crossing \(\Sigma\) from the interior to the exterior. Neumann next takes the \emph{double layer} potential \(v\), solution of the \emph{exterior} Dirichlet problem which tends to \(-w\) on \(\Sigma\). Then the function \(u = v + w\) is harmonic outside \(\Sigma\), and 0 in the exterior of \(V\); as the normal derivative of \(v\) is the same on both sides of \(\Sigma\), it follows at once that \(\pder{u}{n}\) tends to \(g\) from \(\Sigma\) from the interior of \(\)V, and therefore solves the Neumann problem.

C. Neumann had not been able to solve the corresponding problem when the boundary condition is \(\pder{u}{n} + hu = g\) for a an constant \(h > 0\). Poincaré tried to solve the problem by representing \(u\) as a power series in \(h, u = u_0 + h u_1 + \ldots + h^n u_n + \ldots\), and was indeed able to obtain in that way (using Neumann's results) a series convergent for all \(h \geq 0\), uniformly in \(\overline{V}\); however, for the first derivatives, his method could only prove uniform convergence in compact subsets of \(V\), so that it was impossible to give meaning to \(\pder{u}{n}\) on the boundary \(\Sigma\), and to show that u was indeed a solution of the problem. The most interesting result in this attempt is that Poincaré, probably for the first time in history, arrives at the idea of ``weak'' solution of a boundary problem; he shows that his function \(u\) is such that, for \emph{any} function \(v\) which is \(C^2\) in \(\overline{V}\), one has
\begin{equation}
    \label{eq:3.37}
    \iiint_V u\Delta v d\omega + \iint_\Sigma gv d\sigma = \iint_\Sigma \left( \pder{v}{n} + hv \right) u d\sigma
\end{equation}
and adds that ``physically'' this is equivalent to a genuine solution.

The last of the three long papers of Poincaré on partial differential equations was written in 1895 (\cite{177}, vol.~IX, p.~202--272). Although it is the one which contains the smallest number of new results, it probably had a greater influence than the others. From his work both on the Dirichlet problem and on the equation of vibrating membranes, Poincaré had become convinced that there were also ``eigenvalues'' and ``eigenfunctions'' linked to the Dirichlet problem. For us this is completely obvious, for if we look for a solution of \(\Delta u = 0\) taking given values on the boundary \(\Sigma\) of \(V\), we extend the function \(g\) given on \(\Sigma\) to a \(C^2\) function \(h\) in \(\overline{V}\) (when this is possible); replacing \(u\) by \(v = u-h\), we have to find a solution of \(\Delta v + f = O\), with \(f = \Delta h\), which vanishes on \(\Sigma\), and this is just the special case of Schwarz's problem for the equation \eqref{eq:3.31} with \(\xi = 0\).

At that time, however, nobody had yet thought of this simple argument\footnote{It is explicitly mentioned in 1909 by E.E. Levi \cite[vol.~II, p.~302--313]{145}; the first statement and proof of the existence of a continuous function in the whole space \(\R^3\) extending a given function defined and continuous in a closed subset (\emph{i.e.} what we now call the Tietze--Urysohn theorem) is due to Lebesgue in 1907 \cite[vol.~IV, p.~99--100]{138}.}, and Poincaré's reasoning is quite different and much more circuitous. He observes that one can formulate both the interior and exterior Dirichlet problems as special cases of the problem which consists in finding a double layer potential \(W\) (for a density on \(\Sigma\)) such that, for \(s \in \Sigma\),
\begin{equation}
    \label{eq:3.38}
    W(s^-) - W(s^+) - \lambda \left( W(s^-) + W(s^+) \right) = 2\Phi(s)
\end{equation}
where \(W(s^-)\) is the limit of \(W\) at \(s\) along the interior normal, \(W(s^+)\) its limit along the exterior normal, \(\lambda\) is a complex parameter and \(\Phi\) a given function on \(\Sigma\); the values \(\lambda = 1\) and \(\lambda = -1\) correspond respectively to the interior and the exterior Dirichlet problem. To this general problem Poincaré associates a new variational problem: for any simple layer potential \(\Psi\) defined by a density on \(\Sigma\), he considers the ratio \(J/J'\), where \(J\) is the Dirichlet integral \(\iiint (\grad \Psi)^2 d\omega\) extended over \(V\), and \(J'\) the integral of the same function, extended to the exterior of \(V\). The usual non rigorous arguments lead him to conjecture: 1° the existence of an increasing sequence \(0 = \lambda_0 \leq \lambda_1 \leq \ldots \leq \lambda_n \leq \ldots\) of eigenvalues, and: 2° for each \(\lambda_i\), the existence of a simple layer potential \(\Phi_i\), such that, on \(\Sigma\),
\begin{equation}
    \label{eq:3.39}
    \pder{\Phi_i}{n}(s^-) + \lambda_i \pder{\Phi_i}{n}(s^+) = 0;
\end{equation}
in addition, for \(i \neq j, \iiint_V \grad(\Phi_i) \cdot \grad(\Phi_j) d\omega = 0\). Normalizing the \(\Phi_i\) by \(\iiint_V (\grad \Phi_i)^2 d\omega = 1\), he assumes that there is a Fourier expansion \(\Phi = \sum\limits_i c_i \Phi_i\) of the given function \(\Phi\) on \(\Sigma\), and ``solves'' the equation \eqref{eq:3.39} by
\[
    W(s^-) = \sum_i A_i \Phi_i(s), \qquad W(s^+) = -\sum_i \lambda_i A_i \Phi_i(s)
\]
with \(A_i = 2c_i / (1 + \lambda_i - \lambda (1 - \lambda_i))\).

All this is of course presented by Poincaré as purely conjectural, and as a motivation for his detailed study (by methods inspired by those of Schwarz) of the ratio \(J/J'\), which forms the central part of his 1896 paper; but the only positive result he is able to deduce from his study is that the Beer--Neumann series (Chap. \ref{ch:2}, \ref{sec:2.4}) converges, not only for convex domains \(V\), but also for domains \(V \subset \R^3\) having the following property: when \(\R^3\) is imbedded in the 3-dimensional sphere \(\Sph_3\) by adjoining a point at infinity, \(V\) can be transformed into a ball by a homeomorphism of \(\Sph_3\) onto itself, leaving fixed the point at infinity, and which is \(C^2\) in \(\Sph_3\) as well as the inverse homeomorphism\footnote{Without the slightest justification, Poincaré claims as ``clear'' the fact that this property holds for \emph{any} bounded domain \(V\) such that the boundary \(\Sigma\) is a smooth simply connected surface (\cite{177}, vol.~IX, p.~223--224). With the tools of modern Differential Topology, it is now possible to prove that theorem. But in 1895, Poincaré was just beginning to formulate the first notions of that theory, and one wonders if he realized the difficulties which lay in the way of a rigorous proof if he had tried to write it down (when smoothness conditions on \(\Sigma\) and on the homeomorphism are dropped, the result is known to be false, a counterexample being the famous ``Alexander horned sphere'').}.

Almost immediately after the publication of Poincaré's papers, several mathematicians were able to complete and extend his results. In 1898, E. Le Roy \cite{144} proved the existence of the simple layer potentials \(\Phi_i\) conjectured by Poincaré in his 1896 paper; he replaced the ratio \(J/J'\) by \((J+J')/I\), where \(I\) is the surface integral \(\iint_\Sigma \rho^2 d\sigma, \rho\) being the density on \(\Sigma\) corresponding to the simple layer potential \(\Psi\), and adapted the methods of Schwarz and Poincaré to the corresponding variational problem. In 1899, S. Zaremba \cite{232}, by a modification of the method of solution of Neumann's problem used by Poincaré, could complete the latter's solution of the ``cooling off'' problem, proving that the ``weak'' solution of Poincaré was a genuine one. In 1901, Zaremba and W. Stekloff, independently, finally showed that one could drop the global topological property of the domain \(V\) which Poincaré and Le Roy had used, and even weaken the ``smoothness'' conditions on \(\Sigma\); they made essential use of a paper of Liapounoff published 3 years earlier \cite{147}, in which he was able to prove the existence of the normal derivative on \(\Sigma\) of the solution of Dirichlet's problem under these less stringent conditions.

\chapter{The idea of infinite dimension}
\label{ch:4}
\setcounter{equation}{0}

\section{Linear algebra in the XIX\supth{} century}
\label{sec:4.1}

I think that in order to understand the trend of ideas which led to Functional Analysis, it is useful to summarize the evolution of linear algebra during the XIX\supth{} century. Until around 1830, it had consisted in the study of systems of linear equations in any number of variables, with real or complex coefficients, most of the times limited to the case in which the number of equations was equal to the number of variables; the Cramer formulas gave the unique solution when the determinant of the system was not 0, but not much effort was spent on the elucidation of the other cases; the only result which was used occasionally was the fact that a system of \(m\) homogeneous equations in \(n > m\) variables always had a non trivial solution (obvious by induction on \(m\)).

Linear changes of variables
\begin{equation}
    \label{eq:4.1}
    y_j = \sum_{k=1}^n a_{jk} x_k \quad (1 \leq j \leq m)
\end{equation}
had been familiar since the XVITI\supth{} century (mostly for \(m=n \leq 3\)). It naturally led to computations done, not on numbers, but on \emph{rectangular arrays} \((a_{jk})\) of numbers, which ``represented'' these changes of variables. Beginning with Gauss, this trend was systematized in the 1850's by Sylvester and Cayley in the theory of \emph{matrices}.

Ever since the invention of cartesian coordinates (``analytic geometry'', as it came to be called in the XVIII\supth{} century), mathematicians had known how to interpret geometrically computations on systems of 2 or 3 variables, and many had envisioned the possibility of similarly interpreting computations on systems of any number \(n\) of variables in a ``geometry in \(n\) dimensions'', which however would be devoid of ``reality''. After 1840, mainly under the influence of Hamilton and Cayley, this geometrical language was gradually adopted by more and more mathematicians, and had become commonplace at the end of the century. But in the XIX\supth{} century, after 1822 ``geometry'' essentially meant \emph{projective} geometry, and most ``geometric'' interpretations of computations were done, not in the vector space \(\R^n\) or \(\C^n\), but in the complex projective spaces \(\mathbb{P}_n(\C)\); for instance, the relations \eqref{eq:4.1} for \(m = n\) were interpreted as defining also a \emph{projective transformation} in \(\mathbb{P}_{n-1}(\C)\) sending the point of homogeneous coordinates \((x_j)\) to the point of homogeneous coordinates \((y_j)\), and the efforts of Grassmann and Peano to introduce vector spaces in an axiomatic way were persistently ignored until 1900.

Between 1850 and 1880 are proved the main theorems of linear algebra, concerning what are called the ``reductions'' of square matrices. One of these is the problem of finding, for a given square matrix \(U\), an invertible matrix \(P\) such that \(PUP^{-1}\) has a ``reduced'' unique canonical form, which here (for complex matrices \(U\)) means a diagonal array of Jordan matrices; this is the way Jordan himself treats the problem, improving on a previous result of Grassmann, who had proved the existence of a ``reduced'' triangular matrix \(PUP^{-1}\) for any \(U\) (using already the intrinsic notion of endomorphism instead of the notion of square matrix).

Unfortunately, another type of ``reduction'' interfered with the preceding one. To a quadratic form \(\sum\limits_{j \leq k} a_{jk} x_j x_k\) corresponds the \emph{symmetric} matrix \((a_{jk}) = U\), and it was well known since Cauchy that if \(U\) is \emph{real} it is possible to find an invertible real matrix \(P\) (which may even be supposed to be \emph{orthogonal}) such that \(PUP^{-1}\) would be a (real) \emph{diagonal} matrix; this is equivalent to finding an orthogonal change of variables for which the quadratic form became equal to a linear combination \(\sum\limits_{j=1}^n \lambda_j y_j^2\) of squares, the \(\lambda_j\) being the elements of the diagonal matrix \(PUP^{-1}\), or equivalently the roots (with their multiplicity) of the ``characteristic equation'' \(\det{(U - \lambda I)} = 0\). Weierstrass, who was the first to find the ``Jordan normal form'' of a square complex matrix (which Jordan only discovered independently 2 years later\footnote{Jordan was not dealing with matrices having elements in \(\R\) or \(\C\), but with matrices having elements in a \emph{finite field} (\cite{123}, p. 114--126).}), presented it as a generalization of the ``reduction'' of a quadratic form, by considering a \emph{bilinear form} \(\sum\limits_{j,k} a_{jk} x_j y_k\) (with \(U = (a_{jk})\) an \emph{arbitrary} square matrix) and applying to the \(x_j\) and \(y_j\) two ``contragredient'' changes of variables, i.e. such that the bilinear form \(\sum\limits_{j,k} x_j y_k\) remains invariant; this amounts to replacing \(U\) by a matrix \(PUP^{-1}\). When, in 1878, Frobenius gave a systematic account of these results (\cite{078}, vol.~I, p.~343--405), he deliberately abandoned the language of matrices in favor of the language of bilinear forms, defining the ``product'' (\emph{Faltung}) of two bilinear forms \(A(x,y), B(x,y)\), as \(\sum\limits_{k=1}^n \pder{A}{y_k} \pder{B}{x_k}\)!\footnote{In 1896, Pincherle reinterpreted Weierstrass's results in terms of endomorphisms (\cite{173}, vol.~I, p.~358--367).}

Finally, the concept of duality in \emph{vector spaces} was completely foreign to mathematicians until 1900. Duality was well understood in the realm of \emph{projective geometry} (it had been one of the big discoveries of the early XIX\supth{} century), as a bijection of points on planes (in projective space of 3 dimensions) and later as a bijection of points on hyperplanes in any number of dimensions. But linear forms were identified with the systems of their coefficients, ``vectors'' and ``forms'' being thus both ``\(n\)-tuples'' of numbers, which one had to distinguish, according to the way they behaved under changes of variables, by the awkward concepts of ``contragredient'' and ``cogredient'' systems. This identification of a vector space and its dual was reverberated in the identification of endomorphisms with bilinear forms, mentioned above \footnote{In modern linear algebra, the space of endomorphisms of a finite dimensional vector space \(E\) is identified with the tensor product \(E^* \otimes E\), whereas the space of bilinear forms on \(E \times E\) is identified with \(E^* \otimes E^*\).}.

To sum up, at the end of the XIX\supth{} century, the main results of linear and multilinear Algebra had been found but were expressed through insufficiently clarified notions. They could therefore be of no help to the generalizations of linear Algebra to infinite dimensional spaces which were called forth by the development of Functional Analysis; these had to go through the same painful stages, first linear equations, then determinants, later bilinear forms, matrices, and only at the very end vector spaces and linear maps; in other words, the \emph{historical evolution}, just as for finite dimensional linear algebra, was exactly in the reverse order of what we not consider to be the \emph{logical} order!

\section{Infinite determinants}
\label{sec:4.2}

The first appearance of infinite systems of linear equations in infinitely many unknowns seems to occur in Fourier's work on the theory of heat. He has to determine an infinite sequence \((a_m)_{m \geq 1}\) of coefficients such that the relation
\begin{equation}
    \label{eq:4.2}
    1 = \sum_{m=1}^\infty a_m \cos(2m - 1)y
\end{equation}
holds for all \(y\) (\cite{067}, vol.~I, p.~149). Fourier's idea is to take derivatives of all orders of both sides of \eqref{eq:4.2} and identify them for \(y = 0\), which gives him the infinite system of linear equations for the \(a_m\)
\begin{equation}
    \label{eq:4.3}
    \begin{dcases}
        1 = \sum_{m=1}^\infty a_m \\
        0 = \sum_{m=1}^\infty (2m - 1)^2 a_m \\
        0 = \sum_{m=1}^\infty (am-1)^4 a_m \\
        \dotfill
    \end{dcases}
\end{equation}
To solve it, he considers the first \(k\) equations where he replaces the \(a_m\) for \(m > k\) by 0; he then solves that system by Cramer's formulas, which give him a system of \(k\) numbers \(a_1^{(k)}, a_2^{(k)}, \ldots, a^{(k)}_k\), and lets \(k\) tend to infinity in each expression of \(a_m^{(k)}\) for fixed \(m\). Using the formulas giving Vandermonde determinants, he obtains
\[
    a_1^{(k)} = \frac{3^2 \cdot 5^2 \ldots (2k-1)^2}{8 \cdot 24 \ldots (4k^2 - 4k)}
\]
tending to \(a_1 = 4/\pi\), and
\[
    \frac{a_{m+1}^{(k)}}{a_m^{(k)}} = \frac{2m-1}{2m+1} \cdot \frac{m+k}{m-k}
\]
which gives him \(a_m = (-1)^{m-1} 4/\pi(2m-1)\); when later in his book he proves the general formula giving the Fourier coefficients, he can of course check that these values of the \(a_m\) are correct. But he never bothered to give any justification of his procedure, where all questions of convergence are completely disregarded; that procedure could of course be repeated for any infinite system
\begin{equation}
    \label{eq:4.4}
    \sum_{k=1}^\infty a_{jk} x_k = b_j \qquad (j = 1, 2, \ldots)
\end{equation}
but nobody undertook to justify it before 1885\footnote{During that period, Fourier's method was used in two little-known papers, one by Fürstenau in 1860 on the computation of roots of an algebraic equation, and another by Kötteritzsch in 1870, for a system \eqref{eq:4.4} in which the \(a_{jk}\) are 0 for \(j > k\) (see \cite{184}, p.~8--12).}. In that year, P. Appell met such a system with \(a_{jk} = a_k^j\) for a given sequence \((a_k)\), in a question relative to elliptic functions, and used the same method as Fourier; his paper attracted Poincaré's attention, and he showed that for such a ``generalized Vandermonde system'', the procedure was justified provided the infinite product \(F(z) = \prod\limits_{k=1}^\infty \left( 1 - \frac{z}{a_k} \right)\) was convergent for all complex numbers \(z\).

The next year, he returned to the subject, in relation with a paper published in 1877 by the American astronomer and mathematician G.W. Hill on the lunar theory \cite{114}. Hill proposed a new approach which rested on the integration of a second order differential equation
\begin{equation}
    \label{eq:4.5}
    w'' + \left( \sum_{n=-\infty}^{+\infty} \theta_n e^{nit} \right) w = 0
\end{equation}
where the \(\theta_n\) are constants, and one looks for a solution of period \(2\pi\). Hill writes such a solution as a trigonometric series
\begin{equation}
    \label{eq:4.6}
    w = \sum_{n=-\infty}^{+\infty} b_n e^{i(n+c)t}
\end{equation}
and substituting in \eqref{eq:4.5}, obtains for the coefficients \(b_n\) the infinite system of equations
\begin{equation}
    \label{eq:4.7}
    \sum_{k=-\infty}^{+\infty} \theta_{n-k} b_k - (n+c)^2 b_n = 0, \qquad -\infty < n < +\infty.
\end{equation}
He probably was unaware of Fourier's procedure, but used a similar one, keeping this time the equations \eqref{eq:4.7} for \(-p \leq n \leq p\), replacing in these equations the \(b_m\) by 0 for \(m < -p\) or \(m > p\), and letting \(p\) tend to infinity in the solutions of the system thus obtained.

Poincaré considers a general system \eqref{eq:4.4}, where he supposes that \(a_{jj}=1\) for all \(j\) (one can always reduce \eqref{eq:4.4} to such a system by dividing the \(j\)-th equation by \(a_{jj}\), when \(a_{jj} \neq 0\)). His idea is to compare the determinant \(D_n = \det(a_{jk})_{1 \leq j, k \leq n}\) to the product \(P_n = \prod\limits_{j=1}^n \left( \sum_{k=1}^n |a_{jk}| \right)\). It is clear from the definition of a determinant that \(|D_n| \leq P_n\), and from the assumption on the diagonal terms, one has also \(|D_m - D_n| \leq P_m - P_n\); this inequality immediately gives Poincaré's \emph{sufficient} condition for the existence of \(D = \lim\limits_{n \to \infty} D_n\) , namely that the double sum \(\sum\limits_{j \neq k} |a_{jk}|\) be \emph{finite}. Furthermore, Poincaré shows that, when the \(k\)-th column of \(D\) is replaced by a sequence \((b_j)\) which is \emph{bounded}, there is still convergence for the new ``infinite determinant'', and that there is a unique bounded solution \(x_k\) of \eqref{eq:4.4} given by the usual Cramer formulas (with ``infinite determinants'' of course).\footnote{One must beware of the fact that the Fourier method (when no condition is imposed on the \(a_{jk}\)) may very well give convergent ``infinite determinants'', but the values given by the Cramer formulas may be such that the left hand sides of \eqref{eq:4.4} are \emph{divergent} series. An example is given by taking \(a_{jk} = 0\) if \(j > k, a_{jk} = 1\) if \(j \leq k, b_j = (-1)^j\); one finds as a ``solution'' \(x_k = 2(-1)^k\).} Finally he extends his results to doubly infinite systems
\begin{equation}
    \label{eq:4.8}
    \sum_{k=-\infty}^{+\infty} a_{jk} x_k = b_j \qquad (-\infty < j < +\infty)
\end{equation}
with the same restriction \(a_{jj} = 1\) for all \(j\); in particular he shows that Hill's method is justified for the system \eqref{eq:4.7} (\cite{177}, vol.~V, p.~95--107).

Ten years later, H. von Koch \cite{220} refined and generalized Poincaré's results. Instead of making assumptions on the diagonal terms, he writes the coefficients \(\delta_{jk} + c_{jk}\) instead of \(a_{jk}\) (with the Kronecker delta), and uses the expression of a determinant \(\Delta_n = \det(\delta_{jk} + c_{jk})_{1 \leq j, k \leq n}\) as a sum of principal minors
\begin{equation}
    \label{eq:4.9}
    \Delta_n = 1 + \sum_{s=1}^n c_{ss} + \frac{1}{2!} \sum_{s_1, s_2}
    \left|
    \begin{matrix}
        c_{s_1 s_1} & c_{s_1 s_2} \\
        c_{s_2 s_1} & c_{s_2 s_2}
    \end{matrix}
    \right|
    + \frac{1}{3!} \sum_{s_1, s_2, s_3}
    \left|
    \begin{matrix}
        c_{s_1 s_1} & c_{s_1 s_2} & c_{s_1 s_3} \\
        c_{s_2 s_1} & c_{s_2 s_2} & c_{s_2 s_3} \\
        c_{s_3 s_1} & c_{s_3 s_2} & c_{s_3 s_3}
    \end{matrix}
    \right|
    + \ldots
\end{equation}
(an expression which will be the starting point of Fredholm's theorems on integral equations 4 years later (Chap. \ref{ch:5}, \ref{sec:5.1})). He is thus able to replace Poincaré's criterion for convergence by a weaker one: it is enough that the sums \(\sum\limits_j |c_{jj}|\) and \(\sum |c_{i_1 i_2} c_{i_2 i_3} \ldots c_{i_p i_1}|\) (extended to all sequences \((i_1, i_2, \ldots, i_p)\) of distinct indices) be finite. Another convergence criterion is that the sum \(\sum\limits_j |c_{jj}|\) and \(\sum_{j,k} |c_{jk}|^2\) be finite.

\section{Groping towards function spaces}
\label{sec:4.3}

It should not be believed that set-theoretic concepts in mathematics were unknown before Boole (1847) or Cantor; they can be traced at least as far back as Aristotle. The use of the word ``class'' (or, in German, ``Gebiet'', ``Inbegriff'', ``Mannigfaltigkeit'', ``System'') to designate a set of objects having a common property, becomes frequent among mathematicians since the beginning of the XIX\supth{} century. But it is only after Boole, in the second half of the century, that using letters to denote more or less arbitrary sets, and \emph{computing} with these letters, will become a widespread practice.

In particular ``classes'' of functions where very often considered in Analysis, even if their description lacks precision most of the time. Even more widespread was the use, since the XVIII\supth{} century, of sequences of functions, or of functions depending on one or several real parameters (for instance in the Calculus of variations). It was of course dimly realized that such families of functions were ``much smaller'' than the ``class'' of \emph{all} functions under consideration; the first attempt to give a clearer expression to that feeling is probably due to Riemann. In his famous inaugural lecture on the foundations of geometry, after having tried to give an idea of what he means by a ``finite dimensional multiplicity (i.e. manifold)'' where the position of a point is determined by a finite set of numbers, he adds that there are ``multiplicities'' (\emph{Mannigfaltigkeiten}) for which such a determination is not possible, but needs ``an infinite sequence or a continuous multiplicity of numbers'', and gives as an example ``all the possible determinations of a function in a given domain'' (\cite{182}, p.~276).

The extension of the concepts of limit and of continuity to mathematical objects other than numbers or points, such as curves, surfaces or functions, is also very old. However, the applications of that idea dealt with sequences of such objects, or families depending on a finite number of real parameters; again, Riemann seems to have been the first to conceive that a whole ``class'' of functions might be given some kind of ``geometrical'' structure (what we now would call a topology), for when he speaks of the functions for which the Dirichlet integral (Chap. \ref{ch:2}, formula \eqref{eq:2.21}) has a meaning, he says that ``this set of functions constitutes a connected domain, closed in itself'' (\cite{182}, p.~30), and although it is not quite clear what he means by that, we may see in that statement a first glimpse of the notion of compactness, which will emerge in the last part of the century (see below).

The rigorous study of limits of sequences of functions, which began around 1820, brought to light a phenomenon which had no counterpart for sequences of numbers or of points in \(\R^n\): there are \emph{several distinct ways} for a sequence \((f_n)\) of functions to tend to a limit \(f\). The first problem occurred with the distinction between simple and uniform convergence, which was only quite cleared up around 1850. This was followed in the last third of the XIX\supth{} century by a deeper study of these notions, chiefly due to the Italian school (Dini, Ascoli, Arzelà); the most important step taken by that school was the introduction by Ascoli in 1883 of the notion of \emph{equicontinuity}. He discovered that the unpleasant phenomenon of a sequence of continuous functions (in a bounded closed interval \(I\)), converging simply to a discontinuous function, would disappear if one assumed on the sequence the following additional property: for each \(\varepsilon > 0\), there exists a \(\delta > 0\) such that, if \(|x'-x''| \leq \delta\), then \(|f_n(x') - f_n(x'')| \leq \varepsilon\) for \emph{all indices} \(n\) (in other words, the continuity is ``uniform'', not only with respect to \(x\), but also with respect to \(n\)) \cite{008}.

One of the fundamental properties of equicontinuous sequences is that, when in addition the \(f_n\) are uniformly bounded, it is possible to find a \emph{subsequence} \((f_n)\) which converges uniformly, a generalization of the ``Bolzano--Weierstrass'' theorem for sequences of numbers, which was well-known after 1880. This ``compactness'' property (which holds for functions defined in a closed bounded set of \(\R^n\)) was thrust in the limelight by Hilbert, who apparently rediscovered it independently in a special case (he does not quote the Italians) and used it as an essential tool in his famous 1900 paper where he invented the ``direct method'' in the Calculus of variations (\cite{111}, vol.~III, p.~10--14) and thus was able to justify Riemann's use of the ``Dirichlet principle'' (chap. \ref{ch:2}, \ref{sec:2.3}) (\emph{loc.cit.}, p. 15--37).

It is also from the Calculus of variations that another notion of ``neighborhood'' for a function emerged during the last years of the XIX\supth{} century. Already at the end of the XVIII\supth{} century mathematicians investigated the problem of deciding if a solution \(y\) of the Euler equation for an integral \(\int_a^b F(x,y,y')dx\) actually gave a ``relative extremum'' for that integral. Legendre tried to give a solution to that problem by replacing \(y\) in the integral by \(y + \varepsilon u\), where \(u = \delta y\) is an arbitrary ``variation'' of class \(C^1\); he thus obtains a function \(\Phi(\varepsilon)\) of the real parameter \(\varepsilon\) and if \(\Phi''(0) > 0\) (resp. \(\Phi''(0) < 0\)) that function reaches a relative minimum (resp. maximum) for \(\varepsilon = 0\). This yields the condition \(\pder{^2 F}{y'^2} > 0\) (resp. \(< 0\)); but it was soon realized that this condition was not sufficient to guarantee that the integral would actually be smaller (resp. larger) than \emph{all} numbers obtained by replacing \(y\) by \(y + \delta y\) for a ``small'' variation \(\delta y\). Clearly this hinges on the question of what exactly is meant by the word ``small''. Ever since Lagrange, it had been taken for granted that the derivative \((\delta y)' = \delta y'\) is ``small'' whenever \(\delta y\) itself is ``small''; but Weierstrass and his school realized that this was an additional assumption, and this led them to distinguish between ``strong extremum'' and ``weak extremum'': the second corresponds to a notion of ``neighborhood'' of a \(C^1\) function \(y\), where \(z\) is ``close'' to \(y\) when the maximum of \(|z-y|\) is small, whereas for the first \(z\) is only considered as ``close'' to \(y\) if \emph{both} the maximum of \(|z-y|\) and the maximum of \(|z' - y'|\) are small.

Finally, we have noticed earlier that Gram and Poincaré were naturally confronted with the notion of ``convergence in the mean square'' in their study of ``Fourier expansions'' (chap. \ref{ch:3}, \ref{sec:3.2}). We may therefore say that in the last years of the XIX\supth{} century, the idea of ``function spaces'' with various ``topologies'' was so to speak ``in the air'', and ready to blossom forth as soon as it could be expressed in sufficiently general and simple terms.\footnote{It is, however, typical of the unpredictability of mathematical developments that nobody seems to have been able to foresee, even conjecturally, the direction which was taken by Functional Analysis in the fateful years 1900--1910, This is clear in the communication made by Hadamard in the first International Congress of mathematicians in 1897 (\cite{094}, vol.~I, p.~311--312); he was keenly interested in these ``set-theoretical'' ideas, and had great expectations of what was to come; but he could think of no serious applications beyond the rehabilitation of the ``Dirichlet principle'' and some vague ideas on what we now call ``precompactness''.}

The concept of \emph{mapping} of a set of functions into \(\R\), or into another set of functions, is also much older than the general definition of a mapping of an arbitrary set into an arbitrary set, which does not seem to have been formulated be fore Dedekind's famous ``\emph{Was sind und was sollen die Zahlen}'', written in 1872 (although only published in 1888) (\cite{048}, vol.~III, p.~335--391). Ever since the beginning of the Calculus of variations, mathematicians were familiar with the idea of attaching for instance to each \(C^1\) function \(y\) in an interval \([a, b]\) a number \(\int_a^b F(x,y,y')dx\) depending on \(y\); such mappings would receive the name of ``functional'' at the end of the XIX\supth{} century. Similarly, as soon as the concept of function emerged at the end of the XVII\supth{} century together with its use in Calculus, the concept of \emph{operator}, yielding a new function when applied to a given function, was in evidence with the examples of the derivatives \(f \mapsto D^\nu f\) or the translation operator \(f \mapsto \gamma(a)f\) (function \(x \mapsto f(x - a)\)); and from Leibniz to Pincherle (end of the XIX\supth{} century) many analysts were led to ponder on the \emph{algebraic} properties of these operators, and their similarity with results of ordinary algebra (which was originally conceived as applying to numbers only). For instance, the similarity of Leibniz's formula for the iterated differential \(d^n(uv)\) of a product, with the binomial theorem, probably gave him the idea of attempting to introduce differentials \(d^\alpha\) with negative or irrational exponents, a problem to which many mathematicians (such as Liouville, Riemann, Pincherle) later returned, and which has only finally been put to rest with the modern theory of distributions. Other examples are the expression of Taylor's formula given by Lagrange as a relation \(\gamma(-a) = e^{aD}\) between operators, or the factoring of a differential polynomial \(D^n + a_1 D^{n-1} + \ldots + a_n\) on the model of the factoring of an ordinary polynomial \(z^n + a_1 z^{n-1} + \ldots + a_n\).

Such ideas, abundantly developed in the period 1790--1830, had much to do with the new conception of Algebra as dealing with symbols rather than with numbers, and later with the axiomatic and formalist conception of the whole of mathematics (see \cite{054}, chap.~XIII, \S III); but they had no perceptible influence on Analysis, probably because they did not pay much attention to questions of continuity. It is only in the last years of the XIX\supth{} century that such questions appear, in a very episodic way, in papers by Pincherle, Bourlet and Volterra.

The first two of these authors only consider one ``space'' \(E\), the set of all holomorphic functions in a domain \(\Delta\) of the complex plane, and they are exclusively concerned with \emph{linear} operators in that space. In 1886, Pincherle studies operators which, to a holomorphic function\footnote{After Grassmann (1862), Pincherle seems to have been one of the first mathematicians to write a function with a single letter \(\varphi\), when all his contemporaries wrote \(\varphi(x)\). In his later papers, he repeatedly insists on the fact that a function should be considered as a ``point'' in some set.} \(\varphi\), associate the function \(x \mapsto \int_\Gamma A(x, y) \varphi(y) dy\), where \(\Gamma\) is a curve in \(\Delta\) and \(A\) is holomorphic, and he writes that function \(\mathbb{a}\varphi\), but he limits himself to special cases, of the type of the Laplace transform (\cite{173}, vol.~I, p.~92--141). He several times returned later to such questions, but failed to obtain any substantial results\footnote{He should however be credited with what is probably the first conception of a closed hyperplane in \(E\) as the kernel of a continuous linear form, and of closed subspaces of finite codimension as intersections of hyperplanes (\cite{173}, vol.~I, p.~39). In 1897--98, he also has the idea of generalizing Lagrange's ``adjoint'' of an operator (chap. \ref{ch:1}, \ref{sec:1.1}, formula \eqref{eq:1.5}) by considering two vector subspaces \(S, S'\) of \(E\), and a nondegenerate bilinear form \(\varphi, \psi\) on \(S \times S'\); to a linear mapping of \(S\) into \(S'\), he then associates the ``adjoint'' \(A\), a linear mapping of \(S'\) into \(S\) such that \((A \cdot \varphi, \psi) = (\varphi, A \cdot \psi)\), and he observes the relation between the kernel of \(A\) and the image of \(A\) (\cite{173}, vol.~II, p.~77--84).}. In 1897, Bourlet \cite{029}, limiting himself to the case in which \(A\) is a disk \(|z| < r\), explicitly determines the linear operators in \(E\) which are ``continuous'' (by which he means continuity for what we now call the topology of compact convergence), showing that they are integral operators of the form considered by Pincherle.

We must finally mention the first attempts at ``Functional Analysis'' of the young Volterra in 1887 (\cite{219}, vol.~I, p.~294--314), to which, under the influence of Hadamard, has been attributed an exaggerated historical importance. Volterra had in mind a generalization of analytic functions, which may be considered as a prefiguration of Hodge's theory\footnote{See A. Weil, \emph{Oeuvres Scientifiques}, vol.~II, Commentaires sur [1952 e], p.~532 of the correct edition (or vol.~III, p.~450 of the first printing), Springer, Berlin-Heidelberg-New York, 1979.}; for this he needs what he calls ``functions of lines''. Although, from our point of view, his definitions are not very precise\footnote{This can be said of practically all mathematicians before 1906.}, he apparently considers the set \(E\) of \(C^1\) mappings of an interval \(I \subset \R\) into \(\R^3\) (the ``lines''), and the mappings \(y: E \ \R\), continuous for the topology of uniform convergence. For these ``functions of lines'' he immediately wants to generalize the classical notion of derivative; in a manner reminiscent of the Calculus of variations, he considers a ``variation'' \(\delta y = y(\varphi + \theta) - y(\varphi)\), where the increment \(\theta\) is supposed to vanish outside of an interval \([a, b]\), and then the quotient \(\delta y/\sigma\), where \(\sigma = \int_a^b |\theta(t)|dt\); this should tend to a limit when \(b-a\) and the maximum of \(|\theta|\) tend to 0. With our experience of 50 years of Functional Analysis, we cannot help feeling that, without even the barest notions of general topology, these \emph{ad hoc} definitions were decidedly premature. Nevertheless, they caught the fancy of Hadamard, who tried to apply similar ideas to Green's functions and encouraged his students to work in that direction (see \cite[vol.~I, p.~401--404 and 435--453]{094} and \cite{146}). But these ideas have not, up to now, produced anything comparable to the applications of spectral theory and distribution theory, which we will describe in chap. \ref{ch:7} and \ref{ch:9}; it might be worthwhile to reexamine them in the light of recent progress in the theory of infinite dimensional manifolds, which could be their natural setting.

\section{The passage ``from finiteness to infinity''}
\label{sec:4.4}

The urge to deal with ``infinity'' has been present from the very beginnings of Greek mathematics, in spite of all philosophical preconceptions and objections, and has taken various forms. The simplest and most ``natural'' passage ``from finiteness to infinity'' is the ``indefinite repetition'' of the arithmetical operation of addition, on smaller and smaller summands, giving birth to the concept of convergent series, of which one can already find examples in Archimedes. Replace addition by multiplication, and you have the infinite product, born with Calculus in the XVII\supth{} century; and still more sophisticated algebraic manipulations would lead to continued fractions and to the infinite determinants which we have discussed in \ref{sec:4.2}.

Another line of thought goes back at least to Eudoxus's ``method of exhaustion'', and was to lead in the first place to the concept of integral. But in the hands of the mathematicians of the XVII\supth{} and XVIII\supth{} century, this idea of decomposing an object into ``infinitesimal'' parts in which the phenomenon they studied became much easier to describe ``in a first approximation'', was developed into a more and more sophisticated method to discover the differential or partial differential equations which governed the phenomenon ``in the large''. It is in that way that the equation of vibrating strings (chap. \ref{ch:1}, \ref{sec:1.2}, formula \eqref{eq:1.7}) was established, either by considering, as D. Bernoulli, a massive string as a limit (for \(n\) tending to infinity) of a system of \(n\) massive points distributed on a massless string, or by analyzing, as d'Alembert, the forces which are exerted on an ``infinitesimal'' portion of the string by its neighbors.

It is this second method that Fourier applied to obtain the heat equation; he takes for granted that in a system of small ``molecules'', a given molecule \(M\) receives in an ``infinitesimal'' time \(dt\) a quantity of heat from another molecule \(M'\) equal to the difference of temperatures of \(M\) and \(M'\), multiplied by \(dt\) and by a coefficient depending only on the distance \(MM'\); the molecule \(M\), if situated at the surface separating the system of molecules from the external world, also radiates a quantity of heat equal to the difference of its temperature and of the external temperature, multiplied by \(dt\) and another coefficient depending on \(M\). He then derives the equation of the ``cooling off'' process (chap. \ref{ch:3}, \ref{sec:3.2}, equation \eqref{eq:3.21}) by decomposing the solid body \(V\) in ``infinitesimal'' cubes and evaluating the amount of heat received by one of them from its 6 neighbors in time \(dt\), which he takes as proportional (with a constant coefficient) to the variation \(du\) of the temperature of that cube the boundary condition (chap. \ref{ch:3}, \ref{sec:3.2}, equation \eqref{eq:3.22}) is similarly obtained by evaluating the amount of heat lost (by radiation) by an infinitesimal cube at the surface of \(V\).

At the end of his 1890 paper on the cooling off problem (chap. \ref{ch:3}, \ref{sec:3.2}), Poincaré suggests another method reminiscent of D. Bernoulli's procedure. He first considers a large number \(N\) of molecules \(M_i\); following Fourier's physical considerations, and denoting by \(v_i(t)\) the temperature of \(M_i\) at time \(t\), these functions satisfy the system of linear differential equations
\begin{equation}
    \label{eq:4.10}
    \frac{dv_i}{dt} + \sum_{k \neq i} C_{ik} (v_i - v_k) + C_i v_i = 0 \qquad (1 \leq i \leq N),
\end{equation}
\(C_{ik} (v_i - v_k)\) being the quantity of heat received from \(M_k\), and \(C_i v_i\) the quantity of heat radiated by \(M_i\) outside the system. But instead of letting the number of molecules increase to infinity, Poincaré \emph{first} integrates the system \eqref{eq:4.10} by the classical Euler--Lagrange method: he writes \(v_i(t) = u_i e^{-\alpha t}\), and, using the fact that the matrix \((C_{ik})\) is \emph{symmetric}, he recognizes in the equation he obtains for \(\alpha\) the equation giving the eigenvalues of the symmetric matrix corresponding to the non degenerate positive \emph{quadratic form}
\begin{equation}
    \label{eq:4.11}
    \Phi(u_1, u_2, \ldots, u_N) = \sum_{i \neq k} C_{ik} (u_i - u_k)^2 + \sum_i C_i u_i^2.
\end{equation}
Let \(\xi_1 \leq \xi_2 \leq \ldots \leq \xi_N\) be these eigenvalues; the classical theory of quadratic forms shows that one may write
\begin{equation}
    \label{eq:4.12}
    \Phi = \xi_1 \varphi_1^2 + \ldots + \xi_N \varphi_N^2
\end{equation}
where the \(\varphi_i\) are linear forms in the variables \(u_1, \ldots, u_N\) such that \(\varphi_1^2 + \ldots + \varphi_N^2 = u_1^2 + \ldots + u_N^2\); if for two such forms \(f = \alpha_1 u_1 + \ldots + \alpha_N u_N, g = \beta_1 u_1 + \ldots + \beta_N u_N\) one writes \((f|g) = \sum\limits_k \alpha_k \beta_k\), the \(N\) forms \(\varphi_i\) are mutually orthogonal for that scalar product. It is then clear that \(\xi_1\) is the smallest value of the function of \(u_1, \ldots, u_N\)
\begin{equation}
    \label{eq:4.13}
    \frac{\Phi(u_1, \ldots, u_N)}{u_1^2 + u_2^2 + \ldots + u_N^2} = \frac{\Phi}{\varphi_1^2 + \ldots + \varphi_N^2}
\end{equation}
where the \(u_i\) are arbitrary; similarly \(\xi_2\) is the minimum of \eqref{eq:4.13} for \(\varphi_1 = 0, \xi_3\) the minimum for \(\varphi_1 = \varphi_2 = 0\) as relations between the \(u_i\); and so on. This is of course the analogous procedure in \(N\) dimensions to the classical determination of the ``axes'' of an ellipsoid in 3-dimensional space. Poincaré's idea is that the expression \eqref{eq:4.13} corresponds exactly to the quotient
\[
    \frac{\Iiint_V \left(\divergence{v}\right)d\omega + h\Iint_\Sigma v^2 d\sigma}{\Iiint_V v^2 d\omega}
\]
in his (or rather Weber's) procedure for the definition of the eigenvalues in the cooling off problem; these eigenvalues (the poles \(\lambda_m\) of his function \([f, \xi]\) (chap. \ref{ch:3}, \ref{sec:3.2})) correspond to the \(\xi_j\) in \eqref{eq:4.12} and the eigenfunctions \(U_m(M) = P(M,\lambda_m)\) to the \(\varphi_j\), the orthogonality of the \(\varphi_j\) corresponding to the relations
\[
    \Iiint_V U_p U_q d\omega = 0 \quad \text{for } p \neq q
\]
between the \(U_m\). Finally, he realizes that the same ideas apply as well to other problems and gives as an example the theory of elasticity, and he suggests that a rigorous proof of the existence of the \(\lambda_m\) and the \(U_m\), which he had not been able to give, might be obtained by simply letting \(N\) tend to \(+\infty\) in the formula \eqref{eq:4.12}. He never came back to the question; but we cannot fail to see that this is exactly the program which Hilbert in 1904 followed to its successful conclusion for integral equations with symmetric kernel (chap. \ref{ch:5}, \ref{sec:5.2}).

A similar ``passage from finiteness to infinity'' emerged in the first \emph{general} theory of integral equations, beginning with the papers of Le Roux in 1894 and Volterra in 1896. In addition to the \emph{particular} integral equations which had been met by Liouville in the Sturm--Liouville problem (chap. \ref{ch:1}, \ref{sec:1.3}, equation \eqref{eq:1.34} and chap. \ref{ch:2}, \ref{sec:2.1}, equation \ref{eq:2.6}) and by Beer and Neumann in the Dirichlet problem (chap. \ref{ch:2}, \ref{sec:2.4}, equation \eqref{eq:2.23}) (not to speak of what we have called ``crypto-integral'' equations, where the equation is not written down explicitly but the method exactly amounts to solving it), other \emph{particular} equations involving integrals had come up in connection with problems not directly related to differential or partial differential equations. The first one (chronologically) was the ``inversion'' problem for the ``transform'' introduced by Fourier in 1822 (and to which we shall return in chap. \ref{ch:7}, \ref{sec:7.6}); it associates to a function \(f\) in \([0, +\infty[\) the function
\begin{equation}
    \label{eq:4.14}
    \varphi(t) = \int_0^{+\infty} f(x)\cos{tx}dx
\end{equation}
and the problem consisted in finding \(f\) when the transform \(\varphi\) is a given function. It was solved by Fourier's inversion formula (\cite{067}, vol.~I, p.~392)
\begin{equation}
    \label{eq:4.15}
    f(x) = \frac{2}{\pi} \int_0^{+\infty} \varphi(x)\cos{tx}dt
\end{equation}
where, as usual with Fourier, both formulas are obtained by a purely formal calculation. A little later, one of the first published papers of Abel (\cite{001}, vol.~I, p.~11--27 and 97--101) was devoted to a problem of mechanics, which amounted to finding a function \(\varphi\) such that
\begin{equation}
    \label{eq:4.16}
    \int_0^x \frac{\varphi(y)dy}{\sqrt{x-y}} = \psi(x)
\end{equation}
is a given function; he obtains the solution by the formula
\begin{equation}
    \label{eq:4.17}
    \varphi(x) = \frac{1}{\pi} \int_0^x \frac{\psi'(y)dy}{\sqrt{x - y}}
\end{equation}
and extends his result to the case in which \(\sqrt{x - y}\) is replaced by \((x-y)^\alpha\) for \(0 < \alpha < 1\). In a letter to Holmboe, he even hinted at more general results, but nothing was found on the subject in his papers. After Abel, a few papers, giving partial generalizations of his results, were published until 1890\footnote{See the long historical introduction given by Volterra in his 1897 paper on integral equations (\cite{219}, vol.~II, p.~279--287)}; but it was only in 1894 that Le Roux attacked the \emph{general} problem of ``inversion of a definite integral'' (as it was called), \emph{i.e.} finding a \(C^1\) function \(\varphi\) in an interval \([a,b]\) satisfying an equation
\begin{equation}
    \label{eq:4.18}
    \int_a^b \varphi(x) H(x, y)dx = f(y)
\end{equation}
where \(f\) and \(H\) are \(C^1\) (in \([a,b]\) and \([a,b] \times [a,b]\) respectively) and \(f(a) = 0\) \footnote{As these conditions are \emph{not} satisfied for Abel's equation, Le Roux's results (which for him are auxiliary properties which he needs in a study of partial differential equations) do not directly generalize those of Abel.}. In contrast with his predecessors, Le Roux is not trying to find a ``closed formula'' similar to \eqref{eq:4.15} and \eqref{eq:4.17} for the unknown function. He assumes that \(h(y) = H(y,y)\) does not vanish in \([a,b]\), takes the derivative of both sides of \eqref{eq:4.18}, obtaining
\begin{equation}
    \label{eq:4.19}
    h(y) \varphi(y) + \int_a^y \pder{H}{y}(x, y) \varphi(y) dx = f'(y)
\end{equation}
and then applies the method of successive approximations which Picard had popularized a few years earlier:
\[
    u_0(y) = \frac{f'(y)}{h(y)}, \qquad u_n(y) = \frac{f'(y)}{h(y)} - \frac{1}{h(y)} \int_a^y \pder{H}{y}(x, y) u_{n-1}(x)dx \quad \text{for } n \geq 1,
\]
proving easily the convergence of the sequence \((u_n)\) to a solution of \eqref{eq:4.18} (\cite{143}, p.~244--246).

In 1896, Volterra (who apparently was unaware of Le Roux's paper) tackles exactly the same problem by the same method, in a series of 4 notes (\cite{219}, vol.~II, p.~216--262). He goes a little beyond Le Roux, by giving an explicit expression of the solution
\begin{equation}
    \label{eq:4.20}
    \varphi(y) = \frac{f'(y)}{h(y)} - \frac{1}{h(y)} \int_a^y \left( \sum_{i=0}^\infty S_i(x, y) \right) f'(x) dx
\end{equation}
where the \(S_i\) are defined by induction:
\begin{equation}
    \label{eq:4.21}
    S_0(x, y) = \frac{1}{h(x)}\pder{H}{y}(x, y), \qquad S_i(x, y) = \int_x^y S_0(\xi, y)S_{i-1}(x, \xi)d\xi \quad \text{for } i \geq 1.
\end{equation}
In the later notes, he discussed the cases in which \(h(y)\) may vanish at a finite number of points, and the case in which \(H(x,y) = G(x,y)/(x-y)^\alpha\) with \(0 < \alpha < 1\) and \(G\) is continuous (the generalization of Abel's equation). But the most influential part of his notes was the following remark he made immediately after obtaining formula \eqref{eq:4.20}: ``If one considers the system
\begin{equation}
    \label{eq:4.22}
    \begin{cases}
        b_1 = a_{11} x_1 \\
        b_2 = a_{12} x_1 + a_{22} x_2 \\
        \dotfill \\
        b_n = a_{1n} x_1 + \ldots + a_{nn} x_n
    \end{cases}
\end{equation}
the concept of integral easily leads to look at the question of functional Analysis represented by equation \eqref{eq:4.18} as a limiting case of the solution of a system similar to \eqref{eq:4.22}, in which the \(a_{ij}\) and \(a_{ii}\) are the analogous of \(H(x,y)\) and \(H(y,y)\).'' Although he limited himself to that (somewhat vague) statement, it seems obvious that what he had in mind was replacing in \eqref{eq:4.18} the variable \(y\) by its values
\[
    y_k = a + \frac{k}{n}(b-a) \quad \text{for } 1 \leq k \leq n
\]
and replacing the integral by the corresponding ``Riemann sum'' for the subdivision of \([a,b]\) by the points \(y_k\) obtaining the system of type \eqref{eq:4.22}
\begin{equation}
    f(y_j) = \frac{b-a}{n} \sum_{k=1}^n \varphi(y_k) H(y_k, y_j) \quad \text{for } 1 \leq j \leq n.
\end{equation}

Finally, although he does not mention the product of matrices, Volterra develops in these notes the formalism which to two ``kernels'' \(H_1(x,y), H_2(x,y)\) associates the kernel
\[
    H(x, y) = \int_x^y H_1(x, \xi) H_2(\xi, y) d\xi
\]
(which much later he will write \(H = H_1 \ast H_2\)). If, for simplicity, we adopt this notation, he shows that, for an \emph{arbitrary} continuous function \(S_0(x,y)\) if we define for \(i \geq 1\) the ``kernel'' \(S_i\) by \(S_i = S_{i-1} \ast S_0\), one has also \(S_i = S_{i-j} \ast S_{j-1}\) for \(1 \leq j \leq i\), and a majoration
\[
    |S_i(x, y)| \leq \frac{M^{i+1}}{i!} |x-y|^i.
\]
This implies uniform convergence for the series
\begin{equation}
    \label{eq:4.23}
    F_0 = \sum_{i=0}^\infty S_i
\end{equation}
and the relation
\begin{equation}
    \label{eq:4.24}
    F_0 - S_0 = S_0 \ast F_0.
\end{equation}

He observes that one may ``invert'' that relation: if the \(F_i\) are defined for \(i \geq 1\) by \(F_i = F_0 \ast F_{i-1}\), one has
\begin{equation}
    \label{eq:4.25}
    S_0 = \sum_{i=0}^\infty (-1)^i F_i.
\end{equation}

And finally, at the end of his notes, he arrives at the general concept of what Hilbert will call an ``integral equation of the second kind''
\begin{equation}
    \label{eq:4.26}
    \varphi(y) - \int_a^y S_0(x, y) \varphi(x) dx = f(y)
\end{equation}
for which the solution is given by
\begin{equation}
    \label{eq:4.27}
    \varphi(y) = f(y) + \int_a^y F_0(x, y) f(x) dx
\end{equation}
as it follows immediately from \eqref{eq:4.24}, the ``kernel'' and the ``resolvent kernel'' playing completely symmetric parts in these formulas.

\chapter{The crucial years and the definition of Hilbert space}
\label{ch:5}
\setcounter{equation}{0}

Between 1900 and 1910, there was a sudden crystallization of all the ideas and methods which had been slowly accumulating during the XIXth century and which we have described in the previous chapters. This was essentially due to the publication of \emph{four fundamental papers}:

Fredholm's 1900 paper on integral equations;

Lebesgue's thesis of 1902 on integration;

Hilbert's paper of 1906 on spectral theory;

Fréchet's thesis of 1906 on metric spaces.

\section{Fredholm's discovery}
\label{sec:5.1}

The name ``integral equation'' (\emph{Integralgleichung}) was used for the first time by P. du Bois-Reymond in 1888, in a paper on the Dirichlet problem \cite{061}; he has in mind equations of the Beer--Neumann type (chap. \ref{ch:2}, \ref{sec:2.4}) and considers that a general theory of such equations presents ``insuperable difficulties''; he is convinced that much progress would come out of such a theory but acknowledges that ``almost nothing is known on this question''. The later work of Poincaré, which we have discussed above (chap. \ref{ch:3}, \ref{sec:3.2}), and of his immediate followers, did nothing to dispel that impression; their results seemed linked to delicate estimates from potential theory. It therefore came as a complete surprise when, in a short Note published in 1900, Fredholm showed that the general theory of all integral equations (or ``crypto-integral'' equations) considered before him was in fact extremely simple (much simpler than anything known at the time in the theory of partial differential equations).

Ivar Fredholm (1866--1927) was a student of Mittag-Leffler in Stockholm in 1888--1890; he only published a few papers during his lifetime, mostly concerned with partial differential equations (we shall return to his thesis of 1898 in chapter \ref{ch:9}, \ref{sec:9.5}). After a visit to Paris, where he had been in contact with all the French analysts and had become familiar with the recent papers of Poincaré, he communicated in August 1899 his first results on integral equations to his former teacher; they were published in 1900 \cite[p.~61--68]{074} and completed 2 years later in a paper published in \emph{Acta Mathematica} (\cite[p.~81--106]{074} and \cite{075}).

Fredholm's 1900 note is entitled ``On a new method for the solution of Dirichlet's problem'', but it is characteristic that from the start, he brushes aside all the particular features of the Beer--Neumann equation, and (as Le Roux and Volterra had done with Abel's equation (chap. \ref{ch:4}, \ref{sec:4.4})) begins with a \emph{general} ``integral equation of the second kind'' (that name will only be given by Hilbert)
\begin{equation}
    \label{eq:5.1}
    \varphi(s) = f(s) + \lambda \int_a^b K(s, t) f(t) dt
\end{equation}
where \(K\) is supposed to be bounded and piecewise continuous in \((a, b] \times [a, b]\), and \(\varphi\) continuous in \([a, b]\), \(\lambda\) being a complex parameter. He briefly mentions the analogy with systems of linear equations and starts right away with the formulas describing his ``determinants'' (see below). But in a lecture given in 1909 \cite[p.~123--131]{074}, he acknowledges: 1° the inspiration derived from Volterra's idea of a ``passage to the limit'' from a system of linear equations to an integral equation; 2° the help he found in von Koch's work on infinite determinants (chap. \ref{ch:4}, \ref{sec:4.2}). From these indications, sparse as they are, it seems one can reconstruct his procedure, with great probability, as consisting in \emph{putting together three simple ideas}:

I) Replacing the integral in \eqref{eq:5.1} by Riemann sums, one obtains, with the notations of chap. \ref{ch:4}, \ref{sec:4.4}, the system of \(n\) linear equations for the \(f(y_j)\)
\begin{equation}
    \label{eq:5.2}
    f(y_j) + \frac{\lambda (b-a)}{n} \sum_{k=1}^n K(y_k, y_j) f(y_k) = \varphi(y_j) \quad (1 \leq j \leq n).
\end{equation}

II) Writing the determinant of that system according to von Koch's formula (chap. \ref{ch:4}, \ref{sec:4.2}, formula \eqref{eq:4.9})
\[
    1 + \frac{\lambda (b-a)}{n} \sum_{k=1}^n K(y_k, y_j) + \frac{\lambda^2 (b-a)^2}{2!n^2} \sum_{k_1, k_2}
    \left|
    \begin{matrix}
        K(y_{k_1}, y_{k_1}) & K(y_{k_1}, y_{k_2}) \\
        K(y_{k_2}, y_{k_1}) & K(y_{k_2}, y_{k_2})
    \end{matrix}
    \right|
    + \ldots
\]
and then letting \(n\) tend to \(+\infty\), which gives the formula for what Fredholm calls the ``determinant'' of the integral equation \eqref{eq:5.1}
\begin{equation}
    \label{eq:5.3}
    \begin{gathered}
        \Delta(\lambda) = 1 + \lambda \int_a^b K(s, s)ds + \frac{\lambda^2}{2!} \int_a^b \int_a^b K
        \left(
        \begin{matrix}
            s_1 & s_2 \\
            s_1 & s_2
        \end{matrix}
        \right) ds_1 ds_2
        + \ldots \\
        + \frac{\lambda^m}{m!} \int_a^b \ldots \int_a^b K
        \left(
        \begin{matrix}
            s_1 & s_2 & \ldots & s_m \\
            s_1 & s_2 & \ldots & s_m
        \end{matrix}
        \right) ds_1 ds_2 \ldots ds_m
        + \ldots
    \end{gathered}
\end{equation}
where he has written
\begin{equation}
    \label{eq:5.4}
    K
    \left(
    \begin{matrix}
        x_1 & x_2 & \ldots & x_m \\
        y_1 & y_2 & \ldots & y_m
    \end{matrix}
    \right)
    =
    \left|
    \begin{matrix}
        K(x_1, y_1) & K(x_1, y_2) & \ldots & K(x_1, y_m) \\
        K(x_2, y_1) & K(x_2, y_2) & \ldots & K(x_2, y_m) \\
        \hdotsfor{4} \\
        K(x_m, y_1) & K(x_m, y_2) & \ldots & K(x_m, y_m)
    \end{matrix}
    \right|
    .
\end{equation}

III) Proving the uniform convergence of the series \eqref{eq:5.3} in any compact set of the complex plane, for which it is enough to majorize the determinants \eqref{eq:5.4} in a suitable way; in his 1899 letter, Fredholm had given the majoration \(n^{n/2} M^n\), where \(M\) is the upper bound of \(|K|\); he had apparently arrived independently to this result, but was made aware that it was a special case of an inequality published by Hadamard in 1893 (\cite{094}, vol.~I, p.~239--243) for an arbitrary square matrix \(A = (a_{ij})\) of order \(n\):
\begin{equation}
    \label{eq:5.5}
    |\det(A)|^2 \leq \prod_{i=1}^n \left( \sum_{j=1}^n |a_{ij}|^2 \right).
\end{equation}

The next ``natural'' steps are of course to apply Cramer's formulas to the system \eqref{eq:5.2} and let again \(n\) tend to infinity in the numerators; the result is described by Fredholm in the following elegant way: a development of the determinant \eqref{eq:5.4} according to the first row yields the formula
\begin{equation}
    \label{eq:5.6}
    \begin{gathered}
        K
        \left(
        \begin{matrix}
            s & x_1 & \ldots & x_m \\
            t & x_1 & \ldots & x_m
        \end{matrix}
        \right)
        = K(s, t) K
        \left(
        \begin{matrix}
            x_1 & \ldots & x_m \\
            x_1 & \ldots & x_m
        \end{matrix}
        \right)
        - \\
        - K(s, x_1) K
        \left(
        \begin{matrix}
            x_1 & x_2 & \ldots & x_m \\
            t   & x_2 & \ldots & x_m
        \end{matrix}
        \right)
        + K(s, x_2) K
        \left(
        \begin{matrix}
            x_1 & x_2 & x_3 & \ldots & x_m \\
            t   & x_1 & x_3 & \ldots & x_m
        \end{matrix}
        \right)
        - \\
        - \ldots + (-1)^m K(s, x_m) K
        \left(
        \begin{matrix}
            x_1 & x_2 & \ldots & x_m \\
            t   & x_1 & \ldots & x_{m-1}
        \end{matrix}
        \right)
        .
    \end{gathered}
\end{equation}
On the other hand, Fredholm defines the ``minor''
\begin{equation}
    \label{eq:5.7}
    \begin{gathered}
        \Delta(s, t; \lambda) = K(s, t) + \lambda \int_a^b K
        \left(
        \begin{matrix}
            s & x_1 \\
            t & x_1
        \end{matrix}
        \right)
        dx_1 + \ldots + \\
        + \frac{\lambda^m}{m!} \int_a^b \ldots \int_a^b K
        \left(
        \begin{matrix}
            s & x_1 & \ldots & x_m \\
            t & x_1 & \ldots & x_m
        \end{matrix}
        \right)
        dx_1 dx_2 \ldots dx_n + \ldots
    \end{gathered}
\end{equation}
and replaces each integrand by its expression \eqref{eq:5.6}, which gives him the simple relation
\begin{equation}
    \label{eq:5.8}
    \Delta(s, t; \lambda) = K(s, t)\Delta(\lambda) - \lambda \int_a^b K(s, \xi) \Delta(\xi, t; \lambda) d\xi.
\end{equation}
He then introduces the function
\begin{equation}
    \label{eq:5.9}
    \Phi(s) = \varphi(s)\Delta(\lambda) - \lambda \int_a^b \Delta(s, \xi; \lambda) \varphi(\xi) d\xi
\end{equation}
and derives from \eqref{eq:5.8} the equation
\begin{equation}
    \label{eq:5.10}
    \Phi(s) + \lambda \int_a^b K(s, t)\Phi(t)dt = \varphi(s)\Delta(\lambda).
\end{equation}
The conclusion is then immediate: if \(\Delta(\lambda) \neq 0\), the function \(f(s) = \Phi(s)/\Delta(\lambda)\) is a solution of \eqref{eq:5.1}. Furthermore, he shows that one has
\begin{equation}
    \label{eq:5.11}
    \frac{d\Delta(\lambda)}{d\lambda} = \int_a^b \Delta(s, s; \lambda)ds
\end{equation}
and from this he deduces that if \(\lambda_0\) is a zero of order \(\nu\)  of the entire function \(\Delta(\lambda), \Phi(s),\) for a suitable choice of \(\varphi\), cannot be divisible by a power of \(\lambda - \lambda_0\) greater than \((\lambda - \lambda_0)^{\nu-1}\); if \(\Phi(s) = (\lambda - \lambda_0)^k \Phi_1(s)\), one then deduces, from \eqref{eq:5.10}, that
\begin{equation}
    \label{eq:5.12}
    \Phi_1(s) + \lambda_0 \int_a^b K(s, t) \Phi_1(t)dt = 0;
\end{equation}
in other words, if there is no nontrivial solution of the homogeneous equation \eqref{eq:5.12}, necessarily \(\Delta(\lambda_0) \neq 0\), hence the solution of \eqref{eq:5.1} for \(\lambda = \lambda_0\) exists and is unique. However, at that time, he does not yet prove that the existence of a non trivial solution of \eqref{eq:5.12} implies that \(\Delta(\lambda_0) = 0\). But the end of the Note is startling: he considers the Beer--Neumann equation for a bounded plane domain with a \(C^3\) boundary; the kernel of that integral equation is then bounded and continuous, and for \(\lambda_0 = 1\) it is very easy to deduce from the properties of double layer potentials that the homogeneous equation \eqref{eq:5.12} has no nontrivial solution. Therefore the existence and uniqueness of the solution of Dirichlet's problem is proved, doing away, with a single stroke of the pen, so to speak, with all the complications of the Neumann--Poincaré solution!

In his 1903 paper, Fredholm completed his results on some important points. He first defines more general ``minors''
\begin{equation}
    \label{eq:5.13}
    \begin{gathered}
        \Delta
        \left(
        \begin{matrix}
            s_1 & s_2 & \ldots & s_m \\
            t_1 & t_2 & \ldots & t_m
        \end{matrix}
        ; \lambda
        \right)
        = K
        \left(
        \begin{matrix}
            s_1 & s_2 & \ldots & s_m \\
            t_1 & t_2 & \ldots & t_m
        \end{matrix}
        \right)
        + \\
        + \sum_{n=1}^\infty \frac{\lambda^n}{n!} \int_a^b \ldots \int_a^b K
        \left(
        \begin{matrix}
            s_1 & s_2 & \ldots s_m & x_1 & \ldots & x_n \\
            t_1 & t_2 & \ldots t_m & x_1 & \ldots & x_n
        \end{matrix}
        \right)
        dx_1 \ldots dx_n.
    \end{gathered}
\end{equation}
Developing this time the determinants \emph{both} according to the first row and the first column, he obtains the identities
\begin{equation}
    \label{eq:5.14}
    \begin{gathered}
        \Delta
        \left(
        \begin{matrix}
            s_1 & s_2 & \ldots & s_m \\
            t_1 & t_2 & \ldots & t_m
        \end{matrix}
        ; \lambda
        \right)
        + \lambda \int_a^b K(s_1, \xi) \Delta
        \left(
        \begin{matrix}
            \xi & s_2 & \ldots & s_m \\
            t_1 & t_2 & \ldots & t_m
        \end{matrix}
        ; \lambda
        \right)
        d\xi = \\
        = K(s_1, t_1) \Delta
        \left(
        \begin{matrix}
            s_2 & \ldots & s_m \\
            t_2 & \ldots & t_m
        \end{matrix}
        ; \lambda
        \right)
        - K(s_1, t_2) \Delta
        \left(
        \begin{matrix}
            s_2 & s_3 & \ldots & s_m \\
            t_1 & t_3 & \ldots & t_m
        \end{matrix}
        ; \lambda
        \right)
        + \ldots
    \end{gathered}
\end{equation}
and
\begin{equation}
    \label{eq:5.15}
    \begin{gathered}
        \Delta
        \left(
        \begin{matrix}
            s_1 & s_2 & \ldots & s_m \\
            t_1 & t_2 & \ldots & t_m
        \end{matrix}
        ; \lambda
        \right)
        + \lambda \int_a^b K(\xi, t_1) \Delta
        \left(
        \begin{matrix}
            s_1 & s_2 & \ldots & s_m \\
            \xi & t_2 & \ldots & t_m
        \end{matrix}
        ; \lambda
        \right)
        d\xi = \\
        = K(s_1, t_1) \Delta
        \left(
        \begin{matrix}
            s_2 & \ldots & s_m \\
            t_2 & \ldots & t_m
        \end{matrix}
        ; \lambda
        \right)
        - K(s_2, t_1) \Delta
        \left(
        \begin{matrix}
            s_1 & s_3 & \ldots & s_m \\
            t_2 & t_3 & \ldots & t_m
        \end{matrix}
        ; \lambda
        \right)
        + \ldots
    \end{gathered}
\end{equation}
which in particular, for \(m = 1\), reduce to \eqref{eq:5.8} and to
\begin{equation}
    \label{eq:5.16}
    \Delta(s, t; \lambda) = K(s, t)\Delta(\lambda) - \lambda \int_a^b K(\xi, t) \Delta(s, \xi; \lambda) d\xi.
\end{equation}

The use he makes of these formulas is a little more sophisticated than in his first Note. He introduces the operator corresponding to the kernel \(K, f \mapsto S_K f\) such that \(S_K f(s) = f(s) + \int_a^b K(s, t)f(t)dt\), and, for two kernels \(K, K'\), writes the composite \(S_K S_{K'}\) as \(S_{K''}\) with
\begin{equation}
    \label{eq:5.17}
    K''(x,t) = K(x, t) + K'(x, t) + \int_a^b K(s, \xi)K'(\xi, t)d\xi.
\end{equation}
Suppose now that \(\Delta(\lambda) \neq 0\), and write
\begin{equation}
    \label{eq:5.18}
    R(s, t; \lambda) = -\Delta(s, t; \lambda)/\Delta(\lambda)
\end{equation}
(the \emph{resolvent kernel} in the later terminology of Hilbert). It then follows from \eqref{eq:5.17}, \eqref{eq:5.8} and \eqref{eq:5.16} that we have
\begin{equation}
    \label{eq:5.19}
    S_{\lambda K} S_R = S_R S_{\lambda K} = \operatorname{Id}
\end{equation}
and Fredholm has thus shown that the necessary and sufficient condition for the existence and uniqueness of a solution of \eqref{eq:5.1} is \(\Delta(\lambda) \neq 0\), the kernel \(\lambda K\) and the resolvent kernel \(R\) playing completely symmetric parts as in the formulas of Volterra (chap. \ref{ch:4}, \ref{sec:4.4}, formulas \eqref{eq:4.26} and \eqref{eq:4.27}).

Next he examines what happens when \(\Delta(\lambda) = 0\). First he generalizes \eqref{eq:5.11} to
\begin{equation}
    \label{eq:5.20}
    \frac{d^m \Delta(\lambda)}{d\lambda^m} = \int_a^b \ldots \int_a^b \Delta
    \left(
    \begin{matrix}
        s_1 & \ldots & s_m \\
        s_1 & \ldots & s_m
    \end{matrix}
    ; \lambda
    \right)
    ds_1 \ldots ds_m
\end{equation}
and from this he deduces that if \(\Delta(\lambda) = 0\), there is always an integer \(m\) such that \(\Delta
\left(
\begin{matrix}
    s_1 & \ldots & s_m \\
    s_1 & \ldots & s_m
\end{matrix}
; \lambda
\right)\)
is not identically 0. If \(m\) is the smallest integer having that property (which is exactly the order of \(\lambda\) as a zero of \(\Delta\)) he exhibits, using \eqref{eq:5.14}, \(m\) solutions of the homogeneous equation
\begin{equation}
    \label{eq:5.21}
    \Phi_1(s) = \frac{\Delta
    \left(
    \begin{matrix}
        s & s_1 & \ldots & s_m \\
        t & t_1 & \ldots & t_m
    \end{matrix}
    \right)}{\Delta
    \left(
    \begin{matrix}
        s_1 & \ldots & s_m \\
        t_1 & \ldots & t_m
    \end{matrix}
    \right)}, \quad
    \Phi_2(s) = \frac{\Delta
    \left(
    \begin{matrix}
        s_1 & s   & s_3 & \ldots & s_m \\
        t_1 & t_2 & t_3 & \ldots & t_m
    \end{matrix}
    \right)}{\Delta
    \left(
    \begin{matrix}
        s_1 & \ldots & s_m \\
        t_1 & \ldots & t_m
    \end{matrix}
    \right)}, \ldots
\end{equation}
for which he shows that they are linearly independent and that every other solution of the homogeneous equation is a linear combination of the \(\Phi_j (1 \leq j \leq m)\). He concludes the theory (which one often calls the ``Fredholm alternative'') by giving necessary and sufficient conditions on \(\varphi\) for the existence of a solution of \eqref{eq:5.1} when \(\lambda\) is a zero of order \(m\) of \(\Delta\). He observes that the ``transposed equation'' obtained from \eqref{eq:5.1} by replacing \(K(s,t)\) by \(K(t,s)\) has the same ``determinant'', and therefore the corresponding homogeneous equation has exactly \(m\) linearly independent solutions \(\psi_1, \ldots, \psi_m\); the condition \(\varphi\) must satisfy are then
\begin{equation}
    \label{eq:5.22}
    \int_a^b \varphi(x)\psi_j(x)dx = 0 \quad \text{for } 1 \leq j \leq m.
\end{equation}

Finally, Fredholm shows that for any two kernels \(K, K'\), if \(\Delta_K\) and \(\Delta_{K'}\) are the corresponding ``determinants'', then for the ``composed'' kernel \(K''\) defined by \eqref{eq:5.17} one has
\begin{equation}
    \label{eq:5.23}
    \Delta_{K''} = \Delta_K \Delta_{K'}
\end{equation}
which justifies the name ``determinant''. He also points out that his results can be generalized when the kernel \(K\) is not bounded any more, but such that \((x-y)^\alpha K(x, y)\) remains bounded, with \(0 < \alpha < 1\); and he mentions that the extension of his theorems to any number of variables is immediate.

This beautiful paper may be considered as the source from which all further developments of spectral theory are derived. It made a deep and lasting impression on the mathematical world, and almost overnight the theory of integral equations became a favorite topic among analysts (\cite{023}, \cite{175}, \cite{107}).

\section{The contributions of Hilbert}
\label{sec:5.2}

One of the most active proponents of the new theory was David Hilbert. As soon as he heard of Fredholm's results, he started doing himself research work on these questions, made them one of the main subjects discussed in his Seminar at Göttingen\footnote{It is reported (by Hellinger) that Hilbert inaugurated a session of his Seminar by announcing the development of a method which would lead to the proof of the Riemann hypothesis: the problem is to prove that a particular entire function has all its zeroes on the real line, and Hilbert hoped that this function would be expressed as the ``determinant'' of an integral equation with symmetric kernel. However, nobody has yet been able to find such an equation.} and supervised many dissertations on the various aspects of the theory. Between 1904 and 1906, he published six papers on integral equations in the \emph{Göttingen Nachrichten}, later brought together in a single volume entitled ``Grundzüge einer allgemeinen Theorie der Integralgleichungen'' \cite{112}.

In his first paper \cite[p.~1--38]{112}, Hilbert starts by doing explicitly what had only been hinted at by Volterra and Fredholm, the ``passage to the limit'' in the system \eqref{eq:5.2}, restricting himself (as he will do in almost all his results) to the case in which the kernel \(K\) is \emph{symmetric}, i.e. a real continuous function such that \(K(t,s) = K(s,t)\). He soon realized that in that particular case he might obtain much more precise results than Fredholm. In the first place, the symmetric matrix \((K(y_k,y_j))\) is then the matrix of the quadratic form \(\sum\limits_{j,k} K(y_k, y_j) \xi_k \xi_j\), and Hilbert undertook to apply also his ``passage to the limit'' to that form. He thus obtained the results which Poincaré had foreseen in the particular case he had considered (chap. \ref{ch:4}, \ref{sec:4.4}): the roots of the Fredholm determinant are then real; if they are written as a sequence \((\lambda_n)\), each being counted with its multiplicity, then, for each \(n\) there is an eigenfunction \(\varphi_n\) such that \(\int_a^b \varphi_m(t) \varphi_n(t)dt = 0\) for \(m \neq n\). Finally, if one normalizes the \(\varphi_n\) by the condition \(\int_a^b \varphi_n(t)^2 dt = 1\), and if for each a continuous function \(x\) in \([a,b]\), one defines the ``Fourier coefficients'' \((x|\varphi_n) = \int_a^b x(t)\varphi_n(t)dt\), Hilbert proves that
\begin{equation}
    \label{eq:5.24}
    \int_a^b \int_a^b K(s,t)x(s)y(t)dsdt = \sum_n \frac{1}{\lambda_n} (x|\varphi_n)(y|\varphi_n)
\end{equation}
for any two continuous functions \(x, y\), a relation which he rightly considers as the natural generalization of the classical reduction of a quadratic form to its ``axes''. What is particularly interesting in the way Hilbert considers this formula is that he shows that the righthand side of \eqref{eq:5.24} is \emph{uniformly convergent} when the functions \(x\) and \(y\) are allowed to vary arbitrarily, subject only to the conditions \(\int_a^b x(t)^2dt \leq 1\) and \(\int_a^b y(t)^2dt \leq 1\), the first prefiguration of what will become ``the unit ball in Hilbert space'' a few years later. Of course Hilbert also justifies for his integral equation the variational definition of the eigenvalues \(\lambda_n\), first proposed by Weber (chap. \ref{ch:3}, \ref{sec:3.1}). He shows that the set of the \(\lambda_n\) is infinite, except when \(K(x,y)\) is a linear combination of a finite number of functions of type \(u(x)v(y)\). He also proves that the resolvent kernel \(R(s,t; \mu)\) (in the sense of Fredholm) has the eigenvalues \(\lambda_n - \mu\), the corresponding eigenfunctions being \(\varphi_n/(\lambda_n - \mu)\) (\(\mu\) distinct from the \(\lambda_n\)) and writes the identity
\begin{equation}
    \label{eq:5.25}
    R(s, t; \mu) - R(s, t; \nu) = (\mu-\nu) \int_a^b R(s, \xi; \mu) R(\xi, t; \nu)d\xi
\end{equation}
for \(\mu\) and \(\nu\) distinct from the \(\lambda_n\). Finally, he shows that if a function \(f\) can be written in the form
\begin{equation}
    \label{eq:5.26}
    f(s) = \int_a^b K(s, t)g(t)dt
\end{equation}
for a continuous function \(g\), then the corresponding ``Fourier expansion''
\begin{equation}
    \label{eq:5.27}
    f(s) = \sum_n (f|\varphi_n)\varphi_n(s)
\end{equation}
is absolutely and uniformly convergent, and one has the ``Parseval identity''
\begin{equation}
    \label{eq:5.28}
    \int_a^b f(s)^2ds = \sum_n (f|\varphi_n)^2.
\end{equation}
However, he could only give that proof under the restrictive assumption that any continuous function could be approximated (in the sense of mean square value, or, as we would now say, for the topology of Hilbert space!) by functions of the form \eqref{eq:5.26}.

The proof that this last condition is superfluous was given in 1905 in the dissertation of Erhard Schmidt, one of the best students of Hilbert \cite{191}; it contained otherwise no startling new results, but it deserves some comments, since it is the first attempt to do away with the Fredholm ``determinants'', and substitute to them a more conceptual approach\footnote{Some of the results of E. Schmidt were also obtained independently by W. Stekloff \cite{204}.}.

E. Schmidt begins by proving the Bessel identity
\begin{equation}
    \label{eq:5.29}
    \int_a^b \left( f(s) - \sum_{n=1}^N (f|\varphi_n)\varphi_n(s) \right)^2 ds = \int_a^b f(s)^2 ds - \sum_{n=1}^N (f|\varphi_n)^2
\end{equation}
for an \emph{arbitrary} orthonormal system \((\varphi_n)\), from which he deduces that for any continuous function \(f, g\), the series \(\sum\limits_n (f|\varphi_n)(g|\varphi_n)\) is absolutely convergent, and the convergence is uniform when \(f\) is allowed to vary subject to the condition \(\int_a^b f(s)^2 ds \leq A\) for a fixed constant \(A\).

Next he assumes the existence of the eigenvalues \(\lambda_n\) and of the corresponding normalized eigenfunctions \(\varphi_n\), and using the Bessel inequality, he proves that
\begin{equation}
    \label{eq:5.30}
    \sum_n \frac{1}{\lambda_n^2} \leq \int_a^b \int_a^b K(s,t)^2 dsdt
\end{equation}
from which it follows that each \(\lambda_n\) has finite multiplicity and that \(|\lambda_n|\) tends to \(+\infty\) with \(n\) if there is an infinity of eigenvalues.

To prepare for the proof of the existence of the eigenvalues, he introduces, as Fredholm and Volterra had done, the iterated kernels
\begin{equation}
    \label{eq:5.31}
    K_m(s,t) = \int_a^b K_{m-1}(s, \xi)K(\xi, t)d\xi \quad \text{for } m > 1, \text{with } K_1 = K
\end{equation}
and shows that, if \(\varphi\) is an eigenfunction for \(K_m\), it is also an eigenfunction for \(K\) if \(m\) is odd, and is sum of two eigenfunctions for \(K\) if \(m\) is even. This allows him to apply Schwarz's method to prove the existence of at least an eigenvalue when \(K\) is not identically 0, as we have shown in chapter \ref{ch:3}, \ref{sec:3.1}, because what he gets in this way is an eigenvalue of \(K_2\).

Finally, for functions \(f\) given by \eqref{eq:5.26}, he obtains the convergence of the Fourier expansion \eqref{eq:5.27} by applying his initial lemma to the functions \(t \mapsto K(s,t)\) and \(g\); and from that he derives Hilbert's formula \eqref{eq:5.24} by multiplying the formula \(x(s) = \sum_n (x|\varphi_n)\varphi_n(s)\) by \(K(s,t)y(t)\) and integrating.

Hilbert's interest in integral equations with symmetric kernels of course stemmed from the possibility of applying them to questions of Analysis such as the Dirichlet problem; it is to such applications that he devoted the second and third of his papers on integral equations, We shall bypass them for the time being, as well as most results in his two last papers on the subject (see chapters \ref{ch:7} and \ref{ch:9}), to concentrate on his fourth paper, published in 1906, a masterpiece and one of the best papers he ever wrote. By the depth and novelty of its ideas, it is a turning point in the history of Functional Analysis, and indeed deserves to be considered as the very \emph{first} paper published in that discipline.

Hilbert's new departure in that paper is clear from the beginning: he deliberately \emph{abandons} the point of view of integral equations, to return to the older conception of the infinite systems of linear equations (chap. \ref{ch:4}, \ref{sec:4.2}), but with a new twist, This is because he realizes that the theory of integral equations can be \emph{subsumed} as a special case of that older theory: indeed, let \((\omega_n)\) be a complete orthonormal system of continuous functions in \([a,b]\), and suppose the continuous function \(f\) is a solution of \eqref{eq:5.1} for \(\lambda = 1\); then, if we consider the ``Fourier coefficients''
\begin{equation}
    \label{eq:5.32}
    \begin{gathered}
        k_{pq} = \int_a^b \int_a^b K(s,t)\omega_p(s)\omega_q(t)dsdt, \\
        b_p = \int_a^b \varphi(s)\omega_p(s)ds, \qquad x_p = \int_a^b f(s)\omega_p(s)ds
    \end{gathered}
\end{equation}
the \(x_p (p=1,2, \ldots)\) satisfy the infinite system of linear equations
\begin{equation}
    \label{eq:5.33}
    x_p + \sum_{q=1}^\infty k_{pq}x_q = b_p \qquad (p=1,2,\ldots).
\end{equation}
The new twist is that, due to the Bessel identity, one has
\begin{equation}
    \label{eq:5.34}
    \sum_{p,q} k_{pq}^2 < +\infty, \quad
    \sum_p b_p^2 < +\infty, \quad
    \sum_p x_p^2 < +\infty.
\end{equation}
Conversely, suppose we have a solution \((x_p)\) of \eqref{eq:5.33} (with conditions \eqref{eq:5.34}), and observe that if \(k_q(s) = \int_a^b K(s,t)\omega_q(t)dt\), the functions \(k_q\) are continuous, and
\[
    \sum_p k_p(s)^2 \leq \int_a^b K(s,t)^2 dt;
\]
the series \(u(s) = \sum\limits_p x_p k_p(s)\) is then absolutely and uniformly convergent; hence \(u\) is continuous and one has \((u|\omega_p)=b_p - x_p\); therefore, if \(f = \varphi - u, (f|\omega_p) = x_p\) and from the completeness of the system \((\omega_p)\) it follows that \(f\) is a solution of \eqref{eq:5.1} for \(\lambda = 1\).

Hilbert then embarks into completely uncharted territory:

1° He exclusively considers sequences \(x = (x_p)\) (for \(p=1,2,\ldots\)) of real numbers such that \(\sum\limits_p x_p^2 < +\infty\).

2° On the contrary, with regard to the double sequence \((k_{pq})\) of real numbers, he \emph{abandons} at first \emph{any} restrictive condition such as the first condition \eqref{eq:5.34}, and only retains the \emph{symmetry} conditions \(k_{qp} = k_{pq}\).

3° The center of interest is not any more the solution of the system \eqref{eq:5.33}, but the ``symmetric bilinear form''
\begin{equation}
    \label{eq:5.35}
    K(x,y) = \sum_{p,q} k_{pq} x_p y_q
\end{equation}
which he wants to ``reduce'' by a formula which would generalize \eqref{eq:5.24}.

Of course, even under the restrictions \(\sum\limits_p x_p^2 < +\infty, \sum\limits_p y_p^2 < +\infty\), the right hand side of \eqref{eq:5.35} is usually meaningless\footnote{To this rather awkward formulation, Hellinger and Toeplitz \cite{106} substituted the consideration of ``infinite matrices'' \((k_{pq})_{1 \leq p, q < +\infty}\) and of their ``calculus'' inspired from Frobenius, but without associating an endomorphism to a matrix.}; proceeding as Fourier, Poincaré and von Koch (chap. \ref{ch:4}, \ref{sec:4.2}), Hilbert considers, for each integer \(n\), the symmetric bilinear form in \(2n\) variables (``sections'' (Abschnitte) of \(K\))
\begin{equation}
    \label{eq:5.36}
    K_n(x,y) = \sum_{p=1}^n \sum_{q=1}^n k_{pq} x_p y_q,
\end{equation}
but instead of investigating the determinants of these forms, he ``reduces'' each one to its ``axes'' and is confronted with the problem of ``passing to the limit'' for these ``reduced'' forms when \(n\) tends to \(+\infty\). We postpone the detailed examination of the original method by which he was able to solve that problem, to chapter \ref{ch:7}, which is devoted to the history of modern spectral theory, of which this paper of Hilbert is the starting point; we shall only discuss here the various new notions he is led to introduce in that paper.

A) Hilbert is not yet using the geometrical language which will become prevalent among his immediate successors (cf. \ref{sec:5.3}), but it is obvious that everything he does in inspired by the analogy with \(n\)-dimensional FEuclidean space. In particular one of his main tools is the generalization of \emph{orthogonal transformations}: by that he means that, to every sequence \((x_p)\) with \(\sum\limits_p x_p^2 < +\infty\), he associates the sequence \((x'_p)\), where
\begin{equation}
    \label{eq:5.37}
    x'_p = \sum_q a_{pq} x_q \qquad (p=1,2,\ldots)
\end{equation}
and where he imposes on the double sequence \((a_{pq})\) the conditions
\begin{equation}
    \label{eq:5.38}
    \begin{gathered}
        \sum_q a_{pq}^2 = 1, \qquad \sum_p a_{pq} a_{pr} = 0 \quad \text{for } q \neq r \\
        \sum_p a_{pq}^2 = 1, \qquad \sum_q a_{pq} a_{np} = 0 \quad \text{for } n \neq p
    \end{gathered}
\end{equation}
from which he immediately deduces that conversely \((x_p)\) is deduced from \((x'_p)\) by applying the ``inverse'' orthogonal transformation defined by \((a'_{pq})\) with \(a'_{pq} = a_{qp}\).

B) Hilbert restricts himself to forms \eqref{eq:5.35} which he calls \emph{bounded}: they are the (not necessarily symmetric) forms such that there exists an \(M > 0\) for which one has \(|K_n(x,y)| \leq M\) for \(\sum\limits_p x_p^2 \leq 1, \sum_p y_p^2 \leq 1\) and \emph{for all} \(n\); he also introduces \emph{bounded linear forms} \(L(x) = \sum\limits_p a_p x_p\) with \(\sum\limits_p a_p^2 < +\infty\), so that for any \(x\) (resp. \(y\)), and any bounded bilinear form \(K\), the linear forms \(K(x,\cdot): x \mapsto K(x,y)\) and \(K(\cdot,y): x \mapsto K(x,y)\) are bounded. One of the things he wants to do (inspired of course by the ``reduction'' of bilinear forms in a finite number of variables) is to operate an orthogonal transformation on \(x\) and \(y\), substituting the expressions \eqref{eq:5.37} for the \(x_p\) and doing the same for the \(y_p\). Unfortunately, he follows Frobenius in his conception of the ``''\emph{Faltung}'' of bilinear forms (instead of the natural idea of ``composing'' transformations). So, for two bounded bilinear forms \(A, B\), he has to show that the forms \(A_n(x,\cdot) B_n(\cdot,y)\) (\emph{Faltung} of \(A_n(x,y)\) and \(B_n(x,y)\), these forms being defined as in \eqref{eq:5.36}) are the forms \(C_n(x,y)\) corresponding to a bounded form \(C(x,y)\) which he calls again the ``\emph{Faltung}'' of \(A(x,y)\) and \(B(x,y)\) and writes \(A(x,\cdot) B(\cdot,y)\). He can then express the action of an orthogonal transformation on a bounded bilinear form \(K(x,y)\) as a ``\emph{Faltung}''
\begin{equation}
    \label{eq:5.39}
    K'(x', y') = K(\cdot, \cdot) O(\cdot, x')O(\cdot, y')
\end{equation}
where \(O(x,y) = \sum\limits_{p,q} a_{pq} x_p y_q\) is the bounded bilinear form which he associates to the orthogonal transformation \eqref{eq:5.37}.

C) For the development of Functional Analysis, the most important concepts introduced by Hilbert were what he calls ``continuity'' and ``complete continuity'', which correspond to what will later be called the ``strong'' and ``weak'' topologies on Hilbert space. If \(F(x)\) is a complex-valued function defined for all sequences \(x = (x_p)\) such that \(\sum\limits_p x_p^2 < +\infty\), Hilbert says that \(F\) is \emph{continuous} if \(F\left(x^{(n)}\right)\) tends to \(F(x)\) when \(\sum\limits_p \left(x_p - x_p^{(n)}\right)^2\) tends to 0, and that \(F\) is \emph{completely continuous} if \(F\left(x^{(n)}\right)\) tends to \(F(x)\) when \(\sum\limits_p x_p^2 \leq 1, \sum\limits_p \left(x_p^{(n)}\right)^2 \leq 1\) and each coordinate \(x_p^{(n)}\) tends to \(x_p\). He shows that a bounded bilinear form \(K(x,y)\) is continuous, and that \(K_n(x,y)\) tends to \(K(x,y)\) when \(n\) tends to \(\infty\). But he pays special attention to the completely continuous symmetric bilinear forms, and gives a separate proof that an orthogonal transformation can reduce any such form to the type
\begin{equation}
    \label{eq:5.40}
    K(x,y) = \frac{1}{\lambda_1} x_1 y_1 +\frac{1}{\lambda_2} x_2 y_2 + \ldots + \frac{1}{\lambda_n} x_n y_n + \ldots
\end{equation}
where the sequence \((|\lambda_n|)\) is either finite or tends to \(+\infty\). He realizes that this is a genuine generalization of formula \eqref{eq:5.24}, which is the special case in which \(\sum\limits_{p,q} k_{pq}^2 < +\infty\) (corresponding to what will later be called the \emph{Hilbert--Schmidt operators}); he also mentions another special case, the one in which \(K(x,x) > 0\) and \(\sum\limits_p k_{pp} < +\infty\) (corresponding to the positive \emph{nuclear operators} of a later date). This formula \eqref{eq:5.40} enables him to go beyond Fredholm by solving a system \eqref{eq:5.33} which is not derived any more from an integral equation, but in which the \(k_{pq}\) are only supposed to be such that the symmetric bilinear form \eqref{eq:5.35} is completely continuous. A final remark is that he repeatedly uses with great power what he calls a ``principle of choice'', which is equivalent to what will later be called the compactness of the unit ball for the weak topology, and that he extends his results to \emph{hermitian} sesquilinear forms
\begin{equation}
    \label{eq:5.41}
    K(x,y) = \sum_{p,q} k_{pq} x_p \overline{y}_q
\end{equation}
where this time the sequences \((x_p), (y_p)\) and \((k_{pq})\) consist of complex numbers, with \(k_{qp} = \overline{k}_{pq}\).

\section{The confluence of Geometry, Topology and Analysis}
\label{sec:5.3}

It may seem obvious to us that the results of Hilbert are but one step removed from what we now call the theory of Hilbert space; but if, in fact, the birth of that theory almost immediately followed the publication of Hilbert's papers, it seems to me that it is due to the fact that this publication precisely occurred during the emergence of a new concept in mathematics, the concept of \emph{structure}.

Until the middle of the XIX\supth{} century, mathematicians had been dealing with \emph{well determined} mathematical ``objects'': numbers, points, curves, surfaces, volumes, functions, operators. But the fact that algebraic manipulations on different kinds of ``objects'' had a strikingly similar appearance soon attracted attention (cf. chap. \ref{ch:4}, \ref{sec:4.3}), and after 1840 it gradually became clear that the essence of these manipulations did not lie in the \emph{nature} of the objects, but in the \emph{rules} to be followed in handling them, which might be the same for very different types of objects. However, a precise formulation of this idea had to wait for the adoption of the set-theoretic concepts and language; and it is only in 1895 that our definition of a \emph{group}, on an \emph{arbitrary} underlying set, was formulated by Weber \cite{225}. The trend towards the definition of algebraic structures then gained momentum, and around 1920 all fundamental notions of present-day Algebra had been defined.

In Analysis, no similar development had yet occurred in 1900. The extensions of the ideas of limit and continuity which had been formulated always were relative to \emph{special} objects such as curves, surfaces or functions. The possibility of defining such notions in an \emph{arbitrary} set is an idea which undoubtedly was first put forward by Fréchet in 1904 \cite{069}, and developed by him in his famous thesis of 1906 \cite{071}. The simplest and most fruitful method which he proposed for such definitions was the introduction of the notion of \emph{distance} (which he called ``écart'') on a set \(E\), a function \(d(x,y)\) defined for any pair \((x,y)\) of elements of \(E\), with values \(\geq 0\) and such that: 1) the relation \(d(x,y) = 0\) is equivalent to \(x = y\); 2) \(d(y,x) = d(x,y)\); 3) \(d(x,z) \leq d(x,y) + d(y,z)\) for any three elements of \(E\). It is extremely remarkable that with such simple axioms it is possible to extend most notions and arguments relative to neighborhoods, limits and continuity in the space \(\R^n\), which usually are introduced in relation to euclidean distance. But the greatest merit of Fréchet lies in the emphasis he put on three notions which were to play a fundamental part in all later developments of Functional Analysis: compactness, completeness and separability. Moreover, he did not limit himself to deriving general theorems in an abstract setting, but more than half of his thesis is devoted to very ``concrete'' metric spaces (as they came to be called later) closely linked to Analysis: the space of continuous real functions on a compact interval of \(\R\) with the topology of uniform convergence, the space \(\R^{\mathbb{N}}\) of all sequences \(n \mapsto x_n\), with the topology of simple convergence, the space of holomorphic functions in the disc \(|z| < 1\), with the topology of uniform convergence in compact subsets, and finally the space of all continuous ``curves'', images of \([0,1]\) in \(\R^3\) by continuous maps, with a ``distance'' which is a special case of what was later called the Hausdorff distance between two compact sets.

Clearly Hilbert's work immediately lent itself to application of these ideas, and even invited a bodily transfer of euclidean geometry in ``infinite dimension''. This is exactly what was done by Fréchet himself \cite{072} and by E. Schmidt \cite{192} in 1908. In E. Schmidt's paper, we find the definition of what we now call the (complex) space \(l^2\) (or \(l^2_{\C}\)), with the notions of scalar product and of norm (already written \(\norm{A}\)), the definition of orthogonality, of closed sets, and of vector subspaces (called ``\emph{lineares Funktionengebilde}''). The most interesting feature of that paper is the proof of the existence of the orthogonal projection of a point on a closed vector subspace, and the purely geometric way in which Schmidt uses this result to discuss the \emph{most general} system of linear equations in Hilbert space
\begin{equation}
    \label{eq:5.42}
    (x|a_n) = c_n \qquad (n=1,2,\ldots)
\end{equation}
where the \(a_n\) are arbitrary vectors of \(l^2\) and the \(c_n\) arbitrary complex numbers,

For each \(n\), Schmidt considers the closed linear affine varieties \(F_n\) of \(l^2\) defined by the equations \((x|a_j) = c_j\) for \(1 \leq j \leq n\), and the orthogonal projection \(x^{(n)}\) of the origin on \(F_n\); the necessary and sufficient condition of existence of a solution of the system \eqref{eq:5.42} is that the increasing sequence \(\left(\norm{x^{(n)}}\right)\) be bounded; the sequence \(\left(x^{(n)}\right)\) then has a weak limit \(x\) in \(l^2\), which is the solution of \eqref{eq:5.42} of smallest norm. Of course, each \(F_n\) must be different from the empty set, which means that any linear relation \(\sum\limits_{k=1}^n \lambda_k a_k = 0\) between the vectors \(a_n\) must imply \(\sum\limits_{k=1}^n \lambda_k c_k = 0\); one can then assume (by dropping some of the equations \eqref{eq:5.42}) that the \(a_n\) are linearly independent, and in that case, Schmidt easily obtains the explicit expression of \(\norm{x^{(n)}}\):
\begin{equation}
    \label{eq:5.43}
    \norm{x^{(n)}}^2 = \Delta_n / D_n
\end{equation}
with
\begin{equation}
    \label{eq:5.44}
    D_n =
    \left|
    \begin{matrix}
        (a_1 | a_1) & (a_1 | a_2) & \ldots & (a_1 | a_n) \\
        (a_2 | a_1) & (a_2 | a_2) & \ldots & (a_2 | a_n) \\
        \hdotsfor{4} \\
        (a_n | a_1) & (a_n | a_2) & \ldots & (a_n | a_n) \\
    \end{matrix}
    \right|
    , \qquad
    \Delta_n =
    \begin{array}{|c c c c c|}
    0 & c_1 & c_2 & \cdots & c_n \\
    \cline{2-5}
    \multicolumn{1}{|c|}{\overline{c}_1} & & & & \\
    \multicolumn{1}{|c|}{\overline{c}_2} &
    \multicolumn{4}{|c|}{} \\
    \multicolumn{1}{|c|}{\vdots} &
    \multicolumn{4}{|c|}{D_n} \\
    \multicolumn{1}{|c|}{\overline{c}_n} &
    \multicolumn{4}{|c|}{} \\
    \end{array}
\end{equation}

This geometric outlook was already shared in 1906--1907 by two other young mathematicians, E. Fischer and F. Riesz, in the remarkable work which led them (independently) to what is now called the Fischer--Riesz theorem, introducing a hitherto unsuspected link between Hilbert space and the theory of integration (\cite{066}, \cite[vol.~I, p.~378--395]{183}). The latter, from Cauchy to Jordan and Peano, had evolved in a manner completely independent from spectral theory \cite{103}. When Fredholm and E. Schmidt had tried to enlarge the scope of their results on integral equations by weakening the assumptions on the kernel \(K(x,y)\), they had nothing else at their disposal beyond the horrible and useless so-called ``Riemann integral''\footnote{As a function \(K(x,y)\) of two variables may be ``Riemann integrable'' even if the partial functions \(x \mapsto K(x,y)\) are not, Fredholm is compelled to assume that integrability both for the kernel \(K\) and all its partial functions! Although E.Schmidt wrote his dissertation in 1905, he probably had no knowledge of Lebesgue's thesis at that time.}, and it is likely that progress in Functional Analysis might have been appreciably slowed down if the invention of the Lebesgue integral had not appeared, by a happy coincidence, exactly at the beginning of Hilbert's work on integral equations. With the help of this marvelous new tool, Fischer and F. Riesz could define the space \(L^2(I)\) over a compact interval \(I \subset \R\), consisting of square integrable functions, when two functions are identified if they only differ in a set of measure 0. Their fundamental result is that, if to each function \(f \in L^2(I)\) one associates the sequence \((x_p)\) of its Fourier coefficients with respect to a complete orthonormal system (equations \eqref{eq:5.32}), this defines an \emph{isomorphism} of \(L^2(I)\) onto \(l^2\); from that it follows that \(L^2(I)\) is complete and separable. A byproduct was of course that the results of Fredholm and E. Schmidt could be applied without change to any integral equation where the kernel is only supposed to belong to \(L^2(I \times I)\), since it is then equivalent to a system of linear equations corresponding to a ``completely continuous'' bilinear form in the sense of Hilbert.

But the most important consequence of the Fischer--Riesz theorem is that it opened the way to the definition of the \(L^p\) spaces and to the general theory of normed spaces, which will be the subject of the next chapter.

\chapter{Duality and the definition of normed spaces}
\label{ch:6}
\setcounter{equation}{0}

\section{The search for continuous linear functionals}
\label{sec:6.1}

In chap. \ref{ch:4}, \ref{sec:4.3}, we saw that in 1897 C. Bourlet solved for the first time the problem of the determination of a linear map \(U: E \ F\) between ``function spaces'' by conditions of \emph{continuity}. In a short Note published in 1903 (\cite{094}, vol.~I, p.~405--408) Hadamard attacked the same problem with \(E = \Cont([a,b])\), space of real continuous functions in an interval \([a,b], F = \R\), and ``continuity'' means for him that \(U(f_n)\) tends to \(U(f)\) when \(f_n\) tends to \(f\) \emph{uniformly}. He chooses a fixed function \(F\) such that for any continuous function \(f\), one has
\begin{equation}
    \label{eq:6.1}
    f(x) = \lim_{n \to \infty} n \int_a^b f(t) F(n(t-x)) dt
\end{equation}
uniformly in \(x\); one has then
\begin{equation}
    \label{eq:6.2}
    U(f) = \lim_{n \to \infty} \int_a^b f(t) \Phi_n(t) dt
\end{equation}
where \(\Phi_n(t)\) is the value of \(U\) at the function \(x \mapsto nF(n(t-x))\) one may take \(F(x) = e^{-x^2}\), so that \(\Phi_n\) is continuous, but the choice of \(F\) is largely arbitrary (the argument is a typical case of what later will be called a ``regularization'' process).

In two papers published in 1904 and 1905, Fréchet gave another proof of Hadamard's theorem, and, what is more interesting, began to investigate the similar problems when \(\Cont([a,b])\) is replaced by another ``function space''; for instance \cite{070}, he remarked that if one takes for \(E\) the space \(\B([a,b])\) of all bounded integrable functions in \([a,b]\) (continuous or not) with the topology of uniform convergence, there were other continuous linear functionals than those given by Hadamard's formula, for instance the mappings \(f \mapsto c_1 f(x_1) + \ldots + c_m f(x_m)\), where the \(x_j\) are arbitrary points of \([a,b]\) and the \(c_j\) constants. Similarly, if one takes for \(E\) the space of all \(C^r\) functions in \([a,b]\), where convergence means uniform convergence for the function and its derivatives up to order \(r\), Fréchet showed that the continuous linear functionals could then be written
\[
    f \mapsto c_0 f(a) + c_1 f'(a) + \ldots + c_{r-1}f^{(r-1)}(a) + \lim_{n \to \infty} \int_a^b f^{(r)}(t) \Phi_n(t)dt.
\]

As soon as the study of Hilbert space began (chap. \ref{ch:5}, \ref{sec:5.3}), Fréchet \cite{072} and F. Riesz (\cite{183}, vol.~I, p.~386--388) independently showed that continuous ``linear functionals'' on Hilbert space \(l^2_{\R}\) (for the strong topology) could be written uniquely as \(x \mapsto (x|a)\) for a vector \(a \in l^2\).

Finally, in 1909, F. Riesz (\cite{183}, vol.~I, p.~400--402) was able to give a better form to Hadamard's theorem by removing the arbitrariness of the sequence \((\Phi_n)\); his idea was to use the Stieltjes integral, as Hilbert had done in his work on spectral theory (see chap. \ref{ch:8}, \ref{sec:8.2}): he showed that any continuous linear functional \(U: \Cont([a,b]) \ \R\) could be written \emph{uniquely}
\begin{equation}
    \label{eq:6.3}
    U: f \mapsto \int_a^b f(x) d\alpha(x)
\end{equation}
where \(\alpha\) is a function of bounded variation in \([a,b]\), provided one imposed on \(\alpha\) the additional conditions of being continuous on the left and such that \(\alpha(a) = 0\). His method consists in considering, for any \(t \in [a,b]\), the function \(f_t \in \Cont([a,b])\) equal to \(x-a\) for \(a \leq x \leq t\), and to \(t-a\) for \(t \leq x \leq b\), and the function \(A: t \mapsto U(f_t)\); he shows that this function is Lipschitzian, and takes for \(-\alpha(t)\) one of the ``derived numbers'' of \(A\) at the point \(t\); it is then easy to show that \(\alpha\) is a function of bounded variation, and it is a standard procedure to modify it in such a way that it satisfies the additional conditions mentioned above without changing \(U\).

Although the contemporaries did not realize the novelty of F. Riesz's approach, we are justified in seeing in his results (as he himself did) a radical departure from the conceptions of linear algebra prevalent in his time:

1° Whereas, even for the space \(L^2\), it was possible, due to the Fischer--Riesz theorem, to identify the elements of the space with sequences of numbers, generalizing the dominant Cayley concept of linear algebra as a theory of ``\(n\)-tuples'', no such identification was possible for \(\Cont([a,b])\), where one had to work directly on vectors, and not on their ``coordinates''

2° Functions of bounded variation may be discontinuous at a denumerable set of points, and therefore it was impossible to \emph{identify} any more the continuous linear functionals on \(\Cont([a,b])\) to the elements of that space (again in contrast to what happened in \(l^2\) according to the Riesz--Fréchet theorem).

These features would be still more conspicuous in the theory of \(L^p\) and \(l^p\) spaces, which F. Riesz began to investigate in 1910 (\cite{183}, vol.~I, p.~403).

\section{The \protect{\(L^p\)} and \protect{\(l^p\)} spaces}
\label{sec:6.2}

Once the \(L^2\) spaces had been defined, it was a natural generalization to define similarly the function spaces \(L^p(I)\) for any interval \(I \subset \R\), as the set of all complex valued measurable functions \(f\) defined in \(I\) and such that \(|f|^p\) is integrable, for any \(p > 0\) (two functions being identified if they are almost everywhere equal). The study of these spaces was begun by F. Riesz in a fundamental paper (\cite{183}, vol.~I, p.~441--497), second only in importance for the development of Functional Analysis to Hilbert's 1906 paper (chap. \ref{ch:5}, \ref{sec:5.2}).

Riesz limited himself from the start to the case \(p > 1\), in order to be able to use the Hölder and Minkowski inequalities
\begin{equation}
    \label{eq:6.4}
    \left| \sum_{k=1}^n a_k b_k \right| \leq \left( \sum_{k=1}^n |a_k|^p \right)^{1 / p} \left( \sum_{k=1}^n |b_k|^q \right)^{1 / q} \quad \text{for } \frac{1}{p} + \frac{1}{q} = 1
\end{equation}
\begin{equation}
    \label{eq:6.5}
    \left( \sum_{k=1}^n |a_k + b_k|^p \right)^{1 / p} \leq \left( \sum_{k=1}^n |a_k|^p \right)^{1 / p} + \left( \sum_{k=1}^n |b_k|^p \right)^{1 / p}
\end{equation}
which he first extended to measurable functions, showing that if \(f \in L^p(I), g \in L^q(I)\) then \(fg\) is integrable and
\begin{equation}
    \label{eq:6.6}
    \left| \int_I f(x)g(x)dx \right| \leq \left( \int_I |f(x)|^p dx \right)^{1 / p} \left( \int_I |g(x)|^q dx \right)^{1 / q},
\end{equation}
and that if \(f \in L^p(I), g \in L^p(I)\), then \(f + g \in L^p(I)\) and
\begin{equation}
    \label{eq:6.7}
    \left( \int_I |f(x) + g(x)|^p dx \right)^{1 / p} \leq \left( \int_I |f(x)|^p dx \right)^{1 / p} + \left( \int_I |g(x)|^p dx \right)^{1 / p}.
\end{equation}

His central theme is the study of infinite systems of linear equations
\begin{equation}
    \label{eq:6.8}
    \int_I f(x)g_\alpha(x)dx = c_\alpha
\end{equation}
where the \(g_\alpha\) belong to \(L^q(I)\) and one looks for a solution \(f \in L^p(I)\); this may be considered as the generalization of the problem E. Schmidt had treated in \(l^2\), due to the Fischer--Riesz theorem (chap. \ref{ch:5}, \ref{sec:5.3}, equations (42)). In order to adapt Schmidt's method to this problem, F. Riesz begins by extending a number of definitions and results from the theory of Hilbert space: strong convergence of a sequence \((f_n)\) of functions of \(L^p(I)\) to \(f \in L^P(I)\) is defined as meaning that \( \int_I |f(x)-f_n(x)|^p dx\) tends to 0. For weak convergence, he first takes as definition that \(\int_a^x f_n(t)dt\) tends to \(\int_a^x f(t)dt\) for all numbers \(x \in I\); and although he proves a little later that this definition is equivalent to the fact that the integrals \(\int_I (f(x)-f_n(x))g(x)dx\) tend to 0 for \emph{all} \(g \in L^q(I)\), he essentially uses the first definition to prove the generalization of Hilbert's ``principle of choice'' (\emph{i.e.} the weak compactness of the unit ball in \(L^p(I)\)), which will be one of his main ingredients in the solution of \eqref{eq:6.8}. The other ingredient is derived from a result obtained by E. Landau in 1907 \cite{136}; in 1906, Hellinger and Toeplitz had shown that if a sequence \((a_n)\) is such that the series \(\sum\limits_n a_n x_n\) is convergent for \emph{all} sequences \((x_n)\) in \(l^2\), then \((a_n)\) itself be longed to \(l^2\) \cite{106}; Landau proved that, more generally, if \(\sum\limits_n a_n x_n\) is convergent for all sequences \(x_n\) such that \(\sum\limits_n |x_n|^p < +\infty\), then \(\sum\limits_n |a_n|^q < +\infty\). Approximating functions of \(L^p\) by functions having only denumerably many values, F. Riesz deduced from Landau's result that if, for a measurable function \(g\), the product \(fg\) is integrable for all functions \(f \in L^p\), then necessarily \(g \in L^q\).

His solution of \eqref{eq:6.8} then proceeds along the same lines as E. Schmidt; he starts with a finite system \eqref{eq:6.8}, for which, using the standard method of analysis (Lagrange multipliers) he proves the existence and uniqueness of a solution \(f \in L^p\) for which \(\int_I |f(x)|^p dx\) is minimum. The problem is then to find a necessary and sufficient condition on the \(c_\alpha\) such that, when one picks from \eqref{eq:6.8} a finite system corresponding to the indices \(\alpha\) in an arbitrary finite subset \(H\), the corresponding minima \(M_H\) of the integral \(\int_I |f(x)|^p dx\) taken for the ``minimal'' solutions, are \emph{uniformly bounded} (independently of \(H\)); the use of the two ingredients mentioned above then leads to the existence of a solution of \eqref{eq:6.8} by an argument similar to E. Schmidt's. Of course, an explicitly expression of \(M_H\) (similar to formula \eqref{eq:5.43} of chap. \ref{ch:5}, \ref{sec:5.3}) is not available here, and the originality of F. Riesz lies in having found a completely different type of condition, namely the existence of a number \(M > 0\) such that, for \emph{any} finite subset \(H\) of indices, and \emph{any} family \((\lambda_\alpha)_{\alpha in H}\) of scalars, one has the inequality
\begin{equation}
    \label{eq:6.9}
    \left| \sum_{\alpha \in H} \lambda_\alpha c_\alpha \right| \leq M \cdot \left( \int_I \, \sum_{\alpha \in H} |\lambda_\alpha g_\alpha(x)|^q dx \right)^{1 / q}.
\end{equation}

F. Riesz in particular applies this conditions to the special case in which the \(g_\alpha\) are \emph{all} the functions of \(L^q(I)\); \eqref{eq:6.9} is then equivalent to the \emph{continuity} in \(L^q\) of the linear functional \(L\) defined by \(L(g_\alpha) = c_\alpha\), and he has thus generalized his previous results on \(l^2\) and \(\Cont(I)\), proving what we would now express by the statement that \emph{the dual of} \(L^q(I)\) \emph{can be identified with} \(L^p(I)\).

Of course the name ``dual'' is not yet used by F. Riesz, but he explicitly considers, for a ``bounded'' linear mapping \(T\) of \(L^p\) into itself (defined by the condition that \(\int_I |T(f)(x)|^p dx\) remains bounded for all \(f\) such that \(\int_I |f(x)|^p dx < 1\)), the \emph{transposed} mapping \(T'\) defined by the equation
\begin{equation}
    \label{eq:6.10}
    \int_I T(f)(x)g(x)dx = \int_I f(x)T'(g)(x)dx \quad \text{for \emph{all} } f \in L^p(I).
\end{equation}
Indeed, for a function \(g \in L^q(I)\), this defines (up to a null set) a unique function \(T'(g)\), which (by F. Riesz's previous results) also belongs to \(L^q(I)\); furthermore, it is easy to show that the mapping \(T'\) of \(L^q\) into itself is also linear and ``bounded''. F. Riesz then used this concept to obtain a necessary and sufficient condition for the mapping \(T\) to be \emph{bijective}: he showed that such a condition is the existence of a number \(m > 0\) such that \emph{both} inequalities
\begin{equation}
    \label{eq:6.11}
    \begin{gathered}
    \int_I |T(f)(x)|^p dx \geq m \cdot \int_I |f(x)|^p dx \\
    \int_I |T'(g)(x)|^q dx \geq m \cdot \int_I |g(x)|^q dx
    \end{gathered}
\end{equation}
are satisfied for all \(f \in L^p(I)\) and all \(g \in L^q(I)\). F. Riesz had thus given, for the first time, examples of what we now call \emph{reflexive} Banach spaces not isomorphic to their dual\footnote{The dual of \(L^1(I)\) for a compact interval \(I \in \R\) was shown to be isomorphic to \(L^\infty(I)\) by H. Steinhaus \cite{202}; he uses the fact that in that case \(L^2(I) \subset L^1(I)\), and therefore a continuous linear functional on \(L^1(I)\) is also continuous on \(L^2(I)\).}. In his 1913 book on infinite systems of linear equations (\cite{183}, vol.~II, p.~835--1016 and \cite{184}) he treated in a similar way the \(l^p\) spaces for \(p > 1\) (defined as the set of sequences \((x_n)\) of complex numbers such that \(\sum\limits_n |x_n|^p < +\infty\)); in addition he stated without proof that for \(p \neq 2\), no isomorphism of \(l^p\) and \(L^p\) existed any more, in contradistinction to the Fischer--Riesz theorem (\emph{ibid}. Vol.~I, p.~444--445).

\section{The birth of normed spaces and the Hahn--Banach theorem}
\label{sec:6.3}

In 1911, F. Riesz combined his methods for the treatment of the system \eqref{eq:6.8} in \(L^p\) with the Hadamard--Riesz theorem on linear functionals in \(\Cont([a,b])\) in order to study the systems of linear equations
\begin{equation}
    \label{eq:6.12}
    \int_a^b g_\alpha(x)d\xi(x) = c_\alpha
\end{equation}
where \([a,b]\) is a compact interval in \(\R\), the \(g_\alpha\), are continuous in \([a,b]\), the \(c_\alpha\) are given scalars, and one has to determine a function \(\xi\) of bounded variation in \([a,b]\) satisfying the equations \eqref{eq:6.12} for all \(\alpha\). This may be considered as a generalization of a problem which had first been proposed and solved by T. Stieltjes in 1894, the ``moment problem'': it consists in determining an increasing function \(\xi\) in \([0,+\infty[\) such that
\begin{equation}
    \label{eq:6.13}
    \int_0^\infty x^n d\xi(x) = c_n \geq 0 \quad \text{for } n = 0,1,2,\ldots
\end{equation}
(the left hand sides are called the ``moments'' of the function \(\xi\), a terminology stemming from probability theory) \cite{205}; the same problem was later considered when the interval \([0,+\infty[\) is replaced by \(]-\infty, +\infty[\) (the ``Hamburger moment problem'') or by a compact interval \([a,b]\) (the ``Hausdorff moment problem'') \cite{003}.

The solutions to these ``moment problems'' consist in giving explicit conditions on the \(c_n\) involving existence (or existence and uniqueness) of the function \(\xi\) (or rather of the measure \(d\xi\)). The condition given by F. Riesz for the existence of a solution \(\xi\) of the general system \eqref{eq:6.12} is similar to condition \eqref{eq:6.9}, namely the existence of a number \(M > 0\) such that, for any finite family \((\lambda_\alpha)_{\alpha \in H}\) of scalars, one has
\begin{equation}
    \label{eq:6.14}
    \left| \sum_{\alpha \in H} \lambda_\alpha c_\alpha \right| \leq M \cdot \sup_x \left| \sum_{\alpha \in H} \lambda_\alpha g_\alpha(x) \right|;
\end{equation}
(he explicitly observed that the right hand side of this inequality is the limit of the right hand side of \eqref{eq:6.9} when \(q\) tends to \(+\infty\)). His proof is similar to the proof for \eqref{eq:6.8}; he first restricts himself to the case of finite systems \eqref{eq:6.12}, obtains the existence of a ``minimal'' solution of such a system, and then, using a ``principle of choice'' (in our language, the weak compactness of the unit ball in the space of Stieltjes measures), he shows that the condition \eqref{eq:6.14} is sufficient for an arbitrary system  \eqref{eq:6.12}; his procedure is more complicated than for  \eqref{eq:6.8}, because even in the case of a finite system \eqref{eq:6.12}, there is no more uniqueness for the ``minimal'' solutions (\cite{183}, vol.~II, p.~798--827).

We now interpret condition \eqref{eq:6.9} in the following way: first, if \(\sum\limits_\alpha \lambda_\alpha g_\alpha = 0\) in \(L^q(I)\), then \(\sum\limits_\alpha \lambda_\alpha c_\alpha = 0\); this implies that, if \(F\) is the vector subspace of \(L^q(I)\) generated by the \(g_\alpha\), there is a well determined linear form \(L\) defined in \(F\) such that \(L(g_\alpha) = c_\alpha\) for every \(\alpha\). Condition \eqref{eq:6.9} then means that this linear form \(L\) is \emph{continuous} in \(F\); the existence of an \(f \in L^p(I)\) such that \(L(g_\alpha) = \int_I f(x)g_\alpha(x)dx\) for all \(\alpha\) then means that \(L\) can be \emph{extended} to a continuous linear form defined in the \emph{whole} space \(L^q(I)\); in other words, it is a special case of what we now call the Hahn--Banach theorem. There is a similar interpretation of condition \eqref{eq:6.14}, replacing \(L^q(I)\) by \(\Cont([a,b])\).

Such an interpretation of his results was \emph{not} given by F. Riesz; the first mention of that point of view appears in a paper written in 1912 by the Austrian mathematician E. Helly (1884--1943), in which he gives a different proof of F. Riesz's results on the systems \eqref{eq:6.12} \cite{107bis}. After an interval of 9 years (due to the first World War, in which he was a prisoner of war in Russia), Helly returned to his method in a paper of 1921 \cite{108} which again should be considered as a landmark in the history of Functional Analysis, since instead of considering \emph{special} spaces such as the \(l^p, L^p\) or \(\Cont([a,b])\), he for the first time deals with general ``normed sequence spaces'' by methods which do not depend on special features of the space, contrasting with the ones used by E. Schmidt and F. Riesz.\footnote{It seems that during the period 1910--1920, F. Riesz always had in mind possible axiomatic generalizations of his results, although he did not publish anything in that direction (\cite{183}, vol.~I, p.~452).}

Helly considers vector subspaces of the vector space \(\C^{\N}\) of all sequences of complex numbers, and assumes that on such a subspace \(E\) there has been defined a \emph{norm} \(\norm{x}\) (he does not use that name nor the notation) such that: 1) \(\norm{x} \geq 0\) and the relation \(\norm{x} = 0\) is equivalent to \(x = 0\); 2) \(\norm{\lambda x} = |\lambda| \cdot \norm{x}\) for any scalar \(\lambda\); 3) \(\norm{x+y} \leq \norm{x} + \norm{y}\); this defines on \(E\) a distance \(d(x,y) = \norm{x-y}\) in the sense of Fréchet. Of course norms had been defined in the spaces \(l^p, L^p\) and \(\Cont([a,b])\); but Helly seems to be the first to have noticed the relations of that notion with the concepts of \emph{convexity} introduced earlier by Minkowski in his ``Geometry of numbers'' (\cite{161} and \cite{162}). He had shown that the concept of norm on a \emph{finite} dimensional space \(\R^n\) (with the scalars limited to real values) was equivalent to the notion of ``symmetric convex body'', i.e. a closed, symmetric, bounded convex set in which the origin \(O\) is an interior point: such a set \(B\) can be defined by an inequality \(p(x) \leq 1\), where \(p\) is a uniquely determined norm. The boundary of such a set is defined by the equation \(p(x) = 1\), and Minkowski had proved that for each point \(x_0\) of that boundary, there existed at least one \emph{hyperplane of support} \(H\), containing \(x_0\), and such that \(B\) lies entirely \emph{on one side} of \(H\). If, for an \(n\)-tuple of real numbers \(u = (u_1, \ldots, u_n)\) and a point \(x = (x_1, \ldots, x_n)\), one writes \(\ip{u,x} = u_1 x_1 + \ldots + u_n x_n\), the equation of \(H\) has the form \(\ip{u,x} = 1\) for a suitable \(u\), and one has the inequality \(\ip{u,x} \leq p(x)\) for all \(x \in \R^n\), with \(\ip{u, x_0} = p(x_0)\); the \(n\)-tuples \(u\) being identified with the corresponding linear forms \(x \mapsto \ip{u, x}\) on \(\R^n\), Minkowski had also defined the ``support function'' \(q(u) = \sup\limits_{x \neq 0}\ip{u,x}/p(x)\), and shown that it was also a norm on \(\R^n\), ``dual'' to \(p\) and such that the hyperplanes of support of \(B\) are the hyperplanes \(\ip{u,x} = 1\) with \(q(u) = 1\); furthermore, \(p\) is the norm ``dual'' to \(q\), in other words \(p(x) = \sup\limits_{u \neq 0}\ip{u,x}/q(u)\).

To transfer to spaces of sequences these concepts and definitions, Helly associates to \(E\) the subspace \(E'\) of \(\C^{\N}\) consisting of all the sequences \(u = (u_n)\) such that the series \(\sum\limits_n u_n x_n\) converges for \emph{all} \(x = (x_n)\) in \(E\) \footnote{This is not always the dual of \(E\) as we now understand that word.}, and he then considers \(\ip{u,x} = \sum\limits_n u_n x_n\). For any \(u \in E'\), the number \(\norm{u} = \sup\limits_{x \neq 0}\ip{u,x}/\norm{x}\) defines a norm on \(E'\) provided it is not 0 for some elements \(u \neq 0\). Excluding that case, Helly first obtains a weak generalization of Minkowski's result on the hyperplanes of support; if \(B\) is the subset of \(E\) defined by \(\norm{x} \leq 1\), he shows that the hyperplane \(H\) defined by \(\ip{u,x} = 1\) meets \(B\) if \(\norm{u} < 1\), does not meet \(B\) for \(\norm{u} > 1\), but if \(\norm{u} = 1\), the intersection \(H \cap B\) may very well be empty: an example is given by taking \(E = l^1, E' = l^\infty\) and for \(H\) the hyperplane \(\sum\limits_{n=1}^\infty \left(1 - \frac{1}{n}\right) x_n = 1\).

The central problem in Helly's paper is the solution of a system
\begin{equation}
    \label{eq:6.15}
    \ip{u^{(\nu)}, x} = c_\nu \qquad (\nu=1,2,\ldots)
\end{equation}
where the \(u^{(\nu)}\) belong to \(E'\) and one looks for a solution \(x \in E\). The inequality \(|\ip{u,x}| \leq \norm{u} \cdot \norm{x}\) immediately yields the necessary condition similar to Riesz's conditions \eqref{eq:6.9} and \eqref{eq:6.14}, namely the existence of a number \(M > 0\) such that
\begin{equation}
    \label{eq:6.16}
    \left| \sum_{\nu=1}^n \lambda_\nu c_\nu \right| \leq M \cdot \norm*{\sum_{\nu=1}^n \lambda_\nu u^{(\nu)}}
\end{equation}
for any \(n\) and \emph{all} choices of scalars \(\lambda_\nu\); but the example given above (for a single equation) shows that there may well be no solution such that \(\norm{x} = M\), even when condition \eqref{eq:6.16} is satisfied.

Helly, as Schmidt and F. Riesz had done, first considers the case of a \emph{finite} system \eqref{eq:6.15} of \(N\) equations, where as usual the \(u^{(\nu)}\) may be supposed to be linearly independent. The mapping \(f: x \mapsto (\ip{u^{(\nu)},x})_{1 \leq \nu \leq N}\) of \(E\) into \(\C^{\N}\) is then surjective; Helly shows that on \(\C^{\N}, \norm{y} = \inf\limits_{f(x)=y} \norm{x}\) is a norm (for us it is the natural norm on \(E/f^{-1}(0)\) deduced from the norm on \(E\)); condition \eqref{eq:6.16} then guarantees the existence of a solution \(x\) of \eqref{eq:6.15} such that \(\norm{x} < M_1\) for \emph{any} \(M_1 > M\) (if not necessarily for \(M_1 = M\)).

The passage from finite systems \eqref{eq:6.15} to the general case is the most original idea of Helly; he splits the problem in two:

A) Given \(M_1 > M\), find a linear form \(L: E' \ \C\), such that \(|L(u)| \leq M_1 \cdot \norm{u}\) for all \(u \in E'\) and such that \(L(u^{(\nu)}) = c_\nu\) for all \(\nu\).

B) When such a linear form \(L\) has been found, find if possible an element \(p \in E\) such that \(\ip{p,u} = L(u)\) for all \(u \in E'\).

To treat problem A), Helly assumes the additional condition that \(E'\) is \emph{separable} as a metric space; he then proves the existence of a solution (a special case of the Hahn--Banach theorem) in the following way. Let \((p^{(\nu)})\) be a sequence of elements of \(E'\) which is dense in that space. Helly chooses an increasing sequence \(M < M^{(1)} < M^{(2)} < \ldots < M_1\) of numbers, and the main point of his proof consists in showing that there exists a family \((\gamma_\nu)\) of complex numbers such that, for any pair of integers \(m \geq 1, n \geq 1\), and any pair of families \((\lambda_\nu), (\mu_\nu)\) of scalars, one has
\begin{equation}
    \label{eq:6.17}
    \left| \sum_{\nu = 1}^n \lambda_\nu c_\nu + \sum_{\nu = 1}^m \mu_\nu \gamma_\nu \right| \leq M^{(m)} \cdot \norm*{\sum_{\nu = 1}^n \lambda_\nu u^{(\nu)} + \sum_{\nu = 1}^m \mu_\nu p^{(\nu)}}.
\end{equation}

It is then easy to show that there exists a linear form \(L\) on \(E'\) such that \(L(p^{(\nu)}) = \gamma_\nu\) for all indices \(\nu\), and that it is a solution of problem A).

The proof of \eqref{eq:6.17} is done by induction on \(m\), the case \(m = 0\) being the assumption \eqref{eq:6.16}. One has then to prove the existence of a point \(\gamma_{m+1} \in \C\) which, for any integer \(n \geq 1\) and any pair of families of scalars \((\lambda_\nu)_{1 \leq \nu \leq n}, (\mu_\nu)_{1 \leq \nu \leq m}\), belongs to \emph{all} disks defined in \(\C\) by
\begin{equation}
    \label{eq:6.18}
        \left| \sum_{\nu = 1}^n \lambda_\nu c_\nu + \sum_{\nu = 1}^m \mu_\nu \gamma_\nu + \gamma_{m+1} \right| \leq M^{(m+1)} \norm*{\sum_{\nu = 1}^n \lambda_\nu u^{(\nu)} + \sum_{\nu = 1}^m \mu_\nu p^{(\nu)} + p^{(m+1)}}.
\end{equation}
However, a general result on convex sets in a finite dimensional space, proved by Helly himself, reduces that question to proving that \emph{any three} of the disks \eqref{eq:6.18} have a common point; and this is shown by Helly to be a consequence of the result proved before for \emph{finite} systems \eqref{eq:6.15}.

Turning to problem B), Helly discovers that it is quite possible that it has \emph{no} solution; in our language, he gives the first example of \emph{non reflexive} Banach spaces\footnote{F. Riesz had already observed that one could define on the space of functions of bounded variation continuous linear functionals which were not of the form \(\xi \mapsto \int_a^b f(x)d\xi(s)\) for a continuous function \(f\) (for instance one can take for \(f\) an increasing discontinuous function) (\cite{183}, vol.~II, p.~827).}. That example is the space \(E\) of all sequences \((x_n)\) such that the series \(\sum\limits_{k=1}^\infty x_k\) converges, with the norm \(\norm{x} = \sup\limits_n \left| \sum\limits_{k=n}^\infty x_k \right|\); Helly proves that \(E'\) consists of all sequences \((u_k)\) such that \(\norm{u} = |u_1| + \sum\limits_{k=1}^\infty |u_{k+1} - u_k|\) is finite, \(\norm{u}\) being the natural norm on \(E'\); then if one takes \(L(u) = \lim\limits_{k \to \infty} u_k, L\) is continuous on \(E'\) but there is no \(p \in E\) such that \(L(u) = \ip{p,u}\).

Starting from the work of F. Riesz and Helly, it was a natural generalization to define norms on \emph{arbitrary} vector spaces over \(\R\) or \(\C\), and not only on spaces of functions or on subspaces of \(\C^{\N}\). This was done independently by H. Hahn \cite{097} and S. Banach \cite{012}, who restrict themselves to \emph{complete} spaces.

Banach's paper is his thesis, written in 1920: although he does not mention convexity, he is careful to develop and extensively use a geometric language. He is mainly interested in continuous linear operators \(u: E \to F\), where \(E\) and \(F\) are arbitrary normed complete spaces, and in limits of sequences of such operators. Hahn's point of view is similar, although he is only concerned with linear forms; neither he nor Banach are at that moment interested in the problem of extension of linear forms, and we postpone a more detailed discussion of their papers of 1922--23 to \ref{sec:6.4}. We should however mention that in his thesis Banach gives the ``abstract'' formulation of the method of successive approximations (chap. \ref{ch:2}, \ref{sec:2.1}) as a ``contraction principle'': if \(F\) is a mapping of a complete normed space \(E\) into itself such that \(\norm{F(x)-F(y)} \leq k\norm{x-y}\) with \(0 < k < 1\), then the sequence \((x_n)\) defined by induction as \(x_{n+1} = F(x_n)\) (\(x_0\) arbitrary) converges to the unique ``fixed point'' \(x\), such that \(F(x) = x\).

It was only in 1927 that Hahn returned to Helly's paper, in the general context of complete normed spaces, and completely solved the extension problem for such spaces \cite{098}. He proceeds by induction as Helly had done, but at the same time he greatly simplifies and generalizes the method by introducing, for the first time in general problems of Functional Analysis \footnote{Transfinite induction had been used by analysts ever since Cantor, but the application of transfinite induction closest to Hahn's is probably the method by which Banach, in 1923, had proved the existence on \(\R\) of a ``measure'' defined on all subsets of \(\R\) and simply additive \cite{013}.}, \emph{transfinite induction} instead of the ordinary kind. In a complete normed space \(E\), one has a vector subspace \(V\) and there is defined on \(V\) a (real valued) linear form \(f\) such that \(|f(x)| \leq M\norm{x}\) for \(x \in V\); the problem is to extend \(f\) to a linear form \(F\) on \(E\) such that \(|F(y)| < M\norm{y}\) for \(y \in E\). Hahn begins by showing the existence of an ordinal \(\gamma\), and of a mapping \(\xi \mapsto V_\xi\), which, to very ordinal \(\xi < \gamma\) associates a vector subspace \(V_\xi\) of \(E\) such that \(V_0 = V, V_\xi \subset V\) for \(\xi < \eta, V_\xi\) has codimension 1 in \(V_{\xi + 1}\), and \(E\) is the union of the \(V_\xi\) for \(\xi < \gamma\). The problem is then easily reduced to the case in which \(V\) has codimension 1 in \(E\), and then \(E\) is generated by \(V\) and an element \(a \not\in V\); Hahn considers the l.u.b \(B\) of the numbers \(f(x) - M\norm{x-a}\) for \(x \in V\), and the g.l.b. \(A\) of the numbers \(f(x) + M\norm{x-a}\) for \(x \in V\), and, using the assumption \(|f(x)| \leq M\norm{x}\) for \(x \in V\), he easily shows that \(A \leq B\); the extension \(F\) is then defined by \(F(x+\lambda a) = f(x) + \lambda c\) for all \(\lambda \in \R\), where \(c\) is any number such that \(A \leq c \leq B\).

As a particular case of his theorem, Hahn shows that for any vector \(a \neq 0\) in \(E\) there exists a continuous linear form \(L\) on \(E\) such that \(\norm{L} = 1\) and \(L(a) = \norm{a}\); he then formally introduces the \emph{dual} space \(E'\) of \(E\) (``polare Raum'' in his terminology) which is not reduced to 0 due to the preceding result; he writes \(B(u,x)\) instead of \(u(x)\) for \(x \in E, u \in E'\), and considers for any \(x \in E\), the linear form \(c(x): u \mapsto B(u,x)\) on \(E'\), for which he shows that \(\norm{c(x)} = \norm{x}\). In other words, he has defined a linear isometry \(c\) of \(E\) into its second dual \(E''\), and he says a space \(E\) is ``regulär'' if \(c\) is bijective (our reflexive spaces). It may therefore rightly be said that with this paper of Hahn, duality theory at last has come into its own.

Two years later, Banach, who apparently was not aware of Hahn's paper, published the same theorem with the same proof (he later acknowledged Hahn's priority); in addition, he recognized that the argument could be generalized: if \(p\) is a real valued function defined in a vector space \(E\) and such that \(p(x+y) \leq p(x) + p(y)\) and \(p(\lambda x) = \lambda p(x)\) for \(\lambda \geq 0\), and if \(f\) is a linear form defined in a vector subspace \(V\) of \(E\) and such that \(f(x) \leq p(x)\) in \(V\), then it is possible to extend \(f\) to a linear form \(F\) defined in \(E\) and such that \(F(x) \leq p(x)\) in \(E\). This extension was to play later an important role in the development of the theory of locally convex spaces (cf. chapter \ref{ch:8}).

\section{The method of the gliding hump and Baire category}
\label{sec:6.4}

In his 1922 paper \cite{097}, Hahn proved the following theorem: let \(E\) be a complete normed space, \((u_n)\) a sequence of continuous linear forms on \(E\), and suppose that for each \(x \in E\), the sequence of numbers \(|u_n(x)|\) is bounded by a number \emph{depending on \(x\)}; then the sequence of the norms \(\norm{u_n}\) is \emph{bounded}. The proof is by contradiction; assuming that the sequence \((\norm{u_n})\) is unbounded, one determines by induction a sequence \((x_k)\) in \(E\) and a sequence \((n_k)\) of integers such that:
\begin{enumerate}[1°,noitemsep]
    \item the series \(\sum\limits_{k=1}^\infty x_k\) converges to an element \(x \in E\);

    \item \(\sum\limits_{j=k+1}^\infty \left| u_{n_k}(x_j) \right| \leq 1\);

    \item \(\left| u_{n_k}(x_k) \right| \geq k + \sum\limits_{j=1}^{k-1} \left| u_{n_k}(x_j) \right|\).
\end{enumerate}
Then one has for each \(k\),
\[
    \left| u_{n_k}(x) \right| \geq \left| u_{n_k}(x_k) \right| - \sum\limits_{j=1}^{k-1} \left| u_{n_k}(x_j) \right| - \sum\limits_{j=k+1}^{\infty} \left| u_{n_k}(x_j) \right| \geq k-1
\]
which contradicts the assumption. To do this, one assumes the \(u_{n_j}\) have been determined for \(j < k\), and one considers a ball
\[
    B_k: \norm{x} \leq 2^{-k} \cdot \inf_{j < k}\left(\norm{u_{n_j} + 1}\right)^{-1}
\]
in \(E\); the assumption that \((\norm{u_n})\) is unbounded guarantees the existence of an index \(n_k\) and a point \(x_k \in B_k\) for which condition 3° holds; conditions 1° and 2° are then deduced from the choice of the radius of the ball \(B_k\).

This is often called the ``method of the gliding hump'': in the sequence of values \(|u_{n_k}(x_j)|\) when \(j\) varies from 1 to \(+\infty\), the index \(j = k\) corresponds to a ``hump'' much bigger than the sum of the contributions of the other indices.

The result can be put in a different form: if the sequence \(|u_{n_k}(x_j)|\) is unbounded, there exists at least one \(x \in E\) such that the sequence \(|u_{n_k}(x)|\) is unbounded.

In this form, the first example of the method of the gliding hump is probably the way in which Lebesgue, in 1905 (\cite[vol.~III, p.~101]{138} and \cite[p.~86--88]{139}) constructed a continuous periodic function \(F(x)\) in \([0,2\pi]\) whose Fourier series diverges at the point 0. He had proved that, if one writes \(S_n(g)\) for the sum of the first \(n\) terms of the Fourier series of a continuous function \(g\), it is possible to find a sequence \((g_n)\) of continuous periodic functions of bounded variation such that \(|g_n(t)| \leq 1\) in \([0,2\pi]\) and that the sequence of values \(S_n(g_n)(0)\) tends to \(+\infty\). He then defines
\[
    F(x) = \varepsilon_1 f_1(n_1 x) + \varepsilon_2 f_2(n_2 x) + \ldots + \varepsilon_k f_k(n_k x) + \ldots
\]
where the \(\varepsilon_k\) are \(> 0\) and such that \(\sum\limits_{k=1}^\infty \varepsilon_k = 1\), the \(f_k\) are continuous periodic functions of bounded variation such that \(f_k(t)| \leq 1\) in \([0,2\pi]\) and that \(\left|S_{p_k}(f_k)(0)\right| \geq k/\varepsilon_k\), for an increasing sequence \((p_k)\) of integers. Finally the increasing sequence \((n_k)\) of integers is chosen in such a way that \(n_k > n_{k-1} p_{k-1}\) and that, for the continuous function of bounded variation \(F_k(x) = \varepsilon_1 f_1(n_1 x) + \ldots + \varepsilon_{k-1} f_{k-1}(n_{k-1} x)\), all the sums \(S_n(F_k)(0)\) are \(\leq 2\) in absolute value for \(n \geq n_k\), (they converge to \(F_k(0)\)). This choice implies that, for \(j > k\), the sum of the first \(n_k p_k\) terms of the Fourier series of \(f_j(n_j x)\) is reduced to the \emph{first} term of the series, hence is \(\leq 1\) in absolute value; using these definitions it is easy to check that \(\left|S_{n_k p_k}(F)(0)\right| \geq k-3\) for all \(k\).

One year later, Hellinger and Toeplitz, two students of Hilbert, found a rather surprising complement to the definition he had given of a \emph{bounded} bilinear form (chap. \ref{ch:5}, \ref{sec:5.2}); instead of assuming that \(|K_n(x,y)| \leq M\) for all \(n\) and all \(x = (x_p)\) and \(y = (y_p)\) such that \(\sum\limits_p x_p^2 \leq 1\) and \(\sum\limits_p y_p^2 \leq 1\), they showed that it was enough to assume that for each such pair \((x,y)\), one had \(|K_n(x,y)| < M_{x,y}\) for all \(n\), where the number \(M_{x,y}\) \emph{might depend} on \(x, y\) \emph{in an arbitrary way}. Independently of Lebesgue, they proved that result by a ``gliding hump'' method, constructing a pair \((x,y)\) for which the sequence \((|K_n(x,y)|)\) is unbounded if \(K\) is not a \emph{bounded} form in Hilbert's sense \cite{106}.

During the next 20 years, many more examples of the ``gliding hump'' method appeared in the literature: Lebesgue used it repeatedly in a 1909 paper on ``singular integrals'' (\cite{138}, vol.~III, p.~259--351), where one looks for conditions on ``kernels'' \(K_n\) insuring that the integrals \(\int_a^b f(t)K_n(t,x)dt\) tend to \(f(x)\) when \(n\) tends to \(+\infty\), for various kinds of function \(f\). The method was also prominent in the study of ``summation processes'', where one ``transforms'' a sequence \((x_n)\) into a sequence \((y_n)\) by the formulas \(y_n = \sum\limits_{m=1}^\infty a_{nm} x_m\), and has to look for conditions on the \(a_{nm}\) insuring that when \((x_n)\) has a limit, \((y_n)\) tends to the same limit (\cite{193}, vol.~II, p.~389--321). Hahn's paper of 1922 \cite{097} was written to give a general background to all these results, showing that they all were consequences of his general theorem. Independently, Banach, in his thesis, proved a theorem more general than Hahn's, the \(u_n\) being now continuous linear operators from a complete normed space \(E\) into a complete normed space \(F\); he showed that the assumption that the norms \(\norm{u_n(x)}\) are bounded for each \(x\) by a number depending on \(x\), implies that the sequence of the norms \(\norm{u_n}\) is bounded.

Finally, in 1927, Banach and Steinhaus (using an idea of Saks) discovered that this theorem could be proved without using the ``gliding hump'' method, by an application of a theorem Baire had proved in 1899 \cite{011}: he had shown that in \(\R^n\), the intersection of a denumerable family of dense open subsets is itself dense\footnote{For \(n =1\), the same result had been proved two years earlier by W. Osgood \cite{170}.}; this implies that if \(u\) is a real function defined and \emph{lower semi-continuous} in \(\R^n\), and if \(u(x) < +\infty\) for each \(x \in \R^n\), then any non empty open subset \(U\) of \(\R^n\) contains a non empty open subset \(V\) such that \(\sup\limits_{x \in V} u(x) < +\infty\). These results and their proofs immediately generalize when \(\R^n\) is replaced by an arbitrary \emph{complete metric space}. If now \(H\) is a set of linear mappings from a complete normed space \(E\) into a complete normed space \(F\), and if for each \(x \in E, \sup\limits_{u \in H} \norm{u(x)} < +\infty\), the function \(p(x) = \sup\limits_{u \in H} \norm{u(x)}\) is lower semi-continuous, and from the Baire theorem it follows that \(p\) is bounded in a neighborhood of 0, which implies that \(\sup\limits_{u \in H} \norm{u}\) is finite \cite{16}.

\section{Banach's book and beyond}
\label{sec:6.5}

In 1932 S. Banach published a book \cite{015} containing a comprehensive account of all results known at that time in the theory of normed spaces, and in particular the theorems he had published in his papers of 1923 and 1929. A large part was devoted to the concept of weak convergence and its generalizations, which he had begun to study in 1929; we shall postpone to chap. \ref{ch:8}, \ref{sec:8.1} the discussion of these questions. The most remarkable result contained in that book is another consequence of Baire's theorem, discovered by Banach, and much deeper than the Banach--Steinhaus theorem: if \(u\) is a continuous linear mapping from a complete normed space \(E\) into a complete normed space \(F\), then either \(u(E)\) is \emph{meager} in \(F\) (a set ``of first category'' in the terminology of Baire), or \(u(E) = F\). An immediate consequence is the famous \emph{closed graph theorem}: if \(u\) is a linear mapping from \(E\) to \(F\) having a closed graph in \(E \times F\), then \(u\) is continuous. These surprising results have become two of the most powerful tools in all applications of Functional Analysis.

These features, as well as many applications to classical Analysis, gave the book a great appeal, and it had on Functional Analysis the same impact that van der Waerden's book had on Algebra two years earlier. Analysts all over the world began to realize the power of the new methods and to apply them to a great variety of problems; Banach's terminology and notations were universally adopted, complete normed spaces became known as \emph{Banach spaces}, and soon their theory was considered as a compulsory part in most curricula of graduate students. After 1935, the theory of normed spaces became part of the more general theory of \emph{locally convex spaces}, which we shall discuss in chapter \ref{ch:8}; more recently however, there has been a renewed surge of interest in the special properties of normed spaces and their ``geometry''; it is too soon, as yet, to have a clear idea of the scope of these results and of their relation to other parts of mathematics, and we refer the interested reader to \cite{004}, \cite{017}, \cite{047}, \cite{050}, \cite{116}, \cite{134}, \cite{149}, \cite{150} and \cite{185}.

\chapter{Spectral theory after 1900}
\label{ch:7}
\setcounter{equation}{0}

\section{F. Riesz's theory of compact operators}
\label{sec:7.1}

We already mentioned that Fredholm's paper attracted many mathematicians to the theory of integral equations, and also to the theory of infinite systems of linear equations, especially after Hilbert had given it a new impetus. We shall not examine these papers, most of which are concerned with special problems, without much bearing on the progress of Functional Analysis, and we refer the interested reader to \cite{107} (in particular p.~1543--1552 and p.~1574--1575).

The ``Fredholm alternative'' corresponded, in ``infinite dimensional linear algebra'' to the classical relation between kernel and image of an endomorphism of a finite dimensional vector space over \(\C\). But for such endomorphisms, much more was known, namely the Jordan normal form which characterized them up to ``similitude'', and a natural question was to investigate similar properties of the Fredholm operators. However, only partial results in that direction were obtained, before F. Riesz in 1916 (in a paper written in Hungarian (\cite{183}, vol.~II, p.~1017--1052) and only published in German in 1918 (\emph{ibid.}, p.~1053--1080)) gave a complete answer to that question, and found the proper context to Fredholm's results, in what is now known as the \emph{Riesz--Fredholm theory of compact operators}.

F. Riesz never adopted Hilbert's method of dealing with linear equations \emph{via} bilinear forms, but followed Fredholm in using instead \emph{operators}. In his work on \(l^p\) spaces (\cite{183}, vol.~II, p.~876--911 and \cite{184}), he therefore had translated Hilbert's conception of a completely continuous bilinear form (chap. \ref{ch:5}, \ref{sec:5.2}) into the notion of \emph{completely continuous operator}: for him it was a linear mapping of \(l^p\) into itself which transformed weakly convergent sequences into strongly convergent ones. The novelty in his 1918 paper is that he realized that he could give an equivalent definition without mentioning weak convergence, using instead the general concept of compactness introduced by Fréchet: the condition was that the linear operator transformed a bounded set into a \emph{relatively compact} one (for the strong topology). Now this can be defined for an \emph{arbitrary} normed space instead of \(l^p\); in his 1918 paper F. Riesz restricted himself to the space \(\Cont(I)\) for a compact interval \(I \subset \R\), but he explicitly mentioned that he merely considered that case as a ``touchstone'' for more general conceptions (\emph{ibid.}, p.~1053). And indeed, after he has defined the norm on \(\Cont(I)\), he never (except when proving that the Fredholm operator for continuous kernels is completely continuous in \(\Cont(I)\)) uses anything except the axiomatic definition of a norm (remember that this definition only appeared in print 4 years later!).

In my opinion, F. Riesz's 1918 paper is one of the most beautiful ever written; it is entirely geometric in language and spirit, and so perfectly adapted to its goal that it has never been superseded and that Riesz's proofs can still be transcribed almost \emph{verbatim}. He starts from two almost obvious remarks: 1) in a normed space \(E\), if \(V\) is a closed vector subspace not equal to \(E\), there is a vector \(x \in E\) such that \(\norm{x} = 1\) and \(\norm{x-y} \geq \frac{1}{2}\) for all \(y \in V\); 2) a subset \(S \subset E\) cannot be relatively compact if there is in \(S\) an infinite sequence \((x_n)\) such that \(\norm{x_j-x_k} \geq \frac{1}{2}\) for all pairs of distinct indices. The first consequence is the celebrated theorem characterizing finite dimensional normed spaces as the only \emph{locally compact ones}: one has only to cover the ball \(\norm{x} \leq 1\) with a finite number of balls \(\norm{x-a_j} \leq 1/4\), and then there cannot be any point such that \(\norm{z} = 1\) and \(\norm{z - y} \geq 1/2\) for all points of the (necessarily closed) vector subspace \(V\) generated by the \(a_j\).

F. Riesz then considers a completely continuous linear mapping \(u\) of \(E\) into itself (or, as we now say, a \emph{compact} linear mapping), and studies the endomorphism \(v = 1_E - u\) of \(E\). Using the two remarks above and very simple arguments, he proves in succession the following properties:
\begin{enumerate}[a),noitemsep]
    \item the kernel \(V^{-1}(0)\) has finite dimension;

    \item the image \(v(E)\) is closed in \(E\);

    \item the codimension of \(v(E)\) in \(E\) is finite.
\end{enumerate}
The next step is to consider the \emph{iterates} \(v^k\) of \(v\), the kernel \(N_k\) and the image \(F_k\) of \(v^k\) \footnote{In finite dimensional spaces, this method to obtain the the Jordan normal form of an endomorphism had been developed by E. Weyr \cite{228}.}; the \(N_k\) form an increasing sequence of closed subspaces of finite dimension, the \(F_k\) a decreasing sequence of closed subspaces of finite codimension. F. Riesz shows, by contradiction and using remark 2) above, that there is a smallest integer \(n\) such that \(N_{k+1} = N_k\) for \(k \geq n\); it is then an easy matter to prove that \(F_{k+1} = F_k\) for \(k \geq n\), and that \(E\) is the \emph{topological direct sum} of \(F_n\) and \(N_n\); the restriction of \(v\) to \(F_n\) is a linear \emph{homeomorphism} of \(F_n\) onto itself. In particular, if \(N_1 = V^{-1}(0) = {0}, v\) is a linear homeomorphism of \(E\) onto itself, and its inverse \(w = v^{-1}\) is such that \((1_E - u) w = 1_E\), in other words \(w = 1_E + uw\) has the same form as \(v\), since \(uw\) is compact.

These results enable F. Riesz to treat completely the question of \emph{eigenvalues} of a compact operator. There are at most denumerably many eigenvalues \(\lambda_n \neq 0\) in \(\C\), and each of them is \emph{isolated} in \(\C - \{0\}\); their set is bounded and 0 belongs to its closure if it is infinite. For each \(\lambda_n \neq 0, E\) splits into a topological direct sum of two closed subspaces \(F(\lambda_n)\) and \(N(\lambda_n)\), which are stable by \(h\); \(N(\lambda_n)\) has finite dimension, and there is a smallest integer \(k_n\) such that the restriction of \((u - \lambda_n \cdot 1_E)^{k_n}\) to \(N(\lambda_n)\) is 0; the restriction of \(u - \lambda_n \cdot 1_E\) to \(F(\lambda_n)\) is a linear homeomorphism of that subspace onto itself. If \(E\) is complete, the function \(\xi \mapsto (u - \xi \cdot 1_E)^{-1}\) is meromorphic in \(\C-\{0\}\) (with values in the space \(\mathcal{L}(E)\) of continuous endomorphisms of \(E\)); at the points other than the \(\lambda_n\), that function is holomorphic, and at each \(\lambda_n\) it has a pole of order \(k_n\). Finally, if \(m \neq n\), the subspace \(N(\lambda_n)\) is contained in \(F(\lambda_m)\). However, there is in general no global decomposition of \(E\) into a sum of subspaces \(N(\lambda_n)\), similar to what happens for compact \emph{self-adjoint} operators in Hilbert space \cite{099}. As a matter of fact, there may be \emph{no eigenvalues} at all, as for instance for Volterra operators.

In the study of \(u - \xi \cdot 1_E\), the value \(\xi = 0\) is completely exceptional; \(u(E)\) is not closed in general and may have in finite codimension, and \(u^{-1}(0)\) may have infinite dimension. This explains the intractability of integral equations ``of the first kind'', special cases of equations \(u(x)=y\) for \(u\) compact, which had baffled the early mathematicians working on integral equations.

Although there has been much work done on compact operators of special types, the general theory of compact operators has remained pretty much what it was after the publication of F. Riesz's 1918 paper. Among more recent results, one can mention the fact that when \(u\) is a compact operator in a complete normed space, its transposed operator \(\prescript{t}{}u\) in the dual \(E'\) is also a compact operator \cite{188}. It has also been proved that, even when a compact operator \(u\) in \(E\) has \emph{no eigenvalue}, there are always closed vector subspaces \(V\) of \(E\), different from \(E\) and \(\{0\}\), such that \(u(V) \subset V\) \cite{007}.

\section{The spectral theory of Hilbert}
\label{sec:7.2}

We now return to the most original part of Hilbert's 1906 paper (chap. \ref{ch:5}, \ref{sec:5.2}), in which he discovered the entirely new phenomenon of the ``\emph{continuous spectrum}''. In his ``Theory of heat'', Fourier had considered trigonometric series representing functions of period \(2a\) (chap. \ref{ch:1}, \ref{sec:1.2}, formula \eqref{eq:1.13}) when a \emph{tends to} \(+\infty\). The eigenvalues \(\lambda_n = \left( 1 + \frac{1}{2}^2 \pi^2 / a^2 \right)\) corresponding Sturm--Liouville problem for the equation \(y'' + \lambda y = 0\) with boundary conditions \(y(-a) = y(a) = 0\) divide the interval \([0,+\infty[\) in intervals of length tending to \(0\) with \(1/a\), and this had led Fourier to consider that the ``limiting case'' of the trigonometric expansion of a function of period \(2a\) would be, for any function \(f\) defined on \(\R\), the representation by an \emph{integral}
\begin{equation}
    \label{eq:7.1}
    f(x) = \frac{1}{\pi} \int_{-\infty}^{+\infty} dt \int_0^\infty f(t) \cos{u(x-t)}du
\end{equation}
where the ``eigenvalues'' would now \emph{fill the interval} \([0, +\infty[\) (\cite{067}, vol. \ref{ch:1}, p.~392).

In 1897, Wirtinger \cite{230} developed similar ideas for Hill's equation
\begin{equation}
    \label{eq:7.2}
    y'' + \lambda q(x) y = 0
\end{equation}
where \(q\) is a continuous \emph{periodic} function of period 1. The general theory of these equations was well-known at the time: starting with a fundamental system of solutions \(u_1, u_2\) such that
\[
    \begin{aligned}
        &u_1(0, \lambda) = 1, \qquad u'_1(0, \lambda) =  0 \\
        &u_2(0, \lambda) = 0, \qquad u'_2(0, \lambda) = -1
    \end{aligned}
\]
so that \(u_2 u'_1 - u_1 u'_2\) was the constant function 1, one considers complex solutions \(f\) such that \(f(x+1) = \rho f(x)\) for all \(x \in \R\); the constants \(\rho(\lambda)\) having that property are solutions of the equation
\begin{equation}
    \label{eq:7.3}
    \rho^2 + \Phi(\lambda) \rho + 1 = 0
\end{equation}
with \(\Phi(\lambda) = u'_2(1, \lambda) - u_1(1, \lambda)\). The solutions of period 1 correspond to ``eigenvalues'' \(\lambda\) such that \(\Phi(\lambda) = 2\), the solutions \(f\) such that \(|f(x+1)| = |f(x)|\) to values of \(\lambda\) such that \(-2 \leq \Phi(\lambda) \leq 2\), which in general constitute disjoint intervals \(I_k\) of \(\R\). Wirtinger looks for solutions of period \(n\) (an arbitrary integer), which correspond to values of \(\lambda\) such that \((\rho(\lambda))^n = 1\), and he shows that when \(n\) tends to \(+\infty\), these values of \(\lambda\) tend to ``fill up'' the intervals \(I_k\). The similarity with the optical spectra of molecules leads him to speak of the ``Bandesspectrum'' of equation \eqref{eq:7.2} formed by the union of the intervals \(I_k\), and he thinks there should be an integral formula similar to \eqref{eq:7.1}, without being able to guess what that formula could be.

Although Hilbert does not mention Wirtinger's paper, it is probable that he had read it (it is quoted by several of his pupils), and it may be that the name ``\emph{Spectrum}'' which he used came from it; but it is a far cry from the vague ideas of Wirtinger to the extremely general and precise results of Hilbert. On the other hand, the influence of Stieltjes's big paper of 1894 on continued fractions is explicitly acknowledged by Hilbert: Stieltjes had had to take the limit of a sequence of rational functions of a complex variable
\[
    \frac{P_{2n+1}(z)}{Q_{2n+1}(z)} = \sum_{i=1}^n \frac{M_i}{z + x_i}
\]
where the \(M_i\) and the \(x_i\) are real numbers, and he had shown that the limit could be written as a ``Stieltjes transform''
\[
    F(z) = \int_0^\infty \frac{d\Phi(u)}{z+u}
\]
for a function \(\Phi\) of bounded variation (creating for that purpose the concept of ``Stieltjes integral'') \cite{205}. It is a similar problem which confronts Hilbert when he wants to pass from the classical ``reduction'' of the \(n\)-th ``Abschnitt''
\begin{equation}
    \label{eq:7.4}
    K_n(x, x) = \sum_{p=1}^n \sum_{q=1}^n x_{pq} x_p x_q
\end{equation}
of his ``bounded'' quadratic form \(K(x,x)\), to a ``reduction'' of \(K(x,x)\) itself: the classical theory shows that one has
\begin{equation}
    \label{eq:7.5}
    K_n(x, x) = \frac{\left( L_1^{(n)}(x) \right)^2}{\lambda_1^{(n)}} + \cdots + \frac{\left( L_n^{(n)}(x) \right)^2}{\lambda_n^{(n)}}
\end{equation}
where the \(\lambda_j^{(n)}\) are real numbers such that \(\lambda_1^{(n)} \leq \lambda_2^{(n)} \leq \ldots \leq \lambda_n^{(n)}\), and the \(L_j^{(n)}(x)\) are linear forms in \(x_1, x_2, \ldots, x_n\), such that
\begin{equation}
    \label{eq:7.6}
    \left( L_1^{(n)}(x) \right)^2 + \ldots + \left( L_n^{(n)}(x) \right)^2 = x_1^2 + \ldots + x_n^2.
\end{equation}
To each \(x = (x_p)\) of \(l^2\) such that \(x_p=0\) for \(p > n\), Hilbert associates the piecewise linear convex function of the real variable
\begin{equation}
    \label{eq:7.7}
    u^{(n)}(x; \lambda) = \sum_{p=1}^n \left( L_p^{(n)}(x) \right)^2 \left(\lambda - \lambda_p^{(n)}\right)^+.
\end{equation}
His idea is to take (if possible) the limit of each of these functions, for the points of \(l^2\) having only a finite number of coordinates \(\neq 0\). Using his ``principle of choice'' (chap. \ref{ch:5}, \ref{sec:5.2}) and Cantor's diagonal process, he shows that these limits exist at least for a suitable subsequence of the \(u^{(n)}\). In fact, he only uses these limits for the points \(x^{(pp)}\) having the coordinate of index \(p\) equal to 1, all others to 0, and the points \(x^{(pq)}\) for \(p \neq q\), having the coordinates of indices \(p\) and \(q\) equal to \(\frac{1}{\sqrt{2}}\) and all others to 0; he writes the corresponding limits \(u_{pq}(\lambda)\) for all pairs of integers \((p,q)\). They are convex functions of \(\lambda\), hence the right derivative \(v^+_{pq}(\lambda)\) and the left derivative \(v^-_{pq}(\lambda)\) exist for all \(\lambda\), and are equal except for an at most denumerable set of values of \(\lambda\). Hilbert writes \(\lambda_1, \lambda_2, \ldots, \lambda_r, \ldots\) the sequence of these exceptional values of \(\lambda\) for \emph{all} pairs \((p,q)\), and defines quadratic forms
\[
    v^+(x; \lambda) = \sum_{p,q} v^+_{pq}(\lambda)x_p x_q, \qquad v^-(x; \lambda) = \sum_{p,q} v^-_{pq}(\lambda)x_p x_q
\]
and for each index \(r\),
\[
    E_r(x) = v^+(x; \lambda_r) - v^-(x; \lambda_r).
\]
All these are bounded quadratic forms, and more precisely, their values in \(l^2\) are \(\geq 0\) and \(\leq (x|x)\). For the values of \(\lambda\) distinct from the \(\lambda_r\), he writes
\[
    e(x; \lambda) = \sum_{\lambda_r < \lambda} E_r(x)
\]
and
\[
    \sigma(x; \lambda) = v(x; \lambda) - e(x; \lambda)
\]
where \(v\) is the common value of \(v^+\) and \(v^-\). With these notations, his final result is the ``reduction'' of the quadratic form \(K(x,x)\): for each \(x \in l^2\), the function \(\lambda \mapsto sigma(x; \lambda)\) is a continuous function of bounded variation, and one may write
\begin{equation}
    \label{eq:7.8}
    \begin{dcases}
        (x|x) = \sum_r E_r(x) + \int d\sigma(x; \lambda) \\
        K(x,x) = \sum_r \frac{1}{\lambda_r} E_r(x) + \int \frac{1}{\lambda} d\sigma(x; \lambda).
    \end{dcases}
\end{equation}
He says that the set of the \(\lambda_r\) is the \emph{point spectrum} of \(K\) and the complement of the set where all the \(\sigma(x; \lambda)\) are constant the \emph{continuous spectrum} of \(K\).

Of course, when \(K\) is a ``completely continuous'' form, the continuous spectrum is absent, but Hilbert gives examples for which there is no point spectrum. For instance (see \cite{183}, vol.~II, p.~986--989), if \((\varphi_p)\) is a complete orthonormal system in a compact interval \([a,b]\) of \(\R\) and \(f\) a bounded measurable function in \([a,b]\), one defines a bounded quadratic form \(A(x,x) = \sum\limits_{p,q} a_{pq} x_p x_q\) by the formulas
\begin{equation}
    \label{eq:7.9}
    a_{pq} = \int_a^b f(\mu) \varphi_p(\mu) \varphi_q(\mu) d\mu
\end{equation}
and one has the ``reduction'' formulas
\begin{equation}
    \label{eq:7.10}
    \begin{gathered}
        (x|x) = \int_a^b \left( \sum_p x_p \varphi_p(\mu) \right)^2 d\mu \\
        A(x,x) = \int_a^b f(\mu) \left( \sum_p x_p \varphi_p(\mu) \right)^2 d\mu
    \end{gathered}
\end{equation}
from which it is easy to see that there is no point spectrum (unless \(f\) is constant in an interval), and if f is continuous, the continuous spectrum consists of the whole interval \([m,M]\), where \(m\) and \(M\) are the minimum and maximum value taken by \(f\) in \([a,b]\). If one takes \(a = 0, b = \pi, \varphi_p(\mu) = \sqrt\frac{2}{\pi} \sin{p\mu}, f(\mu) = \cos{\mu}\), one obtains the first example given by Hilbert
\begin{equation}
    \label{eq:7.11}
    A(x,x) = x_1 x_2 + x_2 x_3 + x_3 x_4 + \ldots .
\end{equation}
If one takes \(a = -\pi, b = \pi, \varphi_p(\mu) = \frac{1}{\sqrt{2\pi}} (\sin{p\mu} + \cos{p\mu})\) for \(-\infty < p < +\infty\), and \(f(\mu) = -\pi - \mu\) if \(\mu < 0, f(\mu) = \pi - \mu\) if \(\mu > 0\), one gets \(a_{pq} = \frac{1}{p+q}\) if \(p + q \neq 0, a_{pq} = 0\) if \(q = -p\). Making \(x_p = 0\) for \(p \leq 0\), one gets the second example of Hilbert \(A(x,x) = \sum_{p,q=1} \frac{x_p x_q}{p+q}\), and making \(x_p = 0\) for \(p \leq 0, y_p = 0\) for \(p \geq 0\), one gets still another example of Hilbert \(A(x,y) = \sum_{p,q} \frac{x_p y_q}{p-q}\) (where the summation extends to the pairs \((p,q)\) of integers \(>0\) and distinct); both have for continuous spectrum the interval \([-\pi, +\pi]\)\footnote{One also should mention how Stieltjes' results on continued fractions and on the moment problem were soon recognized as belonging to the Hilbert--von Neumann spectral theory. Jacobi, in 1848, had cunsidered the special quadratic forms \(J(x) = \sum\limits_{k=1}^n a_k x_k^2 - 2 \sum\limits_{k=1}^{n-1} b_{k+1} x_k x_{k+1}\), and he had shown that the eigenvalues of that form are the roots of the denominator of the limited continued fraction (\cite{120}, vol.~VI, p.~318--321)
\begin{equation}
    \tag{F}
    \label{eq:7.F}
    \frac{1}{a_0 - z} - \pfrac{b_1^2}{a_1 - z} -  \pfrac{b_2^2}{a_2 - z} - \cdots - \pfrac{b_n^2}{a_n - z}.
\end{equation}
Already in 1878, Heine \cite[vol.~1, p.~421]{104} had hinted at the possibility that an \emph{unlimited} continued fraction of type \eqref{eq:7.F} would be similarly related to a Jacobi quadratic form \(\sum\limits_{n=0}^\infty a_n x_n^2 - 2 \sum\limits_{n=0}^\infty b_{n+1} x_n x_{n+1}\) in an infinite system of variables. What Stieltjes had done was to study directly unlimited continued fractions of type \eqref{eq:7.F}, representing them as ``Stieltjes transforms'' and being led to his ``problem of moments'' by the problem of determining the Stieltjes measure corresponding to a given Stieltjes transform; but he had not considered the relation between the continued fraction and quadratic forms. This was done by Toeplitz in 1910 \cite{214} for the case of Jacobi bounded quadratic forms, and later extended to the general case; it turned out that these forms exactly corresponded to spectra of \emph{multiplicity} 1 \cite{207}.}.

In his 1913 book (\cite{183}, vol.~II, p.~956--989 and \cite{184}), F. Riesz gave an exposition of Hilbert's results based on an entirely different method, and which was to remain standard until around 1950. As we have already pointed out, he replaces the bilinear forms of Hilbert by the much more natural continuous \emph{endomorphisms} of \(E = l^2_{\R}\); to such an endomorphism \(A\) is as sociated the ``bounded'' bilinear form \((x,y) \mapsto (A \cdot x|y)\) and conversely each such form can be uniquely written in that way. F. Riesz's central idea is to define, for such endomorphisms \(A\), ``functions'' \(f(A)\) which would again be continuous endomorphisms of \(E\), for suitable functions \(f\) of a real variable, and to use such ``functions'' to write for a \emph{symmetric} endomorphism \(A\) (i.e. such that \((A \cdot x|y) = (x|A \cdot y)\) for all \(x,y\) in \(E\)) a canonical ``spectral decomposition'' corresponding to formulas \eqref{eq:7.8} of Hilbert.

To develop these ideas, F. Riesz begins by some general results on the algebra \(\mathcal{L}(E)\) of all continuous endomorphisms of \(E\). For the norm \(\norm{A} = \sup\limits_{\norm{x}=1}{\norm{A \cdot x}}\), it is a Banach space, has a unit (the identity mapping \(1_E\)) and is such that \(\norm{AB} \leq \norm{A} \cdot \norm{B}\). However, F. Riesz, for his purpose, is led to use, not the notion of (``uniform'') convergence derived from the norm of \(\mathcal{L}(E)\), but the notion of \emph{strong convergence}: he says a sequence \((A_n)\) converges strongly to \(A\) if, for every \(x \in E\), the sequence \((\norm{A_n \cdot x - A \cdot x})\) tends to 0.

F. Riesz then considers the subspace \(\mathcal{H}(E)\) of all symmetric operators; there is in \(\mathcal{H}(E)\) an \emph{order relation}, \(A \leq B\) meaning that \((A \cdot x|x) < (B \cdot x|x)\) for all \(x \in E\). Suppose \(a \cdot 1_E \leq A \leq b \cdot 1_E\); then, for every polynomial \(P(\xi)\) with real coefficients such that \(P(\xi) \geq 0\) in the interval \([a,b]\), one has \(P(A) \geq 0\). Indeed, one may assume \(a = 0, b = 1\), and then \(P\) is a sum of polynomials of one of the types \(Q(\xi)^2, \xi(Q(\xi))^2, (1-\xi)(Q(xi))^2\) or \(\xi(1-\xi)(Q(\xi))^2\), and it is enough to prove that \(A(1_E - A)\) is \(\geq 0\); but from the Cauchy--Schwarz inequality for the positive quadratic form \((A \cdot x|x)\) it follows that \(\norm{A \cdot x}^4 \leq (A \cdot x|x)(A^2 \cdot x|A \cdot x) < \norm{x}^2 \norm{A \cdot x}^2\); this first implies that \(\norm{A} < 1\), and then \(\norm{A \cdot x}^4 \leq (A \cdot x|x)\norm{A \cdot x}^2\) and finally \(A^2 \cdot x|x) \leq (A \cdot x|x)\). From this it follows at once that if a polynomial \(P(\xi)\) with real coefficients is such that \(m \leq f(\xi) \leq M\) in \([a,b]\), then one has \(m \cdot 1_E \leq P(A) \leq M \cdot 1_E\), and \(\norm{P(A)} \leq \sup(|m|, |M|)\).\footnote{This argument is not the one used by Riesz, who deduces the result by a passage to the limit from the known result for the ``Abschnitte'' \(A_n\) of \(A\).}

These results first imply that if a sequence \((P_n)\) of polynomials converges uniformly to a continuous function \(f\) in \([a,b]\), then the sequence \((P_n(A))\) is a Cauchy sequence in the Banach space \(\mathcal{L}(E)\), and its limit only depends on \(f\), hence can be written \(f(A)\); furthermore the mapping \(f \mapsto f(A)\) is a homomorphism of the algebra \(\Cont([a,b])\) into \(\mathcal{L}(E)\), with values in \(\mathcal{H}(E)\), which justifies the notation; in addition, if \(m \leq f(\xi) \leq M\) in \([a,b]\), one has again \(m \cdot 1_E \leq f(A) \leq M \cdot 1_E\).

But F. Riesz goes further. If \((f_n)\) is an increasing sequence of continuous functions in \([a,b]\), uniformly bounded, then for any \(x \in E\), the sequence of the \((f_n(A) \cdot x|x)\) is increasing and bounded, hence has a limit, from which it follows by linearity that the sequence of the \((f_n(A) \cdot x|y)\) converges for any pair of elements \(x, y\) in \(E\); the Hellinger--Toeplitz theorem (chap. \ref{ch:6}, \ref{sec:6.4}) shows that the limit can be written \((B \cdot x|y)\) where \(B \in \mathcal{H}(E)\); if \(g\) is the (simple) limit of the sequence \((f_n)\), one writes again \(B = g(A)\), and this enables one to define \(g(A)\) for any bounded (upper or lower) semi-continuous function in \([a,b]\) or any linear combination of such functions, which again form an algebra and for which \(g \mapsto g(A)\) is a homomorphism.

F. Riesz then uses these results to obtain the spectral decomposition in the following way; if \(e_\xi\) is the function defined in \(\R\) and such that \(e_\xi(\mu) = 1\) for \(\mu < \xi, e_\xi(\mu) = 0\) for \(\mu \geq \xi, e_\xi(A) = A_\xi\) is defined since \(e_\xi\) is bounded and lower semi-continuous. For any pair of vectors \(x, y\) the function \(\xi \mapsto (A_\xi \cdot x|y)\) is then a function of bounded variation, and for any continuous function \(f\), one has
\begin{equation}
    \label{eq:7.12}
    (f(A) \cdot x | y) = \int_{-\infty}^{+\infty} f(\xi)d(A_\xi \cdot x | y)
\end{equation}
a formula which one also writes
\begin{equation}
    \label{eq:7.13}
    f(A) = \int_{-\infty}^{+\infty} f(\xi)dA_\xi
\end{equation}
and which is justified by the fact that for any \(\varepsilon > 0\), it is possible to divide the interval \([a,b]\) by points \(\xi_k\) in such a way that
\[
    -\varepsilon \cdot 1_E \leq f(A) - \sum_k f(\mu_k)\left( A_{\xi_{k+1}} - A_{\xi_k} \right) \leq \varepsilon \cdot 1_E \qquad (\text{with } \xi_k \leq \mu_k \leq \xi_{k+1})
\]
The \emph{spectrum}\footnote{To pass from the Hilbert notion of ``spectrum'' to the one used by F. Riesz, one must replace the parameter \(\lambda\) of Hilbert by \(1/\xi\).} of \(A\), contained in the interval \([a,b]\), is the complement of the set of points having a neighborhood where \(\xi \mapsto A_\xi\) is constant.

The operator \(A_\xi\) is the \emph{orthogonal projector} of \(E\) onto a closed subspace \(E_\xi\), which is stable under \(A\) and in which \((A \cdot x|x) < \xi(x|x)\) for all \(x \neq 0\); \(1_E - A_\xi\) is the orthogonal projector of \(E\) onto the subspace \(E_\xi^\perp\) orthogonal to \(E_\xi\), and in which \((A_\xi \cdot x|x) \geq \xi(x|x)\). The \emph{point spectrum} of \(A\) is the denumerable set of values of \(\xi\) where at least one of the functions \(\xi \mapsto (A_\xi \cdot x|y)\) is discontinuous; it consists of all the \emph{eigenvalues} of \(A\), but the subspace \(N_\xi\) formed by the corresponding eigenvectors may have infinite dimension. The subspaces \(E_\xi\) are such that \(E_\xi \subset E_\eta\) for \(\xi < \eta, E_\xi = \{0\}\) for \(\xi < a, E_\xi = E\) for \(\xi > b\); the intersection of the subspaces \(E_\eta\) for \(\xi < \eta\) is reduced to \(E_\xi\) if \(\xi\) is not in the point spectrum and is the direct sum of the (orthogonal) subspaces \(E_\xi\) and \(N_\xi\) if \(\xi\) is in the point spectrum.

If \((A \cdot x|y) = \sum\limits_{i,k} a_{ik} x_i y_k\) is the ``bounded'' bilinear form corresponding to the operator \(A\), the eigenvectors \(x = (x_k)\) corresponding to an eigenvalue \(\lambda\) have coordinates which are solutions of the system of linear equations
\begin{equation}
    \label{eq:7.14}
    \lambda x_i - \sum_k a_{ik} x_k = 0 \qquad (i=1,2,\ldots).
\end{equation}
However, Hilbert's theory left unanswered the question of determining ``objects'' which would replace the eigenvectors when a number \(\lambda\) in the spectrum was not in the point spectrum. In some cases this question had a curious answer; for instance, for the form \eqref{eq:7.11}, where there is no point spectrum and the spectrum is the interval \([-1,1]\), for each value \(\lambda = \cos{t}\) in that interval, the system \eqref{eq:7.14} does have a solution, namely \(x_k(t) = \sin{kt}\) for \(k = 1, 2, \ldots\), as it is easily verified; but for such a sequence, the series \(\sum\limits_k x_k(t)^2\) is \emph{not convergent}. The existence of such ``generalized eigenvectors'' was only incorporated in a general theorem much later (see chap. \ref{ch:9}, \ref{sec:9.2}); but in the case of the form \eqref{eq:7.11}, one could observe that if one wrote \(\rho_k(\xi) = \int_0^\xi x_k(t)dt\), then the vector \((\rho_k(\xi))_{k=1,2,\ldots}\) this time belonged to \(E\) for all \(\xi \in \R\), and it followed from the relations \eqref{eq:7.14} satisfied by the \(x_k(t)\) that one could write, for any interval \([\lambda_1, \lambda_2]\) of \(\R\)
\begin{equation}
    \label{eq:7.15}
    \int_{\lambda_1}^{\lambda_2} \xi d\rho_i(\xi) - \sum_k a_{ik}(\rho_k(\lambda_2) - \rho_k(\lambda_1)) = 0 \qquad (i = 1, 2, \ldots).
\end{equation}
This led Hellinger \cite{105} to study systematically the sequences of functions \((\rho_k(\xi))_{k=1,2,\ldots}\) of bounded variation which satisfied \eqref{eq:7.15} for all intervals and for which \(\sum\limits_k \rho_k(\xi)^2\) is finite for all \(\xi\); he called the \(d\rho_k\) ``eigendifferentials'' of the operator \(A\). This study allowed him to attack a problem which naturally arose from Hilbert's spectral theory, by analogy with the finite dimensional case: can one give necessary and sufficient conditions for two operators \(A,B\) with symmetric matrices to be ``similar'', i.e. such that there is an orthogonal transformation \(U\) of \(E\) onto itself such that \(B = UAU^{-1}\)? In the finite dimensional case, the condition is that the eigenvalues of \(A\) and \(B\) be the same, with the same multiplicity for each; the combined efforts of Hellinger and H. Hahn \cite{095} succeeded in obtaining necessary and sufficient conditions for operators in \(l^2_{\R}\), expressed in terms of special systems of ``eigendifferentials''. We shall not give here the detail of these complicated conditions, which we shall formulate in a much simpler way using the Gelfand theory of commutative Banach algebras (\ref{sec:7.5}).

One of the byproducts of F. Riesz's method is that it enabled him to give a \emph{direct} definition of the whole spectrum of \(A\), without any reference to the decomposition \eqref{eq:7.12}: a point \(\lambda \in \R\) is in the spectrum if and only if the operator \(\lambda \cdot 1_E - A\) \emph{does not have a continuous inverse}. Finally, he remarked that his method could (just as Hilbert's method) be extended to \emph{self-adjoint} bounded operators \(A\) in \emph{complex} Hilbert space \(l^2_{\C}\), i.e. those which satisfy the same condition \((A \cdot x|y) = (x|A \cdot y)\) for the hermitian scalar product in that space; the mapping \(\zeta \mapsto (\zeta \cdot 1_E - A)^{-1}\) is then \emph{holomorphic} outside of the spectrum of \(A\). After 1913, almost all papers on spectral theory in Hilbert space dealt exclusively with complex Hilbert space.

\section{The work of Weyl and Carleman}
\label{sec:7.3}

Hilbert's method associating to an integral equation with symmetric kernel \(K(s,t)\) a ``bounded'' bilinear form \(K(x,y)\) (chap. \ref{ch:5}, \ref{sec:5.2}) worked even if \(K(s,t)^2\) was not integrable in \([a,b] \times [a,b]\), but the corresponding bilinear form might not be ``completely continuous''; already in his lectures in 1906 Hilbert had given the example \(K(s,t) = (s+t)^{-1}\) for the interval \([0,1]\) (\cite{227}, vol.~I, p.~83). He also had observed in his Seminar that the Fredholm theorems might fail when the interval \([a,b]\) was unbounded, and had given as example an interpretation of the Fourier inversion formula \eqref{eq:7.1} (\emph{ibid}., p.~2): for \(K(s,t) = \cos{st}\) in the interval \([0,+\infty[, \sqrt\frac{2}{\pi}\) and \(-\sqrt\frac{2}{\pi}\) are the only eigenvalues, but each of them has \emph{infinite multiplicity}\footnote{This is apparently the first appearance of a relation between Fourier transforms and Functional Analysis (see \ref{sec:7.6}).}. He therefore encouraged his most gifted student, Hermann Weyl, to elucidate such ``singular'' integral equations, and in particular to determine conditions on the kernel \(K(s,t)\) implying that the bilinear form \(K(x,y)\) would be ``bounded'' and therefore amenable to his spectral theory. This was the theme of Weyl's dissertation; in it and in a subsequent paper (\emph{ibid}., p.~2--86 and 102--153) he gave the following sufficient conditions for \([a,b[ = [0,+\infty[\): 1° for each \(s \geq 0\), the integral \(\int_0^\infty (K(s,t))^2 dt\) is finite; 2° there is a constant \(M > 0\) such that the inequality \(\left| \int_0^\infty \int_0^\infty K(s,t)u(s)v(t)dsdt \right| \leq M\) for all pairs of continuous functions \(u, v\) such that \(\int_0^\infty (u(s))^2 ds \leq 1\) and \(\int_0^\infty (v(s))^2 ds \leq 1\).

Another direction of research derived from an interpretation of the Sturm--Liouville theory (chap. \ref{ch:1}, \ref{sec:1.3}) in terms of integral equations. Consider a second order differential equation
\begin{equation}
    \label{eq:7.16}
    y'' - q(x)y + \lambda y = f(x)
\end{equation}
where \(q\) and \(f\) are continuous functions in a compact interval \([a,b], q\) having real values, \(f\) real or complex values and \(\lambda\) is a complex parameter; in addition we have two boundary conditions
\begin{equation}
    \label{eq:7.17}
    y(a)\cos\alpha - y'(a)\sin\alpha = 0, \qquad y(b)\cos\beta - y'(b)\sin\beta = 0
\end{equation}
where \(\alpha\) and \(\beta\) are two positive constants. An elementary argument shows that for \(f = 0\), the homogeneous equation \eqref{eq:7.16} with boundary conditions \eqref{eq:7.17} has no solution if \(\lambda \leq -r\), where \(r\) is a number \(>0\) depending only on \(q\). Replacing \(q(x)\) by \(q(x) + r\) and \(\lambda\) by \(\lambda + r\), one may assume that for \(\lambda < 0\), the homogeneous equation \(y'' - q(x)y + \lambda y = 0\) has no solution satisfying the conditions (17).

Now consider first equation \eqref{eq:7.16} for \(\lambda = 0\); there are two solutions \(u_1, u_2\) of the equation \(y''- q(x)y = 0\) such that \(u_1(a)\cos\alpha - u'_1(a)\sin\alpha = 0, u_2(b)\cos\beta - u'_2(b)\sin\beta = 0, u_1\) and \(u_2\) being \emph{linearly independent}. For each \(t \in [a,b]\), define the function
\begin{equation}
    \label{eq:7.18}
    \begin{dcases}
        K(t,x) = -u_2(t)u_1(x)/d & \text{for } a \leq x \leq t \\
        K(t,x) = -u_1(t)u_2(x)/d & \text{for } t \leq x \leq q
    \end{dcases}
\end{equation}
where \(d\) is the constant \(u_1(x)u'_2(x) - u_2(x)u'_1(x)\). The function \(x \mapsto K(t,x)\) is then continuous in \([a,b]\); in each of the semi-open intervals \(a \leq x \leq t, t \leq x \leq b\) (for \(a < t < b\)) it satisfies the equation \(y'' - q(x)y = 0\), and in addition it satisfies conditions \eqref{eq:7.17}; finally, the function \(x \mapsto \pder{}{x} K(x,t)\) has at the point \(x = t\) a discontinuity such hat \(\pder{}{x} K(t,t+) - \pder{}{x} K(t,t-) = -1\). A routine calculation then shows that in order that a function \(y\) be a solution of \(y'- q(x)y = f(x)\) satisfying conditions \eqref{eq:7.17}, it is necessary and sufficient that
\begin{equation}
    \label{eq:7.19}
    y(x) = -\int_a^b K(t,x)f(t)dt
\end{equation}
and therefore the solutions of \eqref{eq:7.16}, satisfying \eqref{eq:7.17}, are exactly the solutions of the integral equation
\begin{equation}
    \label{eq:7.20}
    y(x) = -\lambda \int_a^b K(t,x)y(t)dt = g(x)
\end{equation}
where
\begin{equation}
    \label{eq:7.21}
    g(x) = -\int_a^b K(t,x)f(t)dt.
\end{equation}
Clearly this method was patterned after Schwarz's method for solving the equation of vibrating membranes with the help of the Green function for the Laplacian (chap. \ref{ch:3}, \ref{sec:3.1}), and the function \(K\) was therefore called the \emph{Green function} for the operator \(L(y) = y'' - q(x)y\) and the boundary conditions \eqref{eq:7.17} \cite{035}. As obviously \(K(x,t) = K(t,x)\), the Sturm--Liouville problem was thus reduced to a special case of the Fredholm--Hilbert theory of integral equations with symmetric kernels.

In his second paper (1904) on integral equations, Hilbert developed that method and expanded it to other boundary conditions than \eqref{eq:7.17}. He also was aware that many of the ``special functions'' which had been introduced in Analysis since the XVIII\supth{} century (hypergeometric functions, Bessel functions, Legendre polynomials, Hermite functions, etc.) satisfied equations of type \eqref{eq:7.16} but with \emph{less restrictive} conditions: the interval \([a,b]\) would be replaced by an unbounded interval and the function \(q\) might have \emph{singular points} at the extremities of the interval; in such a case, Hilbert proposed that the corresponding boundary condition should be replaced by the condition that \(y\) remain bounded at such an extremity, or tends to infinity not faster than some given singular function. He showed then, on various examples, that one could even in such ``singular'' cases, define a ``Green function'' \(K\), with the symmetric property \(K(t,x) = K(x,t)\), and the same discontinuity for the partial derivatives for \(x = t\); for instance, if \(L(y) = y'' + y\), one has, for the interval \(]-\infty, +\infty[, K(t,x) = -\frac{1}{2} \sin{|x-t|}\). At that time, he had not yet developed his theory of ``bounded'' bilinear forms, so he limited himself to cases in which the Green function K was a kernel to which the Fredholm theory was applicable (\cite{112}, p.~39--58). But of course he was aware that in examples such as the one above, one would fall on ``singular'' integral equations, and one of his students, E. Hilb, wrote his ``Habilitationschrift'' in 1908 on the application of Hilbert's theory of ``bounded'' forms to two special cases of ``singular'' Sturm--Liouville problems \cite{110}.

Then, in 1909--1910, H. Weyl discovered that he could apply the results of his dissertation to handle the \emph{most general} such problems for second order operators of the type
\begin{equation}
    \label{eq:7.22}
    L(y) = \frac{d}{dx} \left(p(x) \frac{dy}{dx}\right) - q(x)y
\end{equation}
where \(p\) and \(q\) are real continuous functions in an interval \(I \subset \R\) (bounded or not) subject to the only restriction that \(p(x) > 0\) in \(I\). In his ``Habilitationschrift'' (\cite{227}, vol.~I, p.~248--297) he considers the case \(I = [0,+\infty[\); his very original method consists in studying the equation
\begin{equation}
    \label{eq:7.23}
    L(y) + \lambda y = 0
\end{equation}
for \emph{non real} values of \(\lambda\), and in fact, he sees that it is enough to consider the case \(\lambda = i\). Let \(u_1, u_2\) be the two solutions of \(L(y) + iy = 0\) in \([0,+\infty[\) satisfying the initial conditions
\[
    u_1(0) = 1, \quad p(0)u'_1(0) = 0,
    \qquad
    u_2(0) = 0, \quad p(0)u'_2(0) = 1
\]
Let \(\alpha\) be any number \(\geq 0\), and consider the two solutions of \(L(y) + iy = 0\)
\[
    v_\alpha = -\sin{\alpha \cdot u_1} + \cos{\alpha + \cdot u_2},
    \qquad
    w_\alpha = \cos{\alpha \cdot u_1} + \sin{\alpha + \cdot u_2},
\]
so that \(v_\alpha\) is a solution satisfying the boundary condition
\begin{equation}
    \label{eq:7.24}
    \cos{\alpha \cdot y(0)} + \sin{\alpha \cdot p(0)y'(0)} = 0.
\end{equation}
Now take any number \(a > 0\), a number \(\beta \geq 0\), and determine the complex number \(\mu\) by the condition that the solution \(u_\mu = v_\alpha + \mu w_\alpha\) satisfies the boundary condition
\begin{equation}
    \label{eq:7.25}
    \cos{\beta \cdot y(a)} + \sin{\beta \cdot p(a)y'(a)} = 0.
\end{equation}
Weyl shows that the uniquely determined number \(\mu\) describes a \emph{circle} \(\Gamma_a\) in the upper half plane \(\mathcal{I} z > 0\) when \(\beta\) varies from 0 to \(2\pi\).

From the two solutions \(v_\alpha\) and \(u_\mu\) one can form a \emph{Green function} \(K^\mu_\alpha(t,x)\) in the interval \([0,a]\) by the same formulas as \eqref{eq:7.18}, except that now has complex values. Now let \(a\) \emph{tend to} \(+\infty\); Weyl shows that the circles \(\Gamma_a\) form a nested family of decreasing radius, hence have a \emph{limit} \(\Gamma_\infty\) which may be either a circle of radius \(>0\), or a single point. In any case, if one lets the points \(\mu \in \Gamma_a\) tend to a limit \(\mu_0 \in \Gamma_\infty, u_\mu\) tends to a solution \(u_{\mu_0}\) and \(K^\mu_a\) tends to a kernel \(K^{\mu_0}(t,x)\) which always satisfies the two conditions of Weyl's dissertation; this enables one to apply to the corresponding singular integral equation Hilbert's theory of ``bounded'' bilinear forms. Just as for the usual Sturm--Liouville problem the function
\begin{equation}
    \label{eq:7.26}
    y(x) = -\int_0^\infty K^{\mu_0}(t,x)f(t)dt
\end{equation}
is then a solution of
\begin{equation}
    \label{eq:7.27}
    L(y) + iy = f(x)
\end{equation}
satisfying the boundary condition \eqref{eq:7.24} at the extremity 0, and the condition ``at infinity''
\begin{equation}
    \label{eq:7.28}
    \lim_{t \to +\infty} p(t)\left( y(t)u'_{\mu_0}(t) - u_{\mu_0}(t)y'(t) \right) = 0.
\end{equation}

Furthermore, the solution \(u_{\mu_0}\) always \emph{belongs to} \(L^2(0,\infty)\). Weyl then shows that the problem splits in two cases:

I) The ``\emph{limit circle}'' case, \(\Gamma_\infty\) being a circle of radius \(>0\), Then \emph{all} solutions of \(L(y) + iy = 0\) belong to \(L^2(0,\infty)\), and for all \(\mu_0 \in \Gamma_\infty, K^{\mu_0}\) is a Hilbert--Schmidt kernel; the conclusions of the Sturm--Liouville theory are then again valid for equation \eqref{eq:7.23} with the boundary condition \eqref{eq:7.24} at extremity 0, and the condition that the real part of \eqref{eq:7.28} vanishes at extremity \(+\infty\).

II) The ``\emph{limit point}'' case, when \(\Gamma_\infty\) is reduced to a single point \(\mu_0\). Then \(u_{\mu_0}\) is (up to a constant factor) the \emph{only} solution of \(L(y) + iy = 0\) belonging to \(L^2(0,\infty)\), and Weyl's main objective is to set up integral formulas which should be subgituted to the ``Fourier expansions'' of the classical Sturm--Liouville theory, as had been expected by Wirtinger, and obtained by Hilb in special cases; this he is able to do by applying the results of his dissertation to the ``singular'' integral operator having as kernel the imaginary part of \(K^{\mu_0}\), and extensively using Hellinger's ``eigendifferentials''. He also defines the \emph{spectrum} of the differential operator \(L\) as the complement of the set of parameters \(\lambda \in \R\) for which the differential equation \(L(y) + \lambda y = g(x)\) has a solution belonging to \(L^2(0,\infty)\) and satisfying \eqref{eq:7.24}, for every continuous function \(g\) belonging to \(L^2(0,\infty)\); he studies the structure of that spectrum under various assumptions on the functions \(p\) and \(g\), and in particular he gives an example where that spectrum is the \emph{whole real line}. Finally, in a subsequent paper (\cite{227}, vol.~I, p.~222--247), he shows how his theory may be extended when the interval \(I\) is the whole line \(\R\), and the equation belongs to the ``limit point'' type at both extremities.

Viewed from the vantage point of the later von Neumann theory (see \ref{sec:7.4}) these remarkable results of Weyl constitute the first study of an unbounded hermitian operator in Hilbert space, with non zero ``defects'', and of its self-adjoint extensions; the singular integral operator defined by the kernel \(K\) with complex values is probably the first example of a normal operator in Hilbert space which is not self-adjoint\footnote{The concept of a normal square matrix \(A\) with complex elements had been defined in 1877 by Frobenius by the condition that \(AA^* = A^*A\); he had proved that this condition was equivalent to the existence of a unitary matrix \(U\) such that \(UAU^{-1}\) is a \emph{diagonal} matrix \cite[vol.~I, p.~391]{078}.}.

Such phenomena became even more apparent in the work of T. Carleman on singular integral equations (\cite{037} and \cite[p.~313--342]{038}) beginning in 1920. He starts from a kernel which only satisfies the \emph{first} of Weyl's assumptions, namely \(\int_a^b |K(s,t)|^2dt\) is almost everywhere finite (they are now called \emph{Carleman kernels}, and the corresponding operators, which to \(f\) associate \(s \mapsto \int_a^b K(s,t)f(t)dt\), \emph{Carleman operators}). He treats the theory of these operators (for hermitian kernels) by a method similar to H. Weyl's: namely, he considers an increasing sequence \((A_n)\) of measurable subsets of \([a,b]\), the union of which is equal to \([a,b]\) up to a null set, and which are such that he kernel \(K_n(s,t)\), equal to \(K(s,t)\) for \((s,t) \in A \times A\) and to 0 outside, has a finite integral \(\int_a^b \int_a^b |K_n(s,t)|^2 dsdt\). If one now considers the integral equation
\begin{equation}
    \label{eq:7.29}
    \varphi(s) - \lambda \int_a^b K(s,t)\varphi(t)dt = f(s)
\end{equation}
for \emph{non real} \(\lambda, f\) being in \(L^2([a,b])\), one approximates it by the sequence of ordinary Hilbert--Schmidt integral equations
\begin{equation}
    \label{eq:7.30}
    \varphi_s(s) - \lambda \int_a^b K_n(s,t)\varphi_n(t)dt = f(s)
\end{equation}
which (due to the choice of \(\lambda\)) has a unique solution \(\varphi_n \in L^2([a,b])\). Carleman's original procedure is to integrate \eqref{eq:7.30} after multiplication by \(\overline{\varphi_n(s)}\), which gives him the identity between the imaginary parts
\begin{equation}
    \label{eq:7.31}
    \left(\frac{1}{\lambda} - \frac{1}{\overline\lambda}\right) \int_a^b |\varphi_n(s)|^2 ds = \frac{1}{\lambda} \int_a^b \overline{\varphi_n(s)}f(s)ds - \frac{1}{\lambda} \int_a^b \varphi_n(s) \overline{f(s)}ds
\end{equation}
independent of the kernel, and from which he deduces, by the Cauchy--Schwarz inequality,
\begin{equation}
    \label{eq:7.32}
    \int_a^b |\varphi_n(s)|^2 ds \leq \frac{4|\lambda|^2}{|\lambda - \overline\lambda|^2} \int_a^b |f(s)|^2 ds
\end{equation}
Applying the usual ``principle of choice'' and density arguments as Hilbert and F. Riesz had done, Carleman is able (by letting \(n\) tend to \(+\infty\)) to define a linear mapping \(f \mapsto T_\lambda \cdot f\) in \(L^2\) such that \(\varphi = T_\lambda \cdot f\) is a solution of \eqref{eq:7.29} for each \(f \in L^2\), and a passage to the limit in \eqref{eq:7.32} shows that \(T_\lambda\) is continuous. He even goes as far as writing an equation equivalent to
\begin{equation}
    \label{eq:7.33}
    T_\lambda \cdot f = \frac{\overline\lambda}{\lambda - \overline\lambda} (f + U_\lambda \cdot f)
\end{equation}
and showing that \(\norm{U_\lambda \cdot f} \leq \norm{f}\).

\mytodo{there must be a typo in the denominator of his equation (33).}

He then realized that the solution of \eqref{eq:7.29} for non real \(\lambda\) might be \emph{non unique}, and he gave examples of kernels where this phenomenon happens, which he called kernels of class II; the other ones he called kernels of class I, and he showed that they may be more general than the continuous operators of F. Riesz or those considered by H. Weyl.

For any functions \(\varphi, \psi\) in \(L^2\), the functions
\begin{equation}
    \label{eq:7.34}
    S \cdot \varphi: s \mapsto \int_a^b K(s,t)\varphi(t)dt,
    \qquad
    S' \cdot \psi: s \mapsto \int_a^b K(s,t)\psi(t)dt
\end{equation}
are always defined for a hermitian Carleman kernel \(K\), but the set \(D\) (resp. \(D'\)) of functions of \(L^2\) such that \(S \cdot \psi \in L^2\) (resp. \(S' \cdot \psi in L^2\)) is in general a proper subspace of \(L^2\); \(D'\) consists of the complex conjugates of the functions of \(D\). Carleman showed that a necessary and sufficient condition for \(K\) to be ``of class I'' was that the relation
\begin{equation}
    \label{eq:7.35}
    \int_a^b f(s)(S' \cdot \overline{g})(s)ds = \int_a^b \overline{g(s)}(S \cdot f)ds
\end{equation}
should hold for all \(f,g\) in \(D\).

He next proceeded to let also \(n\) tend to \(+\infty\) in the Hilbert formula for Hilbert--Schmidt kernels (chap. \ref{ch:5}, \ref{sec:5.2}, formula \eqref{eq:5.24}) for the kernels \(K_n\), and their conjugates \(\overline{K}_n\) and obtained for the operators \(T_\lambda\) and \(T'_\lambda\) corresponding to \(K\) and \(\overline{K}\) formulas similar to those obtained by Hilbert in his theory of ``bounded'' quadratic forms. For kernels ``of class II'', the study of these formulas led Carleman to single out the case in which the operator \(U_\lambda\), in \eqref{eq:7.33} and the corresponding operator \(U'_\lambda\) for \(T'_\lambda\); are both \emph{unitary}; he shows that this property is independent of the choice of the non real number \(\lambda\) in one of the half planes \(\mathcal{I}\lambda > 0, \mathcal{I}\lambda < 0\), and that it implies that the dimensions of the spaces of solutions of \(\varphi - \lambda S \cdot \varphi = 0\) and \(\psi - \lambda S' \cdot \psi = 0\) in \(L^2\) are the \emph{same}; finally he proved that in this case there are \emph{infinitely} many operators \(T_\lambda\) and \(T'_\lambda\) having the above property (what he calls ``maximal solutions'') for the same kernel \(K\).

All these results were quite surprising, in particular the existence of solutions \(\varphi \neq 0\) for the equation \(\varphi - \lambda S \cdot \varphi = 0\) in \(L^2\) for non real \(\lambda\), which seemed to contradict the classical argument (going back to Poisson (chap. \ref{ch:1}, \ref{sec:1.3}), and even to Lagrange (\cite{135}, vol.~XII, p.~239) in the finite dimensional case) which, from the \emph{reality} of the number \((S \cdot \varphi | varphi)\) for all \(\varphi\), concluded to the impossibility of a non real number \(\lambda\) satisfying \(\varphi = \lambda S \cdot \varphi\) for \(\varphi \neq 0\), since it implied \((\varphi | \varphi) = \lambda(S \cdot \varphi | \varphi)\). We shall see in the next section how this apparent contradiction was resolved in the von Neumann theory, which put the pioneering results of Carleman in their proper context.

\section{The spectral theory of von Neumann}
\label{sec:7.4}

In the fall of 1926, the young J. von Neumann (1903--1957) arrived at Göttingen, to take up his duties as Hilbert's assistant. These were the hectic years during which quantum mechanics was developing at breakneck speed, with a new idea popping up every few weeks from all over the horizon \cite{121}. The theoretical physicists who were developing the new theory were groping for adequate mathematical tools, trying in succession infinite matrices without any consideration of convergence\footnote{As late as 1924, most physicists did not even know that a \emph{finite} matrix was!}, differential operators, ``continuous'' matrices (whatever that might mean) etc. It finally dawned upon them that their ``observables'' had properties which made them look like hermitian operators in Hilbert space, and that, by an extraordinary coincidence, the ``spectrum'' of Hilbert (a name which he had apparently chosen from a superficial analogy) was to be the central conception in the explanation of the ``spectra'' of atoms. It was therefore natural that they should enlist Hilbert's help in trying to put some mathematical sense in their computations. With the assistance of Nordheim and von Neumann, he first tried integral operators in \(L^2\), but that needed the use of the Dirac ``\(\delta\)-function'', a concept which for the mathematics of that time was self-contradictory (cf. chap. \ref{ch:8}, \ref{sec:8.3}); von Neumann therefore resolved to try another approach.

Ever since the discovery of the Fischer--Riesz theorem (chap. \ref{ch:5}, \ref{sec:5.3}) the isomorphism of the space of sequences \(l^2_{\C}\) and of the \(L^2_{\C}\) spaces of classes of quadratically integrable functions in some subset \(\Omega\) of an \(\R^n\) had been familiar to analysts, but by ``Hilbert space'' one always understood one of these ``concrete'' spaces, on which the ``operators'' would there fore be, either ``matrices'' or ``integral operators''s of some kind. Von Neumann was the first to conceive of an ``abstract'' Hilbert space, defined axiomatically as a complex vector space with a hermitian scalar product, separable and complete for the corresponding norm, so that the usual ``concrete'' Hilbert spaces would only be ``incarnations'' so to speak of that ``abstract'' space. Obvious as it now seems to us, this was a momentous step and opened the way to a complete understanding of spectral theory of normal and hermitian operators in Hilbert space, which von Neumann proceeded to develop in 3 fundamental papers published between 1929 and 1932, and which (with the exception of the description of the spectrum, see \ref{sec:7.5}) are still today, in substance, the definitive account of the subject (\cite{221}, vol.~II, p.~1--85, 86--143 and 242--258)\footnote{During the same period, M.H. Stone, independently of von Neumann, obtained the same results concerning self-adjoint (unbounded) operators \cite{206}, and later gave a didactic exposition of the whole theory and of its applications at that time, much clearer than von Neumann's papers, and which remained the reference book on the subject for many years \cite{207}.}.

Abandoning any ``concrete'' presentation of Hilbert space, von Neumann was compelled to work \emph{intrinsically}, using only notions which could directly be defined from the concepts enumerated in the axioms, to the exclusion of anything else. This led him to discover a remarkable series of entirely new ideas and methods.

1) Most operators used in quantum mechanics could \emph{not} be defined in the whole Hilbert space, as for instance in \(L^2(\R)\) multiplication of functions by a fixed function such as \(x \mapsto x\), or derivation of functions. One therefore had to consider, in general, linear mappings \(T\) taking their values in a Hilbert space \(E\), but only defined in a \emph{proper} vector subspace \(\dom(T)\) (the ``domain'' of \(T\)); the most interesting case concerned the operators \(T\) \emph{densely defined}, i.e. those for which \(\dom(T)\) is dense in \(E\) (as in the two examples above).

2) If \(\dom(T)\) is dense in \(E\) and \(T\) is continuous, it can immediately be extended to the whole space \(T\), and one is brought back to the Hilbert theory. But von Neumann had the idea to introduce a weaker substitute for continuity, namely the fact that the \emph{graph} \(\Gamma(T)\) of \(T\) be closed in \(E \times E\); one then says that \(T\) is a \emph{closed} operator. It is obvious that for \(\dom(T) = E\), if \(T\) is continuous, then \(T\) is closed, and the converse follows from the closed graph theorem (chap. \ref{ch:6}, \ref{sec:6.5}). Of the two examples given in 1), the first is closed but the second is not.

3) This last example raises the problem of \emph{extending} (if possible) a densely defined operator \(T\) which is not closed to a closed operator; von Neumann was able to give a beautiful answer to that problem by linking it to a generalization of the notion of \emph{adjoint} operator, well known for bounded operators. In general, for a densely defined operator \(T\) and a vector \(y \in E\), the linear form \(x \mapsto (T \cdot x | y))\), defined in \(\dom(T)\), is not necessarily continuous; if it is, it can be extended uniquely by continuity to \(E\), and then can be written \(x \mapsto (x | y^*)\) for a unique vector \(y^* \in E\); the set of vectors \(v \in E\) having that property is a vector subspace, and if one writes \(y^* = T^* \cdot y\) for those vectors, \(T^*\) is a linear operator defined in that subspace (which is therefore \(\dom(T^*)\)). Now it is easy to show that \(T^*\) is always closed (even if \(T\) is not), and its graph is the subspace of E which is the orthogonal supplement to the closure of \(J(\Gamma(T))\), where \(J\) is the linear automorphism \((x,y) \mapsto (y,-x)\) of \(E \times E\) (``rotation of a right angle''!). This interpretation of \(\dom(T^*)\) gave to von Neumann the proof of the \emph{equivalence} of the two following properties: a) \(T\) can be extended to a closed operator (one says \(T\) is \emph{closable}); b) \(T^*\) is densely defined. One can easily give examples in which \(\dom(T^*) = \{0\}\); if \(T\) is closable, the smallest closed extension of \(T\) is \(T^{**}\), and one has \(T(T^*) = \Gamma(T)\) and \((T^{**})^* = T^*\).

4) The fact that closed densely defined operators are not everywhere defined raises difficulties concerning algebraic combinations of such operators: \(A + B\) is only defined in \(\dom(A) \cap \dom(B), AB\) only in \(\dom(B) \cap B^{-1}(\dom(A))\); one can give examples of closed densely defined operators \(T\) such that \(\dom(T^2) = \{0\}\). However, using the decomposition of \(E \times E\) in the direct sum of the closed orthogonal subspaces \(\Gamma(T)\) and \(J(\Gamma(T^*))\), von Neumann could prove that for any closed densely defined operator \(T, \dom(T^* T)\) is dense, \(T^* T\) is closed and \((T^* T)^* = T^* T\). Furthermore, \(1_E + T^* T\) (closed and defined in \(\dom(T^* T)\)) is a bijection of \(\dom(T^* T)\) onto the whole space \(E\), the inverse \(B(1_E + T^* T)^{-1}\) is a \emph{bounded} self-adjoint and injective operator, the spectrum of which is contained in the interval \([0,1]\).

5) These results enabled von Neumann to completely elucidate the spectral theory of \emph{normal} operators in \(E\). By definition, they are the closed densely defined operators \(N\) such that
\(\dom(N^* N) = \dom(N N^*)\) and \(N^* N = N N^*\). The most important normal operators are the \emph{self-adjoint} operators (which von Neumann called ``hypermaximal''), defined by the condition \(N^* = N\) (implying of course \(\dom(N^*) = \dom(N)\)), and the \emph{unitary} operators, which are bounded and such that \(N^* N = 1_E\) (hence invertible and such that \(N^{-1} = N^*\)).

Now F. Riesz's definition of the spectrum of a bounded operator can be generalized for any \emph{closed} operator \(T\) in \(E\). One says a complex number \(\zeta\) is a \emph{regular value} for \(T\) if the operator \(T - \zeta 1_E\) is a \emph{bijection} of the subspace \(\dom(T)\) onto the whole space \(E\) and if the inverse mapping \(R_T(\zeta)\) (also called the \emph{resolvent} of \(T\)) is a \emph{bounded} operator mapping \(E\) onto \(\dom(T)\); it is enough for that to know that \(T - \zeta 1_E\) is injective, that its image \(L\) is dense in \(E\), and the inverse mapping \((T - \zeta 1_E)^{-1}\) of \(L\) onto \(\dom(T)\) is continuous. The complement \(\spec(T)\) of the set of regular values of \(T\) in \(\C\) is by definition the \emph{spectrum} of \(T\), and the mapping \(\zeta \mapsto R_T(\zeta)\) of \(\C - \spec(T)\) into \(\mathcal{L}(E)\) is \emph{holomorphic}. For a number \(\zeta \in \spec(T)\), there are 3 possibilities:

1° \(T - \zeta 1_E\) is not injective, which means there exists an \(x \in \dom(T)\) such that \(x \neq O\) and \(T \cdot x = \zeta x\), in other words \(\zeta\) is an \emph{eigenvalue} of \(T\); one then says \(\zeta\) belongs to the \emph{point spectrum} of \(T\).

2° \(T - \zeta 1_E\) is injective and its image \(L\) is dense in \(E\), but the inverse mapping \((T - \zeta 1_E)^{-1}\) is not continuous in \(L\); then \(\zeta\) is said to belong to the \emph{continuous spectrum} of \(T\).

3° \(T - \zeta 1_E\) is injective, but its image \(L\) is not dense in \(E\); one says belong to the \emph{residual spectrum} of \(T\).

For \emph{normal} operators, there is \emph{no residual spectrum}; self-adjoint operators are characterized as normal operators for which the spectrum is \emph{contained in} \(\R\), and unitary operators are normal operators for which the spectrum is \emph{contained in the unit circle} \(\mathbb{U}: |\zeta| = 1\).

6) Generalizing F. Riesz's presentation of the Hilbert spectral theory (\ref{sec:7.2}), von Neumann shows that to every \emph{self-adjoint} operator \(A\) in \(E\) is naturally associated a unique \emph{decomposition of unity}. He means by that a family \(\lambda \mapsto E(\lambda)\) of \emph{orthogonal projectors} in \(E\), depending on a real parameter \(\lambda\), and such that:

1° \(E(\lambda) E(\mu) = E(\mu) E(\lambda) = E(\lambda)\) for \(\lambda \leq \mu\);

2° when \(\lambda > \lambda_0\) tends to \(\lambda_0, E\) tends to \(E(\lambda_0)\) strongly; when \(\lambda\) tends to \(+\infty, E(\lambda)\) tends strongly to \(0\), and when \(\lambda\) tends to \(+\infty, E(\lambda)\) tends strongly to \(1_E\);

3° for any \(x \in E\), the mapping \(\lambda \mapsto \norm{E(\lambda) \cdot x}^2\) increases from 0 to \(\norm{x}^2\) in \(\R\); \(\dom(A)\) is exactly the set of \(x \in E\) such that the Stieltjes integral
\[
    \int_{-\infty}^{+\infty} \lambda^2 d\left( \norm{E(\lambda) \cdot x}^2 \right)
\]
is \emph{finite};

4° for any \(x \in \dom(A)\) and any \(y \in E\), the function \(\lambda \mapsto (E(\lambda) \cdot x | y)\) is a function of bounded variation, and one has the expressions
\begin{equation}
    \label{eq:7.36}
    (A \cdot x | y) = \int_{-\infty}^{+\infty} \lambda d((E(\lambda) \cdot x | y)),
    \qquad
    (x|y) = \int_{-\infty}^{+\infty} d((E(\lambda) \cdot x | y))
\end{equation}
as Stieltjes integrals.

Conversely, for any family \(\lambda \mapsto E(\lambda)\) of orthogonal projectors satisfying 1° and 2°, conditions 3° and 4° \emph{define} a self-adjoint operator \(A\) and its domain, to which the given family is its decomposition of unity. The operator \(A\) is bounded if and only if there is a compact interval \([\alpha, \beta]\) such that \(E(\lambda) = 0\) for \(\lambda < \alpha\) and \(E(\lambda) = 1_E\) for \(\lambda > \beta\). The spectrum of \(A\) is the complement of the set of points \(\mu \in \R\) such that \(E(\lambda)\) is constant in a neighborhood of \(\mu\), and the point spectrum is the set of points \(\mu\) such that \(E(\mu-)\) is distinct from \(E(\mu)\).

For \emph{unitary} operators \(U\), there is a similar result: to \(U\) corresponds a unique decomposition of unity \(\lambda \mapsto E(\lambda)\) satisfying conditions 1° and 2° above, with \(E(\lambda) = 0\) for \(\lambda < 0\) and \(E(\lambda) = 1_E\) for \(\lambda  > 1\); condition 3° is then automatically satisfied, and the first relation \eqref{eq:7.36} is replaced by
\begin{equation}
    \label{eq:7.37}
    (U \cdot x | y) = \int_0^1 e^{2i \pi \lambda} d((E(\lambda) \cdot x | y)).
\end{equation}
There is a similar ``decomposition'' for all \emph{normal} operators, but we shall give it in a much simpler equivalent form in \ref{sec:7.5}.

7) The most original part of von Neumann's work on spectral theory is his discovery and study of \emph{hermitian} operators in Hilbert space \(E\), as \emph{distinct} from self-adjoint operators. A hermitian operator \(H\) is a densely defined operator such that \(\dom(H) \subset dom(H^*)\) and that the restriction of \(H^*\) to \(\dom(H)\) is equal to \(H\), in other words
\begin{equation}
    \label{eq:7.38}
    (H \cdot x | y) = (x | H \cdot y) \quad \text{for } x,y \text{ in } \dom{H},
\end{equation}
and in particular \((H \cdot x|x)\) is a \emph{real} number for all \(x \in \dom(H)\). This implies that \(H\) is closable, and its closure \(H^{**}\) is again a hermitian operator; one may therefore restrict the study to \emph{closed} hermitian operators. The new idea of von Neumann is to adapt to Hilbert space a device introduced in 1855 by Cayley to parametrize the orthogonal group: he had shown that, for an \(n \times n\) skew-real symmetric matrix \(S\), such that \(\det(I+S) \neq 0, U = (I-S)(I+S)^{-1}\) was an orthogonal matrix, and any orthogonal matrix \(U\) such that \(\det(I+U) \neq 0\) could be written in that way. Similarly, for a closed hermitian operator \(H\), one has \(\norm{x}^2 \leq \norm{H \cdot x + ix}^2\) for \(x \in \dom(H)\), which implies that the closed operator \(H + iI\) is injective in \(\dom(H)\), and maps \(\dom(H)\) on a \emph{closed} subspace \(F\) in such a way that \((H+iI)^{-1}\) is continuous in \(F\), and \(V: y \mapsto (H - iI)(H + iI)^{-1} \cdot y\) is an \emph{isometry} of \(F\) on a closed subspace \(V(F)\) of \(E\). Conversely, if \(U\) is an isometry of a closed subspace \(F\) of \(E\) onto another closed subspace \(U(F)\) such that the image \(G\) of \(F\) by \(I - U\) is \emph{dense} in \(E\), then \(I-U\) is a \emph{bijection} of \(F\) onto G, and if, for \(y \in G\), one writes \(H \cdot y = i(I+U)(I-U)^{-1} \cdot y, H\) is a closed hermitian operator such that \(\dom(H) = G\) and \(U = (H-iI)(H+iI) = V\) defined above.

Furthermore, if \(E^+_H\) is the orthogonal supplement of \(F\) in \(E\), it is exactly the subspace of \(\dom(H^*)\) consisting of the solutions of \(H^* \cdot x = ix\); similarly, the orthogonal supplement \(E^-_H\) of \(V(F)\) in \(E\) is the subspace of the solutions of \(H^* \cdot x = -ix\) in \(\dom(H^*)\), and \(\dom(H^*)\) is the \emph{direct sum} of the three subspaces \(\dom(H), E^+_H\) and \(E^-_H\).

8) This method enables von Neumann to give a description of all hermitian operators \(H_1\) which \emph{extend} a given hermitian operator \(H\). It is enough to describe the isometry \(V_1\) which is the ``Cayley transform'' of \(H_1\): one takes a closed subspace M of \(E\); and an isometry \(W\) of \(M\) onto a closed subspace \(N\) of \(E_H^+\) is then defined in the Hilbert sum \(F_1 = E \oplus M\), equal to \(V\) in \(F\) and to \(W\) in \(M\); \(E^+_{H_1}\) is the orthogonal supplement of \(M\) in \(E^+_H\) and \(E^-_{H_1}\) the orthogonal supplement of \(N\) in \(E^-_H\). Such a construction is of course only possible if \(\dim(M) \leq \dim(E^-_H)\). The dimensions \(d^+\) of \(E^+_H\) and \(d^-\) of \(E^-_H\) are called the \emph{defects} of \(H\); examples may be given in which they take any integral value or are infinite.

Self-adjoint operators are by definition hermitian operators for which \(H^* = H\), or equivalently those for which the defects are \((0,0)\). It follows at once from the preceding remarks that the closed hermitian operators which can be extended to self-adjoint operators are exactly those for which both defects are \emph{equal} (finite or infinite); unless they are both 0, there are infinitely many such self-adjoint extensions.

To give an example of a closed hermitian operator of defects \((1,0)\), von Neumann takes an orthonormal basis \((e_n)_{n \geq 0}\) of \(E\), and in \(E\) considers the closed hyperplane \(F\) orthogonal to \(e_0\), hence spanned by the \(e_n\), with \(n \geq 1\); he denotes by \(U\) the isometry of \(F\) onto \(E\) defined by \(U \cdot e_n = e_{n-1}\) for \(n \geq 1\); it is easy to show that the image of \(F\) by \(I - U\) is dense in \(E\), and therefore \(U\) is the Cayley transform of a closed hermitian operator \(F\) having the required property. Another (non closed) hermitian operator is given by taking \(E = L^2(I)\), where \(I\) is any interval in \(\R\), and \(H = i\frac{d}{dx}\) which is defined in the subspace of \(E\) consisting of \(C^1\) functions vanishing at both extremities of \(I\) and whose derivative is square integrable (or any subspace of that subspace which is still dense in \(L^2(I)\), for instance the space of \(C^\infty\) functions with compact support in \(I\)); it may then be shown that the defects of \(H^{**}\) are \((1,1)\) if \(I\) is bounded, \((1,0)\) if \(I\) is only bounded from above, \((0,1)\) if \(I\) is bounded from below and \((0,0)\) if \(I = \R\).

As we already mentioned (\ref{sec:7.3}) the results of H. Weyl on linear second order equations with real coefficients can easily be interpreted in the von Neumann theory: the differential operator \(L\) is hermitian; the defects of \(L^{**}\) are \((2,2)\) in the ``limit circle'' case, and \((1,1)\) in the ``limit point'' case. They prefigurated the general spectral theory of formally self-adjoint linear differential equations which developed around 1950 (see chap. \ref{ch:9}, \ref{sec:9.3}).

Similarly, Carleman's results are interpreted in the following way: for a Carleman kernel \(K\), if one writes \(k(s)^2 = \int_a^b |K(s,t)|^2 dt\), in order to get a hermitian operator, one a should \emph{restrict} the operator \(S\) defined in \ref{sec:7.3} (formula \eqref{eq:7.34}) to the subspace (dense in \(L^2\)) of functions \(f\) such that the integral \(\int_a^b k(s)|f(s)|ds\) is finite; \(S\) is then the \emph{adjoint} of that operator, which explains the existence of non trivial solutions of \(S \cdot \varphi = i \varphi\) in \(D = \dom(S)\), and shows that the operator \(U_\lambda\) suitably restricted, coincides with the ``Cayley transform'' which von Neumann later defined in a more general context.

One should finally mention that von Neumann took pains, in a special paper (\cite{221}, vol.~II, p.~141--172), to investigate how hermitian operators might be represented by infinite matrices (to which many mathematicians, and even more physicists, were sentimentally attached); he pointed out that if one wanted to associate to a hermitian operator \(H\) a matrix \((a_{mn})\) by the usual rule \(a = (H \cdot e_m | e_n)\) for an orthonormal basis \((e_n)\) of Hilbert space, one immediately ran into difficulties, since the vectors \(H \cdot e_m\) should be defined, in other words one should have \(E_n \in \dom(H)\) for all \(n\); furthermore, the sums \(\sum\limits_n |a_{mn}|^2\) should all be finite. But if \(H\) is not maximal (i.e. both defects are \(>0\)), any hermitian operator which extends \(H\) obviously has the \emph{same} matrix \((a_{mn})\); and von Neumann showed in great detail how this lack of ``one-to-oneness'' in the correspondence between matrices and operators led to the weirdest pathology, convincing once for all the analysts that matrices were a totally inadequate tool in spectral theory.

\section{Banach algebras}
\label{sec:7.5}

We have seen (\ref{sec:7.2}) that F. Riesz probably was the first mathematician to consider the algebra \(\mathcal{L}(E)\) of all continuous endomorphisms of a separable Hilbert space \(E\), with its norm and what later came to be called its strong topology. In his second paper on spectral theory (\cite{221}, vol.~II, p.~86--143) in which he introduced the concept of normal operator in its most general form, von Neumann began a more detailed study of \(\mathcal{L}(E)\) and its subalgebras. He introduced the weak topology on \(\mathcal{L}(E)\) (see chap. \ref{ch:8}, \ref{sec:8.1}), and (inspired by the work of I. Schur on linear representations of groups) the concept of \emph{commutant} \(M'\) of a subset \(M\) of \(\mathcal{L}(E)\), but with an additional condition: \(M'\) should consist of all operators \(A\) such that, not only \(A\) \emph{but also} \(A^*\), was permutable with all elements of \(M\). He focused his interest on the subalgebras of \(\mathcal{L}(E)\) (later called \emph{involutive} or \(^*\)-subalgebras) which, with every element \(A\), also contained its adjoint \(A^*\); and he proved in that paper the first two non trivial results on such subalgebras: the \emph{double commutant} \(M''\) of \emph{any} involutive subalgebra \(M\) of \(\mathcal{L}(E)\) containing \(1_E\) is the \emph{weak closure} of \(M\), and any weakly closed commutative subalgebra of \(\mathcal{L}(E)\) is generated by a single self-adjoint operator \(A\). A little later he completed this last result by showing that one could define ``functions \(f(A)\)'' of a self-adjoint operator \(A\) for all universally measurable bounded functions \(f\) defined in \(\R\) (and not only for semi-continuous functions, as F. Riesz had done), and he proved that the weakly closed subalgebra generated by \(A\) consisted of all operators \(f(A)\) thus defined (\cite{221}, vol.~II, p.~177--212).

But for von Neumann this was only a beginning. The period 1926--1932 had seen the blossoming forth of the theory of ``hypercomplex numbers'' of Molien, E. Cartan and Wedderburn into the beautiful theory of ``rings with descending chain conditions'' of E. Artin and E. Noether, followed by their applications to linear representations of groups and number theory by R. Brauer, H. Hasse and A. Albert. Von Neumann was very much interested by these developments, and wondered if it could not be possible to build up some similar theory for involutive subalgebras of \(\mathcal{L}(E)\), where of course ``chain conditions'' could not be expected, but suitable topological restrictions would be a substitute, allowing one to obtain a reasonable classification (loc.cit., p. 89). It would take us too far away from our main theme to describe in some detail the series of papers, beginning in 1935, in which, with the partial collaboration of F. Murray, he achieved a great part of this program for what we now call the \emph{von Neumann algebras}, namely the involutive subalgebras equal to their double commutant in \(\mathcal{L}(E)\). By the wealth and novelty of their techniques and their results, these wonderful papers are certainly the most profound and most difficult which von Neumann ever wrote (\cite{221}, vol. III); they revealed a large number of completely unsuspected phenomena, the most conspicuous one being the appearance, in the classification of the von Neumann algebras with trivial center (those called \emph{factors}), of \emph{five} types of algebras labeled I\(_n\), I\(_\infty\), II\(_1\), II\(_\infty\) and III, where type I\(_n\) means algebras of \(n \times n\) matrices, type I\(_\infty\) the algebra \(\mathcal{L}(E)\) itself, but the three other types were entirely unexpected and exhibit new features, such as the attribution to the projectors contained in these algebras of a ``dimension'' which, for algebras of Type II, may be \emph{any real number} (in \([0,1]\) or \([0,+\infty]\)) instead of an integer. The elucidation of the properties of these new algebras, begun by Murray and von Neumann, has engaged many mathematicians during the last 40 years, and it is only recently that some difficult questions, such as the classification of algebras of type III, have begun to be understood (see \cite{057}, \cite{210} and \cite{044}). Furthermore, since 1950 the von Neumann algebras have been an important tool in the theory of linear representations of locally compact groups (see \ref{sec:7.6}); more recently they have been associated to foliations and to generalizations of the Atiyah--Singer index (see \cite{058}, \cite{010}, \cite{031}, \cite{036}, \cite{045}, \cite{113}, \cite{132}, \cite{199}).

Surprisingly enough, the difficult theory of von Neumann algebras was developed 5 years before the elementary concepts of the theory of normed algebras had been defined! The creation of that theory was the work of I. Gelfand in 1941 \cite{083}; a normed algebra \(A\) (over the \emph{complex} field) is an algebra over \(\C\) on which is defined a structure of normed space with the condition that the mapping \((x,y) \mapsto xy\) of \(A \times A\) into \(A\) be continuous. It is then possible to choose on \(A\) a norm compatible with the vector space structure and the topology of \(A\), and such that in addition \(\norm{xy} \leq \norm{x} \cdot \norm{y}\). If \(A\) has a unit element \(e\), one may suppose in addition that \(\norm{e} = 1\). If \(A\) has no unit element, it is always possible to imbed \(A\) into a normed algebra \(\tilde{A}\) with a unit element \(e\), such that \(\tilde{A}\) is the direct sum of \(A\) and \(\C e\). For any normed space \(E\) over \(\C\), the algebra \(\mathcal{L}(E)\) of endomorphisms of \(E\) is a normed algebra for the norm \(\norm{A} = \sup\limits_{\norm{x} \leq 1}{\norm{A \cdot x}}\), but there are many other types of normed algebras, the most elementary one being the algebra \(\Cont(I)\) of complex continuous functions in a compact interval \(I\) of \(\R\), with \(\norm{f} = \sup\limits_{t \in I}{\norm{f(t)}}\).

Gelfand's main idea, which proved extraordinarily fruitful, was to extend spectral theory to elements of normed algebras; if \(A\) is a normed algebra with unit element \(e\), one may apply F. Riesz's definition of the spectrum (\ref{sec:7.2}) to define the \emph{spectrum} of an arbitrary element \(x \in A\): it is the set of complex numbers \(\zeta\) such that \(x - \zeta e\) is \emph{not invertible} in \(A\). Gelfand recognized that to get substantial results one must assume that \(A\) is \emph{complete} as a Banach space, what is called a \emph{Banach algebra}. Then very elementary arguments show that the spectrum \(\spec_A(x)\) of any element \(x \in A\) is a non empty compact subset, contained in the disc \(|\zeta| < \norm{x}\); the invertible elements in \(A\) form an open group \(G\) containing the ball \(\norm{x - e} \leq 1\), and the topology induced on \(G\) is compatible with the group structure; for any \(x \in A\), the map \(\zeta \mapsto (x - \zeta e)^{-1}\) of the complement \(\C - \spec_A(x)\) into \(A\) is holomorphic. Finally, Gelfand obtained a beautiful formula for the radius of the smallest disc of center 0 containing \(\spec_A(x)\); this number, called the \emph{spectral radius} of \(x\), is equal to
\begin{equation}
    \label{eq:7.39}
    \rho(x) = \lim_{n \to \infty} \left( \norm{x^n}^{1/n} \right).
\end{equation}

Next Gelfand undertook the study of general \emph{commutative} Banach algebras by a very original method. Probably inspired by the theory of commutative groups (see \ref{sec:7.6}), he defined a \emph{character} \(\chi\) of a Banach algebra \(A\) as a \emph{homomorphism} of that algebra in the field \(\C\) (considered as \(\C\)-algebra) which is not identically 0. Suppose for simplicity that \(A\) has a unit element \(e\); then any character \(\chi\) is such that \(\chi(e) = 1\), and is a continuous linear form on \(A\), of norm \(\norm{\chi} = 1\); furthermore, for each \(x \in A\), one has \(\chi(x) \in \spec_A(x)\). Gelfand then associates to \(A\) the set \(\Chi(A)\) of \emph{all} characters of \(A\); the map \(\chi \mapsto \chi^{-1}(0)\) is a bijection of \(\Chi(A)\) on the set of all \emph{maximal ideals} in \(A\) (which are automatically closed). Now, in 1937, Stone \cite{208} had already considered the set of maximal ideals of a very special type of ring, a ``boolean ring'' \(B\), which is commutative and such that \(x^2 = x\) and \(2x = 0\) for all \(x \in B\); this kind of ring itself had been suggested to Stone by the set of characteristic functions \(\varphi_M\) of subsets \(M\) of an arbitrary set \(E\), where multiplication is the usual one, and addition \(\dotplus\) is defined by \(\varphi_M \dotplus \varphi_N = \varphi_{M \cup N} - \varphi_{M \cap N}\); furthermore, Stone of course was well aware that for a self-adjoint operator \(A\) in Hilbert space, the orthogonal projectors \(\varphi_M(A)\) for universally measurable subsets \(M\) of \(\R\) form a boolean ring for the same addition. The consideration of the set of maximal ideals of a commutative Banach algebra was therefore not at all foreign to the spirit of spectral theory at that time.

As the set \(\Chi(A)\) is contained in the unit ball \(\norm{x'} \leq 1\) of the dual \(A'\) of the Banach \emph{space} \(A\), the natural embedding of \(A\) into its second dual \(A''\) associates to each element \(x \in A\) the map \(\chi \mapsto \chi(x)\) of \(\Chi(A)\) into \(\C\), which is called the \emph{Gelfand transform} of \(x\) and is written \(\mathcal{G}x\). It is easy to see that \(\Chi(A)\) is \emph{compact} for the weak topology of \(A'\), and that \(\mathcal{G}x\) is a \emph{continuous} function on \(\Chi(A)\) for that topology; one has therefore defined a \emph{continuous homomorphism} \(x \mapsto \mathcal{G}x\) of the Banach algebra \(A\) into the Banach algebra \(\Cont(\Chi(A))\), such that the set of values of \(\mathcal{G}x\) is the spectrum of \(x\), and therefore \(\norm{\mathcal{G}x} = \rho(x) \leq \norm{x}\). The compact space \(\Chi(A)\) is therefore called the \emph{spectrum} of the Banach algebra \(A\).

If one starts from the Banach algebra \(A = \Cont(K)\) of continuous functions on a compact space \(K\), then it is easy to see that \(x \mapsto \mathcal{G}x\) is an isomorphism of \(A\) onto \(\Cont(\Chi(A)), \Chi(A)\) being identified with the space of Dirac measures on \(K\). But in general the homomorphism \(x \mapsto \mathcal{G}x\) of \(A\) into \(\Cont(\Chi(A))\) is neither surjective nor injective. A little later, in collaboration with Naimark \cite{085}, Gelfand began to study Banach algebras in which there is an \emph{involution} \(x \mapsto x^*\) (i.e. such that \((x+y)^* = x^* + y^*, (xy)^* = y^* x^*, (\lambda x)^* = \overline\lambda x^*\) for any scalar \(\lambda\) and \((x^*)^* = x\)) for which in addition \(\norm{x^* x} = \norm{x}^2\); these algebras are now called \(C^*\)-\emph{algebras}. The main result proved by Gelfand and Naimark is that, for a commutative \(C^*\)-algebra \(A\) having a unit element \(e\), the mapping \(x \mapsto \mathcal{G}x\) is an \emph{isometry} of \(A\) onto \(\Cont(\Chi(A))\) such that \(\mathcal{G}x^* = \overline{\mathcal{G}x}\) for all \(x \in A\). Furthermore, if there exists in \(A\) an element \(x_0\) such that the subalgebra of \(A\) generated by \(x_0, x^*_0\) and \(e\) is dense in \(A\), then the map \(\chi \mapsto \chi(x_0)\) is a \emph{homeomorphism} of \(\Chi(A)\) onto \(\spec_A(x_0)\), a compact subset of \(\C\) which one therefore identifies with the spectrum of \(A\).

The Gelfand--Naimark theorem paved the way for a new interpretation of Hilbert's spectral theory. Let \(E\) be a separable Hilbert space, \(N\) a continuous normal operator in \(E\); then the closure \(A\) in \(\mathcal{L}(E)\) (for the normed topology) of the algebra generated by \(1_E, N\) and \(N^*\) is a separable \emph{commutative} \(C^*\)-\emph{algebra} with unit, the mapping \(\eta: \chi \mapsto \chi(N)\) being a homeomorphism of \(\Chi(A)\) onto the spectrum \(\spec(N) \subset \C\). From the Gelfand--Naimark theorem, it follows that the mapping \(f \mapsto \mathcal{G}^{-1}(f \circ \eta)\) is an \emph{isometry} of the algebra \(\Cont(\spec(N))\) on a subalgebra of \(\mathcal(E)\), which one writes \(f \mapsto f(N)\), obtaining in this way a new definition of a ``continuous function of a normal operator'' which had been considered by F. Riesz and von Neumann. Following the method of von Neumann, it is then easy to extend the homomorphism \(f \mapsto f(N)\) to the algebra \(\mathcal{U}(\spec(N))\) of all \emph{universally measurable} bounded functions in \(\spec(N)\).

Finally, by adapting the arguments of von Neumann, Hellinger and Hahn, one arrives at the modern description of the Riesz--von Neumann ``decomposition of unity'' and of the ``multiplicity theory" of Hellinger--Hahn:

1° There is a decomposition of \(E\) into a Hilbert sum (finite or not) \((E_j)_{1 \leq j < \omega}\) (\(\omega\) being an integer or \(+\infty\)) of closed subspaces, each of which is stable by \(N\) and by \(N^*\).

2° There is a positive measure \(\nu\) on the compact space \(\spec(N) \subset \C\) with support \(\spec(N)\), and a decreasing sequence \((S_j)_{1 \leq j < \omega}\) with \(S_1 = \spec(N)\), consisting of universally measurable sets.

3° For each \(j\) such that \(1 \leq j < \omega\), there is an \emph{isometry} \(T_j\) of the Hilbert space \(E_j\) onto the Hilbert space \(F_j = L^2(\spec(N), \varphi_{S_j}, \nu)\) such that the normal operator \(T_j N T^{-1}_j\) in \(F_j\) is the ``multiplication operator'' which, to the class of any function \(u_j\) defined and square integrable (for \(\nu\)) in \(S_j\), associates the class of the function \(\zeta \mapsto \zeta u_j(\zeta)\).

4° In this description, the measure \(\nu\) (considered as a measure on \(\C\)) is determined \emph{up to equivalence}, the sets \(S_j\) are determined \emph{up to a null set} (for \(\nu\)). The set \(M_j = S_j - S_{j+1}\) is the part of \(\spec(N)\) of \emph{multiplicity} \(j\), and (when \(\omega = +\infty\)), \(M_\infty = \bigcap_j S_j\) the part of \(\spec(N)\) of \emph{infinite multiplicity}. If \(P_j\) is the orthogonal projector \(\varphi_{M_j}(N)\), and \(H_{ik} = P_k(E_i)\), the restrictions of \(N\) to the \(k\) orthogonal subspaces \(H_{ik} (1 \leq i \leq k)\) are \emph{equivalent}; the subspaces \(E_i\) are not uniquely determined, but the subspaces \(G_k = P_k(E) = H_{1k} \oplus H_{2k} \oplus \ldots \oplus H_{kk}\) are. The \emph{equivalence class} of \(\nu\) and the \emph{classes of the sets} \(S_j\) are the \emph{unitary invariants} of \(N\) which determine it up to a unitary equivalence \(N \mapsto U N N^{-1}\).

One says that this description is a \emph{diagonalization} of the normal operator \(N\). This name is justified when one considers the classical case in which \(N\) is a normal endomorphism of a finite dimensional space \(E\): \(\spec(N)\) is then a discrete subset of \(\C\) consisting of the eigenvalues of \(N, S_j\) the subset consisting of the eigenvalues of multiplicity \(\geq j\), \(\nu\) the measure having mass +1 at each point of \(\spec(N)\), and \(G_j\) is the subspace of \(E\) which is the direct sum of the eigenspaces of \(N\) corresponding to the eigenvalues of multiplicity \(j\).

It is easy, using von Neumann's results, to extend the preceding description to \emph{unbounded} normal operators \(J\): \(\spec(N)\) is then an arbitrary closed subset of \(\C\) (it may be \(\C\) itself), and the \(S_j\) arbitrary universally measurable subsets of \(\spec(N)\) forming a decreasing sequence; \(N\) is not defined in the whole subspace \(E_j\), but \(\dom(N) \cap E_j\) is the subspace transformed by \(T_j\) into the subspace of \(F_j\) consisting of the \(u_j\) such that the function \(\zeta \mapsto \zeta u_j(\zeta)\) is square integrable for \(\nu\).

Furthermore, for each universally measurable function \(f\) in \(\spec(N)\) (bounded or not), \(f(N)\) is a (generally unbounded) normal operator, which one may define in the following way: \(\dom(f(N)) \cap E_j\) is the subspace transformed by \(T_j\) into the subspace of \(F_j\) consisting of the \(u_j\) such that the function \(\zeta \mapsto f(\zeta) u_j(\zeta)\) is square integrable, and the class of this function is the image of the class of \(u_j\) by \(T_j f(N) T^{-1}_j\).

For self-adjoint operators \(A\) in \(E\), the connection with the ``eigendifferentials'' of Hellinger is made by the following remark, due to F. Riesz: for every \(x = (x_k) \in l^2\), a vector \((\rho_k(\xi))\) is defined by taking \(A_\xi \cdot x\) for every \(\xi \in \R\).

\section{Later developments}
\label{sec:7.6}

Since 1940, an enormous number of papers have been published on Banach algebras, spectral theory and their applications. I think a fair and well organized account of all these developments will have to wait till more time has elapsed and has put them in their proper perspective\footnote{Glaring examples of lack of perspective are given by the Hellinger--Toeplitz article of 1923 in the \emph{Enzyklopädie der math. Wiss.} \cite{107}, which gives undue emphasis to integral equations, and by Hadamard's article on Functional Analysis of 1928 (\cite{094}, vol.~I, p.~435--453), which barely mentions F. Riesz and does not speak of spectral theory at all!}. With the exception of the theory of differential (ordinary or partial) and integral equations, which has a complex background in which more than spectral theory is involved, and which will be considered in chap. \ref{ch:9}, we shall limit our survey to bare indications of the general trends, and to references to a few papers and books.

A) \emph{Structure of Banach algebras}.

After Gelfand and his school had investigated the general properties of all Banach algebras, mathematicians concentrated their efforts on two particular classes of such algebras, the \emph{commutative} and the \emph{involutive} ones.

For a commutative Banach algebra \(A\), a central problem was to define ``functions'' of elements \(x \in A\) more general than polynomials, after the pattern set by F. Riesz and von Neumann. The latter had even shown that it was possible to define functions \(f(N_1, \ldots, N_k)\) of commuting normal operators \(N_j\) in a Hilbert space, for all continuous functions \(f\) defined in \(C^k\). For a general commutative Banach algebra \(A\), such a definition was only possible under some restrictions on \(f\); if \(x_1, \ldots, x_k\) are any elements of \(A\), their \emph{joint spectrum} is the image in \(\C^k\) of the mapping \(\chi \mapsto (\chi(x_1), \chi(x_2), \ldots, \chi(x_k))\) where \(\chi\) runs through \(\Chi(A)\) (for \(k = 1\), it is of course the spectrum of \(x_1\)); one then could prove that if B is the algebra of (germs of) functions \(f\) \emph{holomorphic in a neighborhood} (\emph{depending on} f) \emph{of the joint spectrum} of \(x_1, \ldots, x_k\), there is a homomorphism \(B \to A\) written \(f \mapsto f(x_1, \ldots, x_k)\), which uniquely extends the natural homomorphism of the algebra of polynomials on \(\C^k\) into \(A\) written similarly (\cite{222}, \cite{027}).

On the spectrum \(\Chi(A)\) of a commutative Banach algebra \(A\), one soon was led to consider a topology different from the one induced by the weak topology of the dual \(A'\) of the Banach space \(A\). For Boolean rings, Stone had introduced the idea of defining on the space of \emph{maximal ideals} of such a ring \(B\) a topology, in which the closed sets were defined as the sets of maximal ideals containing a given (arbitrary) ideal of \(B\). As the set \(\Chi(A)\) of characters of a commutative Banach algebra \(A\) corresponds in a one-to-one way to the set of maximal ideals, Stone's topology can be defined in the same way on \(\Chi(A)\); it is in general coarser than the weak topology, and one says \(A\) is a \emph{regular} commutative Banach algebra if these two topologies on \(\Chi(A)\) coincide; for instance, the algebra \(\Cont(K)\) of continuous functions on a compact space \(K\) is a regular algebra.

To any closed ideal \(J\) in \(A\), one attaches the set \(h(J)\) of all characters \(\chi \in \Chi(A)\) which vanish on \(J\); a natural question is to ask if the intersection of all kernels \(\chi^{-1}(0)\) (maximal ideals of \(A\)) such that \(\chi \in h(J)\), which always contains \(J\), is actually equal to \(J\); one then says that the ideal \(J\) admits \emph{spectral synthesis}. Giving conditions for a closed ideal to admit spectral synthesis in a regular commutative Banach algebra is a problem which has been extensively studied (\cite{021}, \cite{058}).

\emph{Involutive} Banach algebras \(A\) (not necessarily commutative) are those equipped with an involution \(x \mapsto x^*\) such that \(\norm{x^*} = \norm{x}\) for all \(x \in A\); \(C^*\)-algebras (\ref{sec:7.5}) are involutive algebras, but there exist involutive Banach algebras which are not \(C^*\)-algebras (see C) below). The central concept is that of \emph{representation} of an involutive Banach algebra \(A\) in a Hilbert space \(E_j\); this means a homomorphism \(f: A \to \mathcal{L}(E)\) of algebras such that in addition \(f(x^*) = f(x)^*\). They have been the subject of a large number of investigations, leading to the elucidation of the structure of several classes of \(C^*\)-algebras; the theory of von Neumann algebras (which are special types of \(C^*\)-algebras) plays a great part in these investigations (\cite{058}, \cite{036}).

B) \emph{Algebras of continuous functions}.

Since 1960, many mathematicians have been interested in the study of \emph{subalgebras} of Banach algebras \(\Cont(K)\) of continuous functions on a compact space \(K\). In classical Analysis, one had much studied the case in which \(K\) is the unit disk \(|z| \leq 1\) in \(\C\); there is then in \(\Cont(K)\) a particularly interesting Banach subalgebra, namely the algebra \(B\) of functions which are \emph{holomorphic} in the interior \(|z| < 1\) of the disc. It can be identified with the algebra \(B_0\) of the \emph{restrictions} of the functions of \(B\) to the unit circle \(\mathbb{U}: |z| = 1\), and \(B_0\) is also the closure in \(\Cont(\mathbb{U})\) of the algebra of \emph{trigonometric polynomials}. It turns out that the study of \(B_0\) is closely linked to the \emph{completions} of the space of trigonometric polynomials in the various spaces \(L^p(\mu)\), where \(\mu\) is Haar measure on \(\mathbb{U}\), and many beautiful properties of these spaces (known as the \emph{Hardy spaces} \(H^p(\mu))\) had been discovered. But in the light of the theory of commutative Banach algebras, it was found that these results could be much better understood if they were generalized to subalgebras of an algebra \(\Cont(K)\) where \(K\) is any compact space, and put in relation with some kinds of measures on \(K\) (see \cite{018}, \cite{033}, \cite{079}, \cite{116}, \cite{140}).

C) \emph{Harmonic Analysis}.

We have already stressed the fact (chap. \ref{ch:1}, \ref{sec:1.2}) that Fourier series provided the starting point of spectral theory when it was realized that they could be generalized to ``expansions'' in series of ``orthogonal''' functions arising from boundary value problems.

It was, however, very soon observed that the ``trigonometric system'' \((e^{inx})_{n \in \mathbb{Z}}\) possessed very peculiar properties not shared by general ``orthogonal systems'', and linked to the functional equation \(e^{i(x + x')} = e^{ix} e^{ix'}\). For instance, if \(f(x) = \sum\limits_{n=-\infty}^{+\infty} a(n) e^{nix}, g(x) = \sum\limits_{n=-\infty}^{+\infty} b(n) e^{nix}\) were two Fourier series, one had for the Fourier series \(f(x)g(x) = \sum\limits_{n=-\infty}^{+\infty} c(n)e^{nix}\) of their product, the very simple formula
\begin{equation}
    \label{eq:7.40}
    c(n) = \sum_{p=-\infty}^{+\infty} a(p)b(n-p).
\end{equation}
Similarly, from the formula \(f(x) = \sum\limits_{n=-\infty}^{+\infty} a(n) e^{nix}\), one obtained
\begin{equation}
    \label{eq:7.41}
    \sum_{n=-\infty}^{+\infty} a(n+1) e^{nix} = e^{-ix}f(x)
\end{equation}
and this property was used by de Moivre and even more by Laplace to solve linear difference equations \(\sum\limits_{k=1}^p \alpha_k a(n+k) = 0\) by associating to the sequence \((a(n))\) the Fourier series \(f(x) = \sum\limits_{n} a(n) e^{nix}\), reducing the difference equation to an algebraic equation for \(f(x)\).

Similar peculiarities were observed for the ``Fourier transform'' associating to an integrable function \(f\) in \(\R\) the function
\begin{equation}
    \label{eq:7.42}
    \mathcal{F}f(x) = \int_{-\infty}^{\infty} e^{-2\pi ixt} f(t)dt.
\end{equation}
Its main virtue, in the eyes of Fourier, Cauchy and Poisson, was that it reduced linear partial differential equations with constant coefficients to algebraic problems, due to the fact that the Fourier transform of the derivative \(f'\) is the function \( x\mapsto 2\pi ix \mathcal{F}f(x)\). Furthermore, in his researches on Probability theory, Tchebycheff had shown that if \(F_1, F_2\) are two independent ``random variables'' with ``probability laws'' \(\alpha_1, \alpha_2\), measures on \(\R\) with densities \(g_1, g_2\), the probability law of \(F_1 + F_2\) had a density given by the convolution \(g = g_1 * g_2\), defined as
\begin{equation}
    \label{eq:7.43}
    g(x) = \int_{-\infty}^{+\infty} g_1(t) g_2(x - t)dt
    \qquad
    (\text{\cite{211}, vol.~II, p.~481--491});
\end{equation}
and Tchebycheff's student Liapounov, who started to use Fourier transforms in Probability theory, observed that
\begin{equation}
    \label{eq:7.44}
    \mathcal{F}(g_1 * g_2) = \mathcal{F}g_1 \cdot \mathcal{F}g_2
    \qquad
    (\text{\cite{148}}).
\end{equation}

One should also mention the Poisson formula (also discovered independently by Cauchy)
\begin{equation}
    \label{eq:7.45}
    \sum_{n \in \mathbb{Z}} f(n) = \sum_{n \in \mathbb{Z}} \mathcal{F}f(n)
\end{equation}
for sufficiently regular functions \(f\) on \(\R\).

It took over 100 years to understand these peculiarities and to connect them with the notion of group, \emph{via} the concepts of \emph{character} and of \emph{group algebra}. Characters were first defined for arbitrary finite commutative groups by H. Weber in 1882, as complex valued functions \(\chi\) on such a group \(G\) with values \(\neq 0\), such that \(\chi(xy) = \chi(x)\chi(y)\) for all \(x, y\) in \(G\); but special cases had long before been considered by Legendre, Gauss and Dirichlet. A meaningful generalization to non commutative finite groups was discovered in 1896 by Frobenius: instead of considering homomorphisms of \(G\) into the multiplicative group \(C^*\), one should consider homomorphisms \(s \mapsto U(s)\) of \(G\) into the general linear group \(\text{GL}(n,\C\) of invertible matrices of order \(n\), for any integer \(n\). This is also called a \emph{linear representation} of degree \(n\) of \(G\) in the vector space \(E = \C^n\): giving such a representation is equivalent to defining an action \((s,x) \mapsto s \cdot x\) of \(G\) on \(E\) such that \(s \cdot (t \cdot x) = (st) \cdot x, e \cdot x = x\) for the neutral element \(e \in G\), and such that each mapping \(x \mapsto s \cdot x\) is linear (with matrix \(U(s)\)).

Now Cayley had defined, for a finite group \(G\), the \emph{group algebra} \(\C[G]\) as the vector space of all formal linear combinations \(\sum\limits_{s \in G} \xi_s s\) with \(\xi_s \in \C\), multiplication being defined by
\begin{equation}
    \label{eq:7.46}
    \left( \sum_{s \in G} \xi_s s \right) \left( \sum_{s \in G} \eta_s s \right)
    =
    \sum_{(s,t) \in G \times G} \xi_s \eta_t st.
\end{equation}

When there is given an action \((s,x) \mapsto s \cdot x\) of \(G\) on \(E\) as above, it defines naturally on \(E\) a structure of \emph{left} \(\C[G]\)-\emph{module} by
\begin{equation}
    \label{eq:7.47}
    \left( \sum_{s \in G} \xi_s s \right) \cdot x
    =
    \sum_{(s,t) \in G \times G} \xi_s (s \cdot x)
\end{equation}
and the study of linear representations of \(G\) is thus equivalent to the study of left \(\C[G]\)-modules.

The fundamental results of Frobenius for finite groups may then be expressed in the following way. An element \(\sum\limits_{s \in G} \xi_s s\) of \(\C[G]\) may be identified with a mapping \(f: s \mapsto x_s\) of \(G\) into \(\C\), so that \(\C[G]\) may be identified with the vector space \(\C^G\) of \emph{all} mappings of \(G\) into \(\C\), with the multiplication written \(f * g\) and defined by
\begin{equation}
    \label{eq:7.48}
    (f * g)(s) = (\Card(G))^{-1} \sum_{t \in G} f(t)g(t^{-1} s).
\end{equation}

Then \(\C[G]\) can be written as a direct sum \(A_1 \oplus A_2 \oplus \ldots \oplus A_h\) of mutually annihilating subalgebras, where \(h\) is the number of classes of conjugate elements in \(G\); each \(A_k (1 \leq k \leq h)\) a \emph{matrix algebra} of dimension \(n^2_k\) over \(\C\), which means that it has a basis \((m^k_{ij})\) of \(n^2_k\) elements belonging to \(\C^G\), with the following properties:
\begin{equation}
    \label{eq:7.49}
    m^k_{pq} * m^k_{rs} = \delta_{qr} m^k_{ps}
    \text{ for }
    1 \leq p, q, r, s \leq n_k
\end{equation}
\begin{equation}
    \label{eq:7.50}
    m^k_{pq} * m^{k'}_{rs} = 0
    \text{ if }
    k \neq k'.
\end{equation}
In addition, one has
\begin{equation}
    \label{eq:7.51}
    m^k_{ji}(s) = \overline{m^k_{ij}(s^-1)}
    \text{ for }
    1 \leq i, j \leq n_k, s \in G
\end{equation}
\begin{equation}
    \label{eq:7.52}
    \sum_{s \in G} m^k_{pq}(s) m^{k'}_{rs}(s) = 0
    \text{ unless }
    p=r, q=s, k=k'
\end{equation}
\begin{equation}
    \label{eq:7.53}
    \sum_{s \in G} m^k_{pq}(s) \overline{m^k_pq(s)} = n_k \Card(G)
    \text{ for }
    1 \leq p, q \leq n_k
\end{equation}
(\emph{orthogonality relations}). One has \(n^2_1 + n^2_2 + \ldots + n^2_h = \Card(G)\) and the expression of any \(f \in \C^G\) with respect to the basis \((m^k_{ij})\) of that space,
\begin{equation}
    \label{eq:7.54}
    f(s) = \sum_{i, j, k} c^k_{ij} m^k_{ij}(s)
\end{equation}
is given explicitly, due to the orthogonality relations, by
\begin{equation}
    \label{eq:7.55}
    c^k_{ij} = (n_k \Card(G))^{-1} \sum_{s \in G} f(s)\overline{m^k_{ij}(s)}.
\end{equation}
If one writes \(M_k(s)\) the \(n_k \times n_k\) matrix \((n^{-1}_k m^k_{ij}(s))\), one has
\begin{equation}
    \label{eq:7.56}
    M_k(st) = M_k(s) M_k(t)
    \text{ and }
    M_k(s^{-1}) = (M_k(s))^*
\end{equation}
In other words, \(s \mapsto M_k(s)\) is a \emph{linear representation} of \(G\) of degree \(n_k\), for \(1 \leq k \leq h\); it is \emph{irreducible}, which means that the corresponding \(\C[G]\)-module is \emph{simple} (i.e. has no non trivial submodule). Furthermore, \emph{every} \(\C[G]\)-module of finite dimension over \(\C\) is a \emph{direct sum} of modules each of which corresponds to one of the linear representations \(s \mapsto M_k(s)\); one says that every linear representation of \(G\) is \emph{completely reducible}, and that it contains the irreducible representation \(s \mapsto M_k(s)\) with \emph{multiplicity} \(d\), if in the direct decomposition of the corresponding module, there are \(d_k\) submodules corresponding to \(s \mapsto M_k(s)\).

Now linear representations of degree \(n\) can be defined in the same way for \emph{any} group \(G\), \emph{finite or not}. Already in 1901, I. Schur, in his dissertation (\cite{193}, vol.~I, p.~1--70), could determine all linear representations of the general linear group \(\text{GL}(N,\C)\) which are such that the elements of \(U(s\)) are \emph{polynomials} in the elements of the matrix \(s \in \text{GL}(N,\C)\); he showed that these representations are again completely reducible and he could determine explicitly the irreducible ones (see \cite{053}); but it is clear that for such infinite groups, all the Frobenius relations described above were meaningless. However, in 1924, I. Schur observed that the restrictions of these representations to the group of \emph{rotations} \(G = \text{SO}(N,\R)\) gave him irreducible representations of that \emph{compact} group, and that these representations were continuous, and could be written \(s \mapsto M_k(s)\) where \(M_k(s)\), as in \eqref{eq:7.56}, was a \emph{unitary} matrix; furthermore, he proved the relations which he rightly considered as the analogues of \eqref{eq:7.52}
\begin{equation}
    \label{eq:7.57}
    \int_G m^k_{pq} \overline{m^{k'}_{rs}(s)}d = 0
    \text{ unless }
    p=r, q=s, k=k'
\end{equation}
where \(ds\) is a left and right invariant measure on \(G\), the existence of which was substantially known since S. Lie, and which had already been used to construct invariants of \(\text{SO}(N,\R)\) by Hurwitz in 1898 (\cite{193}, vol.~II, p.~440--494),

This result attracted the attention of H. Weyl; in a beautiful series of 3 papers published the next year, by a skillful combination of Schur's ideas with the ``infinitesimal'' methods by which E. Cartan in 1913 had obtained all finite dimensional representations of the complex semi-simple Lie groups, he was able to determine explicitly \emph{all} continuous irreducible linear representations of \emph{compact semi-simple Lie groups} (including, in the case of \(\text{SO}(N,\R)\), the ``spinor'' representations which had escaped I. Schur). In all cases, the orthogonality relations \eqref{eq:7.57} still held, and every continuous linear representation of a semi-simple compact Lie group was shown to be completely reducible (\cite{227}, vol.~II, p.~633).

Of course, for compact Lie groups, there is an \emph{infinite} system of irreducible representations \(M_k\), and relations such as \eqref{eq:7.54} were out of the question. But for the group \(\text{SO}(1,\R) = \mathbb{U}\) (the circle group), the irreducible representations were the characters \(\zeta \mapsto \zeta^n\) for \(n \in \mathbb{Z}\), and H. Weyl realized that the formula which corresponded to \eqref{eq:7.54} was just the \emph{Fourier series expansion} of \(f\) (when f is sufficiently regular) \cite[vol.~III, p.~34--37]{227}. He then undertook to generalize this expansion to all semi-simple compact Lie groups \(G\); for such a group, the functions \(m^k_{ij}\) which he had determined formed an orthogonal system in the Hilbert space \(L^2(G)\) (for a left and right invariant measure); the problem was to prove that this system was \emph{complete}.

This is what H. Weyl proved in 1927, in a remarkable paper written in collaboration with his student F. Peter \cite[vol.~III, p.~58--75]{227}, which can be considered as the first application of spectral theory to harmonic analysis. He saw that the notion which could serve as a substitute to the group algebra \(\C[G]\) was the space \(\Cont(G)\) of continuous complex-valued functions on the compact group \(G\), on which an algebra structure is defined by \emph{convolution}, generalizing \eqref{eq:7.40}, \eqref{eq:7.43}, and \eqref{eq:7.48}:
\begin{equation}
    \label{eq:7.58}
    (f * g)(s) = \int_G f(t) g(t^-1 s)dt
\end{equation}
where the integration is for a left and right invariant positive measure on \(G\) with total mass 1. H. Weyl next observed that, given a linear representation \(s \mapsto U(s)\) of \(G\) by unitary matrices of order \(n\), if one wrote
\begin{equation}
    \label{eq:7.59}
    U(f) = \int_G f(s)\overline{U(s)}ds
\end{equation}
one obtained a \emph{homomorphism} of the algebra \(\Cont(G)\) into \(\text{End}(\C^n)\), in other words
\begin{equation}
    \label{eq:7.60}
    U(f * g) = U(f) U(g)
\end{equation}
and furthermore, if one wrote \(\check{\overline{f}}(s) = \overline{f(s^{-1})}\), one had
\begin{equation}
    \label{eq:7.61}
    U(\check{\overline{f}}) = (U(f))^*.
\end{equation}

The crux of his proof is to show that for an \(f \neq 0\) in \(\Cont(G)\), there is at least a continuous representation \(s \mapsto U(s)\) for which \(U(f) \neq 0\); due to the complete reducibility of \(s \mapsto U(s)\), there is then at least one irreducible representation \(s \mapsto M_k(s)\) such that \(M_k(f) \neq 0\), and this hows that \(f\) cannot be orthogonal to all functions \(m^k_{ij}\). However, if the system \((m^k_{ij})\) was not complete, there would exist a non negligible function \(g \in L^2(G)\) orthogonal to all the \(m^k_{ij}\), and as the subspace \(L\) of \(\Cont(G)\) generated by the \(m^k_{ij}\) is a two-sided ideal (one has \((f*U)(s) = U^*(f)U(s)\)), \(h*g\) would also be orthogonal to \(L\) for any function \(h \in \Cont(G)\), and one has \(h*g \in \Cont(G)\) and \(h*g \neq 0\) for suitable functions \(h\).\footnote{This is a slight simplification of Weyl's argument, which consists in obtaining for each function of \(\Cont(G)\) the analogue of the Fischer--Riesz expansion by an inductive application of the Schwarz--E. Schmidt method.}

The proof of the existence of a representation \(H\) such that \(U(f) \neq 0\) is deduced by Weyl from the theory of Hilbert--Schmidt integral equations. He considers the function \(g = f * \check{\overline{f}}\), for which \(U(g) = U(f)U^*(f)\), and it is enough to show that \(U(g) \neq 0\) for some \(U\). One has \(g(e) = \int_G |f(t)|^2 dt > 0\); Weyl forms the sequence of functions \(g_1 = g, g_2 = g*g, \ldots, g_n = g*g_{n-1}, \ldots\) by an adaptation of the method of H.A. Schwarz, as generalized by E. Schmidt to integral equations with symmetric kernels (chap. \ref{ch:3}, \ref{sec:3.1} and chap. \ref{ch:5}, \ref{sec:5.2}), he proves that the sequence of numbers \(\gamma_n = g_n(e)/g_{n-1}(e)\) is increasing and tends to a limit \(\gamma > 0\), and \(g_n/\gamma_n\) tends uniformly to a continuous function \(u\), such that \(g*u = u*g = \gamma u\) and \(u*u = u\); \(\gamma\) is an eigenvalue of the hermitian kernel \(k(s,t) = g(st^{-1})\), and if \(\varphi_j (1 \leq j \leq r)\) form an orthonormal basis of the corresponding eigenspace, one easily proves that \(u(st^{-1}) = \varphi_1(s) \overline{\varphi_1(t)} + \ldots + \varphi_r(s)\overline{\varphi_r(t)}\). Furthermore, for any \(t \in G\), the function \(s \mapsto \varphi_j(st^{-1})\) again is an eigenvector of the same space, hence \(\varphi_j(st^{-1}) = \sum\limits_{k=1}^r \overline{u_{ik}(t)}\varphi_k(s)\), and one shows that if \(U(t) = (u_{jk}(t)), t \mapsto U(t)\) is a linear representation of \(G\) for which \(U(g) \neq 0\).

H. Weyl himself remarked that this method also proved the existence of the irreducible representations \(M_k\) and was applicable to any compact Lie group (not necessarily semi-simple); a little later, when A. Haar had proved in 1933 the existence of a measure invariant by left and right translation on \emph{any} compact subgroup, Weyl's arguments could at once be extended to that general case.

The next year, Pontrjagin, in view of applications to algebraic topology, showed that the Peter--Weyl theory, applied to \emph{commutative metrizable compact groups}, led to a remarkable generalization of the duality between \emph{finite} commutative groups, which had been well-known since Weber. It was of course classical that an irreducible linear representation of a commutative group \(G\) must be of degree 1, in other words it is a character \(\chi\) of \(G\). The Peter--Weyl theory therefore associated to a metrizable compact commutative group \(G\) the set \(\hat{G}\) of all \emph{continuous} characters of \(G\), which of course is itself a denumerable group for ordinary multiplication. Now, for any \(x \in G\), the map \(\chi \mapsto \chi(x)\) is clearly a character of \(\hat{G}\), and Pontrjagin showed that \emph{all} characters of \(\hat{G}\) are of that type. Conversely, ff \(D\) is any denumerable group and \(\hat{D}\) the group of all its characters, one can put on \(\hat{D}\) the topology of simple convergence, for which it becomes a metrizable compact group, and then all \emph{continuous} characters of \(\stackrel{\vee}{G}\) are exactly the maps \(\chi \mapsto \chi(x)\) for all \(x \in D\).

But Pontrjagin went further and could extend this duality to some \emph{locally compact} commutative groups \cite{179}, and in 1936 van Kampen showed, by different methods, that Pontrjagin's results could be generalized to \emph{all} such groups \(G\). The dual \(\hat{G}\), consisting of all continuous characters on \(G\), is given the topology of uniform convergence on compact subsets of \(G\), and it is again locally compact for that topology; to each \(x \in G\) there corresponds the continuous character \(\eta(x): \chi \mapsto \chi(x)\) on \(\hat{G}\), and the duality theorem of Pontrjagin--van Kampen says that \(\eta\) is an \emph{isomorphism of topological groups} of \(G\) onto \(\Hat{\Hat{G}}\) \cite{216}.

This discovery made possible a unified treatment of the Fourier series and of the Fourier integral. In general, for any function \(f \in L^1(G)\), one could define the function \emph{on} \(\hat{G}\)
\begin{equation}
    \label{eq:7.62}
    \mathcal{F}f(\chi) = \int_G f(x)\overline{\chi(x)}dx
\end{equation}
as the \emph{Fourier transform} of \(f\). For \(G = \R\), continuous characters could be uniquely written \(x \mapsto \exp(2\pi ixy)\) for a real number \(y\), so that \(\hat{G}\) could be identified with \(\R\) itself, and \eqref{eq:7.62} was just the definition of the Fourier integral. For \(G = \mathbb{Z}\), all characters are continuous and can be written \(n \mapsto \zeta^n\) for a uniquely determined complex number \(\zeta \in U\); for a function \(n \mapsto c(n)\) on \(\mathbb{Z}\) such that \(\sum\limits_n |c(n)| < +\infty\), the right hand side of \eqref{eq:7.62} becomes the absolutely convergent Fourier series \(\sum\limits_{n \in \mathbb{Z}} c(n)\zeta^n = \sum\limits_{n \in \mathbb{Z}} c(n)e^{ni\theta}\) if \(\zeta = e^{i\theta}\), a function defined on the dual \(\mathbb{U}\) of \(\mathbb{Z}\). Finally, for \(G = \mathbb{U}\), continuous characters are the functions \(\zeta \mapsto \zeta^n\) for a uniquely determined \(n \in \mathbb{Z}\), and the Fourier transform of a function \(f \in L^1(\mathbb{U})\) is the sequence \(n \mapsto c(n)\) of its ``Fourier coefficients''. In 1940, A. Weil \cite{226} showed how most results concerning Fourier series and integrals could be generalized to all locally compact commutative groups; the central theorem was the generalization of the Parseval relation: if a function \(f\) on \(G\) belongs to \(L^1(G) \cap L^2(G)\), its Fourier transform belongs to \(L^2(\hat{G})\), and one has the relation
\begin{equation}
    \label{eq:7.63}
    \int_G |f(x)|^2 dx = \int_{\hat{G}} |\mathcal{F}f(\chi)|^2 d\chi
\end{equation}
for a suitable Haar measure on \(\hat{G}\); this relation, for the case \(G = \R\), had been proved in 1910 by M. Plancherel \cite{174}, and is known as the \emph{Plancherel theorem} for locally compact commutative groups.

We have already mentioned (\ref{sec:7.5}) that Gelfand, when he defined characters on a commutative Banach algebra, had followed the pattern set by H. Weyl and Pontrjagin. In fact, in a joint paper with D. Raikov \cite{086}, he immediately showed how the Pontrjagin--van Kampen--A. Weil theory could be deduced from his general results on Banach algebras in a much simpler way (the earlier proofs relied heavily on detailed information on the structure of locally compact commutative groups). The basic idea is to consider, for a locally compact commutative group \(G\), the space \(L^1(G)\) (for a Haar measure on \(G\)), on which a structure of Banach algebra is defined by the \emph{convolution} product \eqref{eq:7.58}; it is an \emph{involutive} algebra for the involution \(f \mapsto \check{\overline{f}}\), but in general it is not a \(C^*\)-algebra. A character of that algebra (in the sense of Gelfand) can then be uniquely written as \(f \mapsto \mathcal{F}f(\chi)\) for a well-determined character \(\chi\) (in the sense of Pontrjagin), so that the spectrum \(\Chi(L^1(G))\) is identified (with its topology induced by the weak topology of \(L^\infty(G)\)) with the dual group \(\hat{G}\), and then the Fourier transform merely becomes a special case of the Gelfand transform, equation \eqref{eq:7.44} being the expression in that special case of the fact that the Gelfand transform is a homomorphism of algebras!

But this absorption of harmonic Analysis by spectral theory did not stop with commutative groups. One can still define the Banach algebra \(L^1(G)\) for locally compact separable \emph{unimodular} groups (i.e. those for which left invariant Haar measure is also right invariant, for instance compact groups or semi-simple Lie groups), and it is still an involutive algebra. On the other hand, one can define linear representations \(s \mapsto U(s)\) of such a group \(G\) not only when \(U(s)\) is a unitary matrix, but more generally when \(U(s)\) is an \emph{automorphism of a complex Hilbert space} \(E\); one then speaks of \emph{unitary representations} of \(G\) in \(E\), and one adds to the definition the additional condition that for any \(x \in E\), the mapping \(s \mapsto U(s) \cdot x\) of \(G\) into the Hilbert space \(E\) should be \emph{continuous}. For any function \(f \in L^1(G)\), it is then possible to define \(U(f)\) as in \eqref{eq:7.59}, more precisely, one has, for any \(x \in E\) and \(y \in E\)
\begin{equation}
    \label{eq:7.64}
    (U(f) \cdot x | y) = \int_G f(s) (U(s) \cdot x | y) ds
\end{equation}
where \(ds\) is (left and right) invariant Haar measure on \(G\). It is then remarkable that starting from the unitary representations of \(G\) in \(E\), one obtains in this way a bijection of the set of these representations onto the set of all homomorphisms \(V\) of \(L^1(G)\) into the \(C^*\)-algebra \(\mathcal{L}(E)\) which satisfy \eqref{eq:7.61} and are non-degenerate (i.e. such that the \(V(f) \cdot x\) for \(x \in E\) and \(f \in L^1(G)\) generate a dense subspace of E). With convenient modifications, there is still a similar result for \emph{all} locally compact groups, and the general theory of unitary representations of locally compact groups in Hilbert space is thus in a certain sense subordinate to the theory of homomorphisms of involutive Banach algebras in algebras of operators in Hilbert space \cite{058}.

However, most results concerning unitary representations of locally compact groups have up to now been restricted to \emph{Lie groups}, where a large number of more refined and powerful tools (Lie algebras, differential geometry, partial differential equations, etc.) are available. We can only mention here this beautiful and difficult theory (known as \emph{non commutative harmonic Analysis}), which has known an enormous expansion since 1950, and in which many problems are still open; the interested reader is referred to \cite{032}, \cite{043}, \cite{152}, \cite{217} and \cite{223}; for a detailed history of harmonic Analysis (both commutative and non commutative) and its relations with probability theory, quantum mechanics and number theory, see \cite{155}.

D) \emph{Other developments}.

One of the first results on infinite dimensional representations was obtained by M.H. Stone in 1930 \cite{206}: he showed that any unitary representation of the additive group \(\R\) into a separable Hilbert space \(E\) was given by the formula \(t \mapsto e^{itA}\), where \(A\) is an arbitrary (in general unbounded) self-adjoint operator in \(E\), so that one may say that the theory of unitary representations of \(\R\) is \emph{equivalent} to the spectral theory of unbounded self-adjoint operators.

If now \(E\) is an arbitrary \emph{Banach space}, and \(A\) a \emph{bounded} operator in \(E\) it is clear that \(t \mapsto e^{itA}\) is a homomorphism of \(\R\) into the group of invertible elements in \(\mathcal{L}(E)\). Unbounded operators \(A\) may be defined in \(E\) just as in Hilbert space, but for such an operator, \(e^{tA}\) usually has no meaning for \emph{every} real number \(t\). Various questions of Analysis led E. Hille, in a series of papers beginning in 1936, to investigate mappings \(t \mapsto P_t\) into \(\mathcal{L}(E)\), \emph{only defined for} \(t > 0\) and such that, for \(s > 0\) and \(t > 0\)
\begin{equation}
    \label{eq:7.65}
    P_{s+t} = P_s P_t.
\end{equation}

Such mappings are called \emph{semi-groups of operators}. If one assumes that \(\norm{P_t} \leq C\) for all \(t > 0\) and that for every \(x \in E, t \mapsto P_t \cdot x\) is continuous for \(t > 0\), then one may associate to such a semi-group an \emph{unbounded operator} \(A\) in \(E\), defined by
\begin{equation}
    \label{eq:7.66}
    A \cdot x = \lim_{t \to 0} t^{-1} (P_t - 1_E) \cdot x.
\end{equation}

This has been the starting point of an extensive theory with many applications in Analysis \cite{115}.

A large literature has been devoted to various types of operators in Banach spaces. A very general method consists in starting with an operator \(A_0\) whose properties are well-known (for instance a normal operator in a Hilbert space) and to consider operators \(A = A_0 + P\) which differ from \(A_0\) by a ``perturbation'' \(P\) which is ``small'' in some sense; for instance, the norm \(\norm{P}\) is supposed to be small enough, or \(P\) is a compact operator; such assumptions allow in many cases to extend some properties of \(A_0\) to \(A\) (see \cite{124}).

The nice properties of the operators \(1_E + K\), where \(K\) is a compact operator (\ref{sec:7.1}) have inspired the study of generalizations of such operators, for instance the \emph{Fredholm operators} \(U\), which are defined by the properties that \(U^{-1}(0)\) has finite dimension, \(U(E)\) is closed and has finite codimension, but the dimension of \(U^{-1}(0)\) and the codimension of \(U(E)\) are not necessarily equal \cite{133}. Finally, there is an extensive theory of operators for which there is a family of \emph{projectors} having properties similar to the projectors \(E(\lambda)\) associated to a self-adjoint operator in von Neumann's theory (\ref{sec:7.4}); the difficulty is of course to find criteria implying the existence of such a family (\cite{062}, vol.~III).

\chapter{Locally convex spaces and the theory of distributions}
\label{ch:8}
\setcounter{equation}{0}

\section{Weak convergence and weak topology}
\label{sec:8.1}

In his thesis, Fréchet had already noticed that convergence in a metric space could not always correspond to some classical types of ``convergence'' for functions. For instance, if \(\B(\R)\) is the vector space of \emph{all} bounded real functions on \(\R\), it is not possible to define a distance on that space such that \emph{simple} convergence in \(\B(\R)\) would be identical with convergence for that distance. This results from the fact that if \(A\) is a subset of a metric space \(E\), the closure \(\overline{A}\) of \(A\) in \(E\) is identical to the set of limits of all convergent sequences of elements of \(A\). However, if one takes in \(\B(\R)\) the set \(A = \Cont(\R)\) of bounded continuous functions, the limits of sequences of elements of \(A\) for simple convergence are the Baire functions of class 1, and it is known that there are Baire functions of class 2 which are not of class 1, so that \(\overline{A}\) (for the hypothetical distance) could not consist only of functions of class 1 \cite[p.~115]{071}.

There was thus an obvious need for a generalization of the concept of metric space, but none proved adequate for Functional Analysis until Hausdorff, in 1914, created ``General topology'' as we understand it now, based on the concept of neighborhood \cite{100}; but surprisingly enough, it took some time to become aware of that adequacy. Ever since Hilbert, ``weak convergence'' of \emph{sequences} had become a central theme, first in Hilbert spaces, then with F. Riesz and Helly in some types of normed spaces (chap. \ref{ch:6}), and one would have thought that Hausdorff's concept of topology would have been tested on that notion; but until 1934 the only mathematician who seems to have had that idea was von Neumann: he defined weak neighborhoods of a point \(x_0\) in a Hilbert space \(E\) by a finite number of conditions \(|\left(x-x_0|a_j\right)| \leq \varepsilon\) for points \(a_j \in E\), and then went on to define similarly, in the algebra \(\L(E)\) of endomorphisms of \(E\), ``strong neighborhoods'' of an operator \(U_0\) by a finite number of conditions \(\norm{(U-U_0) \cdot x_j} \leq \varepsilon\) for \(x_j \in E\), and ``weak neighborhoods'' by a finite number of conditions \(\norm{\left( (U-U_0) \cdot x_j | y_j \right)} \leq \varepsilon\) \cite[vol.~II, p.~94--104]{221}. But he did not try to extend these ideas to other Banach spaces.

On the contrary, ``weak convergence'' was at the center of Banach's book, and the results he obtained concerning that notion can be considered as some of his deepest work. But to understand what he did, it is probably better first to state the final form which was taken by his 3 main theorems:

I) If \(E\) is a Banach space, and its dual \(E'\) is given the weak topology \(\sigma(E',E)\), the unit ball \(\norm{x} \leq 1\) in \(E'\) is compact for that topology.

II) In order that a vector subspace \(V \subset E'\) be closed for the topology \(\sigma(E',E)\), it is necessary and sufficient that for any closed ball \(B'\) in \(E'\), \(V \cap B'\) be compact for that topology.

III) In order that \(E\) be reflexive, it is necessary and sufficient that the unit ball \(\norm{x} \leq 1\) in \(E\) be compact for the weak topology \(\sigma(E,E')\).

In this form, the theorems were proved by N. Bourbaki in 1938 \cite{025}, and independently a little later by L. Alaoglu \cite{005}. Their proofs use the following ingredients:

a) The weak topology \(\sigma(E',E)\) is defined by taking as neighborhoods of \(x'_0 \in E'\) the sets defined by a finite number of relations \(|\langle x'-'_0, x_j\rangle| \leq \varepsilon\) for arbitrary \(x \in E\); \(\sigma(E,E')\) is defined similarly by exchanging the roles of \(E\) and \(E'\).

b) The word ``compact'' is used in the sense of N. Bourbaki, and means what was defined as ``bicompact sets'' (in Hausdorff spaces) by P. Alexandroff and P. Urysohn in 1924 \cite{006}: for them a space is ``bicompact'' if every open covering of the space contains a finite covering (the ``Borel--Lebesgue axiom'').

c) Compact sets can be characterized equivalently by means of the notion of limit of a ``net''', a notion which generalizes the limit of a sequence and was introduced in 1922 by E.H. Moore and H.L. Smith \cite{164} (N. Bourbaki uses the equivalent concept of limit of a ``filter'').

d) Any product of compact spaces is compact, a theorem proved by A. Tychonoff in 1930.

However, none of these notions or theorems was ever mentioned by Banach or mathematicians of his school until 1940, although they repeatedly quote Hausdorff's book of 1927 \footnote{The bulk of that book \cite{101} is devoted to metric spaces, and general topological spaces are given a very scanty treatment in 5 pages; it seems that Hausdorff had lost faith in his ideas of 1914!}; for them, the word ``compact'' is always taken in the initial sense of Fréchet, meaning a space in which there is no infinite closed discrete set. That notion is equivalent to the notion of ``bicompact'' space when restricted to metrizable spaces; the version of theorem I proved by Banach \cite[p.~123]{015} is therefore limited to \emph{separable} Banach spaces \(E\) (because in that case the ball \(\norm{x'} \leq 1\) is metrizable for the weak topology \(\sigma(E',E)\). It generalizes of course the ``principles of choice'' used by Hilbert, F. Riesz and Helly (chap. \ref{ch:5} and \ref{ch:6})

On the other hand, Banach was able to prove a theorem equivalent to Theorem II for \emph{all} Banach spaces. He starts from the study of vector subspaces of \(E'\) which are \emph{closed} for the topology of the norm, and observes that for such a subspace \(V\), a sequence \((x'_n)\) of points of \(V\) may have a weak limit which does not belong to \(V\); furthermore, if \(V_1\) is the vector space consisting of all these weak limits, it may happen that sequences of points of \(V_1\) have weak limits which do not belong to \(V_1\) and so on \cite[p.~209]{015}. Without speaking of weak topology, he then introduces the weakly closed vector subspaces, under the name of ``regularly closed'' subspaces: he defines such a space \(V\) by the property that for any \(x'_0 \not\in V\), there is an \(x \in E\) such that \(\ip{x',x} \neq 0\) for all \(x' \in V\), but \(\ip{x'_0, x} \neq 0\). Conscious of the fact that weakly convergent sequences are inadequate tools, Banach then introduces an \emph{ad hoc} notion, the ``limits of bounded transfinite sequences'' in the dual \(E'\): for any family \((u'_\xi)_{\xi < \gamma}\) of elements of \(E'\), contained in a ball and indexed by a segment of the ordinals, he shows (by using the Hahn--Banach theorem) that there always exist elements \(u' \in E'\) such that, for every \(x \in E\),
\[
    \Liminf_{\xi \to \gamma} \ip{u'_\xi, x} \leq \ip{u', x} \leq \Limsup_{\xi \to \gamma} \ip{u'_\xi, x}
\]
and calls any such \(u'\) ``a limit'' of the transfinite sequence \((u'_\xi)\) (there may be infinitely many such ``limits''!). He then says \(V\) is ``\emph{transfinitely closed}'' if every bounded transfinite sequence of elements of \(V\) has at least a ``\emph{transfinite limit}'' in \(V\), and what he shows, by a very clever argument, is that ``regularly closed'' and ``transfinitely closed'' are \emph{equivalent} notions \cite[p.~121]{015}. It was an easy matter to replace ``transfinite sequences'' by ``nets'' or ``filters'' in Banach's proof to obtain the equivalent statement of Theorem II.

Finally, by considering \(E\) as naturally imbedded in its second dual \(E''\) and using again ``transfinite limits'', Banach could prove Theorem III, but only when \(E\) is \emph{separable} \footnote{In 1938, Goldstine proved a result which is equivalent to the property that for any Banach space \(E\), the intersection \(E \cap B''\), where \(B''\) is the unit ball \(\norm{x''} \leq 1\) in the second dual \(E''\), is dense in \(B''\) for the weak topology \(\sigma(E'',E')\); from this Theorem III easily follows \cite{088}.}.

It should be mentioned here that these theorems enabled Banach to obtain a series of interesting theorems relating the properties of a continuous linear mapping \(u: E \to F\) (where \(E\) and \(F\) are Banach spaces), those of its transposed mapping \(\prescript{t}{}u: F' \to E'\), and properties of the images \(u(E)\) and \(\prescript{t}{}u(F')\) (see \cite{51}) \footnote{Some of Banach's results had been obtained by Hausdorff in 1931 \cite{102}; it is quite remarkable that he makes no mention of weak topology!}.

\section{Locally convex vector spaces}
\label{sec:8.2}

Although the theory of normed spaces was in the forefront of the development of Functional Analysis after 1906, it was soon realized that they did not exhaust the possibility of applying topological concepts to that discipline; but the various notions belonging to what we now call the general theory of \emph{topological vector spaces} made their appearance in a rather random way and were not the subject of a systematic treatment until 1950.

Already some examples of such spaces are to be found in Fréchet's thesis, where the emphasis is put, not on their algebraic properties but on the possibility of defining their topology by a distance (chap. \ref{ch:5}, \ref{sec:5.3}) and on the fact that the metric spaces thus obtained are complete. The fact that addition and multiplication by a scalar are continuous in such spaces was only explicitly emphasized by Fréchet in 1926 \cite{073}; the idea was picked up by Banach who in his book considered in general these spaces under the name of ``spaces of type \((F)\)'' and showed that the closed graph theorem was also valid for them. A little later, the method which Fréchet had used to define the distance on his examples of 1906 was systematized by S. Mazur and W. Orlicz in what they called the theory of ``spaces of type \((B_0)\)'' \cite{160}: they are what we now call \emph{Fréchet spaces}, where the topology is defined by a \emph{sequence} \((p_n)\) of seminorms with the condition that \(x \neq 0\) implies that \(p_n(x) \neq 0\) for at least one index \(n\); the distance can be defined by the formula
\[
    d(x,y) = \frac{p_1(x-y)}{1 + p_1(x-y)} + \frac{1}{2!}\frac{p_2(x-y)}{1 + p_2(x-y)} + \ldots + \frac{1}{n!}\frac{p_n(x-y)}{1 + p_n(x-y)} + \ldots
\]
and it is supposed that the space is complete for that distance. For the examples of Fréchet (the space \(\R^{\N}\) of all sequences and the space of holomorphic functions in \(|z| < 1\)) it can easily be shown that the topology \emph{cannot} be defined by a single norm.

Other types of spaces were not even metrizable; this was observed by von Neumann in 1929 for the weak topology on Hilbert space (\ref{sec:8.1}). But already in 1910 E.H. Moore had put forward the idea of replacing uniform convergence in \(\R\) by what he called ``relative uniform convergence''; this amounts to consider neighborhoods of 0 defined in the following way: one considers continuous functions \(g\) in \(\R\) such that \(g(x) > 0\) for all \(x \in \R\), and to each such function \(g\), one associates a neighborhood \(V_g\) of 0 consisting of all functions \(f\) such that \(|f| \leq g\); when restricted to functions which are continuous and have compact support, these neighborhoods are exactly those defining what will later be called the \((LF)\)-topology on \(\mathcal{H}(\R)\) \cite{163}.

After 1932, a new notion emerged, that of \emph{boundedness}. It was already realized by Banach that on the same vector space, two norms \(\norm{x}_1\) and \(\norm{x}_2\) such that the ratios \(\norm{x}_1 / \norm{x}_2\) and \(\norm{x}_2 / \norm{x}_1\) are bounded for \(x \neq 0\), define the same topology and therefore, if one defined a \emph{bounded set} in a normed space \(F\) as being contained in some ball, this was a notion independent of the particular norm chosen. However, in an arbitrary metric space, two distances may give rise to the same topology and give quite different notions of ``bounded sets'' when one sticks to the previous definition. But for arbitrary \emph{topological vector spaces} (i.e. those for which addition and multiplication by a scalar are continuous), it turned out that it is possible to give a definition of bounded set which coincides with the previous one for normed spaces, and only depends on the topology: a set \(A\) is \emph{bounded} if, for any sequence \((x_n)\) of points of \(A\), and any sequence \((t_n)\) of scalars tending to 0, the sequence \((t_n x_n)\) converges to 0. An elementary argument shows that this definition is equivalent to the following one: for any neighborhood \(V\) of 0, there is a scalar \(\lambda > 0\) such that \(\lambda A \subset V\). The first general result using this notion was the characterization of (Hausdorff) topological vector spaces for which the topology can be defined by a norm, found by A. Kolmogoroff in 1935 \cite{130}: they are those for which there exists a bounded neighborhood of 0.

Meanwhile, a new kind of topological vector spaces was introduced in 1934 by Köthe and Toeplitz \cite{129}. For any vector subspace \(E\) of the space \(\R^{\N}\) (or \(\C^{\N}\)) of all real (or complex) sequences, they consider the space \(E^*\) of all sequences \((u_n)\) in \(\R^{\N}\) (resp. \(\C^{\N}\)) such that \(\sum\limits_n |u_n x_n|\) converges (nowadays none says that \(E^*\) is the \emph{Köthe dual} of \(E\)). One can then consider the space \(E^{**}\), which obviously contains \(E\); when \(E^{**} = E\), Köthe and Toeplitz say \(E\) is \emph{perfect} (``vollkommen''), and it is this kind of space which is mainly studied in their paper, as well as in many subsequent papers of Köthe and his pupils (\cite{128}, \cite{130}). On such a space \(E\), they first define the weak topology \(\sigma(E,E^*)\) in the same way as von Neumann (\ref{sec:8.1}), neighborhoods of 0 being defined by a finite number of inequalities \(|\ip{x, a^*_j}| \leq 1\) with \(a^*_j \in E^*\) arbitrary. For that topology, they define \emph{bounded sets} as sets \(A \subset E\) such that each function \(x \mapsto \ip{x,a^*}\) (\(a^* \in E*\) arbitrary) is bounded in \(A\), which is of course a special case of the general notion mentioned above, although they introduce it without any reference. But their next step is particularly interesting; as \(E\) and \(E^*\) play symmetrical parts, one can also define the weak topology \(\sigma(E^*,E)\) and bounded sets in \(E^*\); this enables them to define \emph{on} \(E\) a new topology, the \emph{strong topology}, where neighborhoods of 0 are defined in the following way: to each bounded set \(B\) in \(E^*\), one associates the set \(V_B\) of all \(x \in E\) such that \(|\ip{x,y^*}| \leq 1\) for all \(y^* \in B\) (the ``polar set'' of \(B\) in a later terminology), and the \(V_B\) constitute a fundamental system of neighborhoods of 0 for the strong topology. One can then define bounded sets in \(E\) for that topology, and one of the chief results of Köthe and Toeplitz is that bounded sets in \(E\) are \emph{the same} for the weak and strong topology.

All topological vector spaces mentioned above belonged to what we now call \emph{locally convex spaces}, but the general definition of these spaces (under the name ``convex spaces'') was only given in 1935 by von Neumann, in view of a study of almost periodic functions \cite[vol.~II, p.~508--527]{221}. This coincided with a revival of interest in the properties of convex sets in topological vector spaces, which after Helly had been pretty much neglected: in Banach's book, they are only briefly mentioned in a Note at the end of the book \cite[p.~246]{015}. However, in 1933, S. Mazur gave the ``geometric'' version of the Hahn--Banach theorem, generalizing Minkowski's theory by showing that, if \(K\) is an open convex set in a normed space \(E\), there is a closed hyperplane of support of \(K\) through each boundary point of \(K\) \cite{158}. A little later, M. Krein and D. Milman introduced the concept of \emph{extreme point} for a convex set \(K\), i.e. a point of \(K\) such that there is no open line segment containing the point and contained in \(K\); they proved the remarkable fact that there are always ``enough'' extreme points for a \emph{compact convex} set \(K\), more precisely K is the smallest closed convex set containing all the extreme points \cite{131}; a theorem which was to have many important applications in various domains of Functional Analysis.

Once the locally convex spaces had been defined, the Köthe--Toeplitz procedure could be put in a more general context: one starts with \emph{two} vector spaces \(E, F\) and a \emph{bilinear form} \(B\) on \(E \times F\), which is non degenerate, i.e. such that the relation ``\(B(x,y) = 0\) for all \(y \in F\)'' is equivalent to \(x = 0\) and ``\(B(x,y) = 0\) for all \(x \in E\)'' equivalent to \(y = 0\). One then considers on \(E\) all Hausdorff locally convex topologies for which \(F\) is the dual of E; the determination of these topologies was done by G. Mackey \cite{154}, who showed that one gets a fundamental system of neighborhoods of 0 for such a topology by taking the finite intersections of the ``polar'' sets of a family \(\mathfrak{S}\) of subsets of \(F\), which consists of compact symmetric sets for the weak topology \(\sigma(F,E)\) and form a covering of \(F\); in addition, Mackey showed that for all these topologies, the bounded sets in \(E\) are the same.

We shall not try to describe in detail the very numerous papers devoted to topological vector spaces which have been published since 1950, Shortly after that date appeared the first comprehensive treatises on the subject (\cite{026}, \cite{062}, \cite{092}, \cite{125}, \cite{128}, \cite{215}). Most researches have been devoted to the study of particular types of locally convex spaces, such as Fréchet spaces and their direct limits (\cite{055}, \cite{092}), various types of ``vollkommen'' sequence spaces in the sense of Köthe--Toeplitz, which yield a rich harvest of examples and counterexamples, as well as many types of spaces consisting of functions with various properties. The most significant recent results concern the various topologies which one can define on the \emph{tensor product} \(E \otimes F\) of two locally convex spaces; they were studied in depth in a remarkable paper by A. Grothendieck, which deserves to be considered as realizing the greatest progress in Functional Analysis after the work of Banach \cite{091}; this study led its author to the discovery of a new class of locally convex spaces, the \emph{nuclear spaces}, which in a sense are much closer to finite dimensional spaces than even Hilbert spaces (with which they have some surprising connections) \cite{059}. Most spaces occurring in the theory of distributions (\ref{sec:8.3}) are nuclear spaces, and nuclear Fréchet spaces have become quite important in the theory of probability. Finally one should mention a large literature on convex sets in topological vector spaces, taking its origin in a beautiful result of Choquet giving to the Krein--Milman theorem a quantitative interpretation: if \(C\) is the convex hull of the union of \(\{0\}\) and a compact convex set \(K\) contained in a closed hyperplane not containing \(0\), then every point of \(C\) is the barycenter of a positive measure carried by the extreme points of \(C\), and this measure is unique if and only if the order relation defined by the cone \(C'\), union of all the \(\lambda C\) for \(\lambda > 0\), is a lattice \cite{042}. This result has important applications in potential theory.

\section{The theory of distributions}
\label{sec:8.3}

Between 1930 and 1940, several mathematicians began to investigate systematically the concept of ``weak'' solution of a linear partial differential equation, which we have seen appearing episodically (and without a name) in Poincaré's work (chap. \ref{ch:3}, \ref{sec:3.2}). in general, let \(P: f \mapsto \sum\limits_\alpha a_\alpha D^\alpha f\) be any differential operator with \(C^\infty\) coefficients, defined in an open set \(\Omega \subset \R^n\), and write \(\ip{f,g} = \int_\Omega f(x)g(x)dx\) for \(f\) locally integrable in \(\Omega\) and \(g\) continuous with compact support in \(\Omega\); then it is easy to generalize Lagrange's definition of the \emph{adjoint} differential operator \(\prescript{t}{}P\), which in particular satisfies
\begin{equation}
    \label{eq:8.1}
    \ip{P \cdot f, g} = \ip{f, \prescript{t}{}P \cdot g}
\end{equation}
when \(f\) is \(C^\infty\) in \(\Omega\), and \(g\) is \(C^\infty\) in \(\Omega\) with compact support. If \(f\) is a \(C^\infty\) solution of \(P \cdot f = 0\), we have therefore \(\ip{f,\prescript{t}{}P \cdot g} = 0\) for all functions \(g\) which are \(C^\infty\) in \(\Omega\) and have compact support. Conversely, any function \(f\) locally integrable in \(\Omega\) and having that property is called a \emph{weak} solution of the equation \(P \cdot u = 0\), even if it is not differentiable at all, and the problem which had confronted Poincaré was to prove that a weak solution is in fact a genuine \(C^\infty\) solution.

But in fact the same problem, for the simplest differential operator \(D = \frac{d}{dx}\), had already been considered and solved in the affirmative by P. Du Bois-Reymond in 1879 \cite{060}. Prodded by Weierstrass's criticism of the Calculus of variations (chap. \ref{ch:2}, \ref{sec:2.4}), he undertook to prove that if a \(C^1\) \emph{function} \(y\) in an interval \([a,b]\) is an extremum for the integral \(I(y) = \int_a^b F(x,y,y')dx\), where \(F\) is a \(C^1\) function, then the Euler equation
\begin{equation}
    \label{eq:8.2}
    \frac{d}{dx}\left( \pder{F}{y}(x,y,y') \right) - \pder{F}{y}(x,y,y') = 0
\end{equation}
\emph{makes sense}, which certainly is not obvious since nothing guarantees \emph{a priori} that \(\pder{F}{y'}(x, y, y')\) is differentiable! Following the classical procedure of Lagrange, one writes that for any \(C^1\) \emph{function} \(\zeta\) having compact support in \(]a,b[\), the function \(\varepsilon \mapsto I(y + \varepsilon \zeta)\) has an extremum for \(\varepsilon = 0\), which is equivalent to the relation
\begin{equation}
    \label{eq:8.3}
    \int_a^b \left( \pder{F}{y} \zeta + \zeta' \pder{F}{y'} \right) dx = 0;
\end{equation}
but instead of integrating by parts to eliminate \(\zeta'\), one instead eliminates \(\zeta\) by using the fact that \(\zeta(a) = \zeta(b) = 0\): an integration by parts enables indeed to write \eqref{eq:8.3} in the form
\begin{equation}
    \label{eq:8.4}
    \int_a^b \zeta'(x)f(x)dx = 0,
\end{equation}
where \(f(x) = \pder{F}{y'}(x,y(x),y'(x)) - \int_a^b \pder{F}{y}(t,y(t),y'(t))dt\) is only known to be \emph{continuous} in \(]a,b[\). The problem is to prove that \(f\) is a \emph{constant}, for then it will follow that \(\pder{F}{y'}(x,y(x),y'(x)) \) is indeed differentiable and equation \eqref{eq:8.2} is satisfied. However this amounts to showing that any weak solution of \(Du = 0\) is a constant, which is exactly what du Bois-Reymond proves\footnote{It is interesting to remark that in this paper du Bois--Reymond uses (probably for the first time) what we now call ``test functions'', i.e. \(C^\infty\) functions with compact support.}.

Such a result of course could not be expected for any differential operator: for instance, if \(A\) and \(B\) are \emph{any} locally integrable functions in \(\R\), the function \((x,y) \mapsto A(x) + B(y)\) is a weak solution of \(\frac{\partial^2 u}{\partial x \partial y} = 0\), for one has \(\int_{\R} A(x)dx \int_{\R} \frac{\partial^2 g}{\partial x \partial y} dy = 0\) and \(\int_{\R} B(y)dy \int_{\R} \frac{\partial^2 g}{\partial x \partial y} dx = 0\) for any \(C^2\) function \(g\) with compact support.

A step further would lead to defining a ``generalized'' operator \(P\), acting on functions which were not supposed differentiable at all: for a locally integrable function \(f\) in \(\Omega\), \(P \cdot f\) would be (by definition) a locally integrable function \(h\) such that, for \emph{any} \(C^\infty\) function \(g\) in \(\Omega\) with compact support, one has
\begin{equation}
    \label{eq:8.5}
    \ip{h,g} = \ip{f, \prescript{t}{}P \cdot g}.
\end{equation}
In a slightly different context, E. Cartan in 1922 \cite{039} had observed that it was sometimes possible to define an ``exterior derivative'' \(d\omega\) for a differential 2-form \(\omega = Pdy \wedge dz + Qdz \wedge dx + Rdx \wedge dy\), even if \(P, Q, R\) were merely continuous but not necessarily differentiable; one would define \(d\Omega = Sdx \wedge dy \wedge dz\) if \(S\) was a continuous function such that
\begin{equation}
    \label{eq:8.6}
    \iiint_V S dxdydz = \iint_\Sigma (Pdy \wedge dz + Qdz \wedge dx + Rdx \wedge dy)
\end{equation}
for any open set \(V\) with smooth boundary \(\Sigma\). As an example, he gave the form for which \(P = \pder{U}{x}, Q = \pder{U}{y}, R = \pder{U}{z}\), where \(U\) is the potential of a density \(\rho\) which is only supposed to be continuous; then \(P, Q, R\) need not be differentiable, but nevertheless \(S = -4\pi \rho\) satisfies \eqref{eq:8.6}.

The first systematic introduction of such ``generalized'' operators, for \(P = \pder{}{x_j}\) (under the name of ``quasi-dérivées'') is to be found in a paper of J. Leray in 1934 \cite{141} \footnote{Leray's results were rediscovered independently by K. Friedrichs in 1939 \cite{076}.}. In addition, Leray also introduces the process of \emph{regularization} of a locally integrable function \(f\) by a sequence \((\rho_n)\) of \(C^\infty\) functions with compact support tending to 0, such that \(\rho_n \geq 0\) and \(\int \rho_n dx = 1\): he shows that if \(f\) is continuous, \(p_n * f\) is a \(C^\infty\) function which converges uniformly to \(f\) in every compact subset, and if \(h\) is continuous and is the ``generalized derivative'' of \(f\), then \(\rho_n*h\) is the (usual) derivative of \(\rho_n*f\) \footnote{The study of integrals \(\rho_n * f\) for various types of sequences of functions \((\rho_n)\) was a favorite subject of analysts from Weierstrass to Lebesgue, under the name ``singular integrals''. For continuous functions \(\rho_n\) with compact support shrinking to a point, it had been systematically used by H. Weyl on Lie groups, as a substitute for the missing unit element in \(L^1(G)\) (\cite{227}, vol.~III, p.~73).}.

With our present knowledge, we realize that this notion of ``generalized derivative'' was a natural consequence of the use of the Lebesgue integral. Progressively, analysts had become familiar with the idea that two measurable functions which coincided except in a set of measure 0 were not to be distinguished from one another in most operations of Analysis. Not so, however, for differentiation: if \(f\) is a \(C^\infty\) function in \(\R\) and \(\varphi\) is the characteristic function of the set of rationals, \(f + \varphi\) is almost everywhere equal to \(f\), but \(f + \varphi\) has no derivative at any point, in the usual sense! Nevertheless it has of course a ``generalized derivative'' equal to \(f'\), and this could throw doubts on the adequacy of the ``natural'' definition of a derivative in Analysis!

It is easy to see that if a function \(f\) in \(\R\) has a ``generalized derivative'' \(h\) which is locally integrable, then \(f\) is almost everywhere equal to \(\int_0^x h(t)dt + c\), where \(c\) is a constant. This shows that a continuous function may have \emph{almost everywhere} a derivative in the usual sense, without having a ``generalized derivative'', for instance an increasing function \(f\) which is not absolutely continuous, another example of the inadequacy of the concept of derivative in the classical sense.

\emph{A fortiori}, this also shows that a function which is discontinuous at a point of \(\R\) cannot have a ``generalized derivative''. Nevertheless, following Dirac, theoretical physicists did not hesitate to consider that the Heaviside function \(Y\), equal to 0 for \(x < 0\) and to 1 for \(x \geq 0\), had a ``generalized derivative'', the so-called ``Dirac function'' \(\delta\), which would have been equal to 0 for \(x \neq 0\), but such that \(\int \delta(x)dx = 1\); and they even introduced successive ``derivatives'' \(\delta', \delta'', \ldots\) of that ``function'', writing ``equations'' such as
\begin{equation}
    \label{eq:8.7}
    \int g(x)\delta^{(n)}(x-a) = g^{(n)}(a)
\end{equation}
for a \(C^n\) function \(g\), or
\begin{equation}
    \label{eq:8.8}
    \int \delta'(a-x) \delta^{(n)}(x-b) dx = \delta^{(n+1)}(a-b)
    \qquad
    \text{\cite{056}}.
\end{equation}

or some time, mathematicians were puzzled by such manipulations, which eventually led to correct statements on genuine functions. The decisive step was taken in 1936 by S. Sobolev \cite{200}: the outcome of these jugglings with non-existent ``functions'' was finally to define perfectly decent \emph{linear forms} such as \(f \mapsto f^{(n)}(a)\) on the vector space \(\mathcal{D}\) of all \(C^\infty\) functions with compact support defined in an open set \(\Omega \subset \R^n\); Sobolev's idea was therefore to deal \emph{directly} with such linear forms, provided one could characterize them by properties involving only genuine mathematics. As he was led to this idea by a very concrete question, the solution of Cauchy's problem for second order hyperbolic equations with general boundary conditions (see chap. \ref{ch:9}, \ref{sec:9.5}), he could see what kind of properties he needed, and give a general characterization of what he called ``functionals'' on \(\mathcal{D}(\Omega)\), which we now call (after L. Schwartz) \emph{distributions} on \(\Omega\): for each compact subset \(K \subset \Omega\) one considers the subspace \(\mathcal{D}(\Omega; K)\) of \(\mathcal{D}(\Omega)\) consisting of all \(C^\infty\) functions with support in \(K\), and this is a Fréchet space for the semi-norms
\begin{equation}
    \label{eq:8.9}
    p_{m,K}(f) = \sup_{|\alpha| \leq m, x \in K}(D^\alpha f(x));
\end{equation}
distributions are then the linear forms \(T\) on \(\mathcal{D}(\Omega)\), the restriction of which to each subspace \(\mathcal{D}(\Omega; K)\) is \emph{continuous} for the preceding topology\footnote{Sobolev does not speak of topology, but defines convergent sequences in \(\mathcal{D}(\Omega)\) which correspond to these topologies on the spaces \(\mathcal{D}(\Omega; K)\). Another way of expressing the definition is to consider the ``direct limit'' of the topologies of the spaces \(\mathcal{D}(\Omega; K)\); distributions are then the elements of the \emph{dual} of \(\mathcal{D}(\Omega)\) when \(\mathcal{D}(\Omega)\) is given that topology.}. Any locally integrable function \(F\) in \(\Omega\) defined a distribution \(f \mapsto \int_\Omega F(x)f(x)dx\), two almost everywhere functions giving rise to the same distribution (which of course were the measures on \(\Omega\) having a density with respect to Lebesgue measure), so that the space \(L^1_\text{loc}(\Omega)\) of classes of locally integrable functions was identified with a subspace of the space \(\mathcal{D}'(\Omega)\) of all distributions; but more generally all \emph{Radon measures} on \(\Omega\) were particular distributions, and in particular one could define correctly the so-called ``Dirac function'' \(x \mapsto \delta(x-a)\) as the measure \(\varepsilon_a: f \mapsto f(a)\) defined by the mass +1 at the point \(a\). Sobolev pointed out that one can multiply a distribution \(T\) by any \(C^\infty\) function \(g\) in \(\Omega\), by defining \(g \cdot T\) as the distribution \(f \mapsto T(gf)\); more important still, one may define the \emph{derivatives} \(\pder{T}{x_i}\) of \emph{any} distribution \(T\) as the distribution \(f \mapsto -T\left(\pder{f}{x_i}\right)\). Finally, he considered on the space \(\mathcal{D}'(\Omega)\) the weak topology \(\sigma(\mathcal{D}'(\Omega), \mathcal{D}(\Omega))\), and showed that the regularization process could also be applied to distributions: \(\rho_n * T\) is defined as the distribution \(f \mapsto T(\check{\rho}_n * f)\), which turns out to be the class of a \(C^\infty\) function, and \(\rho_n*T\) converges weakly to \(T\) when \(n\) tends to \(+\infty\); the fact that distributions are thus limits (for the weak topology) of \(C^\infty\) functions has led some mathematicians to call them ``generalized functions'' \cite{086}.

During the same period, the need to ``enlarge'' in some way the domain of definition of operators other than differential operators was also felt in different parts of Analysis\footnote{For instance, in the Calculus of variations, one may consider that a smooth \(p\)-dimensional variety \(V\) in an \(\R^n\) defines a linear ``functional'' \(\omega \int_V \omega\) in the vector space of differential \(p\)-forms on \(\R^n\). This leads to the idea of ``generalized varieties'' \cite{231} and of ``currents'' \cite{49}.}, and particularly in classical harmonic Analysis. The definition of the Fourier transform \(\mathcal{F}f\) of a function \(f\) defined in \(\R^n\) only makes sense when \(f \in L^1(\R^n)\); however, as soon as 1910, the Plancherel theorem (chap. \ref{ch:7}, \ref{sec:7.6}) showed that it is possible to define the operator \(f \mapsto \mathcal{F}f\) as an \emph{isometry} of the Hilbert space \(L^2(\R^n)\) onto itself, by extending it by continuity from its original domain of definition \(L^1 \cap L^2\); in other words, for a function \(f \in L^2\) which did not belong to \(L^1\), the Fourier transform \(\mathcal{F}f\) could still be defined, but only by a limit process. Later, efforts were made to define similarly a Fourier transform \(\mathcal{F}f\) for functions \(f\) belonging to other spaces \(L^p\); in his discussion of that problem, A. Weyl observed in 1940 \cite[p.~118]{226} that if \(\Lambda\) is the space of functions \(f \in L^1 \cap L^\infty\) such that \(\mathcal{F}f\) also belongs to \(L^1 \cap L^\infty\), then if two functions \(\Phi, \varphi\) are such that
\begin{equation}
    \label{eq:8.10}
    \int \Phi(x) \cdot \overline{\mathcal{F}f(x)} dx = \int \varphi(x) \cdot \overline{f(x)} dx
\end{equation}
for \emph{all} functions \(f \in \Lambda\), it is legitimate to consider that \(\Phi\) is the Fourier transform of \(\varphi\).

Much earlier, in 1911, H. Weyl, in relation with his work on second order linear differential equations (chap. \ref{ch:7}, \ref{sec:7.6}), had observed that if \(f\) is such that \(f(x)/(1 + |x|)\) is integrable in \(\R\) (for instance, if \(xf(x)\) is bounded) and satisfies an additional regularity condition, then it is possible to write for \(f\) the Fourier inversion formula, provided one replaces \(\mathcal{F}f\) by a Stieltjes measure (\cite{227}, vol.~I, p.~359--360); this amounts to defining the Fourier transform of a bounded Stieltjes measure, a definition which was explicitly given by P. Daniell in 1920 \cite{046}, and which became later a favorite tool of probabilists. Weyl's idea was developed by Hahn \cite{096} and N. Wiener \cite{0229}, and then extended by S. Bochner to functions such that \(f(x)/(1+|x|^k)\) is integrable, where \(k\) is an \emph{arbitrary} integer \cite{024}. Using the fact that by Fourier transformation derivation becomes multiplication by \(x\) (up to a constant), Bochner proceeds as Riemann had done for trigonometric series \cite[p.~245]{182}: in order to obtain an integrable function, he subtracts from \(e^{-ix\xi}\) a function \(L_k(x\xi)\) equal to the first \(k\) terms \(\sum\limits_{k=0}^{m-1} \frac{1}{m!} (-ix\xi)^m\) of the power series expansion of \(e^{-ix\xi}\) in a compact neighborhood of 0, and to 0 outside, and writes
\begin{equation}
    \label{eq:8.11}
    E(\xi, k) = \frac{1}{2\pi} \int_0^\infty f(x) \frac{e^{-ix\xi} -L_k(x\xi)}{(-ix)^k} dx.
\end{equation}
This is of course only defined up to a \emph{polynomial} in \(\xi\) of degree \(\leq k-1\); Bochner's idea would be to take a ``derivative'' in some sense of \(E(\xi, k)\) as ``Fourier transform'' of \(f\), and indeed he writes ``symbolically''
\[
    f(x) \sim{} \int e^{ix\xi} d^k E(\xi, k)
\]
what would be the ``inversion formula''; but as he has no definition of such a ``\(k\)-th derivative'' at his disposal, he is compelled to work only with the \emph{functions} \(E(\xi, k)\) in the applications he gives to difference equations.

It was one of the main contributions of L. Schwartz that he saw, in 1945 \cite{194}, that the concept of distribution introduced by Sobolev (which he had rediscovered independently) could give a satisfactory generalization of the Fourier transform including all the preceding ones. Instead of considering the space \(\Lambda\) introduced by A. Weyl, which is not easy to describe explicitly, he had the idea to take as ``test functions'' the \(C^\infty\) functions \(f\) in \(\R^n\) which are such that \(f\) and \emph{all} its derivatives are ``rapidly decreasing at infinity'', i.e. such that their product with any polynomial is integrable. The essential property of that space \(\mathcal{S}(\R^n)\) of ``declining functions'' is that the Fourier operator \(f \mapsto \mathcal{F}f\) is a \emph{bijection} of \(\mathcal{S}(\R^n)\) with the Fréchet topology defined by the semi-norms
\begin{equation}
    \label{eq:8.12}
    q_{s,m}(f) = \sup_{|\alpha| \leq s, x \in \R^n} (1+|x|)^m |D^\alpha f(x)|.
\end{equation}
It is easy to see that for each compact subset \(K\) of \(\R^n\), the space \(\mathcal{D}(\R^n; K)\) is contained in \(\mathcal{S}(\R^n)\) and the injection \(\mathcal{D}(\R^n; K) \to \mathcal{S}(\R^n)\) is continuous; furthermore the union \(\mathcal{D}(\R^n)\) of all the \(\mathcal{D}(\R^n; K)\) is dense in \(\mathcal{S}(\R^n)\). The continuous linear forms on \(\mathcal{S}(\R^n)\) can thus be considered as special distributions, which Schwartz calls \emph{tempered distributions}. The Fourier transform \(T \mapsto \mathcal{F}T\) in the space \(\mathcal{S}(\R^n)\) of tempered distributions (dual of \(\mathcal{S}(\R^n)\)) is then defined (by a generalization of equation \eqref{eq:8.10}) as the \emph{transposed} automorphism of the Fourier transform in \(\mathcal{S}(\R^n)\), in other words the Fourier transform \(\mathcal{F}T\) of a tempered distribution \(T\) is defined by the relation
\begin{equation}
    \label{eq:8.13}
    \ip{\mathcal{F}T, f} = \ip{T, \mathcal{F}f}
    \text{ for all }
    f \in \mathcal{S}(\R^n).
\end{equation}

To the credit of L. Schwartz must be added his persistent efforts to weld all the previous ideas into a unified and complete theory, which he enriched by many definitions and results (such as those concerning the tensor product and the convolution of distributions) in his now classical treatise \cite{194}. By his own research and those of his numerous students, he began to explore the potentialities of distributions, and gradually succeeded in convincing the world of analysts that this new concept should become central in all linear problems of Analysis, due to the greater freedom and generality it allowed in the fundamental operations of Calculus, doing away with a great many unnecessary restrictions and pathology \footnote{The role of Schwartz in the theory of distributions is very similar to the one played by Newton and Leibniz in the history of Calculus: contrary to popular belief, they of course did not invent it, for derivation and integration were practiced by men such as Cavalieri, Fermat and Roberval when Newton and Leibniz were mere schoolboys. But they were able to systematize the algorithms and notations of Calculus in such a way that it became the versatile and powerful tool which we know, whereas before them it could only be handled \emph{via} complicated arguments and diagrams (see \cite{028}).}. One should reserve a particular mention to what is probably the most original of his contributions, the ``kernel theorem'' \cite{195}. Ever since Hilbert's and F. Riesz's work, it had been realized that integral operators \(f \mapsto K \cdot f\) defined by a ``kernel function'' \(K(x,y)\), as \((K \cdot f)(x) = \int K(x,y)f(y)dy\), were very far from exhausting the general concept of linear operator, since not even the identity could be expressed in that manner! It is therefore very remarkable that if one replaces ``kernel functions'' by ``kernel distributions'' in that definition, one practically obtains \emph{all} linear operators which one meets in problems of Analysis. More precisely, if \(X \subset \R^n\) and \(Y \subset \R^n\) are open sets, any linear mapping \(K\) of \(\mathcal{D}(X)\) into \(\mathcal{D}'(Y)\), which is only supposed to yield partial continuous mappings \(\mathcal{D}(X; L) \to \mathcal{D}'(Y)\) for any compact subset \(L \subset X\) (when \(\mathcal{D}(X; K)\) is given its Fréchet topology and \(\mathcal{D}'(Y)\) the weak topology), can be defined by a uniquely determined ``kernel distribution'' \(K \in \mathcal{D}'(X \times Y)\), in such a way that for any \(u \in \mathcal{D}(X)\), the \emph{distribution} \(K \cdot u\) satisfies
\begin{equation}
    \label{eq:8.14}
    \ip{K \cdot u, v} = \ip{K, u \otimes v}
\end{equation}
for any function \(v \in \mathcal{D}(Y)\). The great interest of this result lies in the fact that most spaces \(E\) of functions defined in \(X\) are such that \(\mathcal{D}(X) \subset X \subset \mathcal{D}'(X)\), the injections \(\mathcal{D}(X; L) \to E\) and \(E \to \mathcal{D}'(X)\) being continuous; if \(A\) (resp. \(B\)) is such a space of functions defined in \(X\), and \(U: A \to B\) a continuous linear mapping, the composed map \(\mathcal{D}(X; L) \to A \xrightarrow{U} B \to \mathcal{D}'(X)\) is continuous, hence is defined by a ``kernel distribution''. For instance, the identity map \(A \to A\) is defined by the distribution \(I\) which is a measure carried by the diagonal \(\Delta_X\) in \(X \times X\) and is such that
\[
    \int_{X \times X} w(x,y) dI(x,y) = \int_X w(x,x) dx.
\]

\chapter{Applications of functional analysis to differential and partial differential equations}
\markboth{Applications of functional analysis}{Applications of functional analysis}
\label{ch:9}
\setcounter{equation}{0}

I will not try to enumerate all the applications which have been made of Functional Analysis in the last 50 years, and which have amply justified the creators of that discipline. But as we have seen in the first 4 chapters how most notions and problems of Functional Analysis had their origin in questions relative to ordinary or partial differential equations. I think it is worthwhile to give a sketchy description of \emph{a few} of the most conspicuous progress in those questions which have been made by an imaginative use of the new tools provided by Functional Analysis, mostly spectral theory and the theory of distributions.

\section{Fixed point theorems}
\label{sec:9.1}

In the first applications which we shall mention, however, little more is used of Banach spaces beyond their definition, and the results primarily concern non linear equations. The main idea is similar to the application of the contraction principle (chap. \ref{ch:6}, \ref{sec:6.3}) to the local existence theorems for differential equations, by writing them in the form \(z = F(z)\) for \(z\) in some Banach space \(E\) of functions; if \(B\) is the closure of an open bounded convex set in \(E\) and \(F\) is a contraction mapping \(B\) into itself, then the contraction principle says there exists a \emph{unique} solution of \(z = F(z)\) in \(B\), what one calls a \emph{fixed point} for \(F\). But after the first years of the XX\supth{} century, new possibilities of obtaining ``fixed point theorems'' appeared with the first results of a new branch of mathematics, Algebraic Topology, created by H. Poincaré in 1895--1900: using the concepts of that theory, L.E.J.Brouwer could show in 1910 that if \(B\) is homeomorphic to a closed ball in some \emph{finite dimensional} vector space \(E\), and \(F\) is any continuous map of \(B\) into itself (which is not supposed any more to be a contraction), then \(F\) has \emph{at least} one fixed point in \(B\).

The problem was to find a similar theorem applicable to \emph{infinite dimensional} Banach spaces \(E\). The first result in that direction was obtained in 1922 by G.D. Birkhoff and O. Kellogg, who considered the case in which \(E = \Cont(I)\) or \(E = l^2\), and showed that Brouwer's theorem could be extended, provided one took for \(B\) a \emph{compact convex set} \cite{022}; their fundamental device consists in using the compactness of \(B\) to ``approximate'' it by a finite dimensional compact convex set \(B_n\), and to similarly ``approximate'' \(F\) by a continuous mapping \(F_n\) of \(B_n\) into itself, to which Brouwer's theorem may be applied. This method was taken up and greatly expanded by J. Schauder, who showed that it could be applied to any Banach space \(E\), and also, for separable Banach spaces, that one could replace compactness of \(B\) and continuity of \(F\) by weak compactness and weak continuity (\cite{186}, \cite{187}). This enabled him to prove, for instance, existence of a solution of the equation
\begin{equation}
    \label{eq:9.1}
    \Delta z = f\left(x, y, z, \pder{z}{x}, \pder{z}{y}\right)
\end{equation}
in a domain \(\Omega \subset \R^2\) with smooth boundary (no connected component of which is reduced to a point), vanishing at the boundary, under the only assumption that \(f\) is bounded and continuous for bounded values of the 5 variables on which it depends; the method consists in transforming \eqref{eq:9.1} into an integro-differential equation
\begin{equation}
    \label{eq:9.2}
    z(x,y) = \iint_\Omega G(x, y, \xi, \eta) f(\xi, \eta, z(\xi, \eta), \pder{z}{\xi}, \pder{z}{\eta}) d\xi d\eta
\end{equation}
where \(G\) is the Green function for \(\Omega\).

A little later, a much more sophisticated approach enabled Schauder to solve Cauchy's problem locally for quasi-linear hyperbolic equations\footnote{The equation is supposed to be ``normal'', which means that the left hand side is such that the quadratic form \(\sum_{i,k} A_{ik} \xi_i \xi_k\) has signature \((1,n-1)\).}
\begin{equation}
    \label{eq:9.3}
    \sum_{i,k} A_{ik} \left( x_1, \ldots, x_n, z, \pder{z}{x_1}, \ldots, \pder{z}{x_n} \right) \frac{\partial^2 z}{\partial x_i \partial x_k} = A \left( x_1, \ldots, x_n, z, \pder{z}{x_1}, \ldots, \pder{z}{x_n} \right).
\end{equation}
His method consists, for a \emph{given} function \(z(x_1, \ldots, x_n)\), to solve the Cauchy problem for the linear hyperbolic equation in the unknown function \(Z\)
\begin{equation}
    \label{eq:9.4}
    \sum_{i,k} A_{ik} \left( x_1, \ldots, x_n, z, \pder{z}{x_1}, \ldots, \pder{z}{x_n} \right) \frac{\partial^2 Z}{\partial x_i \partial x_k} = A \left( x_1, \ldots, x_n, z, \pder{z}{x_1}, \ldots, \pder{z}{x_n} \right).
\end{equation}
When \(z\) is in a suitable set \(K\), the problem has a unique solution \(Z(z)\), and existence of a solution \(z\) of \eqref{eq:9.3} in \(K\) will be obtained if one shows that the equation \(Z(z) = z\) has a solution in \(K\). The problem consists in choosing \(K\) such that the extension of Brouwer's fixed point theorem is applicable; the main point is to obtain \emph{a priori inequalities} for solutions of ``normal'' linear hyperbolic equations
\begin{equation}
    \label{eq:9.5}
    \begin{gathered}
    L(u) \equiv \sum_{i,k} A_{ik}(x_1, \ldots, x_n) \frac{\partial^2 u}{\partial x_i \partial x_k} + \sum_j B_j(x_1, \ldots, x_n) \pder{u}{x_j} \\
    + C(x_1, \ldots, x_n)u = F(x_1, \ldots, x_n)
    \end{gathered}
\end{equation}
defined in a truncated pyramid \(P\) having its larger base \(B\) in \(\R^{n-1}\). Using a method first introduced in 1926 by Friedrichs and H. Lewy, which consists in transforming the integral \(\int_P \pder{u}{x_n} L(u)d\omega\) by integration by parts and Stokes' \(P\) formula, Schauder obtains (for suitable restrictions on \(P\) and the coefficients \(A_{ik}, B_j\) and \(C\) in \eqref{eq:9.5}) an inequality
\begin{equation}
    \label{eq:9.6}
    \begin{gathered}
    \int_P \left( \left(\pder{u}{x_1}\right)^2 + \ldots + \left(\pder{u}{x_n}\right)^2 \right) d\omega \leq \\
    \leq M \left( \int_B \left( u^2 + \left(\pder{u}{x_1}\right)^2 + \ldots + \left(\pder{u}{x_n}\right)^2 \right) d\sigma + \int_P F^2 d\omega \right)
    \end{gathered}
\end{equation}
with a constant \(M\) independent of \(u\); by differentiation he gets similar inequalities for all derivatives (of any order) of \(u\). The space \(E\) is then defined by the norm \(\sup\limits_{x \in P} \left( \sum\limits_{|\alpha| \leq r} |D^\alpha f(x)| \right)\) for a suitable value of \(r\), but the set \(K\) is a ball in \(E\) for \emph{another norm}, namely \(\left( \int_P \left( \sum_{|\alpha| \leq s} |D^\alpha f(x)|^2 \right) dx \right)^{1/2}\) for another value of \(s\); his \emph{a priori} inequalities enable then Schauder to show that \(K\) is compact in \(E\) \footnote{This is probably the first time the norm of the space \(H^s\) makes its appearance.} and \(z \mapsto Z(z)\) continuous in \(K\) \cite{190}.

Another of the famous theorems proved by Brouwer was the \emph{invariance of domain}: if \(\Omega\) is an open subset of \(\R^n\), and \(F\) an \emph{injective} continuous map of \(\Omega\) into \(\R^n\), then the image \(F(\Omega)\) is again an open subset. In 1929, Schauder showed that the theorem was still true in some types of Banach spaces, for maps \(F\) of the type \(x \mapsto x + H(x)\), where \(H\) is \emph{completely continuous} in the sense of F. Riesz, but not necessarily linear. From this he deduced for instance that if one knows that the equation
\begin{equation}
    \label{eq:9.7}
    \Delta z - f \left( x, y, z, \pder{z}{x}, \pder{z}{y} \right) = \psi(x, y)
\end{equation}
has \emph{at most} a solution taking given values \(\varphi(s)\) at the boundary of \(\Omega \subset \R^2\), then, if for given functions \(\varphi_0, \psi_0\) there \emph{exists} such a solution, the same is true for functions \(\varphi, \psi\) sufficiently close to \(\varphi_0, \psi_0\) \cite{189}.

But the most sophisticated application of Algebraic Topology to functional equations was made in the famous 1934 paper by J. Leray and J. Schauder \cite{142}. If \(U\) is an open set in \(\R^n\), such that \(\overline{U}\) is compact, and \(f\) is a continuous mapping of \(\overline{U}\) in \(\R^n\), then, for each point \(z \in \R^n\) which does not be long to the image by \(f\) of the boundary of \(U\), Brouwer had shown that one may attach an integer \(d(f,U,z)\) which only depends on the connected component of \(\R^n - f(\text{Fr}(U))\) to which \(z\) belongs, and varies continuously with \(\lambda\) when \(f(x) = F(x, \lambda)\) where \(F\) is continuous and \(\lambda\) a real parameter; furthermore, when \(d(f,U,z) \neq 0\), the inverse image \(f^{-1}(z)\) is not empty. Leray and Schauder were able, by approximating compact subsets of a Banach space by finite dimensional subsets, to prove, by application of this theorem of Brouwer, the following existence theorem. Let \(E\) be a Banach space, \(I \subset \R\) a compact interval, \(\Omega \subset E \times I\) a bounded open set in \(E \times I, F: \overline{\Omega} \to E\) a mapping which is completely continuous, and in addition uniformly continuous. One assumes that for every \(\lambda \in I\), there are no solution of the equation \(x - F(x, \lambda)) = 0\) in the boundary of \(\Omega\), and that, for \emph{one} value \(\lambda_0 \in I\), the equation \(x - F(x, \lambda_0) = 0\) has exactly one solution in \(\Omega\); then, for \emph{every} \(\lambda \in I\), there exists \emph{at least one} solution of \(x - F(x, \lambda) = 0\) in \(\Omega\).

The application of that theorem to partial differential equations usually necessitates subtle \emph{a priori} inequalities which guarantee that all the assumptions of the theorem are satisfied.

\section{Carleman operators and generalized eigenvectors}
\label{sec:9.2}

For applications to partial differential equations, it is necessary to generalize the notion of Carleman operators defined in chap. \ref{ch:5}, \ref{sec:5.3}. Let \(X\) be a locally compact metrizable and separable space, \(\mu\) a positive measure on \(X\) and \(H\) a separable Hilbert space (finite dimensional or not); let \((a_n)_{n \in J}\) be a Hilbert basis of \(H\), where \(J\) is a finite or denumerable set. One defines a new Hilbert space \(L^2_H(X,\mu)\) as the space of \emph{vector valued} functions \(f: x \mapsto \sum\limits_{n \in J} f_n(x)a_n\), mappings of \(X\) into \(H\) such that each \(f_n\) is a function of \(L^2(X,\mu)\) with complex values, and that \(|f|^2 = \sum\limits_{n \in J} |f_n|^2\) is \(\mu\)-integrable; the scalar product in that Hilbert space is then defined by
\begin{equation}
    \label{eq:9.8}
    (f|g) = \sum_{n \in J} \int_X f_n \overline{g}_n d\mu.
\end{equation}
Let now \(Y\) be another locally compact metrizable and separable space, \(\nu\) a positive measure on \(Y\). A \emph{Carleman kernel} on \(X \times Y\) (for the measure \(\mu \otimes \nu\) and the Hilbert space \(H\)) is then a mapping \(K: (x,y) \mapsto (K_n(x,y))_{n \in J}\) of \(X \times Y\) into \(\C^J\) such that:

1° each complex function \(K_n\) is \((\mu \otimes \nu)\)-measurable;

2° there is a null set \(N \subset Y\) such that, for each \(y \not\in N\), the function \(x \mapsto K_n(x,y)\) is \(\mu\)-measurable and the function \(x \mapsto \sum\limits_{n \in J} |K_n(x, y)|^2\) is \(\mu\)-integrable.

If \(f = \sum\limits_{n \in J} f_n a_n\) is a function of \(L^2(X,\mu)\), the function \(x \mapsto \sum\limits_{n \in J} K_n(x,y) f_n(x)\) is \(\mu\)-integrable for all \(y \in N\), and the function
\begin{equation}
    \label{eq:9.9}
    g(y) = \int_X \sum_{n \in J} K_n(x, y) f_n(x) d\mu(x)
\end{equation}
defined in \(Y-N\), is \(\nu\)-measurable. One writes \(K \cdot f = g\), and \(K\), defined in \(L^2_H(X,\mu)\) is called the \emph{Carleman operator} defined by the Carleman kernel \(K\).

In 1952, F. Mautner discovered that Carleman operators can be characterized by properties which are independent of the definition by a kernel \cite{157}: suppose that there is a null set \(N \subset Y\) and, for each \(y \in Y - N\), a \emph{continuous linear map} \(F_Y: L_H^2(X, \mu) \to \C\) such that for any function \(f \in L^2_H(X,\mu)\), the map \(y \mapsto F_y(f)\) is \(\mu\)-measurable. Then there is a Carleman kernel \(K = (K_n)\) and a null set \(N' \supset N\) such that, for any \(y \not\in N', F_y(f) = (K \cdot f)(y)\) for all functions \(f \in L^2_H(X,\mu)\).

We have seen (chap. \ref{ch:7}, \ref{sec:7.5}) that a continuous normal operator has a ``diagonalization'' which transforms it into multiplication by the function ``identity'' \(\zeta \mapsto \zeta\). More generally, one defines a \emph{diagonalization} of a continuous normal operator \(B\) in a Hilbert space \(E\) for a function \(\zeta mapsto \Phi(\zeta)\) as in chap. \ref{ch:7}, \ref{sec:7.5}, by replacing \(\spec(B)\) by a separable locally compact space \(Y\), the function \(\zeta \mapsto \zeta\) being replaced by a mapping \(\zeta mapsto \Phi(\zeta)\) of Y into \(\C\).

With the help of his characterization of Carleman operators, Mautner was able to get a much more precise description of such a diagonalization when the operator \(B\) is defined in a Hilbert space \(E = L^2(X,\mu)\), and is a \emph{Carleman operator} corresponding to a Carleman kernel \((x,y) \mapsto K(x,y)\) defined in \(X \times X\), relative to the measure \(\mu \otimes \mu\); in addition one assumes that for the isometry \(T = (T_j)_{1 \leq j < \omega}\) defining the diagonalization for the function \(\Phi\), one has \(\Phi(\zeta) \neq 0\) almost every where in \(Y\) (for the measure \(\nu\)), which is equivalent to assuming that \(B\) and \(B^*\) are \emph{injective}. It is then possible to describe the isometries
\[
    T_j: E_j \to F_j = L^2(Y, \varphi_{S_j} \cdot \nu)
\]
performing that diagonalization, in the following way:

1° For each index \(j\) such that \(1 \leq j < \omega\), there is a \((\mu otimes \nu)\)-measurable function \((x, \zeta) \mapsto e_j(x, \zeta)\) such that \(e_j(x, \zeta) = 0\) for \(\zeta \not\in S_j\), and that for almost all \(x \in X\), the function
\[
    \zeta \mapsto |\Phi(\zeta)|^2 \sum_{1 \leq j < \omega} |e_j(x, \zeta)|^2
\]
is \(\nu\)-integrable,

2° Let \(r(x) = \left( \int_X |K(x,y)|^2 d\mu(y) \right)^{1/2}\), which is \(\mu\)-measurable and almost everywhere finite; then, for every function \(f \in E_j\) such that \(\int_X^* r(x)|f(x)|d\mu(x) < +\infty\) one has
\begin{equation}
    \label{eq:9.10}
    (T_j \cdot f)(\zeta) = \int_X f(x)\overline{e_j(x, \zeta)} d\mu(x)
\end{equation}
for almost every \(\zeta \in Y\).

3° There is a null set \(N \subset X\) having the following property: for every \(j\) such that \(1 \leq j < \omega\), let \(u_j\) be a function belonging to \(F_j\), and suppose that the function
\[
    \zeta \mapsto |\Phi(\zeta)|^{-2} \sum_{1 \leq j < \omega} |u_j(\zeta)|^2
\]
is \(\nu\)-integrable. Then, for every \(x \not\in N\), the function \(\zeta \mapsto \sum_{1 \leq j < \omega} e_j(x, \zeta) u_j(\zeta)\) is \(\nu\)-integrable, and
\begin{equation}
    \label{eq:9.11}
    (T^{-1} \cdot (u_j))(x) = \int_Y \left( \sum_{1 \leq j < \omega} e_j(x, \zeta) u_j(\zeta) \right) d\nu(\zeta).
\end{equation}

The nature of this result is better understood when one specializes it to a situation stemming from harmonic Analysis. Take \(X = G\), a separable commutative locally compact group; let \(\mu\) be a Haar measure on \(G\), and consider a complex function \(b \in L^1(G) \cap L^2(G)\); then \(B: f \mapsto b * f\) is a continuous normal operator in \(L^2(G)\), such that \(B^* \cdot f = \check{\overline{b}} * f\); furthermore, from the definition
\[
    (b * f)(x) = \int_G b(x-y) f(y) d\mu(y)
\]
it follows that \(B\) is a \emph{Carleman operator} corresponding to the Carleman kernel \(K(x,y) = b(x-y)\). The Plancherel theorem and the multiplicative property of the Fourier transform (chap. \ref{ch:7}, formula \eqref{eq:7.44}) show that the isometry \(T: f \mapsto \mathcal{F}f\) defines a \emph{diagonalization} of the operator \(B\), with \(Y = \hat{G}, \omega = 2, \Phi(\zeta) = \mathcal{F}b(\zeta)\), and \(\nu\) a Haar measure on \(\hat{G}\). Formulas \eqref{eq:9.10} and \eqref{eq:9.11} then boil down to
\begin{equation}
    \label{eq:9.12}
    (T \cdot f)(\zeta) = \int_G f(x) \overline{e(x, \zeta)} d\mu(x),
    \qquad
    (T^{-1} \cdot u)(x) = \int_{\hat{G}} e(x, \zeta) u(\zeta) d\nu(\zeta)
\end{equation}
with \(e(x,\zeta) = \ip{x, \zeta}\) for \(x \in G\) and \(\zeta \in \hat{G}\), i.e. the definitions of the \emph{Fourier transform} and of its inverse; these formulas are only valid for \(f \in L^1(G) \cap L^2(G)\) and \(u \in L^1(\hat{G}) \cap L^2(\hat{G})\), which shows that the restrictions imposed on the functions \(f\) and \(u_j\), in \eqref{eq:9.10} and \eqref{eq:9.11} cannot be completely suppressed. Finally, for every \(\zeta \in \hat{G}\), one has
\begin{equation}
    \label{eq:9.13}
    b * e (\cdot, \zeta) = \Phi(\zeta) e(\cdot, \zeta)
\end{equation}
and although the functions \(e(\cdot, \zeta)\) \emph{do not} belong to \(L^2(G)\) in general, they are in some sense ``generalized eigenvectors'' for the operator \(B\), exhibiting the same phenomenon already observed by F. Riesz (chap. \ref{ch:7}, \ref{sec:7.2}).

In 1953, it was simultaneously observed by L. Gårding \cite{081} and F. Browder \cite{34} that this phenomenon of ``generalized eigenvectors'' occurs for all self-adjoint operators stemming from \emph{formally self-adjoint elliptic differential operators} (see \ref{sec:9.5}). It is assumed that such an operator \(P\) or order \(m\) in \(L^2(X)\) (where \(X\) is an open bounded subset of \(\R^n\)) possesses a \emph{self-adjoint} extension \(A_P\); then \(L = A_P + iI\) (\(I\) identity) is a \emph{normal} unbounded operator, which is a bijection of \(\dom(A_P)\) onto \(L^2(X)\); the inverse \(L^{-1}\) is therefore a \emph{continuous} normal operator in \(L^2(X)\), and the same is true of course of its iterates \(B = L^{-q}\). It follows from the existence of a parametrix of \(P\) (see \ref{sec:9.5}) that for \(qm > n, B\) is a \emph{Carleman operator}; it then also follows from the hypoellipticity of \(P\) (see \ref{sec:9.5}) and from Mautner's theorem that there is a diagonalization of \(A_P\) with \(Y = \spec(A_P)\) and \(\Phi(\zeta) = \zeta\), for which (with the preceding notations) each function \(e_j(\cdot, \zeta)\), for \(\zeta \not\in \spec(A_P)\) is a \(C^\infty\) function (generally \emph{not} in \(L^2(X)\)) solution of the partial differential equation
\begin{equation}
    \label{eq:9.14}
    (P \cdot e_j(\cdot, \zeta))(x) = \zeta e_j(x, \zeta).
\end{equation}

\section{Boundary problems for ordinary differential equations}
\label{sec:9.3}

The results of H. Weyl on the spectral theory of second order linear differential equations (chap. \ref{ch:7}, \ref{sec:7.3}) naturally raised the question of their generalization to linear differential equations of arbitrary order, but that problem was only attacked by K. Kodaira in 1949 \cite{126}. Surprisingly enough, although Stone had shown in his book \cite{207} how von Neumann's spectral theory could be applied to yield H. Weyl's results, Kodaira elected to follow Weyl's method, suitably extended. Simultaneous work by Glazman and Neumark, and later papers by many authors completed Kodaira's results and also inserted them within von Neumann's theory; we refer the reader to \cite[vol.~II, p.~1588--1592]{062} for more historical details, and we will only describe the main features of the theory.

One considers a differential operator of even order \(2r\)
\begin{equation}
    \label{eq:9.15}
    L: u \mapsto D^r(p_0 D^r u) + D^{r-1}(p_1 D^{r-1}u) + \ldots + D(p_{r-1}Du) + p_r u
\end{equation}
where \(p_0, p_1, \ldots, p_r\) are \emph{real} \(\C^\infty\) functions in an open interval \(J = ]\alpha, \beta[\) (bounded or not) of \(\R\), and \(p_0(t) \neq 0\) for all \(t \in J\); it is \emph{formally self-adjoint}, i.e. for any two functions \(u, v\) of \(\mathcal{D}(J)\) (space of \(C^\infty\) functions in \(J\) with compact support)
\begin{equation}
    \label{eq:9.16}
    (L \cdot u | v) = (u | L \cdot v)
\end{equation}
the scalar product being taken in \(L^2(J)\). We write \(T_L\) the operator \(L\), considered as a hermitian (unbounded) operator in \(L^2(J)\) with \(\dom(T_L) = \mathcal{D}(J)\). Then the adjoint \(T^*_L\) is densely defined; more precisely, \(\dom(T^*_L) = H_L\) is the space of all functions \(u\) of class \(C^{2r-1}\) in \(J\) such that the \emph{distribution} \(L \cdot u\) is a \emph{function} of \(L^2{J}\), with \(T^*_L \cdot u = L \cdot u\). The von Neumann spectral theory (chap. \ref{ch:7}, \ref{sec:7.4}) shows that \(H_L\) is the direct sum of \(\dom(T^{**}_L), E^+_L\) and \(E^-_L\), where \(E^\pm_L\) is the subspace of functions \(u\) of class \(C^{2r-1}\) which are solutions of \(L \cdot u = \pm iu\) and in addition are \emph{square integrable} in \(J\); in fact, they are of class \(C^\infty\). Due to the fact that the \(p_j\) are real functions, both spaces \(E^+_L\) and \(E^-_L\) have the \emph{same} dimension \(p \leq 2r\). The \emph{dual} \(\mathfrak{M}\) of \(E^+_L \oplus E^-_L\) can be identified with the space of linear forms on \(H_L\) which vanish on \(\dom(T^{**}_L)\); it is the direct sum \(\mathfrak{M}_\alpha \oplus \mathfrak{M}_\beta\), where \(\mathfrak{M}_\alpha\) (resp. \(\mathfrak{M}_\beta\)) is the subspace of \(\mathfrak{M}\) consisting of all forms \(\theta\) such that \(\theta(u) = 0\) for all functions \(u \in H_L\) which vanish in a neighborhood of \(\alpha\) (resp. \(\beta\)). If one writes the Lagrange ``adjunction'' formula (chap. \ref{ch:1}, \ref{sec:1.1}, formula \eqref{eq:1.5})
\[
    v(L \cdot u) - u(L \cdot v) = \frac{d}{dt}(C(u,v))
\]
for functions \(u, v\) which are \(C^\infty\) in \(J\), then it can be shown that any linear form \(\theta \in \mathfrak{M}_\alpha\) (resp. \(\mathfrak{M}_\beta\)) can be written
\[
    \theta(u) = \lim_{t \to \alpha} C(u,w)(t)
    \qquad
    (\text{resp. } \theta(u) = \lim_{t \to \beta} C(u,w)(t))
\]
for a \(C^\infty\) function \(w\) in \(J\) for which the limit exists for every \(u \in H_L\), and conversely any such function defines a linear form in \(\mathfrak{M}_\alpha\) (resp. \(\mathfrak{M}_\beta\)). Due to this result, one says that for any \(\theta \in \mathfrak{M}\), the relation \(\theta(u) = 0\) is a \emph{boundary condition} for the differential operator \(L\).

As the defects of \(T^{**}_L\) are both equal to \(p\), there exist \emph{self-adjoint extensions} \(A_L\) of \(T_L\) (infinitely many if \(p > 0\)); for each of them, \(\dom(A_L)\) is a subspace of \(H_L\) of codimension \(p\), defined by \(p\) independent ``limit conditions'' \(\theta_j(u) = 0\) with \(\theta_j \in \mathfrak{M}\) (the linear forms \(\theta_j\) are not arbitrary, since one must have \((T^*_L \cdot u | v) = (u | T^*_L \cdot v)\) for \(u\) and \(v\) in \(\dom(A_L)\)). For a given self-adjoint extension \(A_L\), let \(S\) be its spectrum, which is a closed subset of \(\R\) (it may be the whole line \(\R\), and it is always infinite and unbounded). For any \(\zeta \not\in S, (A_L - \zeta I)^{-1}\) is a continuous normal operator in \(L^2(J)\); one shows that it is a \emph{Carleman operator}, with kernel \((s,t) \mapsto G(\zeta,s,t)\), called the \emph{Green function} of \(A_L - \zeta I\); one has
\setcounter{equation}{15}
\mytodo{either one of the two equations 9.16 is the impostor, or we need to change the numbers from here on.}
\begin{equation}
    \label{eq:9.16b}
    G(\zeta, t, s) = \overline{G(\overline{\zeta}, s, t)}.
\end{equation}

For each \(\zeta \not\in S\) and \(t \in J\), the function \(s \mapsto G(\zeta, t, s)\) is a \(C^{2r-1}\) function, and in each interval \(]\alpha,t[, ]t,\beta[\), it is a \(C^\infty\) solution of the equation \(L \cdot u - \zeta u = 0\), satisfying the boundary conditions which define \(A_L\), and in addition the function \(s \mapsto \pder{^{2r-1}}{s^{2r-1}} G(\zeta, t, s)\) has limits on the left and on the right at the point \(s = t\), such that
\begin{equation}
    \label{eq:9.17}
    \pder{^{2r-1}}{s^{2r-1}} G(\zeta, t, t-) - \pder{^{2r-1}}{s^{2r-1}} G(\zeta, t, t+) = 1/p_0(t).
\end{equation}
These conditions obviously generalize those seen in chap. \ref{ch:7}, \ref{sec:7.3} for second order equations; they completely determine \(G\) once a fundamental system of \(2r\) solutions \(t \mapsto v_j(\zeta, t)\) of \(L \cdot u - \zeta u = 0\) is known, and it is easily seen that one can write in matrix notation
\begin{equation}
    \label{eq:9.18}
    \begin{dcases}
        G(\zeta, s, t) = \vec{v}(\overline{\zeta}, s)^* \, W^-(\zeta)\vec{v}(\zeta, t) \quad \text{for } t < s \\
        G(\zeta, s, t) = \vec{v}(\overline{\zeta}, s)^* \, W^+(\zeta)\vec{v}(\zeta, t) \quad \text{for } t > s
    \end{dcases}
\end{equation}
where \(\vec{\zeta, t}\) is the one column matrix \((v_j(\zeta, t))_{1 \leq j \leq 2r}\), \(W^-(\zeta)\) and \(W^+(\zeta)\) are two square matrices of order \(2r\) which only depend on \(\zeta\). Furthermore, if the \(v_j\) have been chosen so that for each \(t \in J\), the functions \(\zeta \mapsto v_j(\zeta, t)\) are \emph{holomorphic} in an open subset \(H \subset \C\), then \(W^-\) and \(W^+\) are \emph{holomorphic} in \(H \cap (\C - S)\).

As the operator \(L\) is elliptic and formally self-adjoint, one can apply to it the Gårding--Browder theorem (\ref{sec:9.2}). It can be shown that, with the notations introduced in \ref{sec:9.2}, one has \(\omega \leq 2r + 1\), in other words, \(A_L\) has \emph{at most multiplicity} \(2r\) in its spectrum \(S\); for convenience, if \(\omega < 2r + 1\), one defines the function \(e_j(t, \xi)\) to be identically 0 for \(\omega \leq j \leq 2r, t \in J\) and \(\xi in S\); for the other values of \(j, (t,\xi) \mapsto e_j(t,\xi)\) is a (\(\lambda \otimes \nu\))-measurable function in \(J \times S\), such that for each \(\xi \in S, t \mapsto e_j(t,\xi)\) is a solution of the equation \(L \cdot u - \xi u = 0\) and for almost all \(t \in J\), the function \(\xi \mapsto e_j(t, \xi)\) is square integrable (for \(\nu\)) in each compact subset of \(S\); in addition, one has \(e_j(t,g) = 0\) for \(\xi \not\in S\). For \(\xi \in S \subset \R\), one may write (with \(\vec{e}(t,\xi) = (e_j(t, \xi))_{1 \leq j \leq 2r}\), one column matrix)
\begin{equation}
    \label{eq:9.19}
    \vec{e}(t, \xi) = Q(\xi) \cdot \vec{v}(\xi, t)
\end{equation}
and the elements of the matrix \(Q\) are \(\nu\)-measurable and square integrable in every compact subset of \(S\). Let
\begin{equation}
    \label{eq:9.20}
    P(\xi) = Q(\xi)^* Q(\xi)
\end{equation}
which is a positive hermitian matrix for all \(\xi \in S\). The spectral decomposition of the operator \((A_L - \zeta I)^{-1}\) can then be written \emph{explicitly} in the following way: for any function \(f \in \mathcal{D}(J)\), write
\begin{equation}
    \label{eq:9.21}
    (U \cdot f)(\zeta) = \left( \int_J \overline{f(t)} v_j(\zeta, t) dt \right)_{1 \leq j \leq 2r} \quad \text{(one column matrix);}
\end{equation}
then, for \(f, g\) in \(\mathcal{D}(J)\), one has
\begin{equation}
    \label{eq:9.22}
    (f|g) = \int_S ((U \cdot f)(\xi))^* P(\xi)((U \cdot g)(\xi)) d\nu(\eta)
\end{equation}
and
\begin{equation}
    \label{eq:9.23}
    ((A_L - \zeta I)^{-1} \cdot f| g) = \int_S (\xi - \zeta)^{-1} ((U \cdot f)(\xi))^* P(\xi)((U \cdot g)(\xi)) d\nu(\xi)
\end{equation}

The set \(S_j - S_{j+1}\) is then the set of \(\xi \in S\) such that \emph{the matrix} \(P(\xi)\) \emph{has rank equal to} \(j\), and the (vector) measure \(P \cdot \nu\) can be recovered from the knowledge of the matrix \(W^+\) (formula \eqref{eq:9.18}) by the relation
\begin{equation}
    \label{eq:9.24}
    \begin{gathered}
        \int_{[a,b]} P(\xi)d\nu(\xi) - \frac{1}{2} (P(a)\nu(\{a\}) + P(b)\nu(\{b\})) = \\
        = \frac{1}{2i\pi} \lim_{\varepsilon \to 0} \int_a^b (W^+(\sigma + i\varepsilon) - W^+(\sigma - i\varepsilon)) d\sigma.
    \end{gathered}
\end{equation}

These results had previously been obtained for second order operators by Titchmarsh \cite{212}.

Much work has been done to determine the spectrum \(S\), the various subsets \(S_j\) and the measure \(P \cdot \nu\) under various hypotheses on the operator \(L\) (\cite{062}, \cite{166}). It should how ever be stressed that the behavior of the measure \(P \cdot \nu\) on \(\R\) is essentially \emph{arbitrary}: in a remarkable paper, Gelfand and Levitan have shown in 1951 \cite{084} that, given on an arbitrary compact subset \(H\) of \(\R\) an arbitrary measure \(\rho\), it is always possible to find a second order operator \(L\) \emph{with} \(C^\infty\) \emph{coefficients} such that \(P(\xi)\) has the form
\(
    \begin{pmatrix}
        p_{11}(\xi) & 0 \\
        0           & 0
    \end{pmatrix}
\)
and the restriction of the measure \(p_{11} \cdot \nu\) to \(H\) is the given measure \(\rho\).

\section[Sobolev spaces and a priori inequalities]{Sobolev spaces and \emph{a priori} inequalities}
\label{sec:9.4}

Until 1940, there was no general theory of linear partial differential equations (or systems of such equations) of \emph{arbitrary order}. With the exception of a few special types of equations with constant coefficients (such as the ``biharmonic'' equation \(\Delta^2 u = 0\)), the bulk of papers were concerned with \emph{second order} equations in any number of variables, to which must be added a much smaller number of results on equations of arbitrary order in \emph{two} independent variables.

When mathematicians began to be interested in ``weak'' solutions, and later with the arrival of the theory of distributions (chap. \ref{ch:8}), the scope of the theory of linear partial differential equations was greatly widened; if
\begin{equation}
    \label{eq:9.25}
    u \mapsto P \cdot u = \sum_{|\alpha| \leq m} a_\alpha D^\alpha u
\end{equation}
is a linear differential operator with complex \(C^\infty\) coefficients\footnote{The interest shown to equations with \(C^\infty\) coefficients (or \(C^r\) coefficients generally) is chiefly due to the pioneering efforts of Hadamard, who repeatedly emphasized that for applications to Physics it was unreasonable to study exclusively equation with analytic coefficients \cite{093}.} in an open subset \(\Omega\) of \(\R^n\), then for any distribution \(T \in \mathcal{D}'(\Omega)\), each product \(a_\alpha D^\alpha T\) is defined, hence also \(P \cdot T\), and it makes sense to ask for solutions \(T\) of the equation
\begin{equation}
    \label{eq:9.26}
    P \cdot T = S
\end{equation}
where \(S\) is any given distribution in \(\mathcal{D}'(\Omega)\). In particular, one may take for \(S\) a \(C^\infty\) function in \(\Omega\), and then one asks if this imposes conditions on the distributions \(T\) solutions of \eqref{eq:9.26}. in some cases, solutions of \(P \cdot T = 0\) may be distributions of arbitrary order (i.e. as ``irregular'' as possible); this happens for instance for \(P = \pder{^2}{x \partial y}\), where not only arbitrary locally integrable functions \(A(x) + B(y)\) are solutions, but also arbitrary distributions of type \(A \otimes 1 + 1 \otimes B\), where \(A\) and \(B\) are arbitrary distributions of \(\mathcal{D}'(\Omega)\). On the other hand, it may happen that for \emph{all} \(C^\infty\) functions \(f\) in \(\Omega\), \emph{all} solutions of \(P \cdot T = f\) are necessarily \(C^\infty\) functions; such operators \(P\) are now called \emph{hypoelliptic}. This is the case, for instance, when \(n=1\) and \(P = D^p + a_1 D^{p-1} + \ldots + a_p\) is any linear differential operator with leading coefficient 1, an elementary result which follows by induction from the case \(p = 1\), which is du Bois-Reymond's lemma (chap. \ref{ch:8}, \ref{sec:8.3}). In 1927, S. Zaremba \cite{233} proved a result which, in modern language, means that the laplacian \(\Delta\) is hypoelliptic, and H. Weyl in 1940 gave another proof of that result (\cite{227}, vol.~III, p.~758--791).

After 1950, such questions, as well as extensions of the classical boundary problems, began to be studied for operators \eqref{eq:9.25} of \emph{arbitrary} order, heralding a period of unprecedented expansion in the theory of partial differential equations. Among the many methods developed during that period, we shall postpone to \ref{sec:9.5} those linked to the concepts of elementary solution and parametrix, and consider here the applications of the ``\emph{a priori} inequalities'' which were made possible by the appearance of new tools linked to the theory of distributions, the \emph{Sobolev spaces} and their generalizations.

We have already seen (\ref{sec:9.1}) that Schauder had considered the space of functions \(f\) of class \(C^p\) in such that all derivatives \(D^\alpha f\) for \(|\alpha| \leq p\) belong to \(L^2(\Omega)\), and had used the norm \(\norm{f}_p = \left( \int_\Omega \left( \sum_{|\alpha| \leq p} |D^\alpha f(x)|^2 \right) dx \right)^{1/2}\) on that space, which unfortunately was not complete for that norm. In 1936, Sobolev had the idea of considering the functions \(f \in L^2(\Omega)\) which have \emph{weak} (= distributional) derivatives \(D^\alpha f\) belonging also to \(L^2(\Omega)\) for \(|\alpha| \leq p\), and this time this space (with the same norm) \(H^p(\Omega)\) is complete (i.e. a Hilbert space); moreover, Sobolev observed that the larger the number \(p\), the more regular are the functions of \(H^p(\Omega)\); for \(p > \left[ \frac{n}{2} \right] + 1\), they are functions of class \(C^r\) with \(r = p - \left[ \frac{n}{2} \right] - 1\), hence the intersection of all \(H^p(\Omega)\) consists of the functions of class \(C^\infty\) such that all their derivatives are in \(L^2(\Omega)\) \cite{201}.

Later, it was realized that the \(H^s(\R^n)\) could be defined using the Fourier transform, as the space of distributions \(T \in \mathcal{S}'(\R^n)\) such that the Fourier transform \(\mathcal{F}T\) is a locally integrable function for which the function \(\xi \mapsto (1 + |\xi|^2)^s |\mathcal{F}T(\xi)|^2\) is integrable. It is then clear that the definition may be extended to \emph{all real numbers} \(s\).

These properties were the source of what one may call the ``bootstrap'' method to prove that a function is \(C^\infty\): it is enough to show that if it belongs to \emph{some} Sobolev space \(H^r(\Omega)\), it \emph{also} belongs to \(H^{r+1}(\Omega)\). The first idea of this method is apparently due to K. Friedrichs \cite{077}, who applied it to prove that elliptic operators \eqref{eq:9.25} (see \ref{sec:9.5}) are hypoelliptic, a question which we postpone until \ref{sec:9.5}; his tool is a new type of \emph{a priori} inequality, which was the starting point of a very large number of similar results, for elliptic\footnote{For a description of the various methods based on a priori inequalities in 1956, see \cite{168}.} and other types of operators. We shall only mention here one of the most refined ones \cite[p.~207]{117}; it concerns what are called ``principally normal'' operators \(P\), which we shall not attempt to describe more precisely here, but which include operators with real coefficients and operators with constant coefficients; under assumptions on \(\Omega\) in relation with the characteristic hyperplanes of the operator \(P\), too complicated to reproduce here, the fundamental result is that if \(u\) is a distribution on \(\Omega\) having a compact support \(K\), and such that \(P \cdot u\) belongs to some \(H^s(\Omega)\), then \(u\) necessarily belongs to \(H^{s+m-1}(\Omega)\), and there is a constant \(C_{s,K}\) independent of \(u\) and such that
\begin{equation}
    \label{eq:9.27}
    \norm{u}_{s+m-1} \leq C_{s,K}(\norm{P \cdot u}_s + \norm{u}_{s+m-2}).
\end{equation}
The ``bootstrap'' method shows that if \(P \cdot u = f\) where \(f\) is a \(C^\infty\) function, then \(u\) itself is a \(C^\infty\) function. But from inequality \eqref{eq:9.27} one may derive much more information: for instance the space of solutions of \(P \cdot u = 0\) having support in \(K\) is finite dimensional and consists of \(C^\infty\) functions if \(f\) is a \(C^\infty\) function such that \(\int_\Omega f(x)u(x)dx = 0\) for all these functions \(u\), then there is a \(C^\infty\) function \(v\) in \(\Omega\) solution of the adjoint equation \(\prescript{t}{}P \cdot v = f\) (ibid., p.~210). Similar uses of such inequalities have been successful in proving existence and uniqueness of Cauchy problems for hyperbolic (see \ref{sec:9.5}) equations of arbitrary order (\cite{062}, vol.~II, p.~1748--1766).

\section{Elementary solutions, parametrices and pseudo-differential operators}
\label{sec:9.5}

We have seen (chap. \ref{ch:2}, \ref{sec:2.3}) that from the beginning the idea of ``newtonian'' potential was closely related to the laplacian operator; if one considers in \(\R^3\) the integral operator \(U\) defined by the kernel \(-\frac{1}{4\pi|x-y|}\) (where \(|x|\) is the euclidean norm), then Poisson's equation (chap. \ref{ch:2}, formula \eqref{eq:2.12}) can be written \(\Delta \cdot (U \cdot \rho) = \rho\) for \(\rho \in \mathcal{D}(\R^3)\), and one also may write \(U \cdot (\Delta \cdot f)\) if \(f\) is the potential having density \(\rho\), so that this equation is also valid for all \(f \in \mathcal{D}(\R^3)\). Later, when the existence of the Green function \(G\) was proved for a domain \(\Omega\) in \(\R^3\), one had similarly the relations \(\Delta \cdot (U \cdot f) = f\) and \(U \cdot (\Delta cdot f) = f\) for all \(f \in \mathcal{D}(\Omega)\), where now \(U\) is the integral operator defined by the kernel \(\frac{1}{4\pi}G(x,y)\). In both cases, \(y \mapsto \frac{1}{|x-y|}\) and \(y \mapsto G(x,y)\) are solutions of the Laplace equation \(\Delta u = 0\), but with a \emph{singular point} for \(y=x\).

In 1860 Riemann proposed to solve Cauchy's problem for linear hyperbolic equations of second order in 2 variables \(P \cdot u = 0\) by using a particular solution \(y \mapsto R(x, y)\) of the adjoint equation \(\prescript{t}{}P \cdot v = 0\) depending on the point \(x\), in such a way that the integral operator \(U\) defined by the kernel \(R\) satisfies again the relation \(P \cdot (U \cdot f) = f\) for \(f \in \mathcal{D}(\R^2)\), the function \(R\) playing for \(P\) a role similar to the Green function for the laplacian. In that case \(R\) is continuous; but when Volterra and Hadamard undertook to extend the method to second order hyperbolic equations in \(n \geq 3\) variables, they were beset by difficulties stemming from the fact that the function corresponding to \(R\) would now have, not only (as the Green function) a singularity at one point \(x\), but singularities along lines or surfaces.

We cannot here describe the details of these researches, for which we refer to \cite{093}. What we want to emphasize is that, by the end of the XIX\supth{} century, one had the (rather vague) idea that, for a second order linear differential operator \(P\), one should look for a solution \(y \mapsto R(x,y)\) of the adjoint equation \(\prescript{t}{}P \cdot v = 0\) having a suitable singularity at the point \(x\), and that the integral operator \(U\) defined by the kernel \(R\) would be such that \(P \cdot (U \cdot f) = f\) for \(f \in \mathcal{D}(\Omega)\); such a function \(R\) was called an \emph{elementary} (or \emph{fundamental}) solution for \(\prescript{t}{}P\). This worked, not only for \(\Delta\), but also for instance for the heat operator \(\pder{}{t} - \pder{^2}{x^2}\), one has \(R(t,x) = \frac{1}{2\sqrt{\pi t}} e^{-x^2/4t}\) for \(t > 0, R(t,x) = 0\) for \(t < 0\) (\(R(0,x)\) is undefined); for hyperbolic equations, things were not so simple, for one had to apply the integral operator defined by the kernel \(R\), not to \(f\) but to some derivatives of \(f\) \cite{122}.

After 1900, mathematicians began to investigate the possibility of extending the notion of elementary solution and its applications to equations of higher order. To an operator \eqref{eq:9.25}, one associates the polynomial in \(\xi = (\xi_1, \xi_2, \ldots, \xi_n)\)
\begin{equation}
    \label{eq:9.28}
    \sigma_P(x,\xi) = \sum_{|\alpha| \leq m} a_\alpha(x)(2\pi i\xi)^\alpha
\end{equation}
(later called the ``symbol'' of \(P\)) and the homogeneous polynomial in  consisting in the terms of highest degree \(m\) (the ``principal symbol'' of \(P\))
\begin{equation}
    \label{eq:9.29}
    \sigma^0_P(x,\xi) = \sum_{|\alpha| = m} a_\alpha(2\pi i\xi)^\alpha.
\end{equation}
The operator \(P\) is called \emph{elliptic} if \(\sigma^0_P(x,\xi) \neq 0\) for all \(x \in \Omega\) and all \(\xi \neq 0\). In his thesis, Fredholm considered, for \(n = 3\), elliptic operators with constant coefficients, and proved the existence of an elementary solution by writing it explicitly as an abelian integral (\cite{074}, p.~17--57); this was later generalized to elliptic operators with constant coefficients in any number of variables (Holmgren, Herglotz \cite{109}). In 1907, E.E. Levi considered elliptic operators with variable coefficients, and either \(n = 2\) variables, or operators of order 2 in any number of variables; in both cases, using the fact that for constant coefficients elementary solutions were explicitly known, he showed how one could prove the existence of elementary solutions by showing that their determination could be reduced to the solution of a Fredholm integral equation (\cite{145}, vol.~II, p.~28--84).

For operators with constant coefficients, the theory of distributions completely clarified the concept of ``elementary solution'' \cite{194}; such an operator may be written \(u \mapsto A * u\), where \(A = \sum\limits_\alpha a_\alpha D^\alpha \varepsilon_0\), linear combination of derivatives of the Dirac measure \(\varepsilon_0\) at the origin of \(\R^n\). An elementary solution is then, by definition, a \emph{distribution} \(E\) on \(\R^n\) such that
\begin{equation}
    \label{eq:9.30}
    A * E = \varepsilon_0;
\end{equation}
it follows at once from that definition that, for any \emph{distribution} \(T\) with compact support (and not only for a function), one has
\begin{equation}
    \label{eq:9.31}
    A*(E*T) = E*(A*T) = T.
\end{equation}

In 1954, Ehrenpreis \cite{063} and Malgrange \cite{156} independently proved that \emph{any} operator \(P\) of form \eqref{eq:9.25} with constant coefficients has an elementary solution \(E\). Of course, such a solution is only determined up to addition of any distribution \(S\) solution of the homogeneous equation \(P \cdot S =0\). It is not obvious that among these distributions there would exist \emph{tempered} ones; as \(\mathcal{F}A\) is the polynomial \(\sigma_P(\xi)\) such an elementary solution should be such that
\begin{equation}
    \label{eq:9.32}
    \sigma_P(\xi) \cdot \mathcal{F}E = 1;
\end{equation}
it is only in 1958 that independently Hörmander \cite{118} and Łojasiewicz \cite{153} showed that it is always possible, for any polynomial \(Q\), to find a tempered distribution \(T\) such that \(Q \cdot T = 1\).

Elementary solutions proved to be useful to show that an operator \(P\) is hypoelliptic \cite[p.~100]{117}, or to prove uniqueness of the Cauchy problem for some operators \cite[p.~141]{117}. But gradually it was realized that instead of looking for a ``right inverse'' to a differential operator, it could be much simpler to obtain an ``approximate right inverse'' which would be put to the same uses.

Such an idea was first introduced by Hilbert in 1907 under the name of \emph{parametrix}, in a particularly simple context, the study of an elliptic operator \(P\) of order 2 on the sphere \(\Sph_2\); he proves that there is an integral operator \(Q\), defined by a kernel having a singularity similar to the logarithmic singularity of the Green function, and such that \(QP = I + R\), where \(R\) is an integral operator; a solution of \(P \cdot u = f\) is therefore a solution of the \emph{integral equation} \(u + R \cdot u = Q \cdot f\) \cite[p.~233--242]{112}; \(Q\) could therefore be considered as an ``approximate left inverse'' of \(P\).

Two years later, E.E. Levi independently introduced a similar method in a much more general and difficult question, the generalization of the Dirichlet problem for elliptic operators of \emph{arbitrary even order}, completely unexplored until then except for a few special operators such as the iterated laplacian. On these special examples it transpired that what should correspond to the Dirichlet problem for an operator of even order \(2m\) was the boundary condition consisting in fixing the values of the solution and \emph{its first} \(m-1\) \emph{normal derivatives} on the boundary \(I\) of a bounded open set \(\Omega \subset \R^n\). Levi is only concerned with the case of \(n = 2\) variables and first shows that the problem (for smooth boundary \(\Gamma\)) may be reduced to the case in which one has to find a solution of \(P \cdot u = f\) such that \(u\) and its \(m-1\) first normal derivatives take the value 0 on \(\Gamma\). His idea is then to determine two functions, \(\varphi(x,t)\) and \(K(x,y,x_1,y_1)\) defined respectively in \(\overline{\Omega}\) and \(\overline{\Omega} \times \overline{\Omega}\), and such that: 1° for each point \((x_1, y_1) \in \Omega\), the function \((x,y) \mapsto (x,y,x_1,y_1)\) and its \(m-1\) first normal derivatives vanish on \(\Gamma\); 2° the function
\begin{equation}
    \label{eq:9.33}
    u(x,y) = \iint_\Omega K(x,y,x_1,y_1)\varphi(x_1,y_1) dx_1 dy_1
\end{equation}
satisfies the equation \(P \cdot u = f\). In order to obtain that result, he chooses \(K\) in such a way that
\begin{equation}
    \label{eq:9.34}
    (P \cdot u)(x,y) = \varphi(x,y) + \iint_\Omega K_1(x,y,x_1,y_1)\varphi(x_1,y_1) dx_1 dy_1
\end{equation}
where \(K_1\) is a kernel to which Fredholm's theory is applicable, and he has thus reduced the problem to a Fredholm integral equation. If one writes \(Q \cdot \varphi\) the right hand side of \eqref{eq:9.33}, one may say that the operator \(Q\) is such that \(PQ = I + R\) where \(R\) is an integral operator; this time \(Q\) is an ``approximate right inverse'' of \(P\). The determination of the function \(K\) satisfying these conditions is a difficult problem, and it is not surprising that after Levi not much work was done in that direction until around 1960 (\cite{145}, vol.~II, p.~207--343).

At that time progress came from a completely different direction. In his work on integrals of complex functions along paths in \(\C\), Cauchy, in 1814, had observed that if \(f\) is a \(C^1\) function in an interval \([-a,a]\) of \(\R\), the function \(f(x)/x\) is not integrable if \(f(0) \neq 0\), but the sum \(\int_{-a}^{-\varepsilon} \frac{f(x)dx}{x} + \int_\varepsilon^a \frac{f(x)dx}{x}\), when \(\varepsilon\) tends to 0, has a limit which he called the ``principal value'' of the integral. Similarly, if \(L\) is a \(C^1\) curve in \(\C\), the limit
\begin{equation}
    \label{eq:9.35}
    (H \cdot f)(x) = \lim_{\varepsilon \to 0} \int_{L_\varepsilon} \frac{f(z)dz}{z-x}
\end{equation}
where \(x \in L\) and \(L_\varepsilon\) is the part of \(L\) for which the arc of \(L\) joining \(x\) and \(z\) has length \(> \varepsilon\), exists for each \(C^1\) function \(f\) defined on \(L\), and is written \(\text{vp} \int_L \frac{f(z)dz}{z-x}\); if \(L\) is a simple closed curve, the boundary of a bounded open set \(\Omega\), the usual line integral \(\frac{1}{2\pi i} \int_L \frac{f(z)dz}{z-x}\) is defined for \(x \not\in L\), and is in \(\Omega\) a holomorphic function \(F^+(x)\), and in the exterior \(\C - \overline{\Omega}\) another holomorphic function \(F^-(x)\); when \(x\) tends to a point \(t \in L\), these functions have limits respectively equal to
\[
    \begin{aligned}
        F^+(t) &= \frac{1}{2} f(t) + \frac{1}{2\pi i} \text{vp} \int_L \frac{f(z)dz}{z - t}, \\
        F^-(t) &= -\frac{1}{2} f(t) + \frac{1}{2\pi i} \text{vp} \int_L \frac{f(z)dz}{z - t}.
    \end{aligned}
\]
In his third paper on integral equations, Hilbert, using these formulas, showed that one could find two holomorphic functions \(F^+(x)\) in \(\Omega, F^-(x)\) in \(\C - \overline{\Omega}\), such that for \(t \in L\), the limits \(F^+(t)\) and \(F^-(t)\) of these functions exist when \(x\) tends to \(t\), and satisfy a relation \(F^+(t) = g(t)F^-(t)\), where \(g\) is a \(C^1\) function on \(L\) \cite[p.~81--108]{112} this led to calling the function \(H \cdot f\) defined by \eqref{eq:9.35} the \emph{Hilbert transform} of \(f\).

Between 1910 and 1955, many mathematicians studied various generalizations of this operator to functions of any number of variables, and applied them to various problems of Analysis; we cannot describe this evolution in any detail, and refer the reader to \cite{197}. The most general of these ``singular integral operators'' (or ``Calderon--Zygmund operators'' as they were also called) are defined in the following way: \(\Omega\) is an open subset of \(\R^n, (x, \xi) \mapsto K(x, \xi)\), a locally integrable mapping of \(\Omega \times (\R^n - \{0\})\) into \(\C\), which is \emph{positively homogeneous} of degree \(-n\) in \(\xi\) for every \(x \in \Omega\); in addition it is assumed that for any \(x \in \Omega, \int_{\Sph_{n-1}} K(x, \xi) d\sigma(\xi) = 0\) (\(\sigma\) being the invariant measure on \(\Sph_{n-1}\)). Then, for any function \(u \in \mathcal{D}(\Omega)\), the limit
\begin{equation}
    \label{eq:9.36}
    (P \cdot u)(x) = \lim_{\varepsilon \to 0} \int_{|y-x| \geq \varepsilon} K(x,y-x)u(y)dy
\end{equation}
exists for every \(x\) and defines a \emph{singular integral operator} \(P\).

Around 1960, it was realized that the use of Fourier transform (generalized to distributions) enabled one to define a class of linear operators which contained at the same time differential operators of type \eqref{eq:9.25}, singular integral operators and some ordinary integral operators (with locally integrable kernels); several mathematicians independently contributed to this new theory, but again we cannot go into any historical detail, and we shall merely give a short description of its present status (for more references, see \cite{D}, \cite{119} and \cite{068}). The inversion formula for Fourier transforms shows that the operator \eqref{eq:9.25} can be written
\begin{equation}
    \label{eq:9.37}
    (P \cdot u)(x) = \int_{\R^n} \exp(2\pi i(x|\xi))a(x,\xi)\mathcal{F}u(\xi)d\xi
    \quad
    \text{for } u \in \mathcal{D}(\Omega)
\end{equation}
where \(a(x,\xi) = \sigma_P(x,\xi)\) (formula \eqref{eq:9.28}). The generalization consists in replacing in \eqref{eq:9.37} the polynomial (in \(\xi\)) \(\sigma_P(x,\xi)\) by a more general \(C^\infty\) function defined in \(\Omega \times \R^n\) and which is only submitted to conditions concerning its \emph{growth} as \(|\xi|\) tends to \(+\infty\): for a polynomial \(\sigma_P, D^\alpha_x \sigma_P\) has the same behavior for \(|\xi| \to \infty\) as \(\sigma_P\) itself, whereas \(D^\beta_\xi \sigma_P\) is a polynomial in \(\xi\) of degree \(m - |\beta|\). One defines then a \emph{symbol} as a \(C^\infty\) mapping \((x,\xi) \mapsto a(x, \xi)\) of \(\Omega \times \R^n\) into \(\C\), such that one has
\begin{equation}
    \label{eq:9.38}
    |D^\alpha_x D^\beta_\xi a(x, \xi)| \leq C_{\alpha \beta L} (1 + |\xi|)^{m-|\beta|}
\end{equation}
for all multiindices \(\alpha, \beta\) and all compact subsets \(L \subset \Omega\), where \(x \in L\) and \(\xi \in \R^n\) are arbitrary, and \(C_{\alpha \beta L}\) is independent of \(x \in L\) and \(\xi \in \R^n\); the main difference is that here \(m\) (the \emph{order} of \(a\), or of \(P\)) is an \emph{arbitrary real} number. The corresponding operator \(P\) defined by \eqref{eq:9.37} for \(u \in \mathcal{D}(\Omega)\) is called a \emph{pseudo-differential operator} defined by the symbol \(a\); differential operators therefore correspond to symbols of order \(m\) equal to an integer \(m \geq 1\), the newtonian potential to \(m = -2\), and the singular integral operators \eqref{eq:9.36} (for \(K\) of class \(C^\infty\)) to \(m = 0\). The most interesting case is the one in which \(a = a_0 + a_1\), where \(a_0\) is \emph{positively homogeneous of degree} \(m\) in \(\xi\), and \(a_1\) is a symbol of order \(< m\); one then says that \(a\) is the \emph{principal symbol} of the pseudo-differential operator \eqref{eq:9.37}, and one writes \(a_0 = \sigma^0_P\).

The main properties of pseudo-differential operators are the following ones:

I) \(P\) maps the space \(\mathcal{D}(\Omega)\) into the space \(\mathcal{E}(\Omega)\) of all \(C^\infty\) complex functions in \(\Omega\); it has an adjoint \(P^*\), satisfying
\begin{equation}
    \label{eq:9.39}
    (P \cdot u | v) = (u |  P^* \cdot v)
\end{equation}
for all \(u, v\) in \(\mathcal{D}(\Omega)\) (scalar product of \(L^2(\Omega)\)), which is a pseudo-differential operator of same order \(m\); if \(P\) has a principal symbol, \(P^*\) has a principal symbol such that
\begin{equation}
    \label{eq:9.40}
    \sigma^0_{P^*} = \overline{\sigma^0_P}.
\end{equation}

II) One says \(P\) is of \emph{proper type} if both \(P\) and \(P^*\) apply \(\mathcal{D}(\Omega)\) into \emph{itself}; for any pseudo-differential operator \(Q\) of order \(r\), the compositions \(QP\) and \(PQ\) are then defined and are pseudo-differential operators of order \(m+r\); if \(P\) and \(Q\) have principal symbols, so have \(PQ\) and \(QP\) and
\begin{equation}
    \label{eq:9.41}
    \sigma^0_{PQ} = \sigma^0_{QP} = \sigma^0_P \sigma^0_Q.
\end{equation}

III) If \(m < -1, P\) is \emph{an integral operator}, having a kernel which is locally integrable in \(\Omega \times \Omega\), but has \emph{singularities for} \(x = y\); if \(m < -n-k\), the kernel is of class \(C^k\) in the whole of \(\Omega \times \Omega\). One says that the symbol \(a\) (and the corresponding pseudo-differential operator \(P\)) are of \emph{order} \(-\infty\) if a satisfies inequalities \eqref{eq:9.38} for \emph{every} real number \(m\); \(P\) is then an integral operator with a kernel which is of class \(C^\infty\), and conversely any such operator is a pseudo-differential operator of order \(-\infty\), and its principal symbol is 0.

IV) Any pseudo-differential operator \(P\) may be extended by continuity (for the weak topology) from \(\mathcal{D}(\Omega)\) to the space \(\mathcal{E}'(\Omega)\) of all distributions on with compact support. The operators \(P\) of order \(-\infty\) are \emph{characterized} by the property that for any distribution \(T \in \mathcal{E}'(\Omega), P \cdot T\) is a \(C^\infty\) function on \(\Omega\); one says that these operators are \emph{smoothing} operators. Any pseudo-differential operator is the sum of a pseudo-differential operator of proper type and of a smoothing operator. When \(K\) is a smoothing operator, so are the products \(KP\) and \(PK\) for any pseudo-differential operator \(P\), if one of the two operators \(K, P\) is of proper type.

One writes \(P \sim{} Q\) if \(P - Q\) is a smoothing operator.

V) The most remarkable feature of pseudo-differential operators is the possibility of defining a symbol by an \emph{asymptotic expansion}. Suppose given an infinite sequence \(a_0, a_1, \ldots, a_k, \ldots\), of symbols, having orders \(m_0 > m_1 > \ldots > m_k > \ldots\) with \(\lim\limits_{k \to \infty} m_k = -\infty\); then there exists a symbol \(a\) of order \(m_0\) such that, for any \(k, a - (a_0 + a_1 + \ldots + a_k)\) has order \(< m_k\); this is expressed by writing
\[
    a \sim{} a_0 + a_1 + \ldots + a_k + \ldots
\]
and saying that the right hand side is an asymptotic expansion of \(a\). If \(P_0, \ldots, P_k, \ldots\) are the pseudo-differential operators defined by the symbols \(a_0, \ldots, a_k, \ldots\), and \(P\) the pseudo-differential operator defined by \(a\), one also writes
\[
    P \sim{} P_0 + P_1 + \ldots + P_K + \ldots
\]

VI) One says a pseudo-differential operator \(P\) having a principal symbol of order \(m\) is \emph{elliptic} if \(\sigma^0_P(x, \xi) \neq 0\) for \(x \in \Omega\) and \(\xi \neq 0\); for differential operators, this coincides with the previous definition. It is \emph{equivalent} to say that there exists a pseudo-differential operator  of order \(-m\) and of proper type such that \(QP = I+ R\) and \(PQ = I + R'\) where \(R\) and \(R'\) are smoothing operators; in other words, \(P\) has a (left and right) \emph{parametrix} in a very strong sense. The proof is very simple; the necessity follows from the fact that if \(QP \sim{} I\), one must have \(\sigma^0_P = (\sigma^0_Q)^{-1}\) by \eqref{eq:9.41}. Conversely, if \(P\) is elliptic, there is a pseudo-differential operator \(Q_1\) of proper type defined by the symbol \((\sigma^0_P)^{-1}\); one has then \(Q_1 P = I - P_1\), where \(P_1\) has \emph{order} \(\leq -1\), and one is reduced to finding a pseudo-differential operator \(Q_2\) such that \(Q_2(I - P_1) = I + R\), where \(R\) is a smoothing operator; but it is enough to take \(Q_2 \sim{} I + P_1 + P^2_1 + \ldots + P^k_1 + \ldots\) to obtain that result!

VII) An immediate consequence of the existence of a parametrix \(Q\) for an elliptic differential operator \(P\) is that \(P\) is \emph{hypoelliptic}, for if \(T\) is a distribution such that \(P \cdot T = f \in \mathcal{E}(\Omega)\), one has \(Q \cdot f = T + R \cdot T\), and as \(R\) is a smoothing operator, \(R \cdot T\) and \(Q \cdot f\) are both \(C^\infty\) functions, hence also \(T\). Another easy consequence of the use of pseudo-differential operators is that for each point \(x_0 \in \Omega\), there is a small neighborhood \(U \subset \Omega\) of \(x_0\) such that the equation \(P \cdot u = f\) has solutions in \(U\) (in other words, the H. Lewy phenomenon (chap. \ref{ch:2}, \ref{sec:2.2}) cannot occur for an elliptic differential operator \(P\)); one should note, however, that there are examples of elliptic operators \(P\) defined in \(\R^n\), and such that there are equations \(P \cdot u = f\) which have \emph{no solution} in a large ball containing the support of \(f \in \mathcal{D}(\R^n)\) \cite{176}.

VIII) When \(\Omega \subset \R^n\) is bounded, pseudo-differential operators of proper type in \(\Omega\) have simple continuity properties with respect to the Sobolev spaces: if \(P\) is such an operator of order \(r\), defined in a neighborhood of \(\overline{\Omega}\), and \(s\) is any real number, there is a constant \(C\) depending only on \(P\) and \(s\), such that
\begin{equation}
    \label{eq:9.42}
    \norm{P \cdot u}_s \leq C \norm{u}_{r+s}
\end{equation}
for any \(u \in \mathcal{D}(\Omega)\), the norms being those of \(H^s(\R^n)\) and \(H^{r+s}(\R^n)\). If \(r > 0\) and \(P\) is \emph{elliptic}, applying this result to a parametrix of \(P\) immediately yields an \emph{a priori} inequality of Friedrichs type
\begin{equation}
    \label{eq:9.43}
    \norm{u}_r \leq  C(\norm{P \cdot u}_0 + \norm{u}_0).
\end{equation}

Suppose in addition that \(P = P^*\) is a differential operator such that \(\sigma_P(x,\xi) > 0\) for large \(|\xi|\); then an easy inductive argument determines (by an asymptotic expansion of its symbol) a pseudo-differential operator \(S\) of proper type and of order \(r/2\) such that \(P = S^* S + R\), where \(R\) is a smoothing operator. Applying \eqref{eq:9.42} and \eqref{eq:9.43} to \(S\), one obtains the existence of constants \(a > 0, b > 0, c > 0\) such that, for \(u\) and \(v\) in \(\mathcal{D}(\Omega)\), one has
\begin{equation}
    \label{eq:9.44}
    |(P \cdot u | v)_0| \leq c \norm{u}_{r/2} \norm{v}_{r/2}
\end{equation}
\begin{equation}
    \label{eq:9.45}
    (P \cdot u | u)_0 \geq a \norm{u}^2_{r/2} - b \norm{u}^2_0.
\end{equation}
The second one (for an even integer \(r\)) was first proved by Gårding in 1953 \cite{080}. It enabled him to apply the von Neumann spectral theory to the hermitian operator \(T_P\) in the Hilbert space \(L^2(\Omega)\), with \(\dom(T_P) = \mathcal{D}(\Omega)\). In general, the defects of \(T^{**}_P\) are both infinite, and one can define a particular \emph{self-adjoint} extension \(A_P\) of \(T_P\) by the following process: \(\dom(A_P)\) is the dense subspace of \(L^2(\Omega)\) consisting of functions \(u\) such that the \emph{distribution} \(P \cdot u\) is again in \(L^2(\Omega)\) and then \(A_P \cdot u = P \cdot u\); furthermore, \(\dom(A_P)\) is contained in the space \(H^{r/2}_0(\Omega)\), the closure in \(H^{r/2}(\R^n)\) of \(\mathcal{D}(\Omega)\). The spectrum of \(A_P\) is reduced to the point spectrum, consisting of an increasing sequence \((\lambda_n)\) of real eigenvalues of finite multiplicity, tending to \(+\infty\); the corresponding eigenfunctions (suitably normalized) form a \emph{Hilbert basis} of \(L^2(\Omega)\) and are of class \(C^\infty\); for any \(\zeta \in \C\) distinct from the \(\lambda_n, G_\zeta = (A_P - \zeta I)^{-1}\) is a \emph{compact} operator, which one may call the \emph{Green operator} of \(P - \zeta I\). It is easy to see that the restriction of \(Q_\zeta\) to \(\mathcal{D}(\Omega)\) is a pseudo-differential operator of order \(-r\), which in general is not of proper type; however, for every distribution \(T\) with compact support in \(\Omega\), one has \((P-\zeta I) \cdot (G_\zeta \cdot T) = G_\zeta \cdot ((P - \zeta I) \cdot T) = T\) and in particular, for any point \(x \in \Omega\),
\begin{equation}
    \label{eq:9.46}
    (P - \zeta I) \cdot (G_\zeta \cdot \varepsilon_x) = G_\zeta \cdot ((P - \zeta I) \cdot \varepsilon_x) = \varepsilon_x
\end{equation}
so that one may say that the distribution \(G_\zeta \cdot \varepsilon_x\) is an \emph{elementary solution} of \(P - \zeta I\) at the point \(x\).

IX) The results of VIII) apply in particular to a differential operator of even order \(2p \geq 2\)
\begin{equation}
    \label{eq:9.47}
    (P \cdot u)(x) = \sum_{|\alpha| \leq p, |\beta| \leq p} D^\alpha \left( a_{\alpha \beta}(x) D^\beta u(x) \right)
\end{equation}
where the \(a_{\alpha \beta}\) are \emph{bounded} \(C^\infty\) functions in a neighborhood of the \emph{bounded} set \(\overline{\Omega}\), such that:

1° \((-1)^{|\beta|} a_{\beta \alpha} = (-1)^{|\alpha|} \overline{a_{\alpha \beta}}\), which guarantee that \(P^* = P\);

2° there is a constant \(C > 0\) such that, for every \(x \in \Omega\) and every family \((z_\alpha)_{|\alpha|=p}\) of complex numbers, one has
\begin{equation}
    \label{eq:9.48}
    \sum_{|\alpha|=|\beta|=p} (-1)^p a_{\alpha \beta}(x) z_\alpha \overline{z_\beta}
    \geq
    C \left( \sum_{|\alpha|=p} |z_\alpha|^2 \right)
\end{equation}

The Green operator \(G_\zeta\) is then (for \(\zeta \in \spec(A_P)\)) an \emph{integral operator}, with kernel \((x,y) \mapsto \overline{G(\zeta,x,y)}\) which is locally integrable in \(\Omega \times \Omega, C^\infty\) outside of the diagonal and such that \(G(\overline{\zeta}, x, y) = G(\zeta, x, y)\) (the \emph{Green function} of \(P\)); from \eqref{eq:9.46} it follows that \(y \mapsto G(\zeta, x, y)\) is a \emph{solution} of \(P \cdot u = \zeta u\) in the complement \(\Omega - \{x\}\) of the point \(x\).

One may always take \(\zeta = -b\) for a sufficiently large number \(b > 0\), and for every function \(f \in \mathcal{E}(\Omega) \cap L^2(\Omega)\), there is therefore \emph{a unique solution of the equation} \(P \cdot u + bu = f\) \emph{belonging to the space} \(H^p_0(\Omega)\) \emph{and of class} \(C^\infty\) \emph{in} \(\Omega\).

These results were obtained (of course without the theory of pseudo-differential operators) by Gårding and Višik (independently) in 1953 (\cite{080}, \cite{218}); they may be considered as a ``weak'' solution of the generalization of Dirichlet's problem considered by E.E. Levi: \emph{no} assumption is made on the boundary \(\Gamma\) of \(\Omega\), but all which is required of the solution is that it should be arbitrarily close, \emph{for the topology of} \(H^p(\R^n)\), of \(C^\infty\) functions vanishing in a neighborhood of \(\Gamma\); but it (or its derivatives) may have a very pathological behavior at points of \(\Gamma\) if \(\Gamma\) is not smooth.

If one makes the additional assumption that
\[
    (P \cdot u)(x) = \sum_{|\alpha|=p, |\beta|=p} D^\alpha \left( a_{\alpha\beta}(x) D^\beta u(x) \right) + \sum_{|\nu| < p} D^\nu \left( a_\nu(x) D^\nu u(x) \right)
\]
where \((-1)^{|\nu|} a_\nu(x) \geq 0\) for all \(|\nu| < p\), then one may even take \(b = 0\) in the Gårding--Višik theorem (this is the case in particular for \((-Delta)^p\).

X) If \(E\) and \(F\) are two complex vector bundles over a compact differentiable manifold \(X\), and \(\Gamma(E), \Gamma(F)\) are the vector spaces of \(C^\infty\) sections of these bundles over \(X\), one can define pseudo-differential operators \(P: \Gamma(E) \to \Gamma(F)\), which become matrices of ordinary pseudo-differential operators when expressed in local coordinates. It is then possible to define intrinsically a \emph{principal symbol} \(\sigma^0_P\): for each \(x \in X\) and every tangent covector \(\xi\) to \(X\) at the point \(x, \sigma^0_P(x, \xi)\) is a homomorphism \(E_x \to F_x\) of the vector spaces, fibers of \(E\) and \(F\) at \(x\); in local coordinates, \(\sigma^0_P(x, \xi)\) is the matrix of the principal symbols of the elements of the matrix equal to \(P\). It is possible to define on \(\Gamma(E)\) and \(\Gamma(F)\) structures of prehilbert spaces, and to attach to any pseudo-differential operator \(P\) its adjoint \(P^*: \Gamma(F) \to \Gamma(E)\) such that \((P \cdot u|v) = (u | P^* \cdot v)\); properties \eqref{eq:9.40} and \eqref{eq:9.41} still hold.

A pseudo-differential operator \(P: \Gamma(E) \to \Gamma(E)\) is then called \emph{elliptic} if for every \(x \in X\) and every \(\xi \neq 0, \sigma^0_P(x, \xi)\) is a \emph{bijection} of \(E_x\) onto itself, and the existence of a \emph{parametrix} for such an operator can then be proved as in VI). For elliptic operators such that \(P^* = P\), the application of spectral theory to the hermitian operator \(T_P\) (in the Hilbert space, completion of \(\Gamma(E)\)) is here much simpler than in VIII) due to the absence of ``boundary conditions'': there is a Hilbert basis \((u_k)\) of \(\Gamma(E)\) such that \(P \cdot u_k = \mu_k u_k\) where \(\mu_k\) real and \(|\mu_k|\) tends to \(+\infty\) with \(k\); for every \(f \in \Gamma(E)\), one has \(f = \sum\limits_k (f | u_k) u_k\), the series being convergent for the topology of the Fréchet space \(\Gamma(E)\), and \(P \cdot f = \sum\limits_k \mu_k (f | u_k) u_k\) with the same convergence, in particular, \(\Ker(P)\) is the finite dimensional subspace having as a basis the \(u_k\), for which \(\mu_k\), and \(\operatorname{Im}(P)\) is closed and is a topological supplement of \(\Ker(P)\).

If now \(P\) is \emph{any} elliptic operator \(\Gamma(E) \to \Gamma(E), P^* P\) and \(PP^*\) are both elliptic and equal to their adjoints; the study of these operators enable one to evaluate the difference \(\dim(\Ker(P)) - \dim(\Ker(P^*))\), the \emph{index} of \(P\), and to express it by a formula in terms of the principal symbol of \(P\) and of the cohomology of \(X\); this is the famous \emph{Atiyah--Singer formula}, a fundamental result which has many applications and has spurred research in many directions (\cite{010}, \cite{031}, \cite{198}, \cite{199}). It was in fact due to the development of the necessary tools for the proof of that formula that the theory of pseudo-differential operators got started.

XI) The Gårding--Višik theorem of IX) leaves unanswered two questions: 1° Why is it that, in the Dirichlet problem and its generalizations, \emph{half} of the Cauchy data on the boundary are enough to determine the solution? 2° What can be said of the behavior of the unique solution belonging to \(H^p_0(\Omega)\) in the vicinity of a point of the boundary \(\Gamma\) where \(\Gamma\) is \emph{smooth}?

To answer these questions, one starts by investigating the Cauchy problem for an elliptic operator and seeing why in general it has no solution. Suppose that the bounded open set \(\Omega \subset \R^n\) has a smooth boundary \(\Gamma\), and (for simplicity's sake) that \(P\) is a differential operator defined in a neighborhood \(\Omega_0\) of \(\overline{\Omega}\), has even order \(2p \geq 2\), and possesses in \(\Omega_0\) a parametrix \(Q\) of order \(-2p\) such that \(Q \cdot (P \cdot T) = P \cdot (Q \cdot T)\) for any distribution \(T \in \mathcal{E}'(\Omega_0)\) (this is the case for the operator \(P - \zeta I\) in VIII), but here we \emph{do not} suppose that \(P^* = P\)). For any function \(u \in \mathcal{E}(\Omega)\), we note \(\Dch(u)\) the function \((g_0, g_1, \ldots, g_{2p-1})\) defined on the boundary \(\Gamma\) and with values in \(\C^{2p}\), where \(g_j\) is the normal derivative \(\pder{^j u}{n^j}\) at a point of \(\Gamma\). The starting point is an idea due to Sobolev (\cite{200}, p.~63): let \(u\) be the \emph{discontinuous} function equal to \(u\) in \(\overline{\Omega}\) but to 0 in the complement \(\Omega_0 - \overline{\Omega}\). Then \(P \cdot u^0\) is well defined as a \emph{distribution} on \(\Omega_0\), and it is easy to check that one can write
\begin{equation}
    \label{eq:9.49}
    P \cdot u^0 = (P \cdot u)^0 + N
\end{equation}
where \(N\) is a linear operator (independent of the function \(u\)) which to every \(C^\infty\) function in \((\mathcal{E}(\Gamma))^{2p}\) associates a distribution with \emph{support in} \(\Gamma\) (what one now calls a \emph{multilayer} on \(\Gamma\). As both sides of \eqref{eq:9.49} are distributions on \(\Omega_0\) with compact support, the operator \(Q\) may be applied to them, and yields the relation
\begin{equation}
    \label{eq:9.50}
    u^0 = Q \cdot ((P \cdot u)^0) + Q \cdot (N \cdot \Dch(u)).
\end{equation}
This is the general form of Green's formula \eqref{eq:2.17} of chap. \ref{ch:2}, \ref{sec:2.3}; for any function \(f \in \mathcal{E}(\Omega_0)\), the distribution \(Q \cdot f^0\) has a restriction to \(\Omega\) which is a \(C^\infty\) \emph{function} such that all its derivatives have \emph{limits} at every point of \(\Gamma\); one says that it is the \(Q\)-\emph{potential of the mass distribution} of density \(f\) on \(\Omega\). Similarly, for any vector function \(g \in (\mathcal{E}(\Gamma))^{2p}\), the distribution \(Q \cdot (N \cdot g)\) has the same properties, and one says that its restriction to \(\Omega\) is the \(Q\)-\emph{potential of the multilayer} \(N \cdot g\); these properties obviously generalize the classical properties of the newtonian potentials of a mass distribution, of a single layer and of a double layer (chap. \ref{ch:9}, \ref{sec:9.3}) (of course the restriction of \(Q \cdot (N \cdot q)\) to \(\Omega_0 - \overline{\Omega}\) also has similar properties, but the limits at a point of \(\Gamma\) differ in general from the limits at the same point of the restriction of \(Q \cdot (N \cdot g)\) to \(\Omega\)). Therefore, equation \eqref{eq:9.50} shows that if there is an \(u \in \mathcal{E}(\Omega)\) such that \(P \cdot u = f\) and \(\Dch(u) = g\) are given functions, this function \(u\) (restriction of \(u^0\)) is \emph{unique}, which corresponds to what one may expect of the Cauchy problem. But in addition one must have \(\Dch(u^0) = g\), which gives the \emph{necessary condition}
\begin{equation}
    \label{eq:9.51}
    g = \Dch(Q \cdot f^0) + \Dch(Q \cdot (N \cdot g))
\end{equation}
between \(f\) and \(g\). One proves that \(C: g \mapsto \Dch(Q \cdot (N \cdot g))\) is a pseudo-differential operator of \((\mathcal{E}(\Gamma))^{2p}\) into itself, which is called the \emph{Calderon operator} corresponding to the parametrix \(Q\).

A more detailed study shows that \eqref{eq:9.51} is equivalent to \(p\) linear relations between the \(2p\) functions \(g_0, \ldots, g_{2p-1}\) and their derivatives; this explains why one cannot prescribe the \(2p\) functions \(u, \pder{u}{n}, \ldots, \pder{^{2p-1}}{n^{2p-1}}\) on \(\Gamma\), but only \(p\) of them. More generally, one may consider a differential operator \(B\) of \((\mathcal{E}(\Gamma))^{2p}\) into \((\mathcal{E}(\Gamma))^p\), and instead of the Cauchy problem, consider the boundary conditions \(B \cdot (\Dch(u)) = g\) for a given vector function \(g \in (\mathcal{E}(\Gamma))^p\). It is then possible to describe explicitly a set of sufficient conditions (called the \emph{Lopatinski conditions}) linking \(B\) and \(C\) and implying that the problem can be reduced to Fredholm integral equations on \(\Gamma\); more precisely, these conditions imply that the mapping
\begin{equation}
    \label{eq:9.52}
    u \mapsto (P \cdot u, B \cdot (\Dch(h)))
\end{equation}
of \(\mathcal{E}(\overline{\Omega})\) (space of the restrictions to \(\overline{\Omega}\) of functions of \(\mathcal{E}(\Omega_0)\)) into \(\mathcal{E}(\overline{\Omega}) \times (\mathcal{E}(\Gamma))^p\) has \emph{a finite dimensional kernel and a closed image of finite codimension}.

In particular, one checks that the Lopatinski conditions are always satisfied if one takes for \(B \cdot g\) a consecutive sequence \((g_q, g_{q+1}, \ldots, g_{q+p-1})\) of \(p\) of the functions \(g_j\), and the corresponding problem for \(q = 0\) is just the Dirichlet problem as posed by E.E. Levi.

At this point, one might think, from the example \eqref{eq:9.47}, that except for a \emph{denumerable} set of values of \(\zeta \in \C\), the mapping \(u \mapsto ((P - \zeta I) \cdot u, B\cdot (\Dch(u)))\) would in fact be \emph{bijective}. However, this is \emph{not} always the case, and there are examples for which that mapping is injective for \emph{no} \(\zeta \in \C\).

For operators \eqref{eq:9.47} to which the Gårding--Višik theorem applies, to prove that the preceding mapping is bijective for \(\zeta \in \spec(A_P)\) (with \(B \cdot g = (g_0, \ldots, g_{p-1})\)), it is enough to show that, when \(\Gamma\) is smooth, the unique solution \(u \in H^p_0(\Omega)\), and all its derivatives, can be extended by continuity to \(\overline{\Omega} = \Omega \cup \Gamma\) (the second problem mentioned above). Actually, even if \(\Gamma\) is not smooth everywhere, the existence of limits for these functions is guaranteed at \emph{each point} where \(\Gamma\) is smooth; this was first proved by L. Nirenberg in 1955 \cite{167}, and has been proved by Peetre in 1961 using a different method, which however still relies on \emph{a priori} inequalities \cite{171}. These results may be extended to other types of boundary conditions \(B \cdot (\Dch(u)) = g\) satisfying the Lopatinski conditions, the so-called \emph{coercitive} problems for elliptic operators \(P\) for which \(P^* = P\).

\noindent \emph{Further generalizations}. Formula \eqref{eq:9.37} defining a pseudo-differential operator can also be written, replacing \(\mathcal{F}u\) by its definition
\[
    (P \cdot u)(x) = \widetilde{\iint}_{\Omega \times \R^n} \exp(2\pi i(x - y | \xi)) a(x, \xi) u(y) dy d\xi
\]
where the integral is not any more a Lebesgue integral, but an ``improper'' (or ``oscillating'') one, obtained by passage to the limit from the integral of the same function multiplied by a function \(h(\xi/q)\), where \(h \in \mathcal{D}(\R^n)\) is equal to 1 in a neighborhood of 0, and \(q\) tends to \(+\infty\). It turns out that one can define similar integrals when one replaces \(\exp(2\pi i(x-y|\xi))\) by ``phase functions'' \(\varphi(x, y, \xi)\) positively homogeneous in \(\xi\), and \(a(x, \xi)\) by more general ``symbols'' \(a(x,y,\xi)\).

I) Such operators naturally occur in the theory of \emph{strictly hyperbolic operators}, of which the simplest is the wave operator (or \emph{dalembertian})
\begin{equation}
    \label{eq:9.53}
    \square u = \pder{^2 u}{t^2} - \left( \pder{^2 u}{x^2_1} + \ldots + \pder{^2 u}{x^2_n} \right).
\end{equation}
The Cauchy problem for that operator, consisting in finding a solution of \(\square u = 0\) such that \(u(0,x) = g_0(x)\) and \(\pder{u}{t}(0, x) = g_1(x)\) are given functions, had already been solved by Cauchy; the explicit formula he gave for the solution can be written \(u(t,x) = u_+(t,x) + u_-(t,x)\), where
\begin{equation}
    \label{eq:9.54}
    u_\pm(t,x) = \frac{1}{2} \iint \exp(2\pi i((x-y|\xi) \pm (|\xi|t))) \left( g_0(y) \pm \frac{1}{2\pi i} \frac{g_1(y)}{|\xi|} \right) dy d\xi
\end{equation}
the integrals being ``improper'' in a sense easy to describe. In general, one considers an operator of order \(m\) in \(n+1\) variables \(t, x_1, \ldots, x_n\)
\begin{equation}
    \label{eq:9.55}
    P \cdot u = \pder{^m u}{t^m} + \sum_{j=1}^m \sum_{|\alpha| \leq j} c_{j \alpha}(t,x) D^\alpha_x \left( \pder{^{m-j} u}{t^{m-j}} \right)
\end{equation}
and one assumes that its principal symbol \(\sigma_P(\tau, \xi, t, x)\) be written
\begin{equation}
    \label{eq:9.56}
    \sigma^0_P(\tau, \xi, t, x) = \prod_{j=0}^{m-1} \left( \tau - q_j(t, x, \xi) \right)
\end{equation}
where the \(q_j\) are real functions of class \(C^\infty\) in \(I \times \Omega \times (\R^n - \{0\})\) (\(I\) open subset of \(\R\), \(\Omega\) open subset of \(\R^n\)), positively homogeneous or degree 1 in \(\xi\), \emph{and such that for} \(j \neq k, q_j(t,x,\xi) \neq q_k(t,x,\xi)\) \emph{everywhere}. The Cauchy problem to be solved is to find a function \(v(t,x)\) such that \(P \cdot v = f\) and \(\pder{^j}{t^j} v(t_0, x) = g_j(x) (0 \leq j \leq m-1)\) for a \(t_0 \in I\), in a convenient neighborhood of \((t_0,x_0) \in I \times \Omega, f\) and \(g\) being \(C^\infty\) functions. Taking \eqref{eq:9.54} as a model, one introduces \(m^2\) operators \(E_{jh}(s) (0 \leq j, h \leq m-1)\) for s in a neighborhood of \(t_0\) in \(I\), such that, if \(E_h(s) = \sum\limits_{j=0}^{m-1} E_{jh}(s)\) for \(0 \leq h \leq m-1\), one has (locally)
\begin{equation}
    \label{eq:9.57}
    PE_h(s) = R_h(s), \quad \left(\pder{}{t}\right)^k E_h(s) = \delta_{hk}I \text{ for } t=s \text{ and } 0 \leq k \leq m-1,
\end{equation}
where the \(R_h(s)\) are \emph{smoothing} operators. If one writes
\[
    (L \cdot u)(t, x) = \int_{t_0}^t \left( E_{m-1}(s) \cdot u(s, \cdot) \right)(t,x) ds
\]
one has \(\left(\pder{}{t}\right)^k (L \cdot u)(t_0, x) = 0\) for \(0 \leq k \leq m-1\), and \(P \cdot (L \cdot u) = u - V \cdot u\), where \(V\) is a \emph{Volterra integral operator}
\begin{equation}
    \label{eq:9.58}
    (V \cdot u)(t, x) = \int_{t_0}^t ds \int_U K(t,s,x,y)u(s,y)dy
\end{equation}
\(U\) being a neighborhood of \(x_0\) in \(\Omega\) and \(K\) a \(C^\infty\) function. It is easy to see that \(I + V\) is inverted by \(I + W\), where \(W\) is a similar Volterra operator. The Cauchy problem is then solved by taking in a sufficiently small neighborhood of \((t_0, x_0)\)
\begin{equation}
    \label{eq:9.59}
    v = \sum_{j=0}^{m-1} E_j(t_0) \cdot g_j + L \cdot \left( (I+W) \cdot \left( f - \sum_{j=0}^{m-1} R_j(t_0) \cdot g_j \right) \right).
\end{equation}
The construction of the \(E_h\) follows an idea introduced by P. Lax in 1957, and patterned after the known behavior of the solutions of the wave equation, which ``propagate'' along ``rays''. For operators \eqref{eq:9.25} with analytic coefficients, it follows from the Cauchy--Kowalewska theorem that the Cauchy problem for data given on a hypersurface will fail to have a unique solution if the hypersurface is given locally by an equation \(z(x_1, \ldots, x_n) =\) const., where \(z\) is a solution of the partial differential equation of order 1
\begin{equation}
    \label{eq:9.60}
    \Phi\left(x_1, \ldots, x_n, \pder{z}{x_1}, \ldots, \pder{z}{x_n}\right) = 0
\end{equation}
in which \(\Phi\) is obtained from the principal symbol \(\sigma^0_P(x, \xi)\) by replacing the vector \(\xi\) by \(\left(\pder{z}{x_1}, \ldots, \pder{z}{x_n}\right)\); such hypersurfaces are called \emph{characteristic} for the operator \(P\). For strictly hyperbolic operators \eqref{eq:9.55}, it follows from \eqref{eq:9.56} that the equation of characteristic hypersurfaces \eqref{eq:9.60} splits into \(m\) equations
\begin{equation}
    \label{eq:9.61}
    \pder{z}{t} - q_j(t, x, \grad_x{z}) = 0, \qquad 0 \leq j \leq m-1
\end{equation}
(with \(\grad_x{z} = \left(\pder{z}{x_1}, \ldots, \pder{z}{x_n}\right)\)). For the wave operator \eqref{eq:9.53} the equations \eqref{eq:9.61} are
\[
    \pder{z}{t} = \pm |\grad_x{z}|
\]
with solutions
\[
    z = 2\pi ((x | \xi) \pm |\xi|t)
\]
which reduce to \(2\pi(x|\xi)\) for \(t = 0\); they are precisely the ``phase functions'' which enter in Cauchy's formula \eqref{eq:9.54}. For general strictly hyperbolic operators \eqref{eq:9.54}, one therefore introduces the \(m^2\) operators \(F_{jh}(s)\) defined by
\begin{equation}
    \label{eq:9.62}
    (F_{jh}(s) \cdot u)(t,x) =
    \iint_{U \times \R^n} \exp\left( i(\psi_j(t,s,x,\xi) - 2\pi(y|\xi)) \right) a_{jh}(t,s,x,y,\xi) u(y) dy d\xi
\end{equation}
where \(\psi_j(t,s,x,\xi)\) is the unique solution of \eqref{eq:9.61} satisfying the initial condition
\begin{equation}
    \label{eq:9.63}
    \psi_j(s,s,x,\xi) = 2\pi(x | \xi)
\end{equation}
and \(a_{jh}\) is a symbol of order \(-h\) (in the sense defined for pseudo-differential operators). The goal is to determine the \(a_{jh}\) in such a way that, if one writes \(F_h(s) = \sum\limits_{j=0}^{m-1} F_{jh}(s)\), the following conditions are satisfied:

1° \(Q_h(s) = P \circ F_h(s)\) is a smoothing operator;

2° for each \(g \in \mathcal{D}(U)\), the restriction to the hyperplane \(t=s\) of the function \(\pder{^k}{t^k} \left(F_h(s) \cdot g\right)\) has the form \((s,x) \mapsto \delta_{hk} g(x) + \left(Q_{hk}(s) \cdot g\right)(x)\) where \(Q_{hk}(s)\) is a smoothing operator, for \(0 \leq h, k \leq m-1\).

The conditions \eqref{eq:9.57} are then met by taking
\[
    (R_h(s) \cdot g)(t,x) = \sum_{k=0}^{m-1} \frac{(t-s)^k}{k!} (Q_{hk}(s) \cdot g)(x)
\]
and
\[
    E_h(s) = F_h(s) - R_h(s).
\]

To achieve that goal, one defines the \(a_{jh}\) by \emph{asymptotic expansions}
\[
    a_{jh} \sim{} \sum_{l=0}^\infty a^{(l)}_{jh}
\]
where \(a^{(l)}_{jh}\) is a symbol of order \(-h-l\); the \(a^{(l)}_{jh}\) are determined by \emph{induction on} \(l\), in such a way that if one writes
\[
    (F_{jhN}(s) \cdot u)(t,x) =
    \sum_{l=0}^N \widetilde{\iint}_{U \times \R^n} \exp(i(\psi_j(t,s,x,\xi) - 2\pi(y|\xi))) a^{(l)}_{jh}(t,s,x,y,\xi) u(y) dy d\xi
\]
then:

1° Each operator \(P \circ F_{jhN}(s)\) is defined by a symbol of order \(m-h-N-2\).

2° If \(F_{hN}(s) = \sum\limits_{j=0}^{m-1} F_{jhN}(s)\), the restriction to \(t=s\) of the function \(\pder{^k}{t^k}(F_{hN}(s) \cdot g)\) has the form \((s, x) \mapsto \delta_{hk} g(x) + (Q_{hkN}(s) \cdot g)(x)\), where \(Q_{hkN}(s)\) is a pseudo-differential operator of order \(k-h-N-1\), for \(0 \leq h,k \leq m-1\).

It is in this inductive process that the analogs of the ``rays'' enter. The classical Cauchy method for integration of partial differential equations of order 1 consists, for each equation \eqref{eq:9.61}, in considering, in the space \(I \times \Omega \times \R^{n+1}\), the ``characteristic'' curves
\[
    t \mapsto (t, x_1(t), \ldots, x_n(t), p_0(t), p_1(t), \ldots, p_n(t))
\]
solutions of the system of ordinary differential equations
\begin{equation}
    \label{eq:9.65}
    \begin{dcases}
        \der{x_k}{t} = -\pder{q_j}{\xi_k} (t, x_1, \ldots, x_n, p_1, \ldots, p_n) \\
        \der{p_k}{t} = \pder{q_j}{x_k} (t, x_1, \ldots, x_n, p_1, \ldots, p_n)
    \end{dcases}
    \qquad
    1 \leq k \leq n
\end{equation}
and verifying the condition \(p_0 = q_j(t, x_1, \ldots, x_n, p_1, \ldots, p_n)\); one says that their projections \(t \mapsto (t, x_1(t), \ldots, x_n(t))\) on \(I \times \Omega\) constitute the \(j\)-th family of \emph{bicharacteristic curves} for the operator \(P\). For the wave operator, the bicharacteristic curves which are such that \(x_k(s) = y_k\) for \(1 \leq k \leq n\) are in fact the classical ``rays''
\[
    t \mapsto y_k \pm \frac{\xi_k}{|\xi_k|}(t-s) \qquad (1 \leq k \leq n).
\]

In the general theory, each \(a^{(N)}_{jh}(t,s,x,\xi)\) is taken independent of \(y\), and is obtained by \emph{integrating, along each bicharacteristic curve, an ordinary linear differential equation of the first order} (in the variable \(t\)), whose coefficients are determined when the \(a^{(l)}_{jh}\) are known for \(l \leq N-1\); finally the induction starts with the values of the \(a^{(0)}_{jh}(s,s,x,\xi)\), which are given by the linear system
\begin{equation}
    \label{eq:9.66}
    \sum_{j=0}^{m-1} (i q'_j(s,x,\xi))^k a^{0}_{jh}(s,s,x,\xi) = \delta_{jh} \qquad (0 \leq k \leq m-1)
\end{equation}
where \(q'_j(s,x,\xi) = q_j(s,x,\grad_x{\psi_j(s,s,x,\xi)}\)); the determinant of that system is \(\neq 0\) because the \(q_j\) have been supposed to be \emph{distinct}.

One of the consequences of this remarkable construction is that, in the explicit formula \eqref{eq:9.59} for \(f = 0\), one may replace the ``initial conditions'' \(g_j\) by arbitrary \emph{distributions} \(S_j\) on \(\Omega\); \(v\) is then replaced by a distribution \(T\) solution of \(P \cdot T = 0\). For these equations, the ``trace'' \(T_t\) of such a distribution on the hyperplane \(\{t\} \times \Omega\) may be defined (although \(T\) is not a function) as a distribution \emph{on} \(\Omega\), varying with \(t\), and which may be said to ``propagate'' with the ``time'' \(t\), starting from the ``initial'' distribution \(T_{t_0} = S_0\). It is then possible to show that the singular support of \(T_t\) is contained in a set \(M_t\) obtained in the following way: one considers the union \(M_0\) of the singular supports of the distributions \(S_j\) and all the bicharacteristics issued from points of \(M_0\); \(M_t\) is the set of all points on these bicharacteristics at time \(t\). This gives a precise meaning to a phenomenon which had been well known for second order strictly hyperbolic equations, and particular types of ``initial values'': \emph{the singularities propagate along the bicharacteristics}.

Under additional assumptions, it is possible to extend these results when in the decomposition \eqref{eq:9.56}, some of the \(q_j\) are equal \cite{041}.

II) Another important application of operators generalizing the pseudo-differential operators is the problem of \emph{local existence} of solutions for a partial differential equation \(P \cdot u = f\), which has stimulated much research after H. Lewy's discovery (chap. \ref{ch:2}, \ref{sec:2.2}). Over a period of more than 15 years, the combined efforts of Hörmander, Nirenberg and Trèves succeeded in formulating a system of conditions which were finally proved to be necessary and sufficient for local existence by Beals and Fefferman \cite{020}, using new types of operators \cite{019}. We cannot here do more than refer the reader to these papers.

\nocite{*}
\printbibliography[title=References,notkeyword={noprint}]
\markboth{References}{References}
\fancyhead[C]{\textsc{\leftmark}}

\end{document}
